% Options for packages loaded elsewhere
\PassOptionsToPackage{unicode}{hyperref}
\PassOptionsToPackage{hyphens}{url}
\PassOptionsToPackage{dvipsnames,svgnames,x11names}{xcolor}
%
\documentclass[
  letterpaper,
  DIV=11,
  numbers=noendperiod]{scrartcl}

\usepackage{amsmath,amssymb}
\usepackage{iftex}
\ifPDFTeX
  \usepackage[T1]{fontenc}
  \usepackage[utf8]{inputenc}
  \usepackage{textcomp} % provide euro and other symbols
\else % if luatex or xetex
  \usepackage{unicode-math}
  \defaultfontfeatures{Scale=MatchLowercase}
  \defaultfontfeatures[\rmfamily]{Ligatures=TeX,Scale=1}
\fi
\usepackage{lmodern}
\ifPDFTeX\else  
    % xetex/luatex font selection
\fi
% Use upquote if available, for straight quotes in verbatim environments
\IfFileExists{upquote.sty}{\usepackage{upquote}}{}
\IfFileExists{microtype.sty}{% use microtype if available
  \usepackage[]{microtype}
  \UseMicrotypeSet[protrusion]{basicmath} % disable protrusion for tt fonts
}{}
\makeatletter
\@ifundefined{KOMAClassName}{% if non-KOMA class
  \IfFileExists{parskip.sty}{%
    \usepackage{parskip}
  }{% else
    \setlength{\parindent}{0pt}
    \setlength{\parskip}{6pt plus 2pt minus 1pt}}
}{% if KOMA class
  \KOMAoptions{parskip=half}}
\makeatother
\usepackage{xcolor}
\setlength{\emergencystretch}{3em} % prevent overfull lines
\setcounter{secnumdepth}{-\maxdimen} % remove section numbering
% Make \paragraph and \subparagraph free-standing
\makeatletter
\ifx\paragraph\undefined\else
  \let\oldparagraph\paragraph
  \renewcommand{\paragraph}{
    \@ifstar
      \xxxParagraphStar
      \xxxParagraphNoStar
  }
  \newcommand{\xxxParagraphStar}[1]{\oldparagraph*{#1}\mbox{}}
  \newcommand{\xxxParagraphNoStar}[1]{\oldparagraph{#1}\mbox{}}
\fi
\ifx\subparagraph\undefined\else
  \let\oldsubparagraph\subparagraph
  \renewcommand{\subparagraph}{
    \@ifstar
      \xxxSubParagraphStar
      \xxxSubParagraphNoStar
  }
  \newcommand{\xxxSubParagraphStar}[1]{\oldsubparagraph*{#1}\mbox{}}
  \newcommand{\xxxSubParagraphNoStar}[1]{\oldsubparagraph{#1}\mbox{}}
\fi
\makeatother

\usepackage{color}
\usepackage{fancyvrb}
\newcommand{\VerbBar}{|}
\newcommand{\VERB}{\Verb[commandchars=\\\{\}]}
\DefineVerbatimEnvironment{Highlighting}{Verbatim}{commandchars=\\\{\}}
% Add ',fontsize=\small' for more characters per line
\usepackage{framed}
\definecolor{shadecolor}{RGB}{241,243,245}
\newenvironment{Shaded}{\begin{snugshade}}{\end{snugshade}}
\newcommand{\AlertTok}[1]{\textcolor[rgb]{0.68,0.00,0.00}{#1}}
\newcommand{\AnnotationTok}[1]{\textcolor[rgb]{0.37,0.37,0.37}{#1}}
\newcommand{\AttributeTok}[1]{\textcolor[rgb]{0.40,0.45,0.13}{#1}}
\newcommand{\BaseNTok}[1]{\textcolor[rgb]{0.68,0.00,0.00}{#1}}
\newcommand{\BuiltInTok}[1]{\textcolor[rgb]{0.00,0.23,0.31}{#1}}
\newcommand{\CharTok}[1]{\textcolor[rgb]{0.13,0.47,0.30}{#1}}
\newcommand{\CommentTok}[1]{\textcolor[rgb]{0.37,0.37,0.37}{#1}}
\newcommand{\CommentVarTok}[1]{\textcolor[rgb]{0.37,0.37,0.37}{\textit{#1}}}
\newcommand{\ConstantTok}[1]{\textcolor[rgb]{0.56,0.35,0.01}{#1}}
\newcommand{\ControlFlowTok}[1]{\textcolor[rgb]{0.00,0.23,0.31}{\textbf{#1}}}
\newcommand{\DataTypeTok}[1]{\textcolor[rgb]{0.68,0.00,0.00}{#1}}
\newcommand{\DecValTok}[1]{\textcolor[rgb]{0.68,0.00,0.00}{#1}}
\newcommand{\DocumentationTok}[1]{\textcolor[rgb]{0.37,0.37,0.37}{\textit{#1}}}
\newcommand{\ErrorTok}[1]{\textcolor[rgb]{0.68,0.00,0.00}{#1}}
\newcommand{\ExtensionTok}[1]{\textcolor[rgb]{0.00,0.23,0.31}{#1}}
\newcommand{\FloatTok}[1]{\textcolor[rgb]{0.68,0.00,0.00}{#1}}
\newcommand{\FunctionTok}[1]{\textcolor[rgb]{0.28,0.35,0.67}{#1}}
\newcommand{\ImportTok}[1]{\textcolor[rgb]{0.00,0.46,0.62}{#1}}
\newcommand{\InformationTok}[1]{\textcolor[rgb]{0.37,0.37,0.37}{#1}}
\newcommand{\KeywordTok}[1]{\textcolor[rgb]{0.00,0.23,0.31}{\textbf{#1}}}
\newcommand{\NormalTok}[1]{\textcolor[rgb]{0.00,0.23,0.31}{#1}}
\newcommand{\OperatorTok}[1]{\textcolor[rgb]{0.37,0.37,0.37}{#1}}
\newcommand{\OtherTok}[1]{\textcolor[rgb]{0.00,0.23,0.31}{#1}}
\newcommand{\PreprocessorTok}[1]{\textcolor[rgb]{0.68,0.00,0.00}{#1}}
\newcommand{\RegionMarkerTok}[1]{\textcolor[rgb]{0.00,0.23,0.31}{#1}}
\newcommand{\SpecialCharTok}[1]{\textcolor[rgb]{0.37,0.37,0.37}{#1}}
\newcommand{\SpecialStringTok}[1]{\textcolor[rgb]{0.13,0.47,0.30}{#1}}
\newcommand{\StringTok}[1]{\textcolor[rgb]{0.13,0.47,0.30}{#1}}
\newcommand{\VariableTok}[1]{\textcolor[rgb]{0.07,0.07,0.07}{#1}}
\newcommand{\VerbatimStringTok}[1]{\textcolor[rgb]{0.13,0.47,0.30}{#1}}
\newcommand{\WarningTok}[1]{\textcolor[rgb]{0.37,0.37,0.37}{\textit{#1}}}

\providecommand{\tightlist}{%
  \setlength{\itemsep}{0pt}\setlength{\parskip}{0pt}}\usepackage{longtable,booktabs,array}
\usepackage{calc} % for calculating minipage widths
% Correct order of tables after \paragraph or \subparagraph
\usepackage{etoolbox}
\makeatletter
\patchcmd\longtable{\par}{\if@noskipsec\mbox{}\fi\par}{}{}
\makeatother
% Allow footnotes in longtable head/foot
\IfFileExists{footnotehyper.sty}{\usepackage{footnotehyper}}{\usepackage{footnote}}
\makesavenoteenv{longtable}
\usepackage{graphicx}
\makeatletter
\def\maxwidth{\ifdim\Gin@nat@width>\linewidth\linewidth\else\Gin@nat@width\fi}
\def\maxheight{\ifdim\Gin@nat@height>\textheight\textheight\else\Gin@nat@height\fi}
\makeatother
% Scale images if necessary, so that they will not overflow the page
% margins by default, and it is still possible to overwrite the defaults
% using explicit options in \includegraphics[width, height, ...]{}
\setkeys{Gin}{width=\maxwidth,height=\maxheight,keepaspectratio}
% Set default figure placement to htbp
\makeatletter
\def\fps@figure{htbp}
\makeatother

\KOMAoption{captions}{tableheading}
\usepackage{amsmath}
\usepackage{amssymb}
\makeatletter
\@ifpackageloaded{caption}{}{\usepackage{caption}}
\AtBeginDocument{%
\ifdefined\contentsname
  \renewcommand*\contentsname{Table of contents}
\else
  \newcommand\contentsname{Table of contents}
\fi
\ifdefined\listfigurename
  \renewcommand*\listfigurename{List of Figures}
\else
  \newcommand\listfigurename{List of Figures}
\fi
\ifdefined\listtablename
  \renewcommand*\listtablename{List of Tables}
\else
  \newcommand\listtablename{List of Tables}
\fi
\ifdefined\figurename
  \renewcommand*\figurename{Figure}
\else
  \newcommand\figurename{Figure}
\fi
\ifdefined\tablename
  \renewcommand*\tablename{Table}
\else
  \newcommand\tablename{Table}
\fi
}
\@ifpackageloaded{float}{}{\usepackage{float}}
\floatstyle{ruled}
\@ifundefined{c@chapter}{\newfloat{codelisting}{h}{lop}}{\newfloat{codelisting}{h}{lop}[chapter]}
\floatname{codelisting}{Listing}
\newcommand*\listoflistings{\listof{codelisting}{List of Listings}}
\makeatother
\makeatletter
\makeatother
\makeatletter
\@ifpackageloaded{caption}{}{\usepackage{caption}}
\@ifpackageloaded{subcaption}{}{\usepackage{subcaption}}
\makeatother

\ifLuaTeX
  \usepackage{selnolig}  % disable illegal ligatures
\fi
\usepackage{bookmark}

\IfFileExists{xurl.sty}{\usepackage{xurl}}{} % add URL line breaks if available
\urlstyle{same} % disable monospaced font for URLs
\hypersetup{
  pdftitle={Distributions (m.dist)},
  colorlinks=true,
  linkcolor={blue},
  filecolor={Maroon},
  citecolor={Blue},
  urlcolor={Blue},
  pdfcreator={LaTeX via pandoc}}


\title{Distributions (\texttt{m.dist})}
\author{}
\date{}

\begin{document}
\maketitle


Every distribution can define a latent parameter (\texttt{name=}), a
likelihood (\texttt{obs=}), draw samples directly (\texttt{sample=True})
or return the distribution object (\texttt{create\_obj=True}). In R use
\texttt{bf.dist.\textless{}name\textgreater{}()}.

\subsection{\texorpdfstring{\texttt{m.dist.asymmetric\_laplace}}{m.dist.asymmetric\_laplace}}\label{asymmetric_laplace}

\begin{Shaded}
\begin{Highlighting}[]
\NormalTok{m.dist.asymmetric\_laplace(}
\NormalTok{    loc}\OperatorTok{=}\FloatTok{0.0}\NormalTok{,}
\NormalTok{    scale}\OperatorTok{=}\FloatTok{1.0}\NormalTok{,}
\NormalTok{    asymmetry}\OperatorTok{=}\FloatTok{1.0}\NormalTok{,}
\NormalTok{    validate\_args}\OperatorTok{=}\VariableTok{None}\NormalTok{,}
\NormalTok{    name}\OperatorTok{=}\StringTok{\textquotesingle{}x\textquotesingle{}}\NormalTok{,}
\NormalTok{    obs}\OperatorTok{=}\VariableTok{None}\NormalTok{,}
\NormalTok{    mask}\OperatorTok{=}\VariableTok{None}\NormalTok{,}
\NormalTok{    infer}\OperatorTok{=}\VariableTok{None}\NormalTok{,}
\NormalTok{    sample}\OperatorTok{=}\VariableTok{False}\NormalTok{,}
\NormalTok{    seed}\OperatorTok{=}\VariableTok{None}\NormalTok{,}
\NormalTok{    shape}\OperatorTok{=}\NormalTok{(),}
\NormalTok{    event}\OperatorTok{=}\DecValTok{0}\NormalTok{,}
\NormalTok{    create\_obj}\OperatorTok{=}\VariableTok{False}\NormalTok{,}
\NormalTok{    to\_jax}\OperatorTok{=}\VariableTok{True}\NormalTok{,}
\NormalTok{)}
\end{Highlighting}
\end{Shaded}

\textbf{Asymmetric Laplace} Samples from an Asymmetric Laplace
distribution.

The Asymmetric Laplace distribution is a generalization of the Laplace
distribution, where the two sides of the distribution are scaled
differently. It is defined by a location parameter (loc), a scale
parameter (scale), and an asymmetry parameter (asymmetry).

\[
f(x, \kappa) = \begin{cases} \frac{1}{\kappa+\kappa^{-1}}\exp(-x\kappa),\quad x\ge0, \\ \frac{1}{\kappa+\kappa^{-1}}\exp(x/\kappa),\quad x<0 \end{cases}
\]

for \(-\infty < x < \infty\), \(\kappa > 0\).

laplace\_asymmetric takes `kappa' as a shape parameter for \(\kappa\).
For \(\kappa = 1\), it is identical to a Laplace distribution.

\textbf{Args} - \emph{loc} (jnp.ndarray or float): Location parameter of
the distribution.

\begin{itemize}
\item
  \emph{scale} (jnp.ndarray or float): Scale parameter of the
  distribution.
\item
  \emph{asymmetry} (jnp.ndarray or float): Asymmetry parameter of the
  distribution.
\item
  \emph{shape} (tuple): A multi-purpose argument for shaping. When
  \texttt{sample=False} (model building), this is used with
  \texttt{shape} to set the distribution's batch shape. When
  \texttt{sample=True} (direct sampling), this is used as
  \texttt{sample\_shape} to draw a raw JAX array of the given shape.
\item
  \emph{event} (int): The number of batch dimensions to reinterpret as
  event dimensions (used in model building).
\item
  \emph{mask} (jnp.ndarray, bool): Optional boolean array to mask
  observations.
\item
  \emph{create\_obj} (bool): If True, returns the raw BF distribution
  object instead of creating a sample site. This is essential for
  building complex distributions like \texttt{MixtureSameFamily}.
\item
  \emph{sample} (bool, optional): A control-flow argument. If
  \texttt{True}, the function will directly sample a raw JAX array from
  the distribution, bypassing the BF model context. If \texttt{False},
  it will create a \texttt{BF.sample} site within a model. Defaults to
  \texttt{False}.
\item
  \emph{obs} (jnp.ndarray, optional): The observed value for this random
  variable. If provided, the sample site is conditioned on this value,
  and the function returns the observed value. If \texttt{None}, the
  site is treated as a latent (unobserved) random variable. Defaults to
  \texttt{None}.
\item
  \emph{name} (str, optional): The name of the sample site in a BF
  model. This is used to uniquely identify the random variable. Defaults
  to `x'.
\item
  validate\_args (bool, optional): Whether to enable validation of
  distribution parameters. Defaults to \texttt{None}.
\end{itemize}

\textbf{Returns} - When \texttt{sample=False}: A BF AsymmetricLaplace
distribution object (for model building).

\begin{itemize}
\item
  When \texttt{sample=True}: A JAX array of samples drawn from the
  AsymmetricLaplace distribution (for direct sampling).
\item
  When \texttt{create\_obj=True}: The raw BF AsymmetricLaplace
  distribution object (for advanced use cases).
\end{itemize}

\textbf{Example Usage} from BayesForge import bf m = bf(platform=`cpu')
m.dist.asymmetric\_laplace(loc=0.0, scale=1.0, asymmetry=1.0,
sample=True)

\textbf{Wrapper of}
https://num.pyro.ai/en/stable/distributions.html\#asymmetriclaplace

\textbf{References} 1. ``Asymmetric Laplace distribution'', Wikipedia
https://en.wikipedia.org/wiki/Asymmetric\_Laplace\_distribution

\begin{enumerate}
\def\labelenumi{\arabic{enumi}.}
\setcounter{enumi}{1}
\tightlist
\item
  Kozubowski TJ and Podgórski K. A Multivariate and Asymmetric
  Generalization of Laplace Distribution, Computational Statistics 15,
  531--540 (2000).
  doi:\href{https://doi.org/10.1007/PL00022717}{10.1007/PL00022717}
\end{enumerate}

\subsection{\texorpdfstring{\texttt{m.dist.asymmetric\_laplace\_quantile}}{m.dist.asymmetric\_laplace\_quantile}}\label{asymmetric_laplace_quantile}

\begin{Shaded}
\begin{Highlighting}[]
\NormalTok{m.dist.asymmetric\_laplace\_quantile(}
\NormalTok{    loc}\OperatorTok{=}\FloatTok{0.0}\NormalTok{,}
\NormalTok{    scale}\OperatorTok{=}\FloatTok{1.0}\NormalTok{,}
\NormalTok{    quantile}\OperatorTok{=}\FloatTok{0.5}\NormalTok{,}
\NormalTok{    validate\_args}\OperatorTok{=}\VariableTok{None}\NormalTok{,}
\NormalTok{    name}\OperatorTok{=}\StringTok{\textquotesingle{}x\textquotesingle{}}\NormalTok{,}
\NormalTok{    obs}\OperatorTok{=}\VariableTok{None}\NormalTok{,}
\NormalTok{    mask}\OperatorTok{=}\VariableTok{None}\NormalTok{,}
\NormalTok{    infer}\OperatorTok{=}\VariableTok{None}\NormalTok{,}
\NormalTok{    sample}\OperatorTok{=}\VariableTok{False}\NormalTok{,}
\NormalTok{    seed}\OperatorTok{=}\VariableTok{None}\NormalTok{,}
\NormalTok{    shape}\OperatorTok{=}\NormalTok{(),}
\NormalTok{    event}\OperatorTok{=}\DecValTok{0}\NormalTok{,}
\NormalTok{    create\_obj}\OperatorTok{=}\VariableTok{False}\NormalTok{,}
\NormalTok{    to\_jax}\OperatorTok{=}\VariableTok{True}\NormalTok{,}
\NormalTok{)}
\end{Highlighting}
\end{Shaded}

\textbf{Asymmetric Laplace Quantile} Samples from an
AsymmetricLaplaceQuantile distribution.

This distribution is an alternative parameterization of the
AsymmetricLaplace distribution, commonly used in Bayesian quantile
regression. It utilizes a \texttt{quantile} parameter to define the
balance between the left- and right-hand sides of the distribution,
representing the proportion of probability density that falls to the
left-hand side.

\textbf{Args} - \emph{loc} (float): The location parameter of the
distribution.

\begin{itemize}
\item
  \emph{sample} (float): The scale parameter of the distribution.
\item
  \emph{quantile} (float): The quantile parameter, representing the
  proportion of probability density to the left of the median. Must be
  between 0 and 1.
\item
  \emph{shape} (tuple): A multi-purpose argument for shaping. When
  \texttt{sample=False} (model building), this is used with
  \texttt{shape} to set the distribution's batch shape. When
  \texttt{sample=True} (direct sampling), this is used as
  \texttt{sample\_shape} to draw a raw JAX array of the given shape.
\item
  \emph{event} (int): The number of batch dimensions to reinterpret as
  event dimensions (used in model building).
\item
  \emph{mask} (jnp.ndarray, bool): Optional boolean array to mask
  observations.
\item
  \emph{create\_obj} (bool, optional): If True, returns the raw BF
  distribution object instead of creating a sample site. This is
  essential for building complex distributions like
  \texttt{MixtureSameFamily}. Defaults to False.
\item
  \emph{sample} (bool, optional): A control-flow argument. If
  \texttt{True}, the function will directly sample a raw JAX array from
  the distribution, bypassing the BF model context. If \texttt{False},
  it will create a \texttt{BF.sample} site within a model. Defaults to
  \texttt{False}.
\item
  \emph{obs} (jnp.ndarray, optional): The observed value for this random
  variable. If provided, the sample site is conditioned on this value,
  and the function returns the observed value. If \texttt{None}, the
  site is treated as a latent (unobserved) random variable. Defaults to
  \texttt{None}.
\item
  \emph{name} (str, optional): The name of the sample site in a BF
  model. This is used to uniquely identify the random variable. Defaults
  to `x'.
\end{itemize}

\textbf{Returns} - When \texttt{sample=False}: A BF
AsymmetricLaplaceQuantile distribution object (for model building).

\begin{itemize}
\item
  When \texttt{sample=True}: A JAX array of samples drawn from the
  AsymmetricLaplaceQuantile distribution (for direct sampling).
\item
  When \texttt{create\_obj=True}: The raw BF AsymmetricLaplaceQuantile
  distribution object (for advanced use cases).
\end{itemize}

\textbf{Example Usage} from BayesForge import bf m = bf(platform=`cpu')
m.dist.asymmetric\_laplace\_quantile(loc=0.0, scale=1.0, quantile=0.5,
sample=True)

\textbf{Wrapper of}
https://num.pyro.ai/en/stable/distributions.html\#asymmetriclaplacequantile

\subsection{\texorpdfstring{\texttt{m.dist.bernoulli}}{m.dist.bernoulli}}\label{bernoulli}

\begin{Shaded}
\begin{Highlighting}[]
\NormalTok{m.dist.bernoulli(}
\NormalTok{    probs}\OperatorTok{=}\VariableTok{None}\NormalTok{,}
\NormalTok{    logits}\OperatorTok{=}\VariableTok{None}\NormalTok{,}
\NormalTok{    validate\_args}\OperatorTok{=}\VariableTok{None}\NormalTok{,}
\NormalTok{    name}\OperatorTok{=}\StringTok{\textquotesingle{}x\textquotesingle{}}\NormalTok{,}
\NormalTok{    obs}\OperatorTok{=}\VariableTok{None}\NormalTok{,}
\NormalTok{    mask}\OperatorTok{=}\VariableTok{None}\NormalTok{,}
\NormalTok{    infer}\OperatorTok{=}\VariableTok{None}\NormalTok{,}
\NormalTok{    sample}\OperatorTok{=}\VariableTok{False}\NormalTok{,}
\NormalTok{    seed}\OperatorTok{=}\VariableTok{None}\NormalTok{,}
\NormalTok{    shape}\OperatorTok{=}\NormalTok{(),}
\NormalTok{    event}\OperatorTok{=}\DecValTok{0}\NormalTok{,}
\NormalTok{    create\_obj}\OperatorTok{=}\VariableTok{False}\NormalTok{,}
\NormalTok{    to\_jax}\OperatorTok{=}\VariableTok{True}\NormalTok{,}
\NormalTok{)}
\end{Highlighting}
\end{Shaded}

\textbf{Bernoulli} The Bernoulli distribution models a single trial with
two possible outcomes: success or failureIt models a single trial with
exactly two possible outcomes:

\begin{itemize}
\tightlist
\item
  \textbf{Success} (often coded as 1) with probability (p),
\item
  \textbf{Failure} (often coded as 0) with probability (1 - p).
\end{itemize}

It is a building block for many more complex discrete distributions
(e.g.~binomial, geometric).

We write \[
X \sim \mathrm{Bernoulli}(p)
\] with \(0 \le p \le 1\).

\begin{center}\rule{0.5\linewidth}{0.5pt}\end{center}

\textbf{Probability Mass Function (PMF)} \[
P(X = x) =
\begin{cases}
p, & \text{if } x = 1, \\
1 - p, & \text{if } x = 0, \\
0, & \text{otherwise.}
\end{cases}
\]

Alternatively, using an ``indicator'' exponent form:

\[
P(X = x) = p^x (1 - p)^{1 - x}, \quad x \in {0,1}.
\]

\textbf{Args} - \emph{probs} (jnp.ndarray, optional): Probability of
success for each Bernoulli trial. Must be between 0 and 1.

\begin{itemize}
\item
  \emph{logits} (jnp.ndarray, optional): Log-odds of success for each
  Bernoulli trial. \texttt{probs\ =\ sigmoid(logits)}.
\item
  \emph{shape} (tuple): A multi-purpose argument for shaping. When
  \texttt{sample=False} (model building), this is used with
  \texttt{shape} to set the distribution's batch shape. When
  \texttt{sample=True} (direct sampling), this is used as
  \texttt{sample\_shape} to draw a raw JAX array of the given shape.
\item
  \emph{event} (int): The number of batch dimensions to reinterpret as
  event dimensions (used in model building).
\item
  \emph{mask} (jnp.ndarray, bool, optional): Optional boolean array to
  mask observations.
\item
  \emph{create\_obj} (bool, optional): If True, returns the raw BF
  distribution object instead of creating a sample site. This is
  essential for building complex distributions like
  \texttt{MixtureSameFamily}.
\item
  \emph{sample} (bool, optional): A control-flow argument. If
  \texttt{True}, the function will directly sample a raw JAX array from
  the distribution, bypassing the BF model context. If \texttt{False},
  it will create a \texttt{BF.sample} site within a model. Defaults to
  \texttt{False}.
\item
  \emph{seed} (int, optional): An integer seed for this draw when
  \texttt{sample=True}. By default (\texttt{None}) the draw comes from
  the session's random stream, set by \texttt{bf(rand\_seed=...)}, so
  consecutive draws differ but the whole sequence is reproducible. No
  effect when \texttt{sample=False}.
\item
  \emph{obs} (jnp.ndarray, optional): The observed value for this random
  variable. If provided, the sample site is conditioned on this value,
  and the function returns the observed value. If \texttt{None}, the
  site is treated as a latent (unobserved) random variable. Defaults to
  \texttt{None}.
\item
  \emph{name} (str, optional): The name of the sample site in a BF
  model. This is used to uniquely identify the random variable. Defaults
  to `x'.
\end{itemize}

\textbf{Returns} BF Bernoulli distribution object (for model building)
when \texttt{sample=False}. JAX array of samples drawn from the
Bernoulli distribution (for direct sampling) when \texttt{sample=True}.
The raw BF distribution object (for advanced use cases) when
\texttt{create\_obj=True}.

\textbf{Example Usage} from BayesForge import bf m = bf(platform=`cpu')
m.dist.bernoulli(probs=0.7, sample=True)

\textbf{Wrapper of:
https://num.pyro.ai/en/stable/distributions.html\#bernoulli}
\textbf{References} -
\href{https://en.wikipedia.org/wiki/Bernoulli_distribution}{Bernoulli
distribution} -
\href{https://stats.libretexts.org/Courses/Saint_Mary\%27s_College_Notre_Dame/MATH_345__-_Probability_\%28Kuter\%29/3\%3A_Discrete_Random_Variables/3.3\%3A_Bernoulli_and_Binomial_Distributions}{3.3:
Bernoulli and Binomial Distributions - Statistics LibreTexts}

\subsection{\texorpdfstring{\texttt{m.dist.bernoulli\_logits}}{m.dist.bernoulli\_logits}}\label{bernoulli_logits}

\begin{Shaded}
\begin{Highlighting}[]
\NormalTok{m.dist.bernoulli\_logits(}
\NormalTok{    logits}\OperatorTok{=}\VariableTok{None}\NormalTok{,}
\NormalTok{    validate\_args}\OperatorTok{=}\VariableTok{None}\NormalTok{,}
\NormalTok{    name}\OperatorTok{=}\StringTok{\textquotesingle{}x\textquotesingle{}}\NormalTok{,}
\NormalTok{    obs}\OperatorTok{=}\VariableTok{None}\NormalTok{,}
\NormalTok{    mask}\OperatorTok{=}\VariableTok{None}\NormalTok{,}
\NormalTok{    infer}\OperatorTok{=}\VariableTok{None}\NormalTok{,}
\NormalTok{    sample}\OperatorTok{=}\VariableTok{False}\NormalTok{,}
\NormalTok{    seed}\OperatorTok{=}\VariableTok{None}\NormalTok{,}
\NormalTok{    shape}\OperatorTok{=}\NormalTok{(),}
\NormalTok{    event}\OperatorTok{=}\DecValTok{0}\NormalTok{,}
\NormalTok{    create\_obj}\OperatorTok{=}\VariableTok{False}\NormalTok{,}
\NormalTok{    to\_jax}\OperatorTok{=}\VariableTok{True}\NormalTok{,}
\NormalTok{)}
\end{Highlighting}
\end{Shaded}

\textbf{Bernoulli Logits} Samples from a Bernoulli distribution
parameterized by logits.

The Bernoulli distribution models a single binary event (success or
failure), parameterized by the log-odds ratio of success. The
probability of success is given by the sigmoid function applied to the
logit.

\[
  p = \sigma(\eta) = \frac{1}{1 + e^{-\eta}}
  \] where (\eta) is the log-odds (the \emph{logit}).

\textbf{Args} - \emph{logits} (jnp.ndarray, optional): Log-odds ratio of
success. Must be real-valued.

\begin{itemize}
\item
  \emph{shape} (tuple): A multi-purpose argument for shaping. When
  \texttt{sample=False} (model building), this is used with
  \texttt{shape} to set the distribution's batch shape. When
  \texttt{sample=True} (direct sampling), this is used as
  \texttt{sample\_shape} to draw a raw JAX array of the given shape.
\item
  \emph{event} (int, optional): The number of batch dimensions to
  reinterpret as \emph{event} dimensions (used in model building).
  Defaults to 0.
\item
  \emph{mask} (jnp.ndarray, bool, optional): Optional boolean array to
  mask observations. Defaults to None.
\item
  \emph{create\_obj} (bool, optional): If True, returns the raw BF
  distribution object instead of creating a sample site. This is
  essential for building complex distributions like
  \texttt{MixtureSameFamily}. Defaults to False.
\item
  \emph{sample} (bool, optional): A control-flow argument. If
  \texttt{True}, the function will directly sample a raw JAX array from
  the distribution, bypassing the BF model context. If \texttt{False},
  it will create a \texttt{BF.sample} site within a model. Defaults to
  \texttt{False}.
\item
  \emph{obs} (jnp.ndarray, optional): The observed value for this random
  variable. If provided, the sample site is conditioned on this value,
  and the function returns the observed value. If \texttt{None}, the
  site is treated as a latent (unobserved) random variable. Defaults to
  \texttt{None}.
\item
  \emph{name} (str, optional): The name of the sample site in a BF
  model. This is used to uniquely identify the random variable. Defaults
  to `x'.
\end{itemize}

\textbf{Returns} - When \texttt{sample=False}: A BF BernoulliLogits
distribution object (for model building).

\begin{itemize}
\item
  When \texttt{sample=True}: A JAX array of samples drawn from the
  BernoulliLogits distribution (for direct sampling).
\item
  When \texttt{create\_obj=True}: The raw BF distribution object (for
  advanced use cases).

  \textbf{Example Usage} from BayesForge import bf import jax.numpy as
  jnp m = bf(platform=`cpu')
  m.dist.bernoulli\_logits(logits=jnp.array({[}0.2, 1, 2{]}),
  sample=True)

  \textbf{Wrapper of}
  https://num.pyro.ai/en/stable/distributions.html\#bernoulli-logits

  \textbf{References}

  \begin{itemize}
  \tightlist
  \item
    \href{https://en.wikipedia.org/wiki/Bernoulli_distribution}{Bernoulli
    distribution}
  \item
    \href{https://www.tensorflow.org/probability/api_docs/python/tfp/distributions/Bernoulli}{tfp.
    distributions.Bernoulli - Probability}
  \item
    \href{https://mc-stan.org/docs/2_19/functions-reference/bernoulli-logit-distribution.html}{11.2
    Bernoulli Distribution, Logit Parameterization - Stan}
  \end{itemize}
\end{itemize}

\subsection{\texorpdfstring{\texttt{m.dist.bernoulli\_probs}}{m.dist.bernoulli\_probs}}\label{bernoulli_probs}

\begin{Shaded}
\begin{Highlighting}[]
\NormalTok{m.dist.bernoulli\_probs(}
\NormalTok{    probs,}
\NormalTok{    validate\_args}\OperatorTok{=}\VariableTok{None}\NormalTok{,}
\NormalTok{    name}\OperatorTok{=}\StringTok{\textquotesingle{}x\textquotesingle{}}\NormalTok{,}
\NormalTok{    obs}\OperatorTok{=}\VariableTok{None}\NormalTok{,}
\NormalTok{    mask}\OperatorTok{=}\VariableTok{None}\NormalTok{,}
\NormalTok{    infer}\OperatorTok{=}\VariableTok{None}\NormalTok{,}
\NormalTok{    sample}\OperatorTok{=}\VariableTok{False}\NormalTok{,}
\NormalTok{    seed}\OperatorTok{=}\VariableTok{None}\NormalTok{,}
\NormalTok{    shape}\OperatorTok{=}\NormalTok{(),}
\NormalTok{    event}\OperatorTok{=}\DecValTok{0}\NormalTok{,}
\NormalTok{    create\_obj}\OperatorTok{=}\VariableTok{False}\NormalTok{,}
\NormalTok{    to\_jax}\OperatorTok{=}\VariableTok{True}\NormalTok{,}
\NormalTok{)}
\end{Highlighting}
\end{Shaded}

\textbf{Bernoulli Probs} Samples from a Bernoulli distribution
parameterized by probabilities.

The Bernoulli distribution models the probability of success in a single
trial, where the outcome is binary (success or failure). It is
characterized by a single parameter, \texttt{probs}, representing the
probability of success.

\[
Pr(X=1) = p = 1- Pr(X=0) = 1 - q, \\
f(x)=p \text{ if } k = 1 \text{ else } q \text{ if } k = 0
\]

where: p is the probability of success (0 \textless= p \textless= 1)

\textbf{Args} - \emph{probs} (jnp.ndarray): The probability of success
for each Bernoulli trial. Must be between 0 and 1.

\begin{itemize}
\item
  \emph{shape} (tuple): A multi-purpose argument for shaping. When
  \texttt{sample=False} (model building), this is used with
  \texttt{shape} to set the distribution's batch shape. When
  \texttt{sample=True} (direct sampling), this is used as
  \texttt{sample\_shape} to draw a raw JAX array of the given shape.
\item
  \emph{event} (int): The number of batch dimensions to reinterpret as
  event dimensions (used in model building).
\item
  \emph{mask} (jnp.ndarray, bool): Optional boolean array to mask
  observations.
\item
  \emph{create\_obj} (bool, optional): If True, returns the raw BF
  distribution object instead of creating a sample site. This is
  essential for building complex distributions like
  \texttt{MixtureSameFamily}. Defaults to False.
\item
  \emph{sample} (bool, optional): A control-flow argument. If
  \texttt{True}, the function will directly sample a raw JAX array from
  the distribution, bypassing the BF model context. If \texttt{False},
  it will create a \texttt{BF.sample} site within a model. Defaults to
  \texttt{False}.
\item
  \emph{obs} (jnp.ndarray, optional): The observed value for this random
  variable. If provided, the sample site is conditioned on this value,
  and the function returns the observed value. If \texttt{None}, the
  site is treated as a latent (unobserved) random variable. Defaults to
  \texttt{None}.
\item
  \emph{name} (str, optional): The name of the sample site in a BF
  model. This is used to uniquely identify the random variable. Defaults
  to `x'.
\end{itemize}

\textbf{Returns} - When \texttt{sample=False}: A BF BernoulliProbs
distribution object (for model building).

\begin{itemize}
\item
  When \texttt{sample=True}: A JAX array of samples drawn from the
  BernoulliLogits distribution (for direct sampling).
\item
  When \texttt{create\_obj=True}: The raw BF BernoulliProbs object (for
  advanced use cases).
\end{itemize}

\textbf{Example Usage} from BayesForge import bf import jax.numpy as jnp
m = bf(platform=`cpu') m.dist.bernoulli\_probs(probs=jnp.array({[}0.2,
0.7, 0.5{]}), sample=True)

\textbf{Wrapper of}
https://num.pyro.ai/en/stable/distributions.html\#bernoulliprobs

\subsection{\texorpdfstring{\texttt{m.dist.beta}}{m.dist.beta}}\label{beta}

\begin{Shaded}
\begin{Highlighting}[]
\NormalTok{m.dist.beta(}
\NormalTok{    concentration1,}
\NormalTok{    concentration0,}
\NormalTok{    validate\_args}\OperatorTok{=}\VariableTok{None}\NormalTok{,}
\NormalTok{    name}\OperatorTok{=}\StringTok{\textquotesingle{}x\textquotesingle{}}\NormalTok{,}
\NormalTok{    obs}\OperatorTok{=}\VariableTok{None}\NormalTok{,}
\NormalTok{    mask}\OperatorTok{=}\VariableTok{None}\NormalTok{,}
\NormalTok{    infer}\OperatorTok{=}\VariableTok{None}\NormalTok{,}
\NormalTok{    sample}\OperatorTok{=}\VariableTok{False}\NormalTok{,}
\NormalTok{    seed}\OperatorTok{=}\VariableTok{None}\NormalTok{,}
\NormalTok{    shape}\OperatorTok{=}\NormalTok{(),}
\NormalTok{    event}\OperatorTok{=}\DecValTok{0}\NormalTok{,}
\NormalTok{    create\_obj}\OperatorTok{=}\VariableTok{False}\NormalTok{,}
\NormalTok{    to\_jax}\OperatorTok{=}\VariableTok{True}\NormalTok{,}
\NormalTok{)}
\end{Highlighting}
\end{Shaded}

\textbf{Beta} Samples from a Beta distribution, defined on the interval
{[}0, 1{]}. The Beta distribution is a versatile distribution often used
to model probabilities or proportions. It is parameterized by two
positive shape parameters, usually denoted \(\alpha\) and \(\beta>0\),
control the shape of the density (how much mass is pushed toward 0, 1,
or intermediate).

\[
X \sim Beta(\alpha,\beta), \\f(x)=\frac{ x^{\alpha-1}(1-x)^{\beta-1}}{B(\alpha,\beta)}, \\
B(\alpha,\beta)=\frac{\Gamma(\alpha)\Gamma(\beta)}{\Gamma(\alpha+\beta)}, \\F(x)=I_{x}(\alpha+\beta)
\]

where \(B(\alpha, \beta)\) is the Beta function:

\[
B(\alpha, \beta) = \int_0^1 x^{\alpha - 1} (1 - x)^{\beta - 1} , dx = \frac{\Gamma(\alpha),\Gamma (\beta)}{\Gamma(\alpha + \beta)}.
\]

where :num\(\alpha\) and \(\beta\) are the concentration parameters, and
\(B(x, y)\) is the Beta function.

\textbf{Args} - \emph{concentration1} (jnp.ndarray): The first
concentration parameter (shape parameter). Must be positive.

\begin{itemize}
\item
  \emph{concentration0} (jnp.ndarray): The second concentration
  parameter (shape parameter). Must be positive.
\item
  \emph{shape} (tuple): A multi-purpose argument for shaping. When
  \texttt{sample=False} (model building), this is used with
  \texttt{shape} to set the distribution's batch shape. When
  \texttt{sample=True} (direct sampling), this is used as
  \texttt{sample\_shape} to draw a raw JAX array of the given shape.
\item
  \emph{event} (int): The number of batch dimensions to reinterpret as
  event dimensions (used in model building).
\item
  \emph{mask} (jnp.ndarray, bool): Optional boolean array to mask
  observations.
\item
  \emph{create\_obj} (bool): If True, returns the raw BF distribution
  object instead of creating a sample site. This is essential for
  building complex distributions like \texttt{MixtureSameFamily}.
\item
  \emph{sample} (bool, optional): A control-flow argument. If
  \texttt{True}, the function will directly sample a raw JAX array from
  the distribution, bypassing the BF model context. If \texttt{False},
  it will create a \texttt{BF.sample} site within a model. Defaults to
  \texttt{False}.
\item
  \emph{obs} (jnp.ndarray, optional): The observed value for this random
  variable. If provided, the sample site is conditioned on this value,
  and the function returns the observed value. If \texttt{None}, the
  site is treated as a latent (unobserved) random variable. Defaults to
  \texttt{None}.
\item
  \emph{name} (str, optional): The name of the sample site in a BF
  model. This is used to uniquely identify the random variable. Defaults
  to `x'.
\end{itemize}

\textbf{Returns} - When \texttt{sample=False}: A BF Beta distribution
object (for model building).

\begin{itemize}
\item
  When \texttt{sample=True}: A JAX array of samples drawn from the
  BernoulliLogits distribution (for direct sampling).
\item
  When \texttt{create\_obj=True}: The raw BF Beta object (for advanced
  use cases).
\end{itemize}

\textbf{\#\#\#\# Example Usage} from BayesForge import bf m =
bf(platform=`cpu') m.dist.beta(concentration1=1.0, concentration0=1.0,
sample=True)

\textbf{Wrapper of}
https://num.pyro.ai/en/stable/distributions.html\#beta

\textbf{References} -
\href{https://en.wikipedia.org/wiki/Beta_distribution}{Beta distribution
- Wikipedia} -
\href{https://www.statlect.com/probability-distributions/beta-distribution}{Beta
distribution \textbar{} Properties, proofs, exercises - StatLect} -
\href{https://reference.wolfram.com/language/ref/BetaDistribution.html}{BetaDistribution
- Wolfram Language Documentation}

\subsection{\texorpdfstring{\texttt{m.dist.beta\_binomial}}{m.dist.beta\_binomial}}\label{beta_binomial}

\begin{Shaded}
\begin{Highlighting}[]
\NormalTok{m.dist.beta\_binomial(}
\NormalTok{    concentration1,}
\NormalTok{    concentration0,}
\NormalTok{    total\_count}\OperatorTok{=}\DecValTok{1}\NormalTok{,}
\NormalTok{    validate\_args}\OperatorTok{=}\VariableTok{None}\NormalTok{,}
\NormalTok{    name}\OperatorTok{=}\StringTok{\textquotesingle{}x\textquotesingle{}}\NormalTok{,}
\NormalTok{    obs}\OperatorTok{=}\VariableTok{None}\NormalTok{,}
\NormalTok{    mask}\OperatorTok{=}\VariableTok{None}\NormalTok{,}
\NormalTok{    infer}\OperatorTok{=}\VariableTok{None}\NormalTok{,}
\NormalTok{    sample}\OperatorTok{=}\VariableTok{False}\NormalTok{,}
\NormalTok{    seed}\OperatorTok{=}\VariableTok{None}\NormalTok{,}
\NormalTok{    shape}\OperatorTok{=}\NormalTok{(),}
\NormalTok{    event}\OperatorTok{=}\DecValTok{0}\NormalTok{,}
\NormalTok{    create\_obj}\OperatorTok{=}\VariableTok{False}\NormalTok{,}
\NormalTok{    to\_jax}\OperatorTok{=}\VariableTok{True}\NormalTok{,}
\NormalTok{)}
\end{Highlighting}
\end{Shaded}

\textbf{BetaBinomial} Samples from a BetaBinomial distribution, a
compound distribution where the probability of success in a binomial
experiment is drawn from a Beta distribution. This models situations
where the underlying probability of success is not fixed but varies
according to a prior belief represented by the Beta distribution. It is
often used to model over-dispersion relative to the binomial
distribution.

For \(X \sim \mathrm{BetaBinomial}(N,\alpha,\beta)\), for
\(k = 0,1,\dots,N\):

\[
  \Pr(X = k) = \binom{N}{k} ; \frac{B(k + \alpha,; N - k + \beta)}{B(\alpha,\beta)},
  \]

where (\(B(u,v)\)) is the Beta function
\((B(u,v) = \frac{\Gamma(u),\Gamma(v)}{\Gamma(u+v)})\).

\textbf{Args} - \emph{concentration1} (jnp.ndarray): The first
concentration parameter (alpha) of the Beta distribution.

\begin{itemize}
\item
  \emph{concentration0} (jnp.ndarray): The second concentration
  parameter (beta) of the Beta distribution.
\item
  \emph{total\_count} (jnp.ndarray): The number of Bernoulli trials in
  the Binomial part of the distribution.
\item
  \emph{shape} (tuple): A multi-purpose argument for shaping. When
  \texttt{sample=False} (model building), this is used with
  \texttt{shape} to set the distribution's batch shape. When
  \texttt{sample=True} (direct sampling), this is used as
  \texttt{sample\_shape} to draw a raw JAX array of the given shape.
\item
  \emph{event} (int): The number of batch dimensions to reinterpret as
  event dimensions (used in model building).
\item
  \emph{mask} (jnp.ndarray, bool): Optional boolean array to mask
  observations.
\item
  \emph{create\_obj} (bool): If True, returns the raw BF distribution
  object instead of creating a sample site. This is essential for
  building complex distributions like \texttt{MixtureSameFamily}.
\item
  \emph{sample} (bool, optional): A control-flow argument. If
  \texttt{True}, the function will directly sample a raw JAX array from
  the distribution, bypassing the BF model context. If \texttt{False},
  it will create a \texttt{BF.sample} site within a model. Defaults to
  \texttt{False}.
\item
  \emph{obs} (jnp.ndarray, optional): The observed value for this random
  variable. If provided, the sample site is conditioned on this value,
  and the function returns the observed value. If \texttt{None}, the
  site is treated as a latent (unobserved) random variable. Defaults to
  \texttt{None}.
\item
  \emph{name} (str, optional): The name of the sample site in a BF
  model. This is used to uniquely identify the random variable. Defaults
  to `x'.
\end{itemize}

\textbf{Returns} - When \texttt{sample=False}: A BF BetaBinomial
distribution object (for model building).

\begin{itemize}
\item
  When \texttt{sample=True}: A JAX array of samples drawn from the
  BetaBinomial distribution (for direct sampling).
\item
  When \texttt{create\_obj=True}: The raw BF BetaBinomial distribution
  object (for advanced use cases).

  \textbf{Example Usage} from BayesForge import bf m =
  bf(platform=`cpu') m.dist.beta\_binomial(concentration1=1.0,
  concentration0=1.0, total\_count=10, sample=True)

  \textbf{Wrapper of}
  https://num.pyro.ai/en/stable/distributions.html\#betabinomial

  \textbf{References}

  \begin{itemize}
  \tightlist
  \item
    \href{https://en.wikipedia.org/wiki/Beta-binomial_distribution}{Beta-binomial
    distribution}
  \item
    \href{https://www.statisticshowto.com/beta-binomial-distribution/}{Beta-Binomial
    Distribution: Definition - Statistics How To}
  \item
    \href{https://mc-stan.org/docs/2_19/functions-reference/beta-binomial-distribution.html}{12.3
    Beta-Binomial Distribution \textbar{} Stan Functions Reference}
  \end{itemize}
\end{itemize}

\subsection{\texorpdfstring{\texttt{m.dist.beta\_proportion}}{m.dist.beta\_proportion}}\label{beta_proportion}

\begin{Shaded}
\begin{Highlighting}[]
\NormalTok{m.dist.beta\_proportion(}
\NormalTok{    mean,}
\NormalTok{    concentration,}
\NormalTok{    validate\_args}\OperatorTok{=}\VariableTok{None}\NormalTok{,}
\NormalTok{    name}\OperatorTok{=}\StringTok{\textquotesingle{}x\textquotesingle{}}\NormalTok{,}
\NormalTok{    obs}\OperatorTok{=}\VariableTok{None}\NormalTok{,}
\NormalTok{    mask}\OperatorTok{=}\VariableTok{None}\NormalTok{,}
\NormalTok{    infer}\OperatorTok{=}\VariableTok{None}\NormalTok{,}
\NormalTok{    sample}\OperatorTok{=}\VariableTok{False}\NormalTok{,}
\NormalTok{    seed}\OperatorTok{=}\VariableTok{None}\NormalTok{,}
\NormalTok{    shape}\OperatorTok{=}\NormalTok{(),}
\NormalTok{    event}\OperatorTok{=}\DecValTok{0}\NormalTok{,}
\NormalTok{    create\_obj}\OperatorTok{=}\VariableTok{False}\NormalTok{,}
\NormalTok{    to\_jax}\OperatorTok{=}\VariableTok{True}\NormalTok{,}
\NormalTok{)}
\end{Highlighting}
\end{Shaded}

\textbf{Beta Proportion} The Beta Proportion distribution is a
reparameterization of the conventional Beta distribution in terms of a
the variate mean and a precision parameter. It's useful for modeling
rates and proportions. It's essentially the same family as the standard
Beta((\alpha,\beta)), but the mapping is:

\[
\alpha = \mu , \kappa,\quad \beta = (1 - \mu), \kappa.
\]

Thus the pdf is:

\[
f(\theta \mid \mu, \kappa) = \frac{1}{B(\mu\kappa,,(1-\mu)\kappa)} ; \theta^{\mu\kappa - 1} ; (1 - \theta)^{(1 - \mu)\kappa         - 1}, \\
\theta \in (0,1). 
\]

This parameterization is useful in hierarchical models, since one can
put hyperpriors on (\mu) (the ``center'') and (\kappa) (the ``spread'').

\textbf{Args} - \emph{mean} (jnp.ndarray): The mean of the
BetaProportion distribution,must be between 0 and 1.

\begin{itemize}
\item
  \emph{concentration} (jnp.ndarray): The concentration parameter of the
  BetaProportion distribution.
\item
  \emph{shape} (tuple): A multi-purpose argument for shaping. When
  \texttt{sample=False} (model building), this is used with
  \texttt{shape} to set the distribution's batch shape. When
  \texttt{sample=True} (direct sampling), this is used as
  \texttt{sample\_shape} to draw a raw JAX array of the given shape.
\item
  \emph{event} (int): The number of batch dimensions to reinterpret as
  event dimensions (used in model building).
\item
  \emph{mask} (jnp.ndarray, bool): Optional boolean array to mask
  observations.
\item
  \emph{create\_obj} (bool): If True, returns the raw BF distribution
  object instead of creating a sample site. This is essential for
  building complex distributions like \texttt{MixtureSameFamily}.
\item
  \emph{sample} (bool, optional): A control-flow argument. If
  \texttt{True}, the function will directly sample a raw JAX array from
  the distribution, bypassing the BF model context. If \texttt{False},
  it will create a \texttt{BF.sample} site within a model. Defaults to
  \texttt{False}.
\item
  \emph{obs} (jnp.ndarray, optional): The observed value for this random
  variable. If provided, the sample site is conditioned on this value,
  and the function returns the observed value. If \texttt{None}, the
  site is treated as a latent (unobserved) random variable. Defaults to
  \texttt{None}.
\item
  \emph{name} (str, optional): The name of the sample site in a BF
  model. This is used to uniquely identify the random variable. Defaults
  to `x'.
\end{itemize}

\textbf{Returns} - When \texttt{sample=False}: A BF BetaProportion
distribution object (for model building).

\begin{itemize}
\item
  When \texttt{sample=True}: A JAX array of samples drawn from the
  BetaProportion distribution (for direct sampling).
\item
  When \texttt{create\_obj=True}: The raw BF BetaProportion distribution
  object (for advanced use cases).
\end{itemize}

\textbf{Example Usage} from BayesForge import bf m = bf(platform=`cpu')
samples = m.dist.beta\_proportion(mean=0.5, concentration=2.0,
sample=True, shape=(1000,))

\textbf{Wrapper of}
https://num.pyro.ai/en/stable/distributions.html\#beta\_proportion

\textbf{References}
https://mc-stan.org/docs/2\_19/functions-reference/beta-proportion-distribution.html
``19.2 Beta Proportion Distribution \textbar{} Stan Functions
Reference''

\subsection{\texorpdfstring{\texttt{m.dist.binomial}}{m.dist.binomial}}\label{binomial}

\begin{Shaded}
\begin{Highlighting}[]
\NormalTok{m.dist.binomial(}
\NormalTok{    total\_count}\OperatorTok{=}\DecValTok{1}\NormalTok{,}
\NormalTok{    probs}\OperatorTok{=}\VariableTok{None}\NormalTok{,}
\NormalTok{    logits}\OperatorTok{=}\VariableTok{None}\NormalTok{,}
\NormalTok{    validate\_args}\OperatorTok{=}\VariableTok{None}\NormalTok{,}
\NormalTok{    name}\OperatorTok{=}\StringTok{\textquotesingle{}x\textquotesingle{}}\NormalTok{,}
\NormalTok{    obs}\OperatorTok{=}\VariableTok{None}\NormalTok{,}
\NormalTok{    mask}\OperatorTok{=}\VariableTok{None}\NormalTok{,}
\NormalTok{    infer}\OperatorTok{=}\VariableTok{None}\NormalTok{,}
\NormalTok{    sample}\OperatorTok{=}\VariableTok{False}\NormalTok{,}
\NormalTok{    seed}\OperatorTok{=}\VariableTok{None}\NormalTok{,}
\NormalTok{    shape}\OperatorTok{=}\NormalTok{(),}
\NormalTok{    event}\OperatorTok{=}\DecValTok{0}\NormalTok{,}
\NormalTok{    create\_obj}\OperatorTok{=}\VariableTok{False}\NormalTok{,}
\NormalTok{    to\_jax}\OperatorTok{=}\VariableTok{True}\NormalTok{,}
\NormalTok{)}
\end{Highlighting}
\end{Shaded}

\textbf{Binomial} The Binomial distribution models the number of
successes in a sequence of independent Bernoulli trials. It represents
the probability of obtaining exactly \emph{k} successes in \emph{n}
trials, where each trial has a probability \emph{p} of success. It
assumes:

\begin{verbatim}
1. Each trial has exactly two possible outcomes (success or failure). 
2. The probability of success $p is constant across trials.
3. Trials are independent.
\end{verbatim}

We denote: \[
X \sim \mathrm{Binomial}(n, p).
\]

\textbf{Probability Mass Function (PMF)} For \(k = 0,1,2,\dots,n\):

\[
\Pr(X = k) = \binom{n}{k} ; p^k ; (1 - p)^{,n - k},
\]

where \(\binom{n}{k} = \frac{n!}{k!(n-k)!}\).

This formula counts the number of ways to choose which \(k\) out of
\(n\) trials are the successes (the binomial coefficient) and then
multiplies by the probability of each such configuration
\(p^k (1-p)^{n-k}\).

\textbf{Args} - \emph{total\_count} (int): The number of trials
\emph{n}.

\begin{itemize}
\item
  \emph{probs} (jnp.ndarray, optional): The probability of success
  \emph{p} for each trial. Must be between 0 and 1.
\item
  \emph{logits} (jnp.ndarray, optional): The log-odds of success for
  each trial. \texttt{probs\ =\ jax.nn.sigmoid(logits)}.
\item
  \emph{shape} (tuple): A multi-purpose argument for shaping. When
  \texttt{sample=False} (model building), this is used with
  \texttt{shape} to set the distribution's batch shape. When
  \texttt{sample=True} (direct sampling), this is used as
  \texttt{sample\_shape} to draw a raw JAX array of the given shape.
\item
  \emph{event} (int): The number of batch dimensions to reinterpret as
  event dimensions (used in model building).
\item
  \emph{mask} (jnp.ndarray, bool, optional): Optional boolean array to
  mask observations.
\item
  \emph{create\_obj} (bool, optional): If True, returns the raw BF
  distribution object instead of creating a sample site. This is
  essential for building complex distributions like
  \texttt{MixtureSameFamily}.
\item
  \emph{sample} (bool, optional): A control-flow argument. If
  \texttt{True}, the function will directly sample a raw JAX array from
  the distribution, bypassing the BF model context. If \texttt{False},
  it will create a \texttt{BF.sample} site within a model. Defaults to
  \texttt{False}.
\item
  \emph{seed} (int, optional): An integer seed for this draw when
  \texttt{sample=True}. By default (\texttt{None}) the draw comes from
  the session's random stream, set by \texttt{bf(rand\_seed=...)}, so
  consecutive draws differ but the whole sequence is reproducible. No
  effect when \texttt{sample=False}.
\item
  \emph{obs} (jnp.ndarray, optional): The observed value for this random
  variable. If provided, the sample site is conditioned on this value,
  and the function returns the observed value. If \texttt{None}, the
  site is treated as a latent (unobserved) random variable. Defaults to
  \texttt{None}.
\item
  \emph{name} (str, optional): The name of the sample site in a BF
  model. This is used to uniquely identify the random variable. Defaults
  to `x'.
\end{itemize}

\textbf{Returns} Binomial distribution object (for model building) when
\texttt{sample=False}. JAX array of samples drawn from the Binomial
distribution (for direct sampling) when \texttt{sample=True}. The raw BF
distribution object (for advanced use cases) when
\texttt{create\_obj=True}.

\textbf{Example Usage} from BayesForge import bf m = bf(platform=`cpu')
m.dist.binomial(total\_count=10, probs=0.5, sample=True)

\textbf{Wrapper of:
https://num.pyro.ai/en/stable/distributions.html\#binomial}
\textbf{References} -
\href{https://en.wikipedia.org/wiki/Binomial_distribution}{Binomial
distribution} -
\href{https://www.itl.nist.gov/div898/handbook/eda/section3/eda366i.htm}{1.3.6.6.18.
Binomial Distribution - Information Technology Laboratory}

\subsection{\texorpdfstring{\texttt{m.dist.binomial\_logits}}{m.dist.binomial\_logits}}\label{binomial_logits}

\begin{Shaded}
\begin{Highlighting}[]
\NormalTok{m.dist.binomial\_logits(}
\NormalTok{    logits,}
\NormalTok{    total\_count}\OperatorTok{=}\DecValTok{1}\NormalTok{,}
\NormalTok{    validate\_args}\OperatorTok{=}\VariableTok{None}\NormalTok{,}
\NormalTok{    name}\OperatorTok{=}\StringTok{\textquotesingle{}x\textquotesingle{}}\NormalTok{,}
\NormalTok{    obs}\OperatorTok{=}\VariableTok{None}\NormalTok{,}
\NormalTok{    mask}\OperatorTok{=}\VariableTok{None}\NormalTok{,}
\NormalTok{    infer}\OperatorTok{=}\VariableTok{None}\NormalTok{,}
\NormalTok{    sample}\OperatorTok{=}\VariableTok{False}\NormalTok{,}
\NormalTok{    seed}\OperatorTok{=}\VariableTok{None}\NormalTok{,}
\NormalTok{    shape}\OperatorTok{=}\NormalTok{(),}
\NormalTok{    event}\OperatorTok{=}\DecValTok{0}\NormalTok{,}
\NormalTok{    create\_obj}\OperatorTok{=}\VariableTok{False}\NormalTok{,}
\NormalTok{    to\_jax}\OperatorTok{=}\VariableTok{True}\NormalTok{,}
\NormalTok{)}
\end{Highlighting}
\end{Shaded}

\textbf{Binomial Logits} The BinomialLogits distribution represents a
binomial distribution parameterized by logits. It is useful when the
probability of success is not directly known but is instead expressed as
logits, which are the natural logarithm of the odds ratio.

\[
p = \mathrm{logit}^{-1}(\alpha) = \frac{1}{1 + e^{-\alpha}}.
\]

\textbf{Args} - \emph{logits} (jnp.ndarray): Log-odds of each success.

\begin{itemize}
\item
  \emph{total\_count} (int): Number of trials.
\item
  \emph{shape} (tuple): A multi-purpose argument for shaping. When
  \texttt{sample=False} (model building), this is used with
  \texttt{shape} to set the distribution's batch shape. When
  \texttt{sample=True} (direct sampling), this is used as
  \texttt{sample\_shape} to draw a raw JAX array of the given shape.
\item
  \emph{event} (int): The number of batch dimensions to reinterpret as
  event dimensions (used in model building).
\item
  \emph{mask} (jnp.ndarray, bool): Optional boolean array to mask
  observations.
\item
  \emph{create\_obj} (bool, optional): If True, returns the raw BF
  distribution object instead of creating a sample site. This is
  essential for building complex distributions like
  \texttt{MixtureSameFamily}. Defaults to False.
\item
  \emph{sample} (bool, optional): A control-flow argument. If
  \texttt{True}, the function will directly sample a raw JAX array from
  the distribution, bypassing the BF model context. If \texttt{False},
  it will create a \texttt{BF.sample} site within a model. Defaults to
  \texttt{False}.
\item
  \emph{obs} (jnp.ndarray, optional): The observed value for this random
  variable. If provided, the sample site is conditioned on this value,
  and the function returns the observed value. If \texttt{None}, the
  site is treated as a latent (unobserved) random variable. Defaults to
  \texttt{None}.
\item
  \emph{name} (str, optional): The name of the sample site in a BF
  model. This is used to uniquely identify the random variable. Defaults
  to `x'.
\end{itemize}

\textbf{Returns} - When \texttt{sample=False}: A BF BinomialLogits
distribution object (for model building).

\begin{itemize}
\item
  When \texttt{sample=True}: A JAX array of samples drawn from the
  BernoulliLogits distribution (for direct sampling).
\item
  When \texttt{create\_obj=True}: The raw BF BinomialLogits object (for
  advanced use cases).
\end{itemize}

\textbf{Example Usage} from BayesForge import bf import jax.numpy as jnp
m = bf(platform=`cpu') m.dist.binomial\_logits(logits=jnp.zeros(10),
total\_count=5, sample=True)

\textbf{Wrapper of}
https://num.pyro.ai/en/stable/distributions.html\#binomialllogits

\subsection{\texorpdfstring{\texttt{m.dist.binomial\_probs}}{m.dist.binomial\_probs}}\label{binomial_probs}

\begin{Shaded}
\begin{Highlighting}[]
\NormalTok{m.dist.binomial\_probs(}
\NormalTok{    probs,}
\NormalTok{    total\_count}\OperatorTok{=}\DecValTok{1}\NormalTok{,}
\NormalTok{    validate\_args}\OperatorTok{=}\VariableTok{None}\NormalTok{,}
\NormalTok{    name}\OperatorTok{=}\StringTok{\textquotesingle{}x\textquotesingle{}}\NormalTok{,}
\NormalTok{    obs}\OperatorTok{=}\VariableTok{None}\NormalTok{,}
\NormalTok{    mask}\OperatorTok{=}\VariableTok{None}\NormalTok{,}
\NormalTok{    infer}\OperatorTok{=}\VariableTok{None}\NormalTok{,}
\NormalTok{    sample}\OperatorTok{=}\VariableTok{False}\NormalTok{,}
\NormalTok{    seed}\OperatorTok{=}\VariableTok{None}\NormalTok{,}
\NormalTok{    shape}\OperatorTok{=}\NormalTok{(),}
\NormalTok{    event}\OperatorTok{=}\DecValTok{0}\NormalTok{,}
\NormalTok{    create\_obj}\OperatorTok{=}\VariableTok{False}\NormalTok{,}
\NormalTok{    to\_jax}\OperatorTok{=}\VariableTok{True}\NormalTok{,}
\NormalTok{)}
\end{Highlighting}
\end{Shaded}

\textbf{Binomial Probs} Samples from a Binomial distribution with
specified probabilities for each trial.

The Binomial distribution models the number of successes in a sequence
of independent Bernoulli trials, where each trial has the same
probability of success.

For \(k = 0,1,2,\dots,N\): \[   
\Pr(X = k) = \binom{N}{k} ; p^k ; (1 - p)^{,N - k}
\]

where \(\binom{N}{k} = \frac{N!}{k!(N-k)!}\).

\textbf{Args} - \emph{probs} (jnp.ndarray): The probability of success
for each trial. Must be between 0 and 1.

\begin{itemize}
\item
  \emph{total\_count} (int): The number of trials in each sequence.
\item
  \emph{shape} (tuple): A multi-purpose argument for shaping. When
  \texttt{sample=False} (model building), this is used with
  \texttt{shape} to set the distribution's batch shape. When
  \texttt{sample=True} (direct sampling), this is used as
  \texttt{sample\_shape} to draw a raw JAX array of the given shape.
\item
  \emph{event} (int): The number of batch dimensions to reinterpret as
  event dimensions (used in model building).
\item
  \emph{mask} (jnp.ndarray, bool): Optional boolean array to mask
  observations.
\item
  \emph{create\_obj} (bool): If True, returns the raw BF distribution
  object instead of creating a sample site. This is essential for
  building complex distributions like \texttt{MixtureSameFamily}.
\item
  \emph{sample} (bool, optional): A control-flow argument. If
  \texttt{True}, the function will directly sample a raw JAX array from
  the distribution, bypassing the BF model context. If \texttt{False},
  it will create a \texttt{BF.sample} site within a model. Defaults to
  \texttt{False}.
\item
  \emph{obs} (jnp.ndarray, optional): The observed value for this random
  variable. If provided, the sample site is conditioned on this value,
  and the function returns the observed value. If \texttt{None}, the
  site is treated as a latent (unobserved) random variable. Defaults to
  \texttt{None}.
\item
  \emph{name} (str, optional): The name of the sample site in a BF
  model. This is used to uniquely identify the random variable. Defaults
  to `x'.
\end{itemize}

\textbf{Returns} - When \texttt{sample=False}: A BF BinomialProbs
distribution object (for model building).

\begin{itemize}
\item
  When \texttt{sample=True}: A JAX array of samples drawn from the
  BernoulliLogits distribution (for direct sampling).
\item
  When \texttt{create\_obj=True}: The raw BF BinomialLogits object (for
  advanced use cases).
\end{itemize}

\textbf{Example Usage} from BayesForge import bf m = bf(platform=`cpu')
m.dist.binomial\_probs(probs=0.5, total\_count=10, sample=True)

\textbf{Wrapper of}
https://num.pyro.ai/en/stable/distributions.html\#binomialprobs

\subsection{\texorpdfstring{\texttt{m.dist.car}}{m.dist.car}}\label{car}

\begin{Shaded}
\begin{Highlighting}[]
\NormalTok{m.dist.car(}
\NormalTok{    loc,}
\NormalTok{    correlation,}
\NormalTok{    conditional\_precision,}
\NormalTok{    adj\_matrix,}
\NormalTok{    is\_sparse}\OperatorTok{=}\VariableTok{False}\NormalTok{,}
\NormalTok{    validate\_args}\OperatorTok{=}\VariableTok{None}\NormalTok{,}
\NormalTok{    name}\OperatorTok{=}\StringTok{\textquotesingle{}x\textquotesingle{}}\NormalTok{,}
\NormalTok{    obs}\OperatorTok{=}\VariableTok{None}\NormalTok{,}
\NormalTok{    mask}\OperatorTok{=}\VariableTok{None}\NormalTok{,}
\NormalTok{    infer}\OperatorTok{=}\VariableTok{None}\NormalTok{,}
\NormalTok{    sample}\OperatorTok{=}\VariableTok{False}\NormalTok{,}
\NormalTok{    seed}\OperatorTok{=}\VariableTok{None}\NormalTok{,}
\NormalTok{    shape}\OperatorTok{=}\NormalTok{(),}
\NormalTok{    event}\OperatorTok{=}\DecValTok{0}\NormalTok{,}
\NormalTok{    create\_obj}\OperatorTok{=}\VariableTok{False}\NormalTok{,}
\NormalTok{    to\_jax}\OperatorTok{=}\VariableTok{True}\NormalTok{,}
\NormalTok{)}
\end{Highlighting}
\end{Shaded}

\textbf{Conditional Autoregressive (CAR)} The CAR distribution models a
vector of variables where each variable is a linear combination of its
neighbors in a graph. The CAR model captures spatial dependence in areal
data by modeling each observation as conditionally dependent on its
neighbors.It specifies a joint distribution of a vector of random
variables \(\mathbf{y} = (y_1, y_2, \dots, y_N)\) based on their
conditional distributions, where each \(y_i\) is conditionally
independent of all other variables given its neighbors.

\begin{itemize}
\item
  \textbf{Application}: Widely used in disease mapping, environmental
  modeling, and spatial econometrics to account for spatial
  autocorrelation.
\item
  \textbf{Conditional Distribution}: Each \(y_i\) given its neighbors
  \(\mathbf{y}_{-i}\) follows a normal distribution:
\end{itemize}

\[
y_i | \mathbf{y}*{-i} \sim \mathcal{N}\left( \frac{\sum*{j \in \mathcal{N}(i)} w_{ij} y_j}{\sum_{j \in \mathcal{N}(i)} w_{ij}}, \sigma^2 \right)
\]

where:

\begin{itemize}
\tightlist
\item
  \(\mathcal{N}(i)\) denotes the set of neighbors of unit \(i\),
\item
  \(w_{ij}\) are the weights representing the strength of the spatial
  relationship between units \(i\) and \(j\),
\item
  \(\sigma^2\) is the variance term. where \(\mu_i\) is a function of
  the values of the neighbors of site \(i\) and \(\Sigma_i\) is the
  variance of site \(i\).
\end{itemize}

\textbf{Args} - \emph{loc} (Union{[}float, Array{]}): Mean of the
distribution.

\begin{itemize}
\item
  \emph{correlation} (Union{[}float, Array{]}): Correlation between
  variables.
\item
  \emph{conditional\_precision} (Union{[}float, Array{]}): Precision of
  the distribution.
\item
  \emph{adj\_matrix} (Union{[}Array, scipy.sparse.spmatrix{]}):
  Adjacency matrix defining the graph. is\_sparse (bool): Whether the
  adjacency matrix is sparse. Defaults to False.
\item
  *validate\_args/ (bool): Whether to validate arguments. Defaults to
  None.
\item
  \emph{sample} (bool, optional): A control-flow argument. If
  \texttt{True}, the function will directly sample a raw JAX array from
  the distribution, bypassing the BF model context. If \texttt{False},
  it will create a \texttt{BF.sample} site within a model. Defaults to
  \texttt{False}.
\item
  \emph{obs} (jnp.ndarray, optional): The observed value for this random
  variable. If provided, the sample site is conditioned on this value,
  and the function returns the observed value. If \texttt{None}, the
  site is treated as a latent (unobserved) random variable. Defaults to
  \texttt{None}.
\item
  \emph{name} (str, optional): The name of the sample site in a BF
  model. This is used to uniquely identify the random variable. Defaults
  to `x'.
\end{itemize}

\textbf{Returns} - When \texttt{sample=False}: A BF CAR distribution
object (for model building).

\begin{itemize}
\item
  When \texttt{sample=True}: A JAX array of samples drawn from the
  BernoulliLogits distribution (for direct sampling).
\item
  When \texttt{create\_obj=True}: The raw BF CAR object (for advanced
  use cases).
\end{itemize}

\textbf{References} -
\href{https://www.pymc.io/projects/examples/en/latest/spatial/conditional_autoregressive_priors.html}{Conditional
Autoregressive (CAR) Models for Spatial Data} -
\href{https://mc-stan.org/learn-stan/case-studies/icar_stan.html}{Intrinsic
Auto-Regressive Models for Areal Data - Stan} -
\href{https://onlinelibrary.wiley.com/doi/abs/10.1002/9781118445112.stat08048}{Conditional
Autoregressive (CAR) Model - Schmidt}

\subsection{\texorpdfstring{\texttt{m.dist.categorical}}{m.dist.categorical}}\label{categorical}

\begin{Shaded}
\begin{Highlighting}[]
\NormalTok{m.dist.categorical(}
\NormalTok{    probs}\OperatorTok{=}\VariableTok{None}\NormalTok{,}
\NormalTok{    logits}\OperatorTok{=}\VariableTok{None}\NormalTok{,}
\NormalTok{    validate\_args}\OperatorTok{=}\VariableTok{None}\NormalTok{,}
\NormalTok{    name}\OperatorTok{=}\StringTok{\textquotesingle{}x\textquotesingle{}}\NormalTok{,}
\NormalTok{    obs}\OperatorTok{=}\VariableTok{None}\NormalTok{,}
\NormalTok{    mask}\OperatorTok{=}\VariableTok{None}\NormalTok{,}
\NormalTok{    infer}\OperatorTok{=}\VariableTok{None}\NormalTok{,}
\NormalTok{    sample}\OperatorTok{=}\VariableTok{False}\NormalTok{,}
\NormalTok{    seed}\OperatorTok{=}\VariableTok{None}\NormalTok{,}
\NormalTok{    shape}\OperatorTok{=}\NormalTok{(),}
\NormalTok{    event}\OperatorTok{=}\DecValTok{0}\NormalTok{,}
\NormalTok{    create\_obj}\OperatorTok{=}\VariableTok{False}\NormalTok{,}
\NormalTok{    to\_jax}\OperatorTok{=}\VariableTok{True}\NormalTok{,}
\NormalTok{)}
\end{Highlighting}
\end{Shaded}

\textbf{Categorical distribution.} The Categorical distribution, also
known as the multinomial distribution, describes the probability of
\(K\) different outcomes with probabilities
\({p_1, p_2, \dots, p_K}\).There's no inherent ordering required among
the categories.

\begin{verbatim}
  - The probabilities satisfy (p_i \ge 0) for all (i) and $ \sum_{i=1}^K p_i = 1$ . 
  - It’s essentially the distribution of a single draw from a discrete set of categories. 
  - It generalises the Bernoulli distribution (which is the special case $K=2$).
\end{verbatim}

It is commonly used to model discrete choices or classifications.

If \(X\) is a categorical random variable taking values in
\({1, 2, \dots, K}\) with probabilities \(p_1, \dots, p_K\), then:

\[
  \Pr(X = i) = p_i, \quad i = 1,\dots, K.
  \]

where \(p_k\) is the probability of outcome \(k\), and the sum is over
all possible outcomes.

\textbf{Args} - \emph{probs} (jnp.ndarray): A 1D array of probabilities
for each category. Must sum to 1.

\begin{itemize}
\item
  \emph{logits} (jnp.ndarray): A 1D array of unnormalized log
  probabilities for each category.
\item
  \emph{shape} (tuple): A multi-purpose argument for shaping. When
  \texttt{sample=False} (model building), this is used with
  \texttt{.expand(shape)} to set the distribution's batch shape. When
  \texttt{sample=True} (direct sampling), this is used as
  \texttt{sample\_shape} to draw a raw JAX array of the given shape.
\item
  \emph{event} (int): The number of batch dimensions to reinterpret as
  event dimensions (used in model building).
\item
  \emph{mask} (jnp.ndarray, bool): Optional boolean array to mask
  observations.
\item
  \emph{create\_obj} (bool): If True, returns the raw BF distribution
  object instead of creating a sample site. This is essential for
  building complex distributions like \texttt{MixtureSameFamily}.
\item
  \emph{sample} (bool, optional): A control-flow argument. If
  \texttt{True}, the function will directly sample a raw JAX array from
  the distribution, bypassing the BF model context. If \texttt{False},
  it will create a \texttt{BF.sample} site within a model. Defaults to
  \texttt{False}.
\item
  \emph{seed} (int, optional): An integer seed for this draw when
  \texttt{sample=True}. By default (\texttt{None}) the draw comes from
  the session's random stream, set by \texttt{bf(rand\_seed=...)}, so
  consecutive draws differ but the whole sequence is reproducible. No
  effect when \texttt{sample=False}.
\item
  \emph{obs} (jnp.ndarray, optional): The observed value for this random
  variable. If provided, the sample site is conditioned on this value,
  and the function returns the observed value. If \texttt{None}, the
  site is treated as a latent (unobserved) random variable. Defaults to
  \texttt{None}.
\item
  \emph{name} (str, optional): The name of the sample site in a BF
  model. This is used to uniquely identify the random variable. Defaults
  to `x'.
\end{itemize}

\textbf{Returns} - When \texttt{sample=False}: A BF Categorical
distribution object (for model building).

\begin{itemize}
\item
  When \texttt{sample=True}: A JAX array of samples drawn from the
  Categorical distribution (for direct sampling).
\item
  When \texttt{create\_obj=True}: The raw BF distribution object (for
  advanced use cases).

  \textbf{Example Usage} from BayesForge import bf import jax.numpy as
  jnp m = bf(platform=`cpu') m.dist.categorical(probs=jnp.array({[}0.2,
  0.3, 0.5{]}), sample=True)

  \textbf{Wrapper of:
  https://num.pyro.ai/en/stable/distributions.html\#categorical}
  \textbf{References}

  \begin{itemize}
  \tightlist
  \item
    https://en.wikipedia.org/wiki/Categorical\_distribution''Categorical
    distribution''
  \item
    https://distribution-explorer.github.io/discrete/categorical.html''Categorical
    distribution --- Probability Distribution Explorer documentation''
  \item
    https://static.hlt.bme.hu/semantics/external/pages/n-gram/en.wikipedia.org/wiki/Categorical\_distribution.html''Categorical
    distribution - Wikipedia''
  \end{itemize}
\end{itemize}

\subsection{\texorpdfstring{\texttt{m.dist.categorical\_logits}}{m.dist.categorical\_logits}}\label{categorical_logits}

\begin{Shaded}
\begin{Highlighting}[]
\NormalTok{m.dist.categorical\_logits(}
\NormalTok{    logits,}
\NormalTok{    validate\_args}\OperatorTok{=}\VariableTok{None}\NormalTok{,}
\NormalTok{    name}\OperatorTok{=}\StringTok{\textquotesingle{}x\textquotesingle{}}\NormalTok{,}
\NormalTok{    obs}\OperatorTok{=}\VariableTok{None}\NormalTok{,}
\NormalTok{    mask}\OperatorTok{=}\VariableTok{None}\NormalTok{,}
\NormalTok{    infer}\OperatorTok{=}\VariableTok{None}\NormalTok{,}
\NormalTok{    sample}\OperatorTok{=}\VariableTok{False}\NormalTok{,}
\NormalTok{    seed}\OperatorTok{=}\VariableTok{None}\NormalTok{,}
\NormalTok{    shape}\OperatorTok{=}\NormalTok{(),}
\NormalTok{    event}\OperatorTok{=}\DecValTok{0}\NormalTok{,}
\NormalTok{    create\_obj}\OperatorTok{=}\VariableTok{False}\NormalTok{,}
\NormalTok{    to\_jax}\OperatorTok{=}\VariableTok{True}\NormalTok{,}
\NormalTok{)}
\end{Highlighting}
\end{Shaded}

\textbf{Categorical Logits} Samples from a Categorical distribution with
logits. This distribution represents a discrete probability distribution
over a finite set of outcomes, where the probabilities are determined by
the logits. The probability of each outcome is given by the softmax
function applied to the logits.

Given logits \(\gamma = (\gamma_1, \dots, \gamma_K)\), the probability
of each category \$ i \$ is:

\[
p_i = \frac{e^{\gamma_i}}{\sum_{j=1}^{K} e^{\gamma_j}}
\]

where \$ p\_i is the probability of category \$ i , and the sum in the
denominator ensures that the probabilities sum to 1.

\textbf{Args} - \emph{logits} (jnp.ndarray): Log-odds of each category.

\begin{itemize}
\item
  \emph{shape} (tuple): A multi-purpose argument for shaping. When
  \texttt{sample=False} (model building), this is used with
  \texttt{shape} to set the distribution's batch shape. When
  \texttt{sample=True} (direct sampling), this is used as
  \texttt{sample\_shape} to draw a raw JAX array of the given shape.
\item
  \emph{event} (int): The number of batch dimensions to reinterpret as
  event dimensions (used in model building).
\item
  \emph{mask} (jnp.ndarray, bool): Optional boolean array to mask
  observations.
\item
  \emph{create\_obj} (bool): If True, returns the raw BF distribution
  object instead of creating a sample site. This is essential for
  building complex distributions like \texttt{MixtureSameFamily}.
\item
  \emph{sample} (bool, optional): A control-flow argument. If
  \texttt{True}, the function will directly sample a raw JAX array from
  the distribution, bypassing the BF model context. If \texttt{False},
  it will create a \texttt{BF.sample} site within a model. Defaults to
  \texttt{False}.
\item
  \emph{obs} (jnp.ndarray, optional): The observed value for this random
  variable. If provided, the sample site is conditioned on this value,
  and the function returns the observed value. If \texttt{None}, the
  site is treated as a latent (unobserved) random variable. Defaults to
  \texttt{None}.
\item
  \emph{name} (str, optional): The name of the sample site in a BF
  model. This is used to uniquely identify the random variable. Defaults
  to `x'.
\end{itemize}

\textbf{Returns} - When \texttt{sample=False}: A BF CategoricalLogits
distribution object (for model building).

\begin{itemize}
\item
  When \texttt{sample=True}: A JAX array of samples drawn from the
  BernoulliLogits distribution (for direct sampling).
\item
  When \texttt{create\_obj=True}: The raw BF CategoricalLogits object
  (for advanced use cases).
\end{itemize}

\textbf{Example Usage} from BayesForge import bf import jax.numpy as jnp
m = bf(platform=`cpu') m.dist.categorical\_logits(logits=jnp.zeros(5),
sample=True)

\textbf{Wrapper of}
https://num.pyro.ai/en/stable/distributions.html\#categoricallogits

\subsection{\texorpdfstring{\texttt{m.dist.categorical\_probs}}{m.dist.categorical\_probs}}\label{categorical_probs}

\begin{Shaded}
\begin{Highlighting}[]
\NormalTok{m.dist.categorical\_probs(}
\NormalTok{    probs,}
\NormalTok{    validate\_args}\OperatorTok{=}\VariableTok{None}\NormalTok{,}
\NormalTok{    name}\OperatorTok{=}\StringTok{\textquotesingle{}x\textquotesingle{}}\NormalTok{,}
\NormalTok{    obs}\OperatorTok{=}\VariableTok{None}\NormalTok{,}
\NormalTok{    mask}\OperatorTok{=}\VariableTok{None}\NormalTok{,}
\NormalTok{    infer}\OperatorTok{=}\VariableTok{None}\NormalTok{,}
\NormalTok{    sample}\OperatorTok{=}\VariableTok{False}\NormalTok{,}
\NormalTok{    seed}\OperatorTok{=}\VariableTok{None}\NormalTok{,}
\NormalTok{    shape}\OperatorTok{=}\NormalTok{(),}
\NormalTok{    event}\OperatorTok{=}\DecValTok{0}\NormalTok{,}
\NormalTok{    create\_obj}\OperatorTok{=}\VariableTok{False}\NormalTok{,}
\NormalTok{    to\_jax}\OperatorTok{=}\VariableTok{True}\NormalTok{,}
\NormalTok{)}
\end{Highlighting}
\end{Shaded}

\textbf{Categorical Probs distribution.} Samples from a Categorical
distribution.

The Categorical distribution is a discrete probability distribution that
represents the probability of each outcome from a finite set of
possibilities. It is often used to model the outcome of a random process
with a fixed number of possible outcomes, such as the roll of a die or
the selection of an item from a list.

\[   
P(x) = \frac{probs_i}{\sum_{k=1}^{K} probs_k}
\]

\textbf{Args} - \emph{probs} (jnp.ndarray): Probabilities for each
category. Must sum to 1.

\begin{itemize}
\item
  \emph{shape} (tuple): A multi-purpose argument for shaping. When
  \texttt{sample=False} (model building), this is used with
  \texttt{shape} to set the distribution's batch shape. When
  \texttt{sample=True} (direct sampling), this is used as
  \texttt{sample\_shape} to draw a raw JAX array of the given shape.
\item
  \emph{event} (int): The number of batch dimensions to reinterpret as
  event dimensions (used in model building).
\item
  \emph{mask} (jnp.ndarray, bool): Optional boolean array to mask
  observations.
\item
  \emph{create\_obj} (bool, optional): If True, returns the raw BF
  distribution object instead of creating a sample site. This is
  essential for building complex distributions like
  \texttt{MixtureSameFamily}. Defaults to False.
\item
  \emph{sample} (bool, optional): A control-flow argument. If
  \texttt{True}, the function will directly sample a raw JAX array from
  the distribution, bypassing the BF model context. If \texttt{False},
  it will create a \texttt{BF.sample} site within a model. Defaults to
  \texttt{False}.
\item
  \emph{obs} (jnp.ndarray, optional): The observed value for this random
  variable. If provided, the sample site is conditioned on this value,
  and the function returns the observed value. If \texttt{None}, the
  site is treated as a latent (unobserved) random variable. Defaults to
  \texttt{None}.
\item
  \emph{name} (str, optional): The name of the sample site in a BF
  model. This is used to uniquely identify the random variable. Defaults
  to `x'.
\end{itemize}

\textbf{Returns} - When \texttt{sample=False}: A BF CategoricalProbs
distribution object (for model building).

\begin{itemize}
\item
  When \texttt{sample=True}: A JAX array of samples drawn from the
  BernoulliLogits distribution (for direct sampling).
\item
  When \texttt{create\_obj=True}: The raw BF CategoricalProbs object
  (for advanced use cases).
\end{itemize}

\textbf{Example Usage} from BayesForge import bf import jax.numpy as jnp
m = bf(platform=`cpu') m.dist.categorical\_probs(probs=jnp.array({[}0.2,
0.3, 0.5{]}), sample=True)

\textbf{Wrapper of}
https://num.pyro.ai/en/stable/distributions.html\#categoricalprobs

\subsection{\texorpdfstring{\texttt{m.dist.cauchy}}{m.dist.cauchy}}\label{cauchy}

\begin{Shaded}
\begin{Highlighting}[]
\NormalTok{m.dist.cauchy(}
\NormalTok{    loc}\OperatorTok{=}\FloatTok{0.0}\NormalTok{,}
\NormalTok{    scale}\OperatorTok{=}\FloatTok{1.0}\NormalTok{,}
\NormalTok{    validate\_args}\OperatorTok{=}\VariableTok{None}\NormalTok{,}
\NormalTok{    name}\OperatorTok{=}\StringTok{\textquotesingle{}x\textquotesingle{}}\NormalTok{,}
\NormalTok{    obs}\OperatorTok{=}\VariableTok{None}\NormalTok{,}
\NormalTok{    mask}\OperatorTok{=}\VariableTok{None}\NormalTok{,}
\NormalTok{    infer}\OperatorTok{=}\VariableTok{None}\NormalTok{,}
\NormalTok{    sample}\OperatorTok{=}\VariableTok{False}\NormalTok{,}
\NormalTok{    seed}\OperatorTok{=}\VariableTok{None}\NormalTok{,}
\NormalTok{    shape}\OperatorTok{=}\NormalTok{(),}
\NormalTok{    event}\OperatorTok{=}\DecValTok{0}\NormalTok{,}
\NormalTok{    create\_obj}\OperatorTok{=}\VariableTok{False}\NormalTok{,}
\NormalTok{    to\_jax}\OperatorTok{=}\VariableTok{True}\NormalTok{,}
\NormalTok{)}
\end{Highlighting}
\end{Shaded}

\textbf{Cauchy} Samples from a Cauchy

The Cauchy distribution, also known as the Lorentz distribution, is a
continuous probability distribution that arises frequently in various
fields, including physics and statistics. It is characterized by its
heavy tails, which extend indefinitely. This means it has a higher
probability of extreme values compared to the normal distribution.

\begin{itemize}
\tightlist
\item
  \textbf{Parameters}:

  \begin{itemize}
  \tightlist
  \item
    \textbf{Location parameter} (\(\mu\)): Determines the peak of the
    distribution.
  \item
    \textbf{Scale parameter} (\(\gamma\)): Describes the half-width at
    half-maximum (HWHM) of the peak. \[
      f(x; \mu, \gamma) = \frac{1}{\pi \gamma \left[1 + \left(\frac{x - \mu}{\gamma}\right)^2\right]}
      \]
  \end{itemize}
\end{itemize}

where:

\begin{itemize}
\tightlist
\item
  \(x\) is the random variable,
\item
  \(\mu\) is the location parameter,
\item
  \(\gamma\) is the scale parameter.
\end{itemize}

\textbf{Args} - \emph{loc} (jnp.ndarray or float, optional): Location
parameter. Defaults to 0.0.

\begin{itemize}
\item
  \emph{sample} (jnp.ndarray or float, optional): Scale parameter. Must
  be positive. Defaults to 1.0.
\item
  \emph{shape} (tuple): A multi-purpose argument for shaping. When
  \texttt{sample=False} (model building), this is used with
  \texttt{shape} to set the distribution's batch shape. When
  \texttt{sample=True} (direct sampling), this is used as
  \texttt{sample\_shape} to draw a raw JAX array of the given shape.
\item
  \emph{event} (int, optional): The number of batch dimensions to
  reinterpret as event dimensions (used in model building). Defaults to
  None.
\item
  \emph{mask} (jnp.ndarray, bool, optional): Optional boolean array to
  mask observations. Defaults to None.
\item
  \emph{create\_obj} (bool, optional): If True, returns the raw BF
  distribution object instead of creating a sample site. This is
  essential for building complex distributions like
  \texttt{MixtureSameFamily}. Defaults to False.
\item
  \emph{sample} (bool, optional): A control-flow argument. If
  \texttt{True}, the function will directly sample a raw JAX array from
  the distribution, bypassing the BF model context. If \texttt{False},
  it will create a \texttt{BF.sample} site within a model. Defaults to
  \texttt{False}.
\item
  \emph{obs} (jnp.ndarray, optional): The observed value for this random
  variable. If provided, the sample site is conditioned on this value,
  and the function returns the observed value. If \texttt{None}, the
  site is treated as a latent (unobserved) random variable. Defaults to
  \texttt{None}.
\item
  \emph{name} (str, optional): The name of the sample site in a BF
  model. This is used to uniquely identify the random variable. Defaults
  to `x'.
\end{itemize}

\textbf{Returns} - When \texttt{sample=False}: A BF Cauchy distribution
object (for model building).

\begin{itemize}
\item
  When \texttt{sample=True}: A JAX array of samples drawn from the
  Cauchy distribution (for direct sampling).
\item
  When \texttt{create\_obj=True}: The raw BF distribution object (for
  advanced use cases).

  \textbf{Example Usage} from BayesForge import bf m =
  bf(platform=`cpu') m.dist.cauchy(loc=0.0, scale=1.0, sample=True)

  \textbf{Wrapper of}
  https://num.pyro.ai/en/stable/distributions.html\#cauchy
\end{itemize}

\subsection{\texorpdfstring{\texttt{m.dist.chi2}}{m.dist.chi2}}\label{chi2}

\begin{Shaded}
\begin{Highlighting}[]
\NormalTok{m.dist.chi2(}
\NormalTok{    df,}
\NormalTok{    validate\_args}\OperatorTok{=}\VariableTok{None}\NormalTok{,}
\NormalTok{    name}\OperatorTok{=}\StringTok{\textquotesingle{}x\textquotesingle{}}\NormalTok{,}
\NormalTok{    obs}\OperatorTok{=}\VariableTok{None}\NormalTok{,}
\NormalTok{    mask}\OperatorTok{=}\VariableTok{None}\NormalTok{,}
\NormalTok{    infer}\OperatorTok{=}\VariableTok{None}\NormalTok{,}
\NormalTok{    sample}\OperatorTok{=}\VariableTok{False}\NormalTok{,}
\NormalTok{    seed}\OperatorTok{=}\VariableTok{None}\NormalTok{,}
\NormalTok{    shape}\OperatorTok{=}\NormalTok{(),}
\NormalTok{    event}\OperatorTok{=}\DecValTok{0}\NormalTok{,}
\NormalTok{    create\_obj}\OperatorTok{=}\VariableTok{False}\NormalTok{,}
\NormalTok{    to\_jax}\OperatorTok{=}\VariableTok{True}\NormalTok{,}
\NormalTok{)}
\end{Highlighting}
\end{Shaded}

\textbf{Chi-squared} The Chi-squared distribution is a continuous
probability distribution that arises frequently in hypothesis testing,
particularly in ANOVA and chi-squared tests. It is defined by a single
positive parameter, degrees of freedom (df), the number of independent
standard normal variables squared and summed, which determines the shape
of the distribution.

\[   
f(x; k) = \begin{cases}
\frac{x^{\frac{k}{2} - 1} e^{-x/2}}{2^{\frac{k}{2}} \Gamma\left(\frac{k}{2}\right)}, & x > 0 \\
0, & \text{otherwise}
\end{cases}
\]

where: * ( x ) is the random variable, * ( k ) is the degrees of
freedom, * ( \Gamma(\cdot) ) is the gamma function.

\textbf{Args} - \emph{df} (jnp.ndarray): Degrees of freedom. Must be
positive.

\begin{itemize}
\item
  \emph{shape} (tuple): A multi-purpose argument for shaping. When
  \texttt{sample=False} (model building), this is used with
  \texttt{shape} to set the distribution's batch shape. When
  \texttt{sample=True} (direct sampling), this is used as
  \texttt{sample\_shape} to draw a raw JAX array of the given shape.
\item
  \emph{event} (int): The number of batch dimensions to reinterpret as
  event dimensions (used in model building).
\item
  \emph{mask} (jnp.ndarray, bool): Optional boolean array to mask
  observations.
\item
  \emph{create\_obj} (bool, optional): If True, returns the raw BF
  distribution object instead of creating a sample site. This is
  essential for building complex distributions like
  \texttt{MixtureSameFamily}. Defaults to False.
\item
  \emph{sample} (bool, optional): A control-flow argument. If
  \texttt{True}, the function will directly sample a raw JAX array from
  the distribution, bypassing the BF model context. If \texttt{False},
  it will create a \texttt{BF.sample} site within a model. Defaults to
  \texttt{False}.
\item
  \emph{obs} (jnp.ndarray, optional): The observed value for this random
  variable. If provided, the sample site is conditioned on this value,
  and the function returns the observed value. If \texttt{None}, the
  site is treated as a latent (unobserved) random variable. Defaults to
  \texttt{None}.
\item
  \emph{name} (str, optional): The name of the sample site in a BF
  model. This is used to uniquely identify the random variable. Defaults
  to `x'.
\end{itemize}

\textbf{Returns} - When \texttt{sample=False}: A BF Chi2 distribution
object (for model building).

\begin{itemize}
\item
  When \texttt{sample=True}: A JAX array of samples drawn from the
  BernoulliLogits distribution (for direct sampling).
\item
  When \texttt{create\_obj=True}: The raw BF Chi2 object (for advanced
  use cases).
\end{itemize}

\textbf{Example Usage} from BayesForge import bf m = bf(platform=`cpu')
m.dist.chi2(df=3.0, sample = True)

\textbf{Wrapper of}
https://num.pyro.ai/en/stable/distributions.html\#chi2

\subsection{\texorpdfstring{\texttt{m.dist.circulant\_normal}}{m.dist.circulant\_normal}}\label{circulant_normal}

\begin{Shaded}
\begin{Highlighting}[]
\NormalTok{m.dist.circulant\_normal(}
\NormalTok{    loc: jax.Array,}
\NormalTok{    covariance\_row: jax.Array }\OperatorTok{=} \VariableTok{None}\NormalTok{,}
\NormalTok{    covariance\_rfft: jax.Array }\OperatorTok{=} \VariableTok{None}\NormalTok{,}
\NormalTok{    validate\_args}\OperatorTok{=}\VariableTok{None}\NormalTok{,}
\NormalTok{    name}\OperatorTok{=}\StringTok{\textquotesingle{}x\textquotesingle{}}\NormalTok{,}
\NormalTok{    obs}\OperatorTok{=}\VariableTok{None}\NormalTok{,}
\NormalTok{    mask}\OperatorTok{=}\VariableTok{None}\NormalTok{,}
\NormalTok{    infer}\OperatorTok{=}\VariableTok{None}\NormalTok{,}
\NormalTok{    sample}\OperatorTok{=}\VariableTok{False}\NormalTok{,}
\NormalTok{    seed}\OperatorTok{=}\VariableTok{None}\NormalTok{,}
\NormalTok{    shape}\OperatorTok{=}\NormalTok{(),}
\NormalTok{    event}\OperatorTok{=}\DecValTok{0}\NormalTok{,}
\NormalTok{    create\_obj}\OperatorTok{=}\VariableTok{False}\NormalTok{,}
\NormalTok{    to\_jax}\OperatorTok{=}\VariableTok{True}\NormalTok{,}
\NormalTok{)}
\end{Highlighting}
\end{Shaded}

\textbf{Circulant Normal Multivariate normal} Circulant Normal
Multivariate normal distribution with covariance matrix \(\mathbf{C}\)
that is positive-definite and circulant
\href{https://en.wikipedia.org/wiki/Matrix_normal_distribution}{1},
i.e., has periodic boundary conditions. The density of a sample
\(\mathbf{x}\in\mathbb{R}^n\) is the standard multivariate normal
density

\[
    p\left(\mathbf{x}\mid \mu,\mathbf{C}\right) =
    \frac{\left(\mathrm{det}\,\mathbf{C}\right)^{-1/2}}{\left(2\pi\right)^{n / 2}}
    \exp\left(-\frac{1}{2}\left(\mathbf{x}-\mu\right)^\intercal
    \mathbf{C}^{-1}\left(\mathbf{x}-\mu\right)\right),
\]

where \(\mathrm{det}\) denotes the determinant and \(^\intercal\) the
transpose. Circulant matrices can be diagnolized efficiently using the
discrete Fourier transform
\href{https://en.wikipedia.org/wiki/Matrix_normal_distribution}{1},
allowing the log likelihood to be evaluated in \(n \log n\) time for
\(n\) observations
\href{https://statproofbook.github.io/P/matn-pdf.html}{2}.

\begin{itemize}
\item
  \emph{loc:} Mean of the distribution \(\mu\).
\item
  \emph{covariance\_row:} First row of the circulant covariance matrix
  \(C\). Because of periodic boundary conditions, the covariance matrix
  is fully determined by its first row (see
  \texttt{jax.scipy.linalg.toeplitz} for further details).
\item
  \emph{covariance\_rfft:} Real part of the real fast Fourier transform
  of :code:\texttt{covariance\_row}, the first row of the circulant
  covariance matrix \(C\).
\item
  \emph{sample} (bool, optional): A control-flow argument. If
  \texttt{True}, the function will directly sample a raw JAX array from
  the distribution, bypassing the BF model context. If \texttt{False},
  it will create a \texttt{BF.sample} site within a model. Defaults to
  \texttt{False}.
\item
  \emph{obs} (jnp.ndarray, optional): The observed value for this random
  variable. If provided, the sample site is conditioned on this value,
  and the function returns the observed value. If \texttt{None}, the
  site is treated as a latent (unobserved) random variable. Defaults to
  \texttt{None}.
\item
  \emph{name} (str, optional): The name of the sample site in a BF
  model. This is used to uniquely identify the random variable. Defaults
  to `x'.
\end{itemize}

References:

\begin{enumerate}
\def\labelenumi{\arabic{enumi}.}
\item
  Wikipedia. (n.d.). Circulant matrix. Retrieved March 6, 2025, from
  https://en.wikipedia.org/wiki/Circulant\_matrix
\item
  Wood, A. T. A., \& Chan, G. (1994). Simulation of Stationary Gaussian
  Processes in \(\left[0, 1\right]^d\). \emph{Journal of Computational
  and Graphical Statistics}, 3(4), 409--432.
  https://doi.org/10.1080/10618600.1994.10474655
\end{enumerate}

\textbf{Args} - \emph{loc} : jnp.ndarray Mean of the distribution
\(\mu\).

\begin{itemize}
\item
  \emph{covariance\_row} : jnp.ndarray, optional. First row of the
  circulant covariance matrix \(\mathbf{C}\). Defaults to None.
\item
  \emph{covariance\_rfft} : jnp.ndarray, optional Real part of the real
  fast Fourier transform of \texttt{covariance\_row}. Defaults to None.
\end{itemize}

\textbf{Returns} - When \texttt{sample=False}: A BF Circulant Normal
Distribution distribution object (for model building).

\begin{itemize}
\item
  When \texttt{sample=True}: A JAX array of samples drawn from the
  BernoulliLogits distribution (for direct sampling).
\item
  When \texttt{create\_obj=True}: The raw BF Circulant Normal
  Distribution object (for advanced use cases).
\end{itemize}

\textbf{Wrapper of:
https://num.pyro.ai/en/stable/distributions.html\#circulantnormal}
\textbf{References} Wikipedia. (n.d.). Circulant matrix. Retrieved March
6, 2025, from https://en.wikipedia.org/wiki/Circulant\_matrix

Wood, A. T. A., \& Chan, G. (1994). Simulation of Stationary Gaussian
Processes in \([0,1]^d\). Journal of Computational and Graphical
Statistics, 3(4), 409--432.
https://doi.org/10.1080/10618600.1994.10474655

\subsection{\texorpdfstring{\texttt{m.dist.delta}}{m.dist.delta}}\label{delta}

\begin{Shaded}
\begin{Highlighting}[]
\NormalTok{m.dist.delta(}
\NormalTok{    v}\OperatorTok{=}\FloatTok{0.0}\NormalTok{,}
\NormalTok{    log\_density}\OperatorTok{=}\FloatTok{0.0}\NormalTok{,}
\NormalTok{    event\_dim}\OperatorTok{=}\DecValTok{0}\NormalTok{,}
\NormalTok{    validate\_args}\OperatorTok{=}\VariableTok{None}\NormalTok{,}
\NormalTok{    name}\OperatorTok{=}\StringTok{\textquotesingle{}x\textquotesingle{}}\NormalTok{,}
\NormalTok{    obs}\OperatorTok{=}\VariableTok{None}\NormalTok{,}
\NormalTok{    mask}\OperatorTok{=}\VariableTok{None}\NormalTok{,}
\NormalTok{    infer}\OperatorTok{=}\VariableTok{None}\NormalTok{,}
\NormalTok{    sample}\OperatorTok{=}\VariableTok{False}\NormalTok{,}
\NormalTok{    seed}\OperatorTok{=}\VariableTok{None}\NormalTok{,}
\NormalTok{    shape}\OperatorTok{=}\NormalTok{(),}
\NormalTok{    event}\OperatorTok{=}\DecValTok{0}\NormalTok{,}
\NormalTok{    create\_obj}\OperatorTok{=}\VariableTok{False}\NormalTok{,}
\NormalTok{    to\_jax}\OperatorTok{=}\VariableTok{True}\NormalTok{,}
\NormalTok{)}
\end{Highlighting}
\end{Shaded}

\textbf{Delta} The Delta distribution, also known as a point mass
distribution, assigns probability 1 to a single point and 0 elsewhere.
It's useful for representing deterministic variables or as a building
block for more complex distributions.

\[   
F(x) =
\begin{cases}
0, & x < a, \\
1, & x \ge a.
\end{cases}
\]

\textbf{Args} - \emph{v} (jnp.ndarray): The location of the point mass.

\begin{itemize}
\item
  \emph{log\_density} (float, optional): The log probability density of
  the point mass. This is primarily for creating distributions that are
  non-normalized or for specific advanced use cases. For a standard
  delta distribution, this should be 0. Defaults to 0.0.
\item
  \emph{event\_dim} (int, optional): The number of rightmost dimensions
  of \texttt{v} to interpret as event dimensions. Defaults to 0.
\item
  \emph{shape} (tuple): A multi-purpose argument for shaping. When
  \texttt{sample=False} (model building), this is used with
  \texttt{shape} to set the distribution's batch shape. When
  \texttt{sample=True} (direct sampling), this is used as
  \texttt{sample\_shape} to draw a raw JAX array of the given shape.
\item
  \emph{event} (int): The number of batch dimensions to reinterpret as
  event dimensions (used in model building).
\item
  \emph{mask} (jnp.ndarray, bool): Optional boolean array to mask
  observations.
\item
  \emph{create\_obj} (bool): If True, returns the raw BF distribution
  object instead of creating a sample site. This is essential for
  building complex distributions like \texttt{MixtureSameFamily}.
\item
  \emph{sample} (bool, optional): A control-flow argument. If
  \texttt{True}, the function will directly sample a raw JAX array from
  the distribution, bypassing the BF model context. If \texttt{False},
  it will create a \texttt{BF.sample} site within a model. Defaults to
  \texttt{False}.
\item
  \emph{obs} (jnp.ndarray, optional): The observed value for this random
  variable. If provided, the sample site is conditioned on this value,
  and the function returns the observed value. If \texttt{None}, the
  site is treated as a latent (unobserved) random variable. Defaults to
  \texttt{None}.
\item
  \emph{name} (str, optional): The name of the sample site in a BF
  model. This is used to uniquely identify the random variable. Defaults
  to `x'.
\end{itemize}

\textbf{Returns} - When \texttt{sample=False}: A BF Circulant Delta
Distribution distribution object (for model building).

\begin{itemize}
\item
  When \texttt{sample=True}: A JAX array of samples drawn from the
  BernoulliLogits distribution (for direct sampling).
\item
  When \texttt{create\_obj=True}: The raw BF Circulant Delta
  Distribution object (for advanced use cases).
\end{itemize}

\textbf{Example Usage} from BayesForge import bf m = bf(platform=`cpu')
m.dist.delta(v=0.0, sample=True)

\textbf{Wrapper of}
https://num.pyro.ai/en/stable/distributions.html\#delta

\subsection{\texorpdfstring{\texttt{m.dist.dirichlet}}{m.dist.dirichlet}}\label{dirichlet}

\begin{Shaded}
\begin{Highlighting}[]
\NormalTok{m.dist.dirichlet(}
\NormalTok{    concentration,}
\NormalTok{    validate\_args}\OperatorTok{=}\VariableTok{None}\NormalTok{,}
\NormalTok{    name}\OperatorTok{=}\StringTok{\textquotesingle{}x\textquotesingle{}}\NormalTok{,}
\NormalTok{    obs}\OperatorTok{=}\VariableTok{None}\NormalTok{,}
\NormalTok{    mask}\OperatorTok{=}\VariableTok{None}\NormalTok{,}
\NormalTok{    infer}\OperatorTok{=}\VariableTok{None}\NormalTok{,}
\NormalTok{    sample}\OperatorTok{=}\VariableTok{False}\NormalTok{,}
\NormalTok{    seed}\OperatorTok{=}\VariableTok{None}\NormalTok{,}
\NormalTok{    shape}\OperatorTok{=}\NormalTok{(),}
\NormalTok{    event}\OperatorTok{=}\DecValTok{0}\NormalTok{,}
\NormalTok{    create\_obj}\OperatorTok{=}\VariableTok{False}\NormalTok{,}
\NormalTok{    to\_jax}\OperatorTok{=}\VariableTok{True}\NormalTok{,}
\NormalTok{)}
\end{Highlighting}
\end{Shaded}

\textbf{Dirichlet} Samples from a Dirichlet distribution.

The Dirichlet distribution is a multivariate generalization of the Beta
distribution. It is a probability distribution over a simplex, which is
a set of vectors where each element is non-negative and sums to one. It
is often used as a prior distribution for categorical distributions.

For ( X = (X\_1, \dots, X\_K) \sim \mathrm{Dir}(\alpha)), the density is

\[
f( x; \alpha) = \frac{1}{B(\alpha)} \prod_{i=1}^K x_i^{\alpha_i - 1}
\]

for ( x) in the simplex, and 0 otherwise. Here (B(\alpha)) is the
multivariate Beta function (the normalization constant):

\[
B(\alpha) = \frac{\prod_{i=1}^K \Gamma(\alpha_i)}{\Gamma\left(\sum_{i=1}^K \alpha_i\right)}.
\]

\textbf{Args} - \emph{concentration} (jnp.ndarray): The concentration
parameter(s) of the Dirichlet distribution. Must be a positive array.

\begin{itemize}
\item
  \emph{shape} (tuple): A multi-purpose argument for shaping. When
  \texttt{sample=False} (model building), this is used with
  \texttt{shape} to set the distribution's batch shape. When
  \texttt{sample=True} (direct sampling), this is used as
  \texttt{sample\_shape} to draw a raw JAX array of the given shape.
\item
  \emph{event} (int): The number of batch dimensions to reinterpret as
  event dimensions (used in model building).
\item
  \emph{mask} (jnp.ndarray, bool): Optional boolean array to mask
  observations.
\item
  \emph{create\_obj} (bool): If True, returns the raw BF distribution
  object instead of creating a sample site. This is essential for
  building complex distributions like \texttt{MixtureSameFamily}.
\item
  \emph{sample} (bool, optional): A control-flow argument. If
  \texttt{True}, the function will directly sample a raw JAX array from
  the distribution, bypassing the BF model context. If \texttt{False},
  it will create a \texttt{BF.sample} site within a model. Defaults to
  \texttt{False}.
\item
  \emph{obs} (jnp.ndarray, optional): The observed value for this random
  variable. If provided, the sample site is conditioned on this value,
  and the function returns the observed value. If \texttt{None}, the
  site is treated as a latent (unobserved) random variable. Defaults to
  \texttt{None}.
\item
  \emph{name} (str, optional): The name of the sample site in a BF
  model. This is used to uniquely identify the random variable. Defaults
  to `x'.
\end{itemize}

\textbf{Returns} - When \texttt{sample=False}: A BF Dirichlet
Distribution distribution object (for model building).

\begin{itemize}
\item
  When \texttt{sample=True}: A JAX array of samples drawn from the
  Dirichlet distribution (for direct sampling).
\item
  When \texttt{create\_obj=True}: The raw BF Dirichlet Distribution
  object (for advanced use cases).
\end{itemize}

\textbf{Example Usage} from BayesForge import bf import jax.numpy as jnp
m = bf(platform=`cpu') m.dist.dirichlet(concentration=jnp.array({[}1.0,
1.0, 1.0{]}), sample=True)

\textbf{Wrapper of:
https://num.pyro.ai/en/stable/distributions.html\#dirichlet}
\textbf{References}
https://distribution-explorer.github.io/multivariate\_continuous/dirichlet.html
https://en.wikipedia.org/wiki/Dirichlet\_distribution

\subsection{\texorpdfstring{\texttt{m.dist.dirichlet\_multinomial}}{m.dist.dirichlet\_multinomial}}\label{dirichlet_multinomial}

\begin{Shaded}
\begin{Highlighting}[]
\NormalTok{m.dist.dirichlet\_multinomial(}
\NormalTok{    concentration,}
\NormalTok{    total\_count}\OperatorTok{=}\DecValTok{1}\NormalTok{,}
\NormalTok{    validate\_args}\OperatorTok{=}\VariableTok{None}\NormalTok{,}
\NormalTok{    name}\OperatorTok{=}\StringTok{\textquotesingle{}x\textquotesingle{}}\NormalTok{,}
\NormalTok{    obs}\OperatorTok{=}\VariableTok{None}\NormalTok{,}
\NormalTok{    mask}\OperatorTok{=}\VariableTok{None}\NormalTok{,}
\NormalTok{    infer}\OperatorTok{=}\VariableTok{None}\NormalTok{,}
\NormalTok{    sample}\OperatorTok{=}\VariableTok{False}\NormalTok{,}
\NormalTok{    seed}\OperatorTok{=}\VariableTok{None}\NormalTok{,}
\NormalTok{    shape}\OperatorTok{=}\NormalTok{(),}
\NormalTok{    event}\OperatorTok{=}\DecValTok{0}\NormalTok{,}
\NormalTok{    create\_obj}\OperatorTok{=}\VariableTok{False}\NormalTok{,}
\NormalTok{    to\_jax}\OperatorTok{=}\VariableTok{True}\NormalTok{,}
\NormalTok{)}
\end{Highlighting}
\end{Shaded}

\textbf{Dirichlet-Multinomial} Creates a Dirichlet-Multinomial compound
distribution, which is a Multinomial distribution with a Dirichlet prior
on its probabilities. It is often used in Bayesian statistics to model
count data where the proportions of categories are uncertain.

The \textbf{Dirichlet-Multinomial} distribution describes a model in
which you first draw a probability vector \[
\mathbf{p} = (p_1, \dots, p_K) \sim \mathrm{Dirichlet}(\alpha)
\] and then you draw counts \[
  \mathbf{x} = (x_1, \dots, x_K) \sim \mathrm{Multinomial}(n, \mathbf{p}).
\]

Integrating out \(\mathbf{p}\) gives you the marginal distribution of
\$\mathbf{x}). * It generalizes the \textbf{Beta-Binomial} model to
multiple categories (i.e.~(K \textgreater{} 2)). * A key feature is that
it introduces \textbf{overdispersion} relative to a standard
multinomial: counts are more ``clumped'' because the underlying
probabilities vary according to the Dirichlet prior.

\textbf{Args} - \emph{concentration} (jnp.ndarray): The concentration
parameter (alpha) for the Dirichlet prior. Values must be positive. The
last dimension is interpreted as the number of categories.

\begin{itemize}
\item
  \emph{total\_count} (int, jnp.ndarray, optional): The total number of
  trials (n). This must be a non-negative integer. Defaults to 1.
\item
  \emph{validate\_args} (bool, optional): Whether to enable validation
  of distribution parameters. Defaults to \texttt{None}.
\item
  \emph{name} (str, optional): The name of the sample site in a BF
  model. This is used to uniquely identify the random variable. Defaults
  to `x'.
\item
  \emph{obs} (jnp.ndarray, optional): The observed value for this random
  variable. If provided, the sample site is conditioned on this value,
  and the function returns the observed value. If \texttt{None}, the
  site is treated as a latent (unobserved) random variable. Defaults to
  \texttt{None}.
\item
  \emph{mask} (jnp.ndarray, bool, optional): Optional boolean array to
  mask observations. If provided, events with a \texttt{True} mask will
  be conditioned on \texttt{obs}, while the remaining events will be
  treated as latent variables. Defaults to \texttt{None}.
\item
  \emph{sample} (bool, optional): A control-flow argument. If
  \texttt{True}, the function will directly sample a raw JAX array from
  the distribution, bypassing the BF model context. If \texttt{False},
  it will create a \texttt{BF.sample} site within a model. Defaults to
  \texttt{False}.
\item
  \emph{seed} (int, optional): An integer seed for this draw when
  \texttt{sample=True}. By default (\texttt{None}) the draw comes from
  the session's random stream, set by \texttt{bf(rand\_seed=...)}, so
  consecutive draws differ but the whole sequence is reproducible. No
  effect when \texttt{sample=False}.
\item
  \emph{shape} (tuple, optional): A multi-purpose argument for shaping.
  When \texttt{sample=False} (model building), this is used with
  \texttt{.expand(shape)} to set the distribution's batch shape. When
  \texttt{sample=True} (direct sampling), this is used as
  \texttt{sample\_shape} to draw a raw JAX array of the given shape.
\item
  \emph{event} (int, optional): The number of batch dimensions to
  reinterpret as event dimensions (used in model building).
\item
  \emph{create\_obj} (bool, optional): If True, returns the raw BF
  distribution object instead of creating a sample site. This is
  essential for building complex distributions like
  \texttt{MixtureSameFamily}. Defaults to \texttt{False}.
\end{itemize}

\textbf{Returns} - When \texttt{sample=False}: A BF
dirichlet\_multinomial Distribution distribution object (for model
building).

\begin{itemize}
\item
  When \texttt{sample=True}: A JAX array of samples drawn from the
  dirichlet\_multinomial distribution (for direct sampling).
\item
  When \texttt{create\_obj=True}: The raw BF dirichlet\_multinomial
  Distribution object (for advanced use cases).
\end{itemize}

\textbf{Example Usage} from BayesForge import bf import jax.numpy as jnp
m = bf(platform=`cpu')

\begin{verbatim}
**Direct sampling**
**Sample a single vector of counts for 10 trials from 3 categories**
counts = m.dist.dirichlet_multinomial(concentration=jnp.array([1.0, 1.0, 1.0]),total_count=10,sample=True)

**Usage within a model**
def my_model(obs_data=None):
    **Define a prior on the concentration parameter**
    alpha = m.dist.half_cauchy(scale=jnp.ones(5), name='alpha', shape=(5,))

    **Model observed counts**
    with m.plate('data', len(obs_data)):
        y = m.dist.dirichlet_multinomial(
            concentration=alpha,
            total_count=100,
            name='y',
            obs=obs_data
        )
\end{verbatim}

\textbf{Wrapper of}
https://num.pyro.ai/en/stable/distributions.html\#dirichletmultinomial

\textbf{References}
https://en.wikipedia.org/wiki/Dirichlet-multinomial\_distribution
https://docs.scipy.org/doc/scipy/reference/generated/scipy.stats.dirichlet\_multinomial.html

\subsection{\texorpdfstring{\texttt{m.dist.discrete\_uniform}}{m.dist.discrete\_uniform}}\label{discrete_uniform}

\begin{Shaded}
\begin{Highlighting}[]
\NormalTok{m.dist.discrete\_uniform(}
\NormalTok{    low}\OperatorTok{=}\DecValTok{0}\NormalTok{,}
\NormalTok{    high}\OperatorTok{=}\DecValTok{1}\NormalTok{,}
\NormalTok{    validate\_args}\OperatorTok{=}\VariableTok{None}\NormalTok{,}
\NormalTok{    name}\OperatorTok{=}\StringTok{\textquotesingle{}x\textquotesingle{}}\NormalTok{,}
\NormalTok{    obs}\OperatorTok{=}\VariableTok{None}\NormalTok{,}
\NormalTok{    mask}\OperatorTok{=}\VariableTok{None}\NormalTok{,}
\NormalTok{    infer}\OperatorTok{=}\VariableTok{None}\NormalTok{,}
\NormalTok{    sample}\OperatorTok{=}\VariableTok{False}\NormalTok{,}
\NormalTok{    seed}\OperatorTok{=}\VariableTok{None}\NormalTok{,}
\NormalTok{    shape}\OperatorTok{=}\NormalTok{(),}
\NormalTok{    event}\OperatorTok{=}\DecValTok{0}\NormalTok{,}
\NormalTok{    create\_obj}\OperatorTok{=}\VariableTok{False}\NormalTok{,}
\NormalTok{    to\_jax}\OperatorTok{=}\VariableTok{True}\NormalTok{,}
\NormalTok{)}
\end{Highlighting}
\end{Shaded}

\textbf{Discrete Uniform} Samples from a Discrete Uniform distribution.

The Discrete Uniform distribution defines a uniform distribution over a
range of integers. It is characterized by a lower bound (\texttt{low})
and an upper bound (\texttt{high}), inclusive.

\[   
P(X = k) = \frac{1}{high - low + 1}, \quad k \in \{low, low+1, ..., high\}
\]

Otherwise (if \(k\) is outside the range), \(P(X = k) = 0\).

\textbf{Args} - \emph{low} (jnp.ndarray): The lower bound of the uniform
range, inclusive.

\begin{itemize}
\item
  \emph{high} (jnp.ndarray): The upper bound of the uniform range,
  inclusive.
\item
  \emph{shape} (tuple): A multi-purpose argument for shaping. When
  \texttt{sample=False} (model building), this is used with
  \texttt{shape} to set the distribution's batch shape. When
  \texttt{sample=True} (direct sampling), this is used as
  \texttt{sample\_shape} to draw a raw JAX array of the given shape.
\item
  \emph{event} (int): The number of batch dimensions to reinterpret as
  event dimensions (used in model building).
\item
  \emph{mask} (jnp.ndarray, bool): Optional boolean array to mask
  observations.
\item
  \emph{create\_obj} (bool): If True, returns the raw BF distribution
  object instead of creating a sample site. This is essential for
  building complex distributions like \texttt{MixtureSameFamily}.
\item
  \emph{sample} (bool, optional): A control-flow argument. If
  \texttt{True}, the function will directly sample a raw JAX array from
  the distribution, bypassing the BF model context. If \texttt{False},
  it will create a \texttt{BF.sample} site within a model. Defaults to
  \texttt{False}.
\item
  \emph{obs} (jnp.ndarray, optional): The observed value for this random
  variable. If provided, the sample site is conditioned on this value,
  and the function returns the observed value. If \texttt{None}, the
  site is treated as a latent (unobserved) random variable. Defaults to
  \texttt{None}.
\item
  \emph{name} (str, optional): The name of the sample site in a BF
  model. This is used to uniquely identify the random variable. Defaults
  to `x'.
\end{itemize}

\textbf{Returns} - When \texttt{sample=False}: A BF DiscreteUniform
Distribution distribution object (for model building).

\begin{itemize}
\item
  When \texttt{sample=True}: A JAX array of samples drawn from the
  DiscreteUniform distribution (for direct sampling).
\item
  When \texttt{create\_obj=True}: The raw BF DiscreteUniform
  Distribution object (for advanced use cases).
\end{itemize}

\textbf{Example Usage} from BayesForge import bf m = bf(platform=`cpu')
m.dist.discrete\_uniform(low=0, high=5, sample=True)

\textbf{Wrapper of}
https://num.pyro.ai/en/stable/distributions.html\#discreteuniform

\textbf{References}

\subsection{\texorpdfstring{\texttt{m.dist.doubly\_truncated\_power\_law}}{m.dist.doubly\_truncated\_power\_law}}\label{doubly_truncated_power_law}

\begin{Shaded}
\begin{Highlighting}[]
\NormalTok{m.dist.doubly\_truncated\_power\_law(}
\NormalTok{    alpha,}
\NormalTok{    low,}
\NormalTok{    high,}
\NormalTok{    validate\_args}\OperatorTok{=}\VariableTok{None}\NormalTok{,}
\NormalTok{    name}\OperatorTok{=}\StringTok{\textquotesingle{}x\textquotesingle{}}\NormalTok{,}
\NormalTok{    obs}\OperatorTok{=}\VariableTok{None}\NormalTok{,}
\NormalTok{    mask}\OperatorTok{=}\VariableTok{None}\NormalTok{,}
\NormalTok{    infer}\OperatorTok{=}\VariableTok{None}\NormalTok{,}
\NormalTok{    sample}\OperatorTok{=}\VariableTok{False}\NormalTok{,}
\NormalTok{    seed}\OperatorTok{=}\VariableTok{None}\NormalTok{,}
\NormalTok{    shape}\OperatorTok{=}\NormalTok{(),}
\NormalTok{    event}\OperatorTok{=}\DecValTok{0}\NormalTok{,}
\NormalTok{    create\_obj}\OperatorTok{=}\VariableTok{False}\NormalTok{,}
\NormalTok{    to\_jax}\OperatorTok{=}\VariableTok{True}\NormalTok{,}
\NormalTok{)}
\end{Highlighting}
\end{Shaded}

\textbf{Doubly Truncated Power Law} This distribution represents a
continuous power law with a finite support bounded between \texttt{low}
and \texttt{high}, and with an exponent \texttt{alpha}. It is normalized
over the interval \texttt{{[}low,\ high{]}} to ensure the area under the
density function is 1.

Key parameters:

\begin{itemize}
\tightlist
\item
  (\alpha \textgreater{} 0): shape (exponent) of the power law
\item
  (x\_\{\min\}): lower cutoff (minimum possible value)
\item
  (x\_\{\max\}): upper cutoff (maximum possible value)
\end{itemize}

The probability density function (PDF) is defined as:

\[
 f(x) = \begin{cases}
 \displaystyle \frac{(1-\alpha), x^{-\alpha}}{x_{\max}^{1-\alpha} - x_{\min}^{1-\alpha}}, & \alpha \neq 1, ; x_{\min} \le x         \le x_{\max},\\
 \displaystyle \frac{1}{x, \ln\left(x_{\max}/x_{\min}\right)}, & \alpha = 1, ; x_{\min} \le x \le x_{\max},\\
 0, & \text{otherwise}.
 \end{cases}
 \]

where the normalization constant \(Z(\alpha,a,b)\) is given by

\[
 Z(\alpha,a,b) =
 \begin{cases}
 \log(b) - \log(a), & \text{if } \alpha = -1, \\
 \dfrac{b^{1+\alpha} - a^{1+\alpha}}{1+\alpha}, & \text{otherwise}.
 \end{cases}
 \]

This distribution is useful for modeling data that follows a power-law
behavior but is naturally bounded due to measurement or theoretical
constraints (e.g., finite-size systems).

\textbf{Args} - \emph{alpha} (float or array-like): Power-law exponent.

\begin{itemize}
\item
  \emph{low} (float or array-like): Lower bound of the distribution
  (must be \(≥\) 0).
\item
  \emph{high} (float or array-like): Upper bound of the distribution
  (must be \(>\) 0).
\item
  \emph{shape} (tuple): A multi-purpose argument for shaping. When
  \texttt{sample=False} (model building), this is used with
  \texttt{shape} to set the distribution's batch shape. When
  \texttt{sample=True} (direct sampling), this is used as
  \texttt{sample\_shape} to draw a raw JAX array of the given shape.
\item
  *validate\_args (bool, optional): Whether to validate the arguments.
  Defaults to True.
\item
  \emph{sample} (bool, optional): A control-flow argument. If
  \texttt{True}, the function will directly sample a raw JAX array from
  the distribution, bypassing the BF model context. If \texttt{False},
  it will create a \texttt{BF.sample} site within a model. Defaults to
  \texttt{False}.
\item
  \emph{obs} (jnp.ndarray, optional): The observed value for this random
  variable. If provided, the sample site is conditioned on this value,
  and the function returns the observed value. If \texttt{None}, the
  site is treated as a latent (unobserved) random variable. Defaults to
  \texttt{None}.
\item
  \emph{name} (str, optional): The name of the sample site in a BF
  model. This is used to uniquely identify the random variable. Defaults
  to `x'.
\end{itemize}

\textbf{Returns} - When \texttt{sample=False}: A BF
doubly\_truncated\_power\_law Distribution distribution object (for
model building).

\begin{itemize}
\item
  When \texttt{sample=True}: A JAX array of samples drawn from the
  doubly\_truncated\_power\_law distribution (for direct sampling).
\item
  When \texttt{create\_obj=True}: The raw BF
  doubly\_truncated\_power\_law Distribution object (for advanced use
  cases).
\end{itemize}

\textbf{Example Usage} from BayesForge import bf m =
bf(platform=`cpu')\\
m.dist.doubly\_truncated\_power\_law(low=0.1, high=10.0, alpha=2.0,
sample=True)

\textbf{References}
https://reference.wolframcloud.com/language/ref/TruncatedDistribution.html

\subsection{\texorpdfstring{\texttt{m.dist.euler\_maruyama}}{m.dist.euler\_maruyama}}\label{euler_maruyama}

\begin{Shaded}
\begin{Highlighting}[]
\NormalTok{m.dist.euler\_maruyama(}
\NormalTok{    t,}
\NormalTok{    sde\_fn,}
\NormalTok{    init\_dist,}
\NormalTok{    validate\_args}\OperatorTok{=}\VariableTok{None}\NormalTok{,}
\NormalTok{    name}\OperatorTok{=}\StringTok{\textquotesingle{}x\textquotesingle{}}\NormalTok{,}
\NormalTok{    obs}\OperatorTok{=}\VariableTok{None}\NormalTok{,}
\NormalTok{    mask}\OperatorTok{=}\VariableTok{None}\NormalTok{,}
\NormalTok{    infer}\OperatorTok{=}\VariableTok{None}\NormalTok{,}
\NormalTok{    sample}\OperatorTok{=}\VariableTok{False}\NormalTok{,}
\NormalTok{    seed}\OperatorTok{=}\VariableTok{None}\NormalTok{,}
\NormalTok{    shape}\OperatorTok{=}\NormalTok{(),}
\NormalTok{    event}\OperatorTok{=}\DecValTok{0}\NormalTok{,}
\NormalTok{    create\_obj}\OperatorTok{=}\VariableTok{False}\NormalTok{,}
\NormalTok{    to\_jax}\OperatorTok{=}\VariableTok{True}\NormalTok{,}
\NormalTok{)}
\end{Highlighting}
\end{Shaded}

\textbf{Euler--Maruyama} Euler--Maruyama methode is a method for the
approximate numerical solution of a stochastic differential equation
(SDE). It simulates the solution to an SDE by iteratively applying the
Euler method to each time step, incorporating a random perturbation to
account for the diffusion term.

\[
   dX_t = f(X_t, t) dt + g(X_t, t) dW_t
\]

where: - \(X_t\) is the state of the system at time \(t\). -
\(f(X_t, t)\) is the drift coefficient. - \(g(X_t, t)\) is the diffusion
coefficient. - \(dW_t\) is a Wiener process (Brownian motion).

\textbf{Args} - \emph{t} (jnp.ndarray): Discretized time steps.

\begin{itemize}
\item
  \emph{sde\_fn} (callable): A function that takes the current state and
  time as input and returns the drift and diffusion coefficients.
\item
  \emph{init\_dist} (Distribution): The initial distribution of the
  system.
\item
  \emph{shape} (tuple): A multi-purpose argument for shaping. When
  \texttt{sample=False} (model building), this is used with
  \texttt{shape} to set the distribution's batch shape. When
  \texttt{sample=True} (direct sampling), this is used as
  \texttt{sample\_shape} to draw a raw JAX array of the given shape.
\item
  *sample\_shape (tuple, optional): The shape of the samples to draw.
  Defaults to None.
\item
  *validate\_args (bool, optional): Whether to validate the arguments.
  Defaults to True.
\item
  \emph{sample} (bool, optional): A control-flow argument. If
  \texttt{True}, the function will directly sample a raw JAX array from
  the distribution, bypassing the BF model context. If \texttt{False},
  it will create a \texttt{BF.sample} site within a model. Defaults to
  \texttt{False}.
\item
  \emph{obs} (jnp.ndarray, optional): The observed value for this random
  variable. If provided, the sample site is conditioned on this value,
  and the function returns the observed value. If \texttt{None}, the
  site is treated as a latent (unobserved) random variable. Defaults to
  \texttt{None}.
\item
  \emph{name} (str, optional): The name of the sample site in a BF
  model. This is used to uniquely identify the random variable. Defaults
  to `x'.
\end{itemize}

\textbf{Returns} jnp.ndarray: Samples drawn from the Euler--Maruyama
distribution.

\textbf{Example Usage} from BayesForge import bf import jax.numpy as jnp
m = bf(platform=`cpu') m.dist.euler\_maruyama(t=jnp.array({[}0.0, 0.1,
0.2{]}), sde\_fn=lambda x, t: (x, 1.0), init\_dist=m.dist.normal(0.0,
1.0, create\_obj=True), sample = True)

\subsection{\texorpdfstring{\texttt{m.dist.exponential}}{m.dist.exponential}}\label{exponential}

\begin{Shaded}
\begin{Highlighting}[]
\NormalTok{m.dist.exponential(}
\NormalTok{    rate}\OperatorTok{=}\FloatTok{1.0}\NormalTok{,}
\NormalTok{    validate\_args}\OperatorTok{=}\VariableTok{None}\NormalTok{,}
\NormalTok{    name}\OperatorTok{=}\StringTok{\textquotesingle{}x\textquotesingle{}}\NormalTok{,}
\NormalTok{    obs}\OperatorTok{=}\VariableTok{None}\NormalTok{,}
\NormalTok{    mask}\OperatorTok{=}\VariableTok{None}\NormalTok{,}
\NormalTok{    infer}\OperatorTok{=}\VariableTok{None}\NormalTok{,}
\NormalTok{    sample}\OperatorTok{=}\VariableTok{False}\NormalTok{,}
\NormalTok{    seed}\OperatorTok{=}\VariableTok{None}\NormalTok{,}
\NormalTok{    shape}\OperatorTok{=}\NormalTok{(),}
\NormalTok{    event}\OperatorTok{=}\DecValTok{0}\NormalTok{,}
\NormalTok{    create\_obj}\OperatorTok{=}\VariableTok{False}\NormalTok{,}
\NormalTok{    to\_jax}\OperatorTok{=}\VariableTok{True}\NormalTok{,}
\NormalTok{)}
\end{Highlighting}
\end{Shaded}

\textbf{Exponential} The Exponential distribution is a continuous
probability distribution that models the time until an event occurs in a
Poisson process, where events occur continuously and independently at a
constant average rate. It is often used to model the duration of events,
such as the time until a machine fails or the length of a phone call.

For ( x \ge 0 ) \[   
  F(x; \lambda) = 1 - e^{-\lambda x}
  \]

Zhere ( \lambda \textgreater{} 0 ) is the rate parameter (the mean is (
1/\lambda )).

\textbf{Args} - \emph{rate} (jnp.ndarray): The rate parameter,
\(\lambda\). Must be positive.

\begin{itemize}
\item
  \emph{shape} (tuple): A multi-purpose argument for shaping. When
  \texttt{sample=False} (model building), this is used with
  \texttt{shape} to set the distribution's batch shape. When
  \texttt{sample=True} (direct sampling), this is used as
  \texttt{sample\_shape} to draw a raw JAX array of the given shape.
\item
  \emph{event} (int): The number of batch dimensions to reinterpret as
  event dimensions (used in model building).
\item
  \emph{mask} (jnp.ndarray, bool): Optional boolean array to mask
  observations.
\item
  \emph{create\_obj} (bool): If True, returns the raw BF distribution
  object instead of creating a sample site. This is essential for
  building complex distributions like \texttt{MixtureSameFamily}.
\item
  \emph{sample} (bool, optional): A control-flow argument. If
  \texttt{True}, the function will directly sample a raw JAX array from
  the distribution, bypassing the BF model context. If \texttt{False},
  it will create a \texttt{BF.sample} site within a model. Defaults to
  \texttt{False}.
\item
  \emph{obs} (jnp.ndarray, optional): The observed value for this random
  variable. If provided, the sample site is conditioned on this value,
  and the function returns the observed value. If \texttt{None}, the
  site is treated as a latent (unobserved) random variable. Defaults to
  \texttt{None}.
\item
  \emph{name} (str, optional): The name of the sample site in a BF
  model. This is used to uniquely identify the random variable. Defaults
  to `x'.
\end{itemize}

\textbf{Returns} - When \texttt{sample=False}: A BF Exponential
distribution object (for model building). - When \texttt{sample=True}: A
JAX array of samples drawn from the Exponential distribution (for direct
sampling). - When \texttt{create\_obj=True}: The raw BF distribution
object (for advanced use cases).

\textbf{Example Usage} from BayesForge import bf m = bf(platform=`cpu')
m.dist.exponential(rate=1.0, sample=True)

\textbf{Wrapper of}
https://num.pyro.ai/en/stable/distributions.html\#exponential

\textbf{References}
https://reference.wolfram.com/language/ref/ExponentialDistribution.html
https://en.wikipedia.org/wiki/Exponential\_distribution

\subsection{\texorpdfstring{\texttt{m.dist.factor}}{m.dist.factor}}\label{factor}

\begin{Shaded}
\begin{Highlighting}[]
\NormalTok{m.dist.factor(name, log\_factor)}
\end{Highlighting}
\end{Shaded}

\textbf{Factor} Adds an arbitrary log-probability (log-likelihood)
factor to the probabilistic model.

This primitive is used to incorporate custom log-likelihood terms,
marginalization over discrete/latent variables, Jacobian volume
adjustments, or unnormalized potentials that cannot be expressed as
standard distribution sample sites.

\textbf{Args} name (str): Unique identifier for this factor statement
site in the model trace. log\_factor (jnp.ndarray or float): A scalar or
batched log-probability value. If batched, it should typically be
evaluated inside a corresponding \texttt{plate} context matching its
batch shape.

\textbf{Returns} None

\textbf{Example Usage} from BayesForge import bf import jax.numpy as jnp
m = bf(platform=`cpu') y = jnp.array({[}0.3, -1.2, 0.8{]})

\begin{verbatim}
def model(y):
    mu = m.dist.normal(0, 1, name='mu')
    **custom (here: Laplace) log-likelihood, one term per observation**
    log_lik_per_obs = -jnp.abs(y - mu)
    with m.dist.plate("obs", y.shape[0]):
        m.dist.factor("custom_log_lik", log_lik_per_obs)

m.fit(model, obs=dict(y=y), num_samples=100, num_warmup=100, num_chains=1, progress_bar=False)
\end{verbatim}

\textbf{Wrapper of}
https://num.pyro.ai/en/stable/primitives.html\#numpyro.primitives.factor

\subsection{\texorpdfstring{\texttt{m.dist.folded\_distribution}}{m.dist.folded\_distribution}}\label{folded_distribution}

\begin{Shaded}
\begin{Highlighting}[]
\NormalTok{m.dist.folded\_distribution(}
\NormalTok{    base\_dist,}
\NormalTok{    validate\_args}\OperatorTok{=}\VariableTok{None}\NormalTok{,}
\NormalTok{    name}\OperatorTok{=}\StringTok{\textquotesingle{}x\textquotesingle{}}\NormalTok{,}
\NormalTok{    obs}\OperatorTok{=}\VariableTok{None}\NormalTok{,}
\NormalTok{    mask}\OperatorTok{=}\VariableTok{None}\NormalTok{,}
\NormalTok{    infer}\OperatorTok{=}\VariableTok{None}\NormalTok{,}
\NormalTok{    sample}\OperatorTok{=}\VariableTok{False}\NormalTok{,}
\NormalTok{    seed}\OperatorTok{=}\VariableTok{None}\NormalTok{,}
\NormalTok{    shape}\OperatorTok{=}\NormalTok{(),}
\NormalTok{    event}\OperatorTok{=}\DecValTok{0}\NormalTok{,}
\NormalTok{    create\_obj}\OperatorTok{=}\VariableTok{False}\NormalTok{,}
\NormalTok{    to\_jax}\OperatorTok{=}\VariableTok{True}\NormalTok{,}
\NormalTok{)}
\end{Highlighting}
\end{Shaded}

\textbf{Folded} Samples from a Folded distribution, which is the
absolute value of a base univariate distribution. This distribution
reflects the base distribution across the origin, effectively taking the
absolute value of each sample.

\[
   p(x) = \sum_{k=-\infty}^{\infty} p(x - 2k)
\]

\textbf{Args} - \emph{loc} (float, optional): Location parameter of the
base distribution. Defaults to 0.0.

\begin{itemize}
\item
  \emph{sample} (float, optional): Scale parameter of the base
  distribution. Defaults to 1.0.
\item
  \emph{shape} (tuple): A multi-purpose argument for shaping. When
  \texttt{sample=False} (model building), this is used with
  \texttt{shape} to set the distribution's batch shape. When
  \texttt{sample=True} (direct sampling), this is used as
  \texttt{sample\_shape} to draw a raw JAX array of the given shape.
\item
  \emph{event} (int): The number of batch dimensions to reinterpret as
  event dimensions (used in model building).
\item
  \emph{mask} (jnp.ndarray, bool): Optional boolean array to mask
  observations.
\item
  \emph{create\_obj} (bool): If True, returns the raw BF distribution
  object instead of creating a sample site. This is essential for
  building complex distributions like \texttt{MixtureSameFamily}.
\item
  \emph{sample} (bool, optional): A control-flow argument. If
  \texttt{True}, the function will directly sample a raw JAX array from
  the distribution, bypassing the BF model context. If \texttt{False},
  it will create a \texttt{BF.sample} site within a model. Defaults to
  \texttt{False}.
\item
  \emph{obs} (jnp.ndarray, optional): The observed value for this random
  variable. If provided, the sample site is conditioned on this value,
  and the function returns the observed value. If \texttt{None}, the
  site is treated as a latent (unobserved) random variable. Defaults to
  \texttt{None}.
\item
  \emph{name} (str, optional): The name of the sample site in a BF
  model. This is used to uniquely identify the random variable. Defaults
  to `x'.
\end{itemize}

\textbf{Returns} BF FoldedDistribution distribution object (for model
building) when \texttt{sample=False}. JAX array of samples drawn from
the FoldedDistribution distribution (for direct sampling) when
\texttt{sample=True}. The raw BF distribution object (for advanced use
cases) when \texttt{create\_obj=True}.

\textbf{Example Usage} from BayesForge import bf m = bf(platform=`cpu')
m.dist.folded\_distribution(m.dist.normal(loc=0.0, scale=1.0,
create\_obj = True), sample=True)

\textbf{Wrapper of}
https://num.pyro.ai/en/stable/distributions.html\#foldeddistribution

\subsection{\texorpdfstring{\texttt{m.dist.gamma}}{m.dist.gamma}}\label{gamma}

\begin{Shaded}
\begin{Highlighting}[]
\NormalTok{m.dist.gamma(}
\NormalTok{    concentration,}
\NormalTok{    rate}\OperatorTok{=}\FloatTok{1.0}\NormalTok{,}
\NormalTok{    validate\_args}\OperatorTok{=}\VariableTok{None}\NormalTok{,}
\NormalTok{    name}\OperatorTok{=}\StringTok{\textquotesingle{}x\textquotesingle{}}\NormalTok{,}
\NormalTok{    obs}\OperatorTok{=}\VariableTok{None}\NormalTok{,}
\NormalTok{    mask}\OperatorTok{=}\VariableTok{None}\NormalTok{,}
\NormalTok{    infer}\OperatorTok{=}\VariableTok{None}\NormalTok{,}
\NormalTok{    sample}\OperatorTok{=}\VariableTok{False}\NormalTok{,}
\NormalTok{    seed}\OperatorTok{=}\VariableTok{None}\NormalTok{,}
\NormalTok{    shape}\OperatorTok{=}\NormalTok{(),}
\NormalTok{    event}\OperatorTok{=}\DecValTok{0}\NormalTok{,}
\NormalTok{    create\_obj}\OperatorTok{=}\VariableTok{False}\NormalTok{,}
\NormalTok{    to\_jax}\OperatorTok{=}\VariableTok{True}\NormalTok{,}
\NormalTok{)}
\end{Highlighting}
\end{Shaded}

\textbf{Gamma} Samples from a Gamma distribution.

The Gamma distribution is a continuous probability distribution
frequently used in Bayesian statistics, particularly as a prior
distribution for variance parameters. It is defined by two positive
shape parameters: concentration ((k)) and rate ((\lambda)).

There are two common parameterizations:

\begin{enumerate}
\def\labelenumi{\arabic{enumi}.}
\item
  \textbf{Shape--scale}: parameters ((\alpha, \theta)), with density {[}
  f(x) = \frac{1}{\Gamma(\alpha) \theta^\alpha} x\^{}\{\alpha - 1\}
  e\^{}\{-x/\theta\}, \quad x \ge 0{]}
\item
  \textbf{Shape--rate}: parameters ((\alpha, \lambda)) where (\lambda =
  1/\theta), with density {[} f(x) =
  \frac{\lambda^\alpha}{\Gamma(\alpha)} x\^{}\{\alpha - 1\}
  e\^{}\{-\lambda x\}, \quad x \ge 0{]}
\end{enumerate}

\textbf{Args} - \emph{concentration} (jnp.ndarray): The shape parameter
of the Gamma distribution (k \textgreater{} 0).

\begin{itemize}
\item
  \emph{rate} (jnp.ndarray): The rate parameter of the Gamma
  distribution (theta \textgreater{} 0).
\item
  \emph{shape} (tuple): A multi-purpose argument for shaping. When
  \texttt{sample=False} (model building), this is used with
  \texttt{shape} to set the distribution's batch shape. When
  \texttt{sample=True} (direct sampling), this is used as
  \texttt{sample\_shape} to draw a raw JAX array of the given shape.
\item
  \emph{event} (int): The number of batch dimensions to reinterpret as
  event dimensions (used in model building).
\item
  \emph{mask} (jnp.ndarray, bool): Optional boolean array to mask
  observations.
\item
  \emph{create\_obj} (bool): If True, returns the raw BF distribution
  object instead of creating a sample site. This is essential for
  building complex distributions like \texttt{MixtureSameFamily}.
\item
  \emph{sample} (bool, optional): A control-flow argument. If
  \texttt{True}, the function will directly sample a raw JAX array from
  the distribution, bypassing the BF model context. If \texttt{False},
  it will create a \texttt{BF.sample} site within a model. Defaults to
  \texttt{False}.
\item
  \emph{obs} (jnp.ndarray, optional): The observed value for this random
  variable. If provided, the sample site is conditioned on this value,
  and the function returns the observed value. If \texttt{None}, the
  site is treated as a latent (unobserved) random variable. Defaults to
  \texttt{None}.
\item
  \emph{name} (str, optional): The name of the sample site in a BF
  model. This is used to uniquely identify the random variable. Defaults
  to `x'.
\end{itemize}

\textbf{Returns} Gamma: A BF Gamma distribution object (for model
building).

\begin{verbatim}
jnp.ndarray: A JAX array of samples drawn from the Gamma distribution (for direct sampling).

Gamma: The raw BF distribution object (for advanced use cases).
\end{verbatim}

\textbf{Example Usage} from BayesForge import bf m = bf(platform=`cpu')
m.dist.gamma(concentration=2.0, rate=0.5, sample=True)

\textbf{Wrapper of}
https://num.pyro.ai/en/stable/distributions.html\#gamma

\textbf{References}
https://reference.wolfram.com/language/ref/GammaDistribution.html
https://en.wikipedia.org/wiki/Gamma\_distribution

\subsection{\texorpdfstring{\texttt{m.dist.gamma\_poisson}}{m.dist.gamma\_poisson}}\label{gamma_poisson}

\begin{Shaded}
\begin{Highlighting}[]
\NormalTok{m.dist.gamma\_poisson(}
\NormalTok{    concentration,}
\NormalTok{    rate}\OperatorTok{=}\FloatTok{1.0}\NormalTok{,}
\NormalTok{    validate\_args}\OperatorTok{=}\VariableTok{None}\NormalTok{,}
\NormalTok{    name}\OperatorTok{=}\StringTok{\textquotesingle{}x\textquotesingle{}}\NormalTok{,}
\NormalTok{    obs}\OperatorTok{=}\VariableTok{None}\NormalTok{,}
\NormalTok{    mask}\OperatorTok{=}\VariableTok{None}\NormalTok{,}
\NormalTok{    infer}\OperatorTok{=}\VariableTok{None}\NormalTok{,}
\NormalTok{    sample}\OperatorTok{=}\VariableTok{False}\NormalTok{,}
\NormalTok{    seed}\OperatorTok{=}\VariableTok{None}\NormalTok{,}
\NormalTok{    shape}\OperatorTok{=}\NormalTok{(),}
\NormalTok{    event}\OperatorTok{=}\DecValTok{0}\NormalTok{,}
\NormalTok{    create\_obj}\OperatorTok{=}\VariableTok{False}\NormalTok{,}
\NormalTok{    to\_jax}\OperatorTok{=}\VariableTok{True}\NormalTok{,}
\NormalTok{)}
\end{Highlighting}
\end{Shaded}

\textbf{Gamma Poisson} The Gamma-Poisson distribution, also known as the
Negative Binomial distribution, models overdispersed count data. It
arises from a hierarchical process where the rate parameter of a Poisson
distribution is itself a random variable following a Gamma distribution.
This structure allows the model to capture variability in count data
that exceeds what is predicted by the Poisson distribution, making it
suitable for applications like modeling RNA-sequencing data and
microbial count.

\[
  f(x) = \frac{\Gamma(x + \beta) \alpha^x}{\Gamma(\beta)(1 + \alpha)^{\beta + x} x!}, \quad x = 0, 1, 2, \ldots
  \] Where \(\Gamma\) denotes the gamma function.

This CDF expresses the probability that the count variable ( X ) takes a
value less than or equal to ( x ), incorporating the overdispersion
introduced by the Gamma distribution.

\textbf{Args} - \emph{concentration} (jnp.ndarray): Shape parameter
(alpha) of the Gamma distribution.

\begin{itemize}
\item
  \emph{rate} (jnp.ndarray): Rate parameter (beta) for the Gamma
  distribution.
\item
  \emph{shape} (tuple): A multi-purpose argument for shaping. When
  \texttt{sample=False} (model building), this is used with
  \texttt{shape} to set the distribution's batch shape. When
  \texttt{sample=True} (direct sampling), this is used as
  \texttt{sample\_shape} to draw a raw JAX array of the given shape.
\item
  \emph{event} (int): The number of batch dimensions to reinterpret as
  event dimensions (used in model building).
\item
  \emph{mask} (jnp.ndarray, bool): Optional boolean array to mask
  observations.
\item
  \emph{create\_obj} (bool, optional): If True, returns the raw BF
  distribution object instead of creating a sample site. This is
  essential for building complex distributions like
  \texttt{MixtureSameFamily}. Defaults to False.
\item
  \emph{sample} (bool, optional): A control-flow argument. If
  \texttt{True}, the function will directly sample a raw JAX array from
  the distribution, bypassing the BF model context. If \texttt{False},
  it will create a \texttt{BF.sample} site within a model. Defaults to
  \texttt{False}.
\item
  \emph{obs} (jnp.ndarray, optional): The observed value for this random
  variable. If provided, the sample site is conditioned on this value,
  and the function returns the observed value. If \texttt{None}, the
  site is treated as a latent (unobserved) random variable. Defaults to
  \texttt{None}.
\item
  \emph{name} (str, optional): The name of the sample site in a BF
  model. This is used to uniquely identify the random variable. Defaults
  to `x'.
\end{itemize}

\textbf{Returns} - When \texttt{sample=False}: A BF GammaPoisson
distribution object (for model building).

\begin{itemize}
\item
  When \texttt{sample=True}: A JAX array of samples drawn from the
  GammaPoisson distribution (for direct sampling).
\item
  When \texttt{create\_obj=True}: The raw BF distribution object (for
  advanced use cases).

  \textbf{Example Usage} from BayesForge import bf m =
  bf(platform=`cpu') m.dist.gamma\_poisson(concentration=1.0, rate=2.0,
  sample=True)

  \textbf{Wrapper of}
  https://num.pyro.ai/en/stable/distributions.html\#gammapoisson

  \textbf{References}
  \href{https://en.wikipedia.org/wiki/Matrix_normal_distribution}{1}
  https://www.statisticshowto.com/gamma-poisson-distribution
  https://en.wikipedia.org/wiki/Negative\_binomial\_distribution
\end{itemize}

\subsection{\texorpdfstring{\texttt{m.dist.gaussian\_copula}}{m.dist.gaussian\_copula}}\label{gaussian_copula}

\begin{Shaded}
\begin{Highlighting}[]
\NormalTok{m.dist.gaussian\_copula(}
\NormalTok{    marginal\_dist,}
\NormalTok{    correlation\_matrix}\OperatorTok{=}\VariableTok{None}\NormalTok{,}
\NormalTok{    correlation\_cholesky}\OperatorTok{=}\VariableTok{None}\NormalTok{,}
\NormalTok{    validate\_args}\OperatorTok{=}\VariableTok{None}\NormalTok{,}
\NormalTok{    name}\OperatorTok{=}\StringTok{\textquotesingle{}x\textquotesingle{}}\NormalTok{,}
\NormalTok{    obs}\OperatorTok{=}\VariableTok{None}\NormalTok{,}
\NormalTok{    mask}\OperatorTok{=}\VariableTok{None}\NormalTok{,}
\NormalTok{    infer}\OperatorTok{=}\VariableTok{None}\NormalTok{,}
\NormalTok{    sample}\OperatorTok{=}\VariableTok{False}\NormalTok{,}
\NormalTok{    seed}\OperatorTok{=}\VariableTok{None}\NormalTok{,}
\NormalTok{    shape}\OperatorTok{=}\NormalTok{(),}
\NormalTok{    event}\OperatorTok{=}\DecValTok{0}\NormalTok{,}
\NormalTok{    create\_obj}\OperatorTok{=}\VariableTok{False}\NormalTok{,}
\NormalTok{    to\_jax}\OperatorTok{=}\VariableTok{True}\NormalTok{,}
\NormalTok{)}
\end{Highlighting}
\end{Shaded}

\textbf{Gaussian Copula} A distribution that links the
\texttt{batch\_shape{[}:-1{]}} of a marginal distribution with a
multivariate Gaussian copula, modelling the correlation between the
axes. A copula is a multivariate distribution over the uniform
distribution on {[}0, 1{]}. The Gaussian copula links the marginal
distributions through a multivariate normal distribution.

\textbf{Args} - \emph{marginal\_dist} (Distribution): Distribution whose
last batch axis is to be coupled.

\begin{itemize}
\item
  \emph{correlation\_matrix} (array\_like, optional): Correlation matrix
  of the coupling multivariate normal distribution. Defaults to None.
\item
  \emph{correlation\_cholesky} (array\_like, optional): Correlation
  Cholesky factor of the coupling multivariate normal distribution.
  Defaults to None.
\item
  \emph{shape} (tuple): A multi-purpose argument for shaping. When
  \texttt{sample=False} (model building), this is used with
  \texttt{shape} to set the distribution's batch shape. When
  \texttt{sample=True} (direct sampling), this is used as
  \texttt{sample\_shape} to draw a raw JAX array of the given shape.
\item
  \emph{event} (int): The number of batch dimensions to reinterpret as
  event dimensions (used in model building).
\item
  \emph{mask} (jnp.ndarray, bool, optional): Optional boolean array to
  mask observations. Defaults to None.
\item
  \emph{create\_obj} (bool, optional): If True, returns the raw BF
  distribution object instead of creating a sample site. This is
  essential for building complex distributions like
  \texttt{Mi\ \ xtureSameFamily}. Defaults to False.
\item
  \emph{sample} (bool, optional): A control-flow argument. If
  \texttt{True}, the function will directly sample a raw JAX array from
  the distribution, bypassing the BF model context. If \texttt{False},
  it will create a \texttt{BF.sample} site within a model. Defaults to
  \texttt{False}.
\item
  \emph{obs} (jnp.ndarray, optional): The observed value for this random
  variable. If provided, the sample site is conditioned on this value,
  and the function returns the observed value. If \texttt{None}, the
  site is treated as a latent (unobserved) random variable. Defaults to
  \texttt{None}.
\item
  \emph{name} (str, optional): The name of the sample site in a BF
  model. This is used to uniquely identify the random variable. Defaults
  to `x'.
\end{itemize}

\textbf{Returns} BF GaussianCopula distribution object: When
\texttt{sample=False} (for model building). JAX array: When
\texttt{sample=True} (for direct sampling). BF distribution object: When
\texttt{create\_obj=True} (for advanced use cases).

\textbf{Example Usage} from BayesForge import bf import jax.numpy as jnp
m = bf(platform=`cpu') m.dist.gaussian\_copula( marginal\_dist =
m.dist.beta(2.0, 5.0, create\_obj = True), correlation\_matrix =
jnp.array({[}{[}1.0, 0.7{]},{[}0.7, 1.0{]}{]}), sample = True )

\textbf{Wrapper of}
https://num.pyro.ai/en/stable/distributions.html\#gaussiancopula

\subsection{\texorpdfstring{\texttt{m.dist.gaussian\_copula\_beta}}{m.dist.gaussian\_copula\_beta}}\label{gaussian_copula_beta}

\begin{Shaded}
\begin{Highlighting}[]
\NormalTok{m.dist.gaussian\_copula\_beta(}
\NormalTok{    concentration1,}
\NormalTok{    concentration0,}
\NormalTok{    correlation\_matrix}\OperatorTok{=}\VariableTok{None}\NormalTok{,}
\NormalTok{    correlation\_cholesky}\OperatorTok{=}\VariableTok{None}\NormalTok{,}
\NormalTok{    validate\_args}\OperatorTok{=}\VariableTok{False}\NormalTok{,}
\NormalTok{    name}\OperatorTok{=}\StringTok{\textquotesingle{}x\textquotesingle{}}\NormalTok{,}
\NormalTok{    obs}\OperatorTok{=}\VariableTok{None}\NormalTok{,}
\NormalTok{    mask}\OperatorTok{=}\VariableTok{None}\NormalTok{,}
\NormalTok{    infer}\OperatorTok{=}\VariableTok{None}\NormalTok{,}
\NormalTok{    sample}\OperatorTok{=}\VariableTok{False}\NormalTok{,}
\NormalTok{    seed}\OperatorTok{=}\VariableTok{None}\NormalTok{,}
\NormalTok{    shape}\OperatorTok{=}\NormalTok{(),}
\NormalTok{    event}\OperatorTok{=}\DecValTok{0}\NormalTok{,}
\NormalTok{    create\_obj}\OperatorTok{=}\VariableTok{False}\NormalTok{,}
\NormalTok{    to\_jax}\OperatorTok{=}\VariableTok{True}\NormalTok{,}
\NormalTok{)}
\end{Highlighting}
\end{Shaded}

\textbf{Gaussian Copula Beta} This distribution combines a Gaussian
copula with a Beta distribution. The Gaussian copula models the
dependence structure between random variables, while the Beta
distribution defines the marginal distributions of each variable.

\textbf{Args} - \emph{concentration1} (jnp.ndarray): The first shape
parameter of the Beta distribution.

\begin{itemize}
\item
  \emph{concentration0} (jnp.ndarray): The second shape parameter of the
  Beta distribution.
\item
  \emph{correlation\_matrix} (array\_like, optional): Correlation matrix
  of the coupling multivariate normal distribution. Defaults to None.
\item
  \emph{correlation\_cholesky} (jnp.ndarray): The Cholesky decomposition
  of the correlation matrix.

  \begin{itemize}
  \tightlist
  \item
    \emph{shape} (tuple): A multi-purpose argument for shaping. When
    \texttt{sample=False} (model building), this is used with
    \texttt{shape} to set the distribution's batch shape. When
    \texttt{sample=True} (direct sampling), this is used as
    \texttt{sample\_shape} to draw a raw JAX array of the given shape.
  \end{itemize}
\item
  \emph{event} (int): The number of batch dimensions to reinterpret as
  event dimensions (used in model building).
\item
  \emph{mask} (jnp.ndarray, bool): Optional boolean array to mask
  observations.
\item
  \emph{create\_obj} (bool): If True, returns the raw BF distribution
  object instead of creating a sample site. This is essential for
  building complex distributions like \texttt{MixtureSameFamily}.
\item
  \emph{sample} (bool, optional): A control-flow argument. If
  \texttt{True}, the function will directly sample a raw JAX array from
  the distribution, bypassing the BF model context. If \texttt{False},
  it will create a \texttt{BF.sample} site within a model. Defaults to
  \texttt{False}.
\item
  \emph{obs} (jnp.ndarray, optional): The observed value for this random
  variable. If provided, the sample site is conditioned on this value,
  and the function returns the observed value. If \texttt{None}, the
  site is treated as a latent (unobserved) random variable. Defaults to
  \texttt{None}.
\item
  \emph{name} (str, optional): The name of the sample site in a BF
  model. This is used to uniquely identify the random variable. Defaults
  to `x'.
\end{itemize}

\textbf{Returns} GaussianCopulaBeta: A BF GaussianCopulaBeta
distribution object (for model building). jnp.ndarray: A JAX array of
samples drawn from the GaussianCopulaBeta distribution (for direct
sampling). Distribution: The raw BF distribution object (for advanced
use cases).

\textbf{Example Usage} from BayesForge import bf import jax.numpy as jnp
m = bf(platform=`cpu') m.dist.gaussian\_copula\_beta( concentration1 =
jnp.array({[}2.0, 3.0{]}), concentration0 = jnp.array({[}5.0, 3.0{]}),
correlation\_cholesky = jnp.linalg.cholesky(jnp.array({[}{[}1.0,
0.7{]},{[}0.7, 1.0{]}{]})), sample = True )

\textbf{Wrapper of}
https://num.pyro.ai/en/stable/distributions.html\#gaussiancopulabetadistribution

\subsection{\texorpdfstring{\texttt{m.dist.gaussian\_random\_walk}}{m.dist.gaussian\_random\_walk}}\label{gaussian_random_walk}

\begin{Shaded}
\begin{Highlighting}[]
\NormalTok{m.dist.gaussian\_random\_walk(}
\NormalTok{    scale}\OperatorTok{=}\FloatTok{1.0}\NormalTok{,}
\NormalTok{    num\_steps}\OperatorTok{=}\DecValTok{1}\NormalTok{,}
\NormalTok{    validate\_args}\OperatorTok{=}\VariableTok{None}\NormalTok{,}
\NormalTok{    name}\OperatorTok{=}\StringTok{\textquotesingle{}x\textquotesingle{}}\NormalTok{,}
\NormalTok{    obs}\OperatorTok{=}\VariableTok{None}\NormalTok{,}
\NormalTok{    mask}\OperatorTok{=}\VariableTok{None}\NormalTok{,}
\NormalTok{    infer}\OperatorTok{=}\VariableTok{None}\NormalTok{,}
\NormalTok{    sample}\OperatorTok{=}\VariableTok{False}\NormalTok{,}
\NormalTok{    seed}\OperatorTok{=}\VariableTok{None}\NormalTok{,}
\NormalTok{    shape}\OperatorTok{=}\NormalTok{(),}
\NormalTok{    event}\OperatorTok{=}\DecValTok{0}\NormalTok{,}
\NormalTok{    create\_obj}\OperatorTok{=}\VariableTok{False}\NormalTok{,}
\NormalTok{    to\_jax}\OperatorTok{=}\VariableTok{True}\NormalTok{,}
\NormalTok{)}
\end{Highlighting}
\end{Shaded}

\textbf{Gaussian Random Walk} Creates a distribution over a Gaussian
random walk of a specified number of steps. This is a discrete-time
stochastic process where the value at each step is the previous value
plus a Gaussian-distributed increment. The distribution is over the
entire path.

\textbf{Args} - \emph{scale} (float, jnp.ndarray, optional): The
standard deviation (\(\sigma\)) of the Gaussian increments. Must be
positive. Defaults to 1.0.

\begin{itemize}
\item
  *num\_steps (int, optional): The number of steps in the random walk,
  which determines the event shape of the distribution. Must be
  positive. Defaults to 1.
\item
  *validate\_args (bool, optional): Whether to enable validation of
  distribution parameters. Defaults to \texttt{None}.
\item
  \emph{name} (str, optional): The name of the sample site in a BF
  model. This is used to uniquely identify the random variable. Defaults
  to `x'.
\item
  \emph{obs} (jnp.ndarray, optional): The observed value for this random
  variable. If provided, the sample site is conditioned on this value,
  and the function returns the observed value. If \texttt{None}, the
  site is treated as a latent (unobserved) random variable. Defaults to
  \texttt{None}.
\item
  \emph{mask} (jnp.ndarray, bool, optional): Optional boolean array to
  mask observations. If provided, events with a \texttt{True} mask will
  be conditioned on \texttt{obs}, while the remaining events will be
  treated as latent variables. Defaults to \texttt{None}.
\item
  \emph{sample} (bool, optional): A control-flow argument. If
  \texttt{True}, the function will directly sample a raw JAX array from
  the distribution, bypassing the BF model context. If \texttt{False},
  it will create a \texttt{BF.sample} site within a model. Defaults to
  \texttt{False}.
\item
  \emph{seed} (int, optional): An integer seed for this draw when
  \texttt{sample=True}. By default (\texttt{None}) the draw comes from
  the session's random stream, set by \texttt{bf(rand\_seed=...)}, so
  consecutive draws differ but the whole sequence is reproducible. No
  effect when \texttt{sample=False}.
\item
  \emph{shape} (tuple, optional): A multi-purpose argument for shaping.
  When \texttt{sample=False} (model building), this is used with
  \texttt{.expand(shape)} to set the distribution's batch shape. When
  \texttt{sample=True} (direct sampling), this is used as
  \texttt{sample\_shape} to draw a raw JAX array of the given shape.
\item
  \emph{event} (int, optional): The number of batch dimensions to
  reinterpret as event dimensions (used in model building).
\item
  \emph{create\_obj} (bool, optional): If True, returns the raw BF
  distribution object instead of creating a sample site. This is
  essential for building complex distributions like
  \texttt{MixtureSameFamily}. Defaults to \texttt{False}.
\end{itemize}

\textbf{Returns} BF.primitives.Messenger: A BF sample site object when
used in a model context (\texttt{sample=False}). jnp.ndarray: A JAX
array of samples drawn from the GaussianRandomWalk distribution (for
direct sampling, \texttt{sample=True}).
numpyro.distributions.Distribution: The raw BF distribution object (if
\texttt{create\_obj=True}).

\textbf{Example Usage} from BayesForge import bf import jax.numpy as jnp
m = bf(platform=`cpu')

\begin{verbatim}
**Direct sampling of a random walk with 100 steps**
path = m.dist.gaussian_random_walk(scale=0.5, num_steps=100, sample=True)

**Usage within a model for a latent time series**
def my_model(data=None):
    **Prior on the volatility of the random walk**
    volatility = m.dist.half_cauchy(scale=1.0, name='volatility')
    
    **The latent random walk**
    latent_process = m.dist.gaussian_random_walk(
        scale=volatility,
        num_steps=len(data) if data is not None else 10,
        name='latent_process'
    )
    
    **Observation model**
    **Assumes the observed data is the latent process plus some noise**
    obs_noise = m.dist.half_cauchy(scale=1.0, name='obs_noise')
    with m.plate('time', len(data) if data is not None else 10):
        return m.dist.normal(loc=latent_process, scale=obs_noise, obs=data, name='obs')
\end{verbatim}

\textbf{Wrapper of}
https://num.pyro.ai/en/stable/distributions.html\#gaussianrandomwalk

\subsection{\texorpdfstring{\texttt{m.dist.gaussian\_state\_space}}{m.dist.gaussian\_state\_space}}\label{gaussian_state_space}

\begin{Shaded}
\begin{Highlighting}[]
\NormalTok{m.dist.gaussian\_state\_space(}
\NormalTok{    num\_steps,}
\NormalTok{    transition\_matrix,}
\NormalTok{    covariance\_matrix}\OperatorTok{=}\VariableTok{None}\NormalTok{,}
\NormalTok{    precision\_matrix}\OperatorTok{=}\VariableTok{None}\NormalTok{,}
\NormalTok{    scale\_tril}\OperatorTok{=}\VariableTok{None}\NormalTok{,}
\NormalTok{    validate\_args}\OperatorTok{=}\VariableTok{None}\NormalTok{,}
\NormalTok{    name}\OperatorTok{=}\StringTok{\textquotesingle{}x\textquotesingle{}}\NormalTok{,}
\NormalTok{    obs}\OperatorTok{=}\VariableTok{None}\NormalTok{,}
\NormalTok{    mask}\OperatorTok{=}\VariableTok{None}\NormalTok{,}
\NormalTok{    infer}\OperatorTok{=}\VariableTok{None}\NormalTok{,}
\NormalTok{    sample}\OperatorTok{=}\VariableTok{False}\NormalTok{,}
\NormalTok{    seed}\OperatorTok{=}\VariableTok{None}\NormalTok{,}
\NormalTok{    shape}\OperatorTok{=}\NormalTok{(),}
\NormalTok{    event}\OperatorTok{=}\DecValTok{0}\NormalTok{,}
\NormalTok{    create\_obj}\OperatorTok{=}\VariableTok{False}\NormalTok{,}
\NormalTok{    to\_jax}\OperatorTok{=}\VariableTok{True}\NormalTok{,}
\NormalTok{)}
\end{Highlighting}
\end{Shaded}

\textbf{Gaussian State Space} Samples from a Gaussian state space model.

\textbf{Args} - \emph{num\_steps} (int): Number of steps.

\begin{itemize}
\item
  \emph{transition\_matrix} (jnp.ndarray): State space transition matrix
  \(\mathbf{A}\).
\item
  \emph{covariance\_matrix} (jnp.ndarray, optional): Covariance of the
  innovation noise \({\epsilon}\). Defaults to None.
\item
  *precision\_matrix (jnp.ndarray, optional): Precision matrix of the
  innovation noise \(\epsilon\). Defaults to None.
\item
  *scale\_tril (jnp.ndarray, optional): Scale matrix of the innovation
  noise \(\epsilon\). Defaults to None.
\item
  \emph{shape} (tuple): A multi-purpose argument for shaping. When
  \texttt{sample=False} (model building), this is used with
  \texttt{.expand(shape)} to set the distribution's batch shape. When
  \texttt{sample=True} (direct sampling), this is used as
  \texttt{sample\_shape} to draw a raw JAX array of the given shape.
\item
  \emph{event} (int): The number of batch dimensions to reinterpret as
  event dimensions (used in model building).
\item
  \emph{mask} (jnp.ndarray, bool, optional): Optional boolean array to
  mask observations. Defaults to None.
\item
  \emph{create\_obj} (bool, optional): If True, returns the raw BF
  distribution object instead of creating a sample site. Defaults to
  False.
\item
  \emph{sample} (bool, optional): A control-flow argument. If
  \texttt{True}, the function will directly sample a raw JAX array from
  the distribution, bypassing the BF model context. If \texttt{False},
  it will create a \texttt{BF.sample} site within a model. Defaults to
  \texttt{False}.
\item
  \emph{obs} (jnp.ndarray, optional): The observed value for this random
  variable. If provided, the sample site is conditioned on this value,
  and the function returns the observed value. If \texttt{None}, the
  site is treated as a latent (unobserved) random variable. Defaults to
  \texttt{None}.
\item
  \emph{name} (str, optional): The name of the sample site in a BF
  model. This is used to uniquely identify the random variable. Defaults
  to `x'.
\end{itemize}

\textbf{Returns} - When \texttt{sample=False}: A BF GaussianStateSpace
distribution object (for model building). - When \texttt{sample=True}: A
JAX array of samples drawn from the GaussianStateSpace distribution (for
direct sampling). - When \texttt{create\_obj=True}: The raw BF
distribution object (for advanced use cases).

\textbf{Example Usage} from BayesForge import bf import jax.numpy as jnp
m = bf(platform=`cpu') m.dist.gaussian\_state\_space(num\_steps=5,
transition\_matrix=jnp.array({[}{[}0.5, 0.1{]}, {[}0.0, 0.8{]}{]}),
covariance\_matrix=jnp.array({[}{[}1.0, 0.6{]}, {[}0.6, 1.0{]}{]}),
sample=True)

\textbf{Wrapper of}
https://num.pyro.ai/en/stable/distributions.html\#gaussianstate

\subsection{\texorpdfstring{\texttt{m.dist.geometric}}{m.dist.geometric}}\label{geometric}

\begin{Shaded}
\begin{Highlighting}[]
\NormalTok{m.dist.geometric(}
\NormalTok{    probs}\OperatorTok{=}\VariableTok{None}\NormalTok{,}
\NormalTok{    logits}\OperatorTok{=}\VariableTok{None}\NormalTok{,}
\NormalTok{    validate\_args}\OperatorTok{=}\VariableTok{None}\NormalTok{,}
\NormalTok{    name}\OperatorTok{=}\StringTok{\textquotesingle{}x\textquotesingle{}}\NormalTok{,}
\NormalTok{    obs}\OperatorTok{=}\VariableTok{None}\NormalTok{,}
\NormalTok{    mask}\OperatorTok{=}\VariableTok{None}\NormalTok{,}
\NormalTok{    infer}\OperatorTok{=}\VariableTok{None}\NormalTok{,}
\NormalTok{    sample}\OperatorTok{=}\VariableTok{False}\NormalTok{,}
\NormalTok{    seed}\OperatorTok{=}\VariableTok{None}\NormalTok{,}
\NormalTok{    shape}\OperatorTok{=}\NormalTok{(),}
\NormalTok{    event}\OperatorTok{=}\DecValTok{0}\NormalTok{,}
\NormalTok{    create\_obj}\OperatorTok{=}\VariableTok{False}\NormalTok{,}
\NormalTok{    to\_jax}\OperatorTok{=}\VariableTok{True}\NormalTok{,}
\NormalTok{)}
\end{Highlighting}
\end{Shaded}

\textbf{Geometric distribution.} The Geometric distribution models the
number of independent Bernoulli trials needed to obtain the first
success (or equivalently the number of failures before the first
success). Each trial has the same probability of success \(p\), and
trials are independent.

\begin{itemize}
\tightlist
\item
  \(X\) = number of trials until the first success (so support
  \(1,2,3,\dots\)).
\end{itemize}

\[
P(X = k) = (1 - p)^{,k-1} ; p,\quad k = 1,2,\dots
\]

\textbf{Args} - \emph{probs} (jnp.ndarray, optional): Probability of
success on each trial. Must be between 0 and 1.

\begin{itemize}
\item
  \emph{logits} (jnp.ndarray, optional): Log-odds of success on each
  trial. \texttt{probs\ =\ jax.nn.sigmoid(logits)}.
\item
  \emph{shape} (tuple): A multi-purpose argument for shaping. When
  \texttt{sample=False} (model building), this is used with
  \texttt{shape} to set the distribution's batch shape. When
  \texttt{sample=True} (direct sampling), this is used as
  \texttt{sample\_shape} to draw a raw JAX array of the given shape.
\item
  \emph{event} (int): The number of batch dimensions to reinterpret as
  event dimensions (used in model building).
\item
  \emph{mask} (jnp.ndarray, bool, optional): Optional boolean array to
  mask observations.
\item
  \emph{create\_obj} (bool, optional): If True, returns the raw BF
  distribution object instead of creating a sample site. This is
  essential for building complex distributions like
  \texttt{MixtureSameFamily}.
\item
  \emph{sample} (bool, optional): A control-flow argument. If
  \texttt{True}, the function will directly sample a raw JAX array from
  the distribution, bypassing the BF model context. If \texttt{False},
  it will create a \texttt{BF.sample} site within a model. Defaults to
  \texttt{False}.
\item
  \emph{seed} (int, optional): An integer seed for this draw when
  \texttt{sample=True}. By default (\texttt{None}) the draw comes from
  the session's random stream, set by \texttt{bf(rand\_seed=...)}, so
  consecutive draws differ but the whole sequence is reproducible. No
  effect when \texttt{sample=False}.
\item
  \emph{obs} (jnp.ndarray, optional): The observed value for this random
  variable. If provided, the sample site is conditioned on this value,
  and the function returns the observed value. If \texttt{None}, the
  site is treated as a latent (unobserved) random variable. Defaults to
  \texttt{None}.
\item
  \emph{name} (str, optional): The name of the sample site in a BF
  model. This is used to uniquely identify the random variable. Defaults
  to `x'.
\end{itemize}

\textbf{Returns} - When \texttt{sample=False}: A BF Geometric
distribution object (for model building). - When \texttt{sample=True}: A
JAX array of samples drawn from the Geometric distribution (for direct
sampling). - When \texttt{create\_obj=True}: The raw BF distribution
object (for advanced use cases).

\textbf{Example Usage} from BayesForge import bf m = bf(platform=`cpu')
m.dist.geometric(probs=0.5, sample=True)

\textbf{Wrapper of:
https://num.pyro.ai/en/stable/distributions.html\#geometric}
\textbf{References} *
https://www.britannica.com/topic/geometric-distribution ``Geometric
distribution \textbar{} Definition, Formula, Examples, Illustration, \&
Applications \textbar{} Britannica'' *
https://en.wikipedia.org/wiki/Geometric\_distribution ``Geometric
distribution'' *
https://distribution-explorer.github.io/discrete/geometric.html
``Geometric distribution --- Probability Distribution Explorer
documentation''

\subsection{\texorpdfstring{\texttt{m.dist.geometric\_logits}}{m.dist.geometric\_logits}}\label{geometric_logits}

\begin{Shaded}
\begin{Highlighting}[]
\NormalTok{m.dist.geometric\_logits(}
\NormalTok{    logits,}
\NormalTok{    validate\_args}\OperatorTok{=}\VariableTok{None}\NormalTok{,}
\NormalTok{    name}\OperatorTok{=}\StringTok{\textquotesingle{}x\textquotesingle{}}\NormalTok{,}
\NormalTok{    obs}\OperatorTok{=}\VariableTok{None}\NormalTok{,}
\NormalTok{    mask}\OperatorTok{=}\VariableTok{None}\NormalTok{,}
\NormalTok{    infer}\OperatorTok{=}\VariableTok{None}\NormalTok{,}
\NormalTok{    sample}\OperatorTok{=}\VariableTok{False}\NormalTok{,}
\NormalTok{    seed}\OperatorTok{=}\VariableTok{None}\NormalTok{,}
\NormalTok{    shape}\OperatorTok{=}\NormalTok{(),}
\NormalTok{    event}\OperatorTok{=}\DecValTok{0}\NormalTok{,}
\NormalTok{    create\_obj}\OperatorTok{=}\VariableTok{False}\NormalTok{,}
\NormalTok{    to\_jax}\OperatorTok{=}\VariableTok{True}\NormalTok{,}
\NormalTok{)}
\end{Highlighting}
\end{Shaded}

\textbf{GeometricLogits} Samples from a GeometricLogits distribution,
which models the number of failures before the first success in a
sequence of independent Bernoulli trials. It is parameterized by logits,
which are transformed into probabilities using the sigmoid function.

\[
P(X = k) = (1 - p)^k , p,
\]

where:

\begin{itemize}
\tightlist
\item
  X is the number of failures before the first success.
\item
  k is the number of failures.
\item
  p is the probability of success on each trial (derived from the
  logits).
\end{itemize}

\textbf{Args} - \emph{logits} (jnp.ndarray): Log-odds parameterization
of the probability of success.

\begin{itemize}
\item
  \emph{shape} (tuple): A multi-purpose argument for shaping. When
  \texttt{sample=False} (model building), this is used with
  \texttt{shape} to set the distribution's batch shape. When
  \texttt{sample=True} (direct sampling), this is used as
  \texttt{sample\_shape} to draw a raw JAX array of the given shape.
\item
  \emph{event} (int): The number of batch dimensions to reinterpret as
  event dimensions (used in model building).
\item
  \emph{mask} (jnp.ndarray, bool): Optional boolean array to mask
  observations.
\item
  \emph{create\_obj} (bool): If True, returns the raw BF distribution
  object instead of creating a sample site. This is essential for
  building complex distributions like \texttt{MixtureSameFamily}.
\item
  \emph{sample} (bool, optional): A control-flow argument. If
  \texttt{True}, the function will directly sample a raw JAX array from
  the distribution, bypassing the BF model context. If \texttt{False},
  it will create a \texttt{BF.sample} site within a model. Defaults to
  \texttt{False}.
\item
  \emph{obs} (jnp.ndarray, optional): The observed value for this random
  variable. If provided, the sample site is conditioned on this value,
  and the function returns the observed value. If \texttt{None}, the
  site is treated as a latent (unobserved) random variable. Defaults to
  \texttt{None}.
\item
  \emph{name} (str, optional): The name of the sample site in a BF
  model. This is used to uniquely identify the random variable. Defaults
  to `x'.
\end{itemize}

\textbf{Returns} BF GeometricLogits distribution object (for model
building) when \texttt{sample=False}. JAX array of samples drawn from
the GeometricLogits distribution (for direct sampling) when
\texttt{sample=True}. The raw BF distribution object (for advanced use
cases) when \texttt{create\_obj=True}.

\textbf{Example Usage} from BayesForge import bf import jax.numpy as jnp
m = bf(platform=`cpu') m.dist.geometric\_logits(logits=jnp.zeros(10),
sample=True)

\textbf{Wrapper of}
https://num.pyro.ai/en/stable/distributions.html\#geometriclogits

\textbf{References}
https://www.tensorflow.org/probability/api\_docs/python/tfp/distributions/Geometric

\subsection{\texorpdfstring{\texttt{m.dist.geometric\_probs}}{m.dist.geometric\_probs}}\label{geometric_probs}

\begin{Shaded}
\begin{Highlighting}[]
\NormalTok{m.dist.geometric\_probs(}
\NormalTok{    probs,}
\NormalTok{    validate\_args}\OperatorTok{=}\VariableTok{None}\NormalTok{,}
\NormalTok{    name}\OperatorTok{=}\StringTok{\textquotesingle{}x\textquotesingle{}}\NormalTok{,}
\NormalTok{    obs}\OperatorTok{=}\VariableTok{None}\NormalTok{,}
\NormalTok{    mask}\OperatorTok{=}\VariableTok{None}\NormalTok{,}
\NormalTok{    infer}\OperatorTok{=}\VariableTok{None}\NormalTok{,}
\NormalTok{    sample}\OperatorTok{=}\VariableTok{False}\NormalTok{,}
\NormalTok{    seed}\OperatorTok{=}\VariableTok{None}\NormalTok{,}
\NormalTok{    shape}\OperatorTok{=}\NormalTok{(),}
\NormalTok{    event}\OperatorTok{=}\DecValTok{0}\NormalTok{,}
\NormalTok{    create\_obj}\OperatorTok{=}\VariableTok{False}\NormalTok{,}
\NormalTok{    to\_jax}\OperatorTok{=}\VariableTok{True}\NormalTok{,}
\NormalTok{)}
\end{Highlighting}
\end{Shaded}

\textbf{GeometricProbs} Samples from a Geometric

The Geometric distribution models the number of trials until the first
success in a sequence of independent Bernoulli trials, where each trial
has the same probability of success.

\[
   P(X = k) = (1 - p)^k p
\]

\textbf{Args} - \emph{probs} (jnp.ndarray): Probability of success on
each trial. Must be between 0 and 1.

\begin{itemize}
\item
  \emph{shape} (tuple): A multi-purpose argument for shaping. When
  \texttt{sample=False} (model building), this is used with
  \texttt{shape} to set the distribution's batch shape. When
  \texttt{sample=True} (direct sampling), this is used as
  \texttt{sample\_shape} to draw a raw JAX array of the given shape.
\item
  \emph{event} (int): The number of batch dimensions to reinterpret as
  event dimensions (used in model building).
\item
  \emph{mask} (jnp.ndarray, bool): Optional boolean array to mask
  observations.
\item
  \emph{create\_obj} (bool): If True, returns the raw BF distribution
  object instead of creating a sample site. This is essential for
  building complex distributions like \texttt{MixtureSameFamily}.
\item
  \emph{sample} (bool, optional): A control-flow argument. If
  \texttt{True}, the function will directly sample a raw JAX array from
  the distribution, bypassing the BF model context. If \texttt{False},
  it will create a \texttt{BF.sample} site within a model. Defaults to
  \texttt{False}.
\item
  \emph{obs} (jnp.ndarray, optional): The observed value for this random
  variable. If provided, the sample site is conditioned on this value,
  and the function returns the observed value. If \texttt{None}, the
  site is treated as a latent (unobserved) random variable. Defaults to
  \texttt{None}.
\item
  \emph{name} (str, optional): The name of the sample site in a BF
  model. This is used to uniquely identify the random variable. Defaults
  to `x'.
\end{itemize}

\textbf{Returns} BF GeometricProbs distribution object (for model
building). JAX array of samples drawn from the GeometricProbs
distribution (for direct sampling). The raw BF distribution object (for
advanced use cases).

\textbf{Example Usage} from BayesForge import bf m = bf(platform=`cpu')
m.dist.geometric\_probs(probs=0.5, sample=True)

\textbf{Wrapper of}
https://num.pyro.ai/en/stable/distributions.html\#geometricprobs

\subsection{\texorpdfstring{\texttt{m.dist.gompertz}}{m.dist.gompertz}}\label{gompertz}

\begin{Shaded}
\begin{Highlighting}[]
\NormalTok{m.dist.gompertz(}
\NormalTok{    concentration,}
\NormalTok{    rate}\OperatorTok{=}\FloatTok{1.0}\NormalTok{,}
\NormalTok{    validate\_args}\OperatorTok{=}\VariableTok{None}\NormalTok{,}
\NormalTok{    name}\OperatorTok{=}\StringTok{\textquotesingle{}x\textquotesingle{}}\NormalTok{,}
\NormalTok{    obs}\OperatorTok{=}\VariableTok{None}\NormalTok{,}
\NormalTok{    mask}\OperatorTok{=}\VariableTok{None}\NormalTok{,}
\NormalTok{    infer}\OperatorTok{=}\VariableTok{None}\NormalTok{,}
\NormalTok{    sample}\OperatorTok{=}\VariableTok{False}\NormalTok{,}
\NormalTok{    seed}\OperatorTok{=}\VariableTok{None}\NormalTok{,}
\NormalTok{    shape}\OperatorTok{=}\NormalTok{(),}
\NormalTok{    event}\OperatorTok{=}\DecValTok{0}\NormalTok{,}
\NormalTok{    create\_obj}\OperatorTok{=}\VariableTok{False}\NormalTok{,}
\NormalTok{    to\_jax}\OperatorTok{=}\VariableTok{True}\NormalTok{,}
\NormalTok{)}
\end{Highlighting}
\end{Shaded}

\textbf{Gompertz} The Gompertz distribution is a distribution with
support on the positive real line that is closely related to the Gumbel
distribution. This implementation follows the notation used in the
Wikipedia entry for the Gompertz distribution. See
https://en.wikipedia.org/wiki/Gompertz\_distribution.

The probability density function (PDF) is:

\[
F(x) = 1 - \exp!\bigl[-a (e^x - 1)\bigr], \quad x \ge 0.
\]

\textbf{Args} - \emph{concentration} (jnp.ndarray): The concentration
parameter. Must be positive.

\begin{itemize}
\item
  \emph{rate} (jnp.ndarray): The rate parameter. Must be positive.
\item
  \emph{shape} (tuple): A multi-purpose argument for shaping. When
  \texttt{sample=False} (model building), this is used with
  \texttt{shape} to set the distribution's batch shape. When
  \texttt{sample=True} (direct sampling), this is used as
  \texttt{sample\_shape} to draw a raw JAX array of the given shape.
\item
  \emph{event} (int): The number of batch dimensions to reinterpret as
  event dimensions (used in model building).
\item
  \emph{mask} (jnp.ndarray, bool): Optional boolean array to mask
  observations.
\item
  \emph{create\_obj} (bool): If True, returns the raw BF distribution
  object instead of creating a sample site. This is essential for
  building complex distributions like \texttt{MixtureSameFamily}.
\item
  \emph{sample} (bool, optional): A control-flow argument. If
  \texttt{True}, the function will directly sample a raw JAX array from
  the distribution, bypassing the BF model context. If \texttt{False},
  it will create a \texttt{BF.sample} site within a model. Defaults to
  \texttt{False}.
\item
  \emph{obs} (jnp.ndarray, optional): The observed value for this random
  variable. If provided, the sample site is conditioned on this value,
  and the function returns the observed value. If \texttt{None}, the
  site is treated as a latent (unobserved) random variable. Defaults to
  \texttt{None}.
\item
  \emph{name} (str, optional): The name of the sample site in a BF
  model. This is used to uniquely identify the random variable. Defaults
  to `x'.
\end{itemize}

\textbf{Returns} BF Gompertz distribution object: When
\texttt{sample=False} (for model building). JAX array: When
\texttt{sample=True} (for direct sampling). BF distribution object: When
\texttt{create\_obj=True} (for advanced use cases).

\textbf{Example Usage} from BayesForge import bf m = bf(platform=`cpu')
m.dist.gompertz(concentration=1.0, rate=1.0, sample=True)

\textbf{Wrapper of}
https://num.pyro.ai/en/stable/distributions.html\#gompertz
\textbf{References} https://en.wikipedia.org/wiki/Gompertz\_distribution

\subsection{\texorpdfstring{\texttt{m.dist.gumbel}}{m.dist.gumbel}}\label{gumbel}

\begin{Shaded}
\begin{Highlighting}[]
\NormalTok{m.dist.gumbel(}
\NormalTok{    loc}\OperatorTok{=}\FloatTok{0.0}\NormalTok{,}
\NormalTok{    scale}\OperatorTok{=}\FloatTok{1.0}\NormalTok{,}
\NormalTok{    validate\_args}\OperatorTok{=}\VariableTok{None}\NormalTok{,}
\NormalTok{    name}\OperatorTok{=}\StringTok{\textquotesingle{}x\textquotesingle{}}\NormalTok{,}
\NormalTok{    obs}\OperatorTok{=}\VariableTok{None}\NormalTok{,}
\NormalTok{    mask}\OperatorTok{=}\VariableTok{None}\NormalTok{,}
\NormalTok{    infer}\OperatorTok{=}\VariableTok{None}\NormalTok{,}
\NormalTok{    sample}\OperatorTok{=}\VariableTok{False}\NormalTok{,}
\NormalTok{    seed}\OperatorTok{=}\VariableTok{None}\NormalTok{,}
\NormalTok{    shape}\OperatorTok{=}\NormalTok{(),}
\NormalTok{    event}\OperatorTok{=}\DecValTok{0}\NormalTok{,}
\NormalTok{    create\_obj}\OperatorTok{=}\VariableTok{False}\NormalTok{,}
\NormalTok{    to\_jax}\OperatorTok{=}\VariableTok{True}\NormalTok{,}
\NormalTok{)}
\end{Highlighting}
\end{Shaded}

\textbf{Gumbel} Samples from a Gumbel (or Extreme Value) distribution.

The Gumbel distribution is a continuous probability distribution named
after German mathematician Carl Gumbel. It is often used to model the
distribution of maximum values in a sequence of independent random
variables.

\[
  F(x;\mu,\beta) = \exp\bigl( -\exp\bigl( -\tfrac{x - \mu}{\beta} \bigr) \bigr), \quad x \in \mathbb{R}.
  \]

\textbf{Args} - \emph{loc} (jnp.ndarray or float, optional): Location
parameter. Defaults to 0.0.

\begin{itemize}
\item
  \emph{scale} (jnp.ndarray or float, optional): Scale parameter. Must
  be positive. Defaults to 1.0.
\item
  \emph{shape} (tuple): A multi-purpose argument for shaping. When
  \texttt{sample=False} (model building), this is used with
  \texttt{shape} to set the distribution's batch shape. When
  \texttt{sample=True} (direct sampling), this is used as
  \texttt{sample\_shape} to draw a raw JAX array of the given shape.
\item
  \emph{event} (int, optional): The number of batch dimensions to
  reinterpret as event dimensions (used in model building). Defaults to
  None.
\item
  \emph{mask} (jnp.ndarray, bool, optional): Optional boolean array to
  mask observations. Defaults to None.
\item
  \emph{create\_obj} (bool, optional): If True, returns the raw BF
  distribution object instead of creating a sample site. This is
  essential for building complex distributions like
  \texttt{MixtureSameFamily}. Defaults to False.
\item
  \emph{sample} (bool, optional): A control-flow argument. If
  \texttt{True}, the function will directly sample a raw JAX array from
  the distribution, bypassing the BF model context. If \texttt{False},
  it will create a \texttt{BF.sample} site within a model. Defaults to
  \texttt{False}.
\item
  \emph{obs} (jnp.ndarray, optional): The observed value for this random
  variable. If provided, the sample site is conditioned on this value,
  and the function returns the observed value. If \texttt{None}, the
  site is treated as a latent (unobserved) random variable. Defaults to
  \texttt{None}.
\item
  \emph{name} (str, optional): The name of the sample site in a BF
  model. This is used to uniquely identify the random variable. Defaults
  to `x'.
\end{itemize}

\textbf{Returns} - When \texttt{sample=False}: A BF Gumbel distribution
object (for model building). - When \texttt{sample=True}: A JAX array of
samples drawn from the Gumbel distribution (for direct sampling). - When
\texttt{create\_obj=True}: The raw BF distribution object (for advanced
use cases).

\textbf{Example Usage} from BayesForge import bf m = bf(platform=`cpu')
m.dist.gumbel(loc=0.0, scale=1.0, sample=True)

\textbf{Wrapper of}
https://num.pyro.ai/en/stable/distributions.html\#gumbel
\textbf{References} https://en.wikipedia.org/wiki/Gumbel\_distribution

\subsection{\texorpdfstring{\texttt{m.dist.half\_cauchy}}{m.dist.half\_cauchy}}\label{half_cauchy}

\begin{Shaded}
\begin{Highlighting}[]
\NormalTok{m.dist.half\_cauchy(}
\NormalTok{    scale}\OperatorTok{=}\FloatTok{1.0}\NormalTok{,}
\NormalTok{    validate\_args}\OperatorTok{=}\VariableTok{None}\NormalTok{,}
\NormalTok{    name}\OperatorTok{=}\StringTok{\textquotesingle{}x\textquotesingle{}}\NormalTok{,}
\NormalTok{    obs}\OperatorTok{=}\VariableTok{None}\NormalTok{,}
\NormalTok{    mask}\OperatorTok{=}\VariableTok{None}\NormalTok{,}
\NormalTok{    infer}\OperatorTok{=}\VariableTok{None}\NormalTok{,}
\NormalTok{    sample}\OperatorTok{=}\VariableTok{False}\NormalTok{,}
\NormalTok{    seed}\OperatorTok{=}\VariableTok{None}\NormalTok{,}
\NormalTok{    shape}\OperatorTok{=}\NormalTok{(),}
\NormalTok{    event}\OperatorTok{=}\DecValTok{0}\NormalTok{,}
\NormalTok{    create\_obj}\OperatorTok{=}\VariableTok{False}\NormalTok{,}
\NormalTok{    to\_jax}\OperatorTok{=}\VariableTok{True}\NormalTok{,}
\NormalTok{)}
\end{Highlighting}
\end{Shaded}

\textbf{HalfCauchy} The HalfCauchy distribution is a probability
distribution that is half of the Cauchy distribution. It is defined on
the positive real numbers and is often used in situations where only
positive values are relevant.

\[   
  f(x) = \frac{1}{2} \cdot \frac{1}{\pi \cdot \frac{1}{scale} \cdot (x^2 + \frac{1}{scale^2})}
  \]

\textbf{Args} - \emph{scale} (jnp.ndarray): The scale parameter of the
Cauchy distribution. Must be positive.

\begin{itemize}
\item
  \emph{shape} (tuple): A multi-purpose argument for shaping. When
  \texttt{sample=False} (model building), this is used with
  \texttt{shape} to set the distribution's batch shape. When
  \texttt{sample=True} (direct sampling), this is used as
  \texttt{sample\_shape} to draw a raw JAX array of the given shape.
\item
  \emph{event} (int): The number of batch dimensions to reinterpret as
  event dimensions (used in model building).
\item
  \emph{mask} (jnp.ndarray, bool): Optional boolean array to mask
  observations.
\item
  \emph{create\_obj} (bool): If True, returns the raw BF distribution
  object instead of creating a sample site. This is essential for
  building complex distributions like \texttt{MixtureSameFamily}.
\item
  \emph{sample} (bool, optional): A control-flow argument. If
  \texttt{True}, the function will directly sample a raw JAX array from
  the distribution, bypassing the BF model context. If \texttt{False},
  it will create a \texttt{BF.sample} site within a model. Defaults to
  \texttt{False}.
\item
  \emph{obs} (jnp.ndarray, optional): The observed value for this random
  variable. If provided, the sample site is conditioned on this value,
  and the function returns the observed value. If \texttt{None}, the
  site is treated as a latent (unobserved) random variable. Defaults to
  \texttt{None}.
\item
  \emph{name} (str, optional): The name of the sample site in a BF
  model. This is used to uniquely identify the random variable. Defaults
  to `x'.
\end{itemize}

\textbf{Returns} - When \texttt{sample=False}: A BF HalfCauchy
distribution object (for model building). - When \texttt{sample=True}: A
JAX array of samples drawn from the HalfCauchy distribution (for direct
sampling). - When \texttt{create\_obj=True}: The raw BF distribution
object (for advanced use cases).

\textbf{Example Usage} from BayesForge import bf m = bf(platform=`cpu')
m.dist.half\_cauchy(scale=1.0, sample=True)

\textbf{Wrapper of:
https://num.pyro.ai/en/stable/distributions.html\#halfcauchy}

\subsection{\texorpdfstring{\texttt{m.dist.half\_normal}}{m.dist.half\_normal}}\label{half_normal}

\begin{Shaded}
\begin{Highlighting}[]
\NormalTok{m.dist.half\_normal(}
\NormalTok{    scale}\OperatorTok{=}\FloatTok{1.0}\NormalTok{,}
\NormalTok{    validate\_args}\OperatorTok{=}\VariableTok{None}\NormalTok{,}
\NormalTok{    name}\OperatorTok{=}\StringTok{\textquotesingle{}x\textquotesingle{}}\NormalTok{,}
\NormalTok{    obs}\OperatorTok{=}\VariableTok{None}\NormalTok{,}
\NormalTok{    mask}\OperatorTok{=}\VariableTok{None}\NormalTok{,}
\NormalTok{    infer}\OperatorTok{=}\VariableTok{None}\NormalTok{,}
\NormalTok{    sample}\OperatorTok{=}\VariableTok{False}\NormalTok{,}
\NormalTok{    seed}\OperatorTok{=}\VariableTok{None}\NormalTok{,}
\NormalTok{    shape}\OperatorTok{=}\NormalTok{(),}
\NormalTok{    event}\OperatorTok{=}\DecValTok{0}\NormalTok{,}
\NormalTok{    create\_obj}\OperatorTok{=}\VariableTok{False}\NormalTok{,}
\NormalTok{    to\_jax}\OperatorTok{=}\VariableTok{True}\NormalTok{,}
\NormalTok{)}
\end{Highlighting}
\end{Shaded}

\textbf{HalfNormal} Samples from a HalfNormal distribution.

The HalfNormal distribution is a distribution of the absolute value of a
normal random variable. It is defined by a location parameter
(implicitly 0) and a scale parameter.

\[
  f(x; \sigma) = \frac{2}{\pi,\sigma} ; \frac{1}{1 + \bigl(x/\sigma\bigr)^2}, \qquad x \ge 0.
  \]

\textbf{Args} - \emph{scale} (float, array): The scale parameter of the
distribution. Must be positive.

\begin{itemize}
\item
  \emph{shape} (tuple): A multi-purpose argument for shaping. When
  \texttt{sample=False} (model building), this is used with
  \texttt{shape} to set the distribution's batch shape. When
  \texttt{sample=True} (direct sampling), this is used as
  \texttt{sample\_shape} to draw a raw JAX array of the given shape.
\item
  \emph{event} (int): The number of batch dimensions to reinterpret as
  event dimensions (used in model building).
\item
  \emph{mask} (jnp.ndarray, bool): Optional boolean array to mask
  observations.
\item
  \emph{create\_obj} (bool): If True, returns the raw BF distribution
  object instead of creating a sample site. This is essential for
  building complex distributions like \texttt{MixtureSameFamily}.
\item
  \emph{sample} (bool, optional): A control-flow argument. If
  \texttt{True}, the function will directly sample a raw JAX array from
  the distribution, bypassing the BF model context. If \texttt{False},
  it will create a \texttt{BF.sample} site within a model. Defaults to
  \texttt{False}.
\item
  \emph{obs} (jnp.ndarray, optional): The observed value for this random
  variable. If provided, the sample site is conditioned on this value,
  and the function returns the observed value. If \texttt{None}, the
  site is treated as a latent (unobserved) random variable. Defaults to
  \texttt{None}.
\item
  \emph{name} (str, optional): The name of the sample site in a BF
  model. This is used to uniquely identify the random variable. Defaults
  to `x'.
\end{itemize}

\textbf{Returns} - When \texttt{sample=False}: A BF HalfNormal
distribution object (for model building). - When \texttt{sample=True}: A
JAX array of samples drawn from the HalfNormal distribution (for direct
sampling). - When \texttt{create\_obj=True}: The raw BF distribution
object (for advanced use cases).

\textbf{Example Usage} from BayesForge import bf m = bf(platform=`cpu')
m.dist.half\_normal(scale=1.0, sample=True)

\textbf{Wrapper of}
https://num.pyro.ai/en/stable/distributions.html\#halfnormal

\textbf{References}
https://distribution-explorer.github.io/continuous/halfcauchy.html

\subsection{\texorpdfstring{\texttt{m.dist.improper\_uniform}}{m.dist.improper\_uniform}}\label{improper_uniform}

\begin{Shaded}
\begin{Highlighting}[]
\NormalTok{m.dist.improper\_uniform(}
\NormalTok{    support,}
\NormalTok{    batch\_shape,}
\NormalTok{    event\_shape,}
\NormalTok{    validate\_args}\OperatorTok{=}\VariableTok{None}\NormalTok{,}
\NormalTok{    name}\OperatorTok{=}\StringTok{\textquotesingle{}x\textquotesingle{}}\NormalTok{,}
\NormalTok{    obs}\OperatorTok{=}\VariableTok{None}\NormalTok{,}
\NormalTok{    mask}\OperatorTok{=}\VariableTok{None}\NormalTok{,}
\NormalTok{    infer}\OperatorTok{=}\VariableTok{None}\NormalTok{,}
\NormalTok{    sample}\OperatorTok{=}\VariableTok{False}\NormalTok{,}
\NormalTok{    seed}\OperatorTok{=}\VariableTok{None}\NormalTok{,}
\NormalTok{    shape}\OperatorTok{=}\NormalTok{(),}
\NormalTok{    event}\OperatorTok{=}\DecValTok{0}\NormalTok{,}
\NormalTok{    create\_obj}\OperatorTok{=}\VariableTok{False}\NormalTok{,}
\NormalTok{    to\_jax}\OperatorTok{=}\VariableTok{True}\NormalTok{,}
\NormalTok{)}
\end{Highlighting}
\end{Shaded}

\textbf{Improper Uniform} A helper distribution with zero
\texttt{log\_prob} over the \texttt{support} domain.

\textbf{Args} - \emph{support}
(numpyro.distributions.constraints.Constraint): The support of this
distribution.

\begin{itemize}
\item
  \emph{batch\_shape} (tuple): Batch shape of this distribution. It is
  usually safe to set \texttt{batch\_shape=()}.
\item
  \emph{event\_shape} (tuple): Event shape of this distribution.
\item
  \emph{shape} (tuple): A multi-purpose argument for shaping. When
  \texttt{sample=False} (model building), this is used with
  \texttt{.expand(shape)} to set the distribution's batch shape. When
  \texttt{sample=True} (direct sampling), this is used as
  \texttt{sample\_shape} to draw a raw JAX array of the given shape.
\item
  \emph{event} (int): The number of batch dimensions to reinterpret as
  event dimensions (used in model building).
\item
  \emph{mask} (jnp.ndarray, bool): Optional boolean array to mask
  observations.
\item
  \emph{create\_obj} (bool, optional): If True, returns the raw BF
  distribution object instead of creating a sample site. This is
  essential for building complex distributions like
  \texttt{MixtureSameFamily}. Defaults to False.
\item
  \emph{sample} (bool, optional): A control-flow argument. If
  \texttt{True}, the function will directly sample a raw JAX array from
  the distribution, bypassing the BF model context. If \texttt{False},
  it will create a \texttt{BF.sample} site within a model. Defaults to
  \texttt{False}.
\item
  \emph{obs} (jnp.ndarray, optional): The observed value for this random
  variable. If provided, the sample site is conditioned on this value,
  and the function returns the observed value. If \texttt{None}, the
  site is treated as a latent (unobserved) random variable. Defaults to
  \texttt{None}.
\item
  \emph{name} (str, optional): The name of the sample site in a BF
  model. This is used to uniquely identify the random variable. Defaults
  to `x'.
\end{itemize}

\textbf{Returns} - When \texttt{sample=False}: A BF ImproperUniform
distribution object (for model building). - When \texttt{sample=True}: A
JAX array of samples drawn from the ImproperUniform distribution (for
direct sampling). - When \texttt{create\_obj=True}: The raw BF
distribution object (for advanced use cases).

\textbf{Example Usage} from BayesForge import bf from
numpyro.distributions import constraints m = bf(platform=`cpu')

\begin{verbatim}
  def model():
      **flat prior over ordered vectors of length 10**
      x = m.dist.improper_uniform(constraints.ordered_vector, (), (10,), name='x')


  **Wrapper of**
  https://num.pyro.ai/en/stable/distributions.html#improperuniform
\end{verbatim}

\subsection{\texorpdfstring{\texttt{m.dist.inverse\_gamma}}{m.dist.inverse\_gamma}}\label{inverse_gamma}

\begin{Shaded}
\begin{Highlighting}[]
\NormalTok{m.dist.inverse\_gamma(}
\NormalTok{    concentration,}
\NormalTok{    rate}\OperatorTok{=}\FloatTok{1.0}\NormalTok{,}
\NormalTok{    validate\_args}\OperatorTok{=}\VariableTok{None}\NormalTok{,}
\NormalTok{    name}\OperatorTok{=}\StringTok{\textquotesingle{}x\textquotesingle{}}\NormalTok{,}
\NormalTok{    obs}\OperatorTok{=}\VariableTok{None}\NormalTok{,}
\NormalTok{    mask}\OperatorTok{=}\VariableTok{None}\NormalTok{,}
\NormalTok{    infer}\OperatorTok{=}\VariableTok{None}\NormalTok{,}
\NormalTok{    sample}\OperatorTok{=}\VariableTok{False}\NormalTok{,}
\NormalTok{    seed}\OperatorTok{=}\VariableTok{None}\NormalTok{,}
\NormalTok{    shape}\OperatorTok{=}\NormalTok{(),}
\NormalTok{    event}\OperatorTok{=}\DecValTok{0}\NormalTok{,}
\NormalTok{    create\_obj}\OperatorTok{=}\VariableTok{False}\NormalTok{,}
\NormalTok{    to\_jax}\OperatorTok{=}\VariableTok{True}\NormalTok{,}
\NormalTok{)}
\end{Highlighting}
\end{Shaded}

\textbf{Inverse Gamma} The InverseGamma distribution is a two-parameter
family of continuous probability distributions. It is defined by its
shape \(\alpha\) and rate \(\beta\) parameters. It is often used as a
prior distribution for precision parameters (inverse variance) in
Bayesian statistics.

\[
  f(x; \alpha, \beta) = \frac{\beta^\alpha}{\Gamma(\alpha)} x^{-(\alpha + 1)} \exp\left(-\frac{\beta}{x}\right), \quad x > 0
  \]

where: * (\Gamma(\alpha)) is the Gamma function.

\textbf{Args} - \emph{concentration} (jnp.ndarray): The concentration
parameter \(\alpha\) of the InverseGamma distribution. Must be positive.

\begin{itemize}
\item
  \emph{rate} (jnp.ndarray): The rate parameter \(\beta\) of the
  InverseGamma distribution. Must be positive.
\item
  \emph{shape} (tuple): A multi-purpose argument for shaping. When
  \texttt{sample=False} (model building), this is used with
  \texttt{shape} to set the distribution's batch shape. When
  \texttt{sample=True} (direct sampling), this is used as
  \texttt{sample\_shape} to draw a raw JAX array of the given shape.
\item
  \emph{event} (int): The number of batch dimensions to reinterpret as
  event dimensions (used in model building).
\item
  \emph{mask} (jnp.ndarray, bool): Optional boolean array to mask
  observations.
\item
  \emph{create\_obj} (bool): If True, returns the raw BF distribution
  object instead of creating a sample site. This is essential for
  building complex distributions like \texttt{MixtureSameFamily}.
\item
  \emph{sample} (bool, optional): A control-flow argument. If
  \texttt{True}, the function will directly sample a raw JAX array from
  the distribution, bypassing the BF model context. If \texttt{False},
  it will create a \texttt{BF.sample} site within a model. Defaults to
  \texttt{False}.
\item
  \emph{obs} (jnp.ndarray, optional): The observed value for this random
  variable. If provided, the sample site is conditioned on this value,
  and the function returns the observed value. If \texttt{None}, the
  site is treated as a latent (unobserved) random variable. Defaults to
  \texttt{None}.
\item
  \emph{name} (str, optional): The name of the sample site in a BF
  model. This is used to uniquely identify the random variable. Defaults
  to `x'.
\end{itemize}

\textbf{Returns} - When \texttt{sample=False}: A BF InverseGamma
distribution object (for model building). - When \texttt{sample=True}: A
JAX array of samples drawn from the InverseGamma distribution (for
direct sampling). - When \texttt{create\_obj=True}: The raw BF
distribution object (for advanced use cases).

\textbf{Example Usage} from BayesForge import bf m = bf(platform=`cpu')
m.dist.inverse\_gamma(concentration=2.0, rate=1.0, sample=True)

\textbf{Wrapper of}
https://num.pyro.ai/en/stable/distributions.html\#inversegamma

\subsection{\texorpdfstring{\texttt{m.dist.kumaraswamy}}{m.dist.kumaraswamy}}\label{kumaraswamy}

\begin{Shaded}
\begin{Highlighting}[]
\NormalTok{m.dist.kumaraswamy(}
\NormalTok{    concentration1,}
\NormalTok{    concentration0,}
\NormalTok{    validate\_args}\OperatorTok{=}\VariableTok{None}\NormalTok{,}
\NormalTok{    name}\OperatorTok{=}\StringTok{\textquotesingle{}x\textquotesingle{}}\NormalTok{,}
\NormalTok{    obs}\OperatorTok{=}\VariableTok{None}\NormalTok{,}
\NormalTok{    mask}\OperatorTok{=}\VariableTok{None}\NormalTok{,}
\NormalTok{    infer}\OperatorTok{=}\VariableTok{None}\NormalTok{,}
\NormalTok{    sample}\OperatorTok{=}\VariableTok{False}\NormalTok{,}
\NormalTok{    seed}\OperatorTok{=}\VariableTok{None}\NormalTok{,}
\NormalTok{    shape}\OperatorTok{=}\NormalTok{(),}
\NormalTok{    event}\OperatorTok{=}\DecValTok{0}\NormalTok{,}
\NormalTok{    create\_obj}\OperatorTok{=}\VariableTok{False}\NormalTok{,}
\NormalTok{    to\_jax}\OperatorTok{=}\VariableTok{True}\NormalTok{,}
\NormalTok{)}
\end{Highlighting}
\end{Shaded}

\textbf{Kumaraswamy} The Kumaraswamy distribution is a continuous
probability distribution defined on the interval {[}0, 1{]}. It is a
flexible distribution that can take on various shapes depending on its
parameters.

\[
  \] f(x; a, b) = a b x\^{}\{a-1\} (1 - x\textsuperscript{a)}\{b-1\},
\quad x \in (0, 1) \$\$

where: * ( a ) and ( b ) are shape parameters. * ( x ) is the random
variable.

\textbf{Args} - \emph{concentration1} (jnp.ndarray): The first shape
parameter. Must be positive.

\begin{itemize}
\item
  \emph{concentration0} (jnp.ndarray): The second shape parameter. Must
  be positive.
\item
  \emph{shape} (tuple): A multi-purpose argument for shaping. When
  \texttt{sample=False} (model building), this is used with
  \texttt{shape} to set the distribution's batch shape. When
  \texttt{sample=True} (direct sampling), this is used as
  \texttt{sample\_shape} to draw a raw JAX array of the given shape.
\item
  \emph{event} (int): The number of batch dimensions to reinterpret as
  event dimensions (used in model building).
\item
  \emph{mask} (jnp.ndarray, bool): Optional boolean array to mask
  observations.
\item
  \emph{create\_obj} (bool): If True, returns the raw BF distribution
  object instead of creating a sample site. This is essential for
  building complex distributions like \texttt{MixtureSameFamily}.
\item
  \emph{sample} (bool, optional): A control-flow argument. If
  \texttt{True}, the function will directly sample a raw JAX array from
  the distribution, bypassing the BF model context. If \texttt{False},
  it will create a \texttt{BF.sample} site within a model. Defaults to
  \texttt{False}.
\item
  \emph{obs} (jnp.ndarray, optional): The observed value for this random
  variable. If provided, the sample site is conditioned on this value,
  and the function returns the observed value. If \texttt{None}, the
  site is treated as a latent (unobserved) random variable. Defaults to
  \texttt{None}.
\item
  \emph{name} (str, optional): The name of the sample site in a BF
  model. This is used to uniquely identify the random variable. Defaults
  to `x'.
\end{itemize}

\textbf{Returns} - When \texttt{sample=False}: A BF Kumaraswamy
distribution object (for model building). - When \texttt{sample=True}: A
JAX array of samples drawn from the Kumaraswamy distribution (for direct
sampling). - When \texttt{create\_obj=True}: The raw BF distribution
object (for advanced use cases).

\textbf{Example Usage} from BayesForge import bf m = bf(platform=`cpu')
m.dist.kumaraswamy(concentration1=2.0, concentration0=3.0, sample=True)

\textbf{Wrapper of}
https://num.pyro.ai/en/stable/distributions.html\#kumaraswamy
\textbf{Rerferences}
https://en.wikipedia.org/wiki/Kumaraswamy\_distribution

\subsection{\texorpdfstring{\texttt{m.dist.laplace}}{m.dist.laplace}}\label{laplace}

\begin{Shaded}
\begin{Highlighting}[]
\NormalTok{m.dist.laplace(}
\NormalTok{    loc}\OperatorTok{=}\FloatTok{0.0}\NormalTok{,}
\NormalTok{    scale}\OperatorTok{=}\FloatTok{1.0}\NormalTok{,}
\NormalTok{    validate\_args}\OperatorTok{=}\VariableTok{None}\NormalTok{,}
\NormalTok{    name}\OperatorTok{=}\StringTok{\textquotesingle{}x\textquotesingle{}}\NormalTok{,}
\NormalTok{    obs}\OperatorTok{=}\VariableTok{None}\NormalTok{,}
\NormalTok{    mask}\OperatorTok{=}\VariableTok{None}\NormalTok{,}
\NormalTok{    infer}\OperatorTok{=}\VariableTok{None}\NormalTok{,}
\NormalTok{    sample}\OperatorTok{=}\VariableTok{False}\NormalTok{,}
\NormalTok{    seed}\OperatorTok{=}\VariableTok{None}\NormalTok{,}
\NormalTok{    shape}\OperatorTok{=}\NormalTok{(),}
\NormalTok{    event}\OperatorTok{=}\DecValTok{0}\NormalTok{,}
\NormalTok{    create\_obj}\OperatorTok{=}\VariableTok{False}\NormalTok{,}
\NormalTok{    to\_jax}\OperatorTok{=}\VariableTok{True}\NormalTok{,}
\NormalTok{)}
\end{Highlighting}
\end{Shaded}

\textbf{Laplace} Samples from a Laplace distribution, also known as the
double exponential distribution. The Laplace distribution is defined by
its location parameter (loc) and scale parameter (scale).

\[
  f(x \mid \mu, b) = \frac{1}{2b} \exp\left( -\frac{|x - \mu|}{b} \right)
  \]

Where: * (\(\mu\)) is the \textbf{location parameter}, indicating the
peak of the distribution. * (b \textgreater{} 0) is the \textbf{scale
parameter}, controlling the spread of the distribution.

\textbf{Args} - \emph{loc} (jnp.ndarray): Location parameter of the
Laplace distribution.

\begin{itemize}
\item
  \emph{sample} (jnp.ndarray): Scale parameter of the Laplace
  distribution. Must be positive.
\item
  \emph{shape} (tuple): A multi-purpose argument for shaping. When
  \texttt{sample=False} (model building), this is used with
  \texttt{shape} to set the distribution's batch shape. When
  \texttt{sample=True} (direct sampling), this is used as
  \texttt{sample\_shape} to draw a raw JAX array of the given shape.
\item
  \emph{event} (int): The number of batch dimensions to reinterpret as
  event dimensions (used in model building).
\item
  \emph{mask} (jnp.ndarray, bool): Optional boolean array to mask
  observations.
\item
  \emph{create\_obj} (bool): If True, returns the raw BF distribution
  object instead of creating a sample site. This is essential for
  building complex distributions like \texttt{MixtureSameFamily}.
\item
  \emph{sample} (bool, optional): A control-flow argument. If
  \texttt{True}, the function will directly sample a raw JAX array from
  the distribution, bypassing the BF model context. If \texttt{False},
  it will create a \texttt{BF.sample} site within a model. Defaults to
  \texttt{False}.
\item
  \emph{seed} (int, optional): An integer seed for this draw when
  \texttt{sample=True}. By default (\texttt{None}) the draw comes from
  the session's random stream, set by \texttt{bf(rand\_seed=...)}, so
  consecutive draws differ but the whole sequence is reproducible. No
  effect when \texttt{sample=False}.
\item
  \emph{obs} (jnp.ndarray, optional): The observed value for this random
  variable. If provided, the sample site is conditioned on this value,
  and the function returns the observed value. If \texttt{None}, the
  site is treated as a latent (unobserved) random variable. Defaults to
  \texttt{None}.
\item
  \emph{name} (str, optional): The name of the sample site in a BF
  model. This is used to uniquely identify the random variable. Defaults
  to `x'.
\end{itemize}

\textbf{Returns} - When \texttt{sample=False}: A BF Laplace distribution
object (for model building). - When \texttt{sample=True}: A JAX array of
samples drawn from the Laplace distribution (for direct sampling). -
When \texttt{create\_obj=True}: The raw BF distribution object (for
advanced use cases).

\textbf{Example Usage} from BayesForge import bf m = bf(platform=`cpu')
m.dist.laplace(loc=0.0, scale=1.0, sample=True)

\textbf{Wrapper of}
https://num.pyro.ai/en/stable/distributions.html\#laplace

\subsection{\texorpdfstring{\texttt{m.dist.left\_truncated\_distribution}}{m.dist.left\_truncated\_distribution}}\label{left_truncated_distribution}

\begin{Shaded}
\begin{Highlighting}[]
\NormalTok{m.dist.left\_truncated\_distribution(}
\NormalTok{    base\_dist,}
\NormalTok{    low}\OperatorTok{=}\FloatTok{0.0}\NormalTok{,}
\NormalTok{    validate\_args}\OperatorTok{=}\VariableTok{None}\NormalTok{,}
\NormalTok{    name}\OperatorTok{=}\StringTok{\textquotesingle{}x\textquotesingle{}}\NormalTok{,}
\NormalTok{    obs}\OperatorTok{=}\VariableTok{None}\NormalTok{,}
\NormalTok{    mask}\OperatorTok{=}\VariableTok{None}\NormalTok{,}
\NormalTok{    infer}\OperatorTok{=}\VariableTok{None}\NormalTok{,}
\NormalTok{    sample}\OperatorTok{=}\VariableTok{False}\NormalTok{,}
\NormalTok{    seed}\OperatorTok{=}\VariableTok{None}\NormalTok{,}
\NormalTok{    shape}\OperatorTok{=}\NormalTok{(),}
\NormalTok{    event}\OperatorTok{=}\DecValTok{0}\NormalTok{,}
\NormalTok{    create\_obj}\OperatorTok{=}\VariableTok{False}\NormalTok{,}
\NormalTok{    to\_jax}\OperatorTok{=}\VariableTok{True}\NormalTok{,}
\NormalTok{)}
\end{Highlighting}
\end{Shaded}

\textbf{Left Truncated} Samples from a left-truncated distribution.

A left-truncated distribution is a probability distribution obtained by
restricting the support of another distribution to values greater than a
specified lower bound. This is useful when dealing with data that is
known to be greater than a certain value. All the ``mass'' below (or
equal to) (a) is \textbf{excluded} (not just unobserved, but removed
from the sample/analysis).

Let (X) be a random variable with PDF (f\_X(x)) and CDF (F\_X(x)).
Choose a truncation point (a). Then define a new random variable (Y = X
\mid (X \textgreater{} a)). The PDF and CDF of (Y) are:

\begin{itemize}
\item
  \textbf{PDF} for ( y \textgreater{} a ): \[
    f_Y(y) = \frac{f_X(y)}{1 - F_X(a)} .
    \] (Here (1 - F\_X(a) = P(X \textgreater{} a)) is the normalising
  constant)
  (\href{https://en.wikipedia.org/wiki/Matrix_normal_distribution}{Wikipédia})
\item
  \textbf{CDF} for (y \textgreater{} a): \[
    F_Y(y) = P(Y \le y) = \frac{F_X(y) - F_X(a)}{1 - F_X(a)} .
    \]
\end{itemize}

\textbf{Args} - \emph{base\_dist}: The base distribution to truncate.
Must be univariate and have real support.

\begin{itemize}
\item
  \emph{low}: The lower truncation bound. Values less than this are
  excluded from the distribution.
\item
  \emph{shape} (tuple): A multi-purpose argument for shaping. When
  \texttt{sample=False} (model building), this is used with
  \texttt{shape} to set the distribution's batch shape. When
  \texttt{sample=True} (direct sampling), this is used as
  \texttt{sample\_shape} to draw a raw JAX array of the given shape.
\item
  \emph{event} (int): The number of batch dimensions to reinterpret as
  event dimensions (used in model building).
\item
  \emph{mask} (jnp.ndarray, bool): Optional boolean array to mask
  observations.
\item
  \emph{create\_obj} (bool): If True, returns the raw BF distribution
  object instead of creating a sample site. This is essential for
  building complex distributions like \texttt{MixtureSameFamily}.
\item
  \emph{sample} (bool, optional): A control-flow argument. If
  \texttt{True}, the function will directly sample a raw JAX array from
  the distribution, bypassing the BF model context. If \texttt{False},
  it will create a \texttt{BF.sample} site within a model. Defaults to
  \texttt{False}.
\item
  \emph{seed} (int, optional): An integer seed for this draw when
  \texttt{sample=True}. By default (\texttt{None}) the draw comes from
  the session's random stream, set by \texttt{bf(rand\_seed=...)}, so
  consecutive draws differ but the whole sequence is reproducible. No
  effect when \texttt{sample=False}.
\item
  \emph{obs} (jnp.ndarray, optional): The observed value for this random
  variable. If provided, the sample site is conditioned on this value,
  and the function returns the observed value. If \texttt{None}, the
  site is treated as a latent (unobserved) random variable. Defaults to
  \texttt{None}.
\item
  \emph{name} (str, optional): The name of the sample site in a BF
  model. This is used to uniquely identify the random variable. Defaults
  to `x'.
\end{itemize}

\textbf{Returns} - When \texttt{sample=False}: A BF
LeftTruncatedDistribution distribution object (for model building).

\begin{itemize}
\item
  When \texttt{sample=True}: A JAX array of samples drawn from the
  LeftTruncatedDistribution distribution (for direct sampling).
\item
  When \texttt{create\_obj=True}: The raw BF distribution object (for
  advanced use cases).

  \textbf{Wrapper of:
  https://num.pyro.ai/en/stable/distributions.html\#lefttruncateddistribution}
  \textbf{References}
  https://en.wikipedia.org/wiki/Truncated\_distribution
  https://encyclopediaofmath.org/wiki/Truncated\_distribution
\end{itemize}

\subsection{\texorpdfstring{\texttt{m.dist.levy}}{m.dist.levy}}\label{levy}

\begin{Shaded}
\begin{Highlighting}[]
\NormalTok{m.dist.levy(}
\NormalTok{    loc,}
\NormalTok{    scale,}
\NormalTok{    validate\_args}\OperatorTok{=}\VariableTok{None}\NormalTok{,}
\NormalTok{    name}\OperatorTok{=}\StringTok{\textquotesingle{}x\textquotesingle{}}\NormalTok{,}
\NormalTok{    obs}\OperatorTok{=}\VariableTok{None}\NormalTok{,}
\NormalTok{    mask}\OperatorTok{=}\VariableTok{None}\NormalTok{,}
\NormalTok{    infer}\OperatorTok{=}\VariableTok{None}\NormalTok{,}
\NormalTok{    sample}\OperatorTok{=}\VariableTok{False}\NormalTok{,}
\NormalTok{    seed}\OperatorTok{=}\VariableTok{None}\NormalTok{,}
\NormalTok{    shape}\OperatorTok{=}\NormalTok{(),}
\NormalTok{    event}\OperatorTok{=}\DecValTok{0}\NormalTok{,}
\NormalTok{    create\_obj}\OperatorTok{=}\VariableTok{False}\NormalTok{,}
\NormalTok{    to\_jax}\OperatorTok{=}\VariableTok{True}\NormalTok{,}
\NormalTok{)}
\end{Highlighting}
\end{Shaded}

\textbf{Levy} Samples from a Levy distribution. The Lévy distribution is
a continuous probability distribution on the positive real line (or
shifted positive line) that is \emph{heavy-tailed}, \emph{skewed}, and
arises naturally in connection with stable distributions specifically
the case with stability index \(\alpha = \tfrac12\). It is often used in
contexts such as hitting‐time problems for Brownian motion, physics
(e.g., van der Waals line‐shapes), and modelling very heavy‐tailed
phenomena. Let (X) be a Lévy‐distributed random variable with
\textbf{location} parameter \(\mu\) and \textbf{scale} parameter (c
\textgreater{} 0). The support is \(x \ge \mu\).

\textbf{PDF} \[
f(x; \mu, c) = \sqrt{\frac{c}{2\pi}} ; \frac{1}{(x - \mu)^{3/2}} ; \exp!\Bigl( -\frac{c}{2 (x - \mu)} \Bigr),
\quad x \ge \mu.
\]

\textbf{CDF} \[
F(x; \mu, c) = \Pr(X \le x)
= \operatorname{erfc}!\Bigl(\sqrt{\tfrac{c}{2 (x - \mu)}}\Bigr)
= 2 - 2,\Phi!\Bigl(\sqrt{\tfrac{c}{(x - \mu)}}\Bigr),
\quad x \ge \mu.
\] Here \(\operatorname{erfc}\) is the complementary error function and
\(\Phi\) is the standard normal CDF.

In the \emph{standard} case with \(\mu = 0\) and \(c = 1\), one often
simply writes \(X \sim \text{Lévy}(0,1)\).

\textbf{Args} - \emph{loc} (jnp.ndarray): Location parameter.

\begin{itemize}
\item
  \emph{sample} (jnp.ndarray): Scale parameter.
\item
  \emph{shape} (tuple): A multi-purpose argument for shaping. When
  \texttt{sample=False} (model building), this is used with
  \texttt{shape} to set the distribution's batch shape. When
  \texttt{sample=True} (direct sampling), this is used as
  \texttt{sample\_shape} to draw a raw JAX array of the given shape.
\item
  \emph{event} (int): The number of batch dimensions to reinterpret as
  event dimensions (used in model building).
\item
  \emph{mask} (jnp.ndarray, bool): Optional boolean array to mask
  observations.
\item
  \emph{create\_obj} (bool): If True, returns the raw BF distribution
  object instead of creating a sample site.
\item
  \emph{sample} (bool, optional): A control-flow argument. If
  \texttt{True}, the function will directly sample a raw JAX array from
  the distribution, bypassing the BF model context. If \texttt{False},
  it will create a \texttt{BF.sample} site within a model. Defaults to
  \texttt{False}.
\item
  \emph{seed} (int, optional): An integer seed for this draw when
  \texttt{sample=True}. By default (\texttt{None}) the draw comes from
  the session's random stream, set by \texttt{bf(rand\_seed=...)}, so
  consecutive draws differ but the whole sequence is reproducible. No
  effect when \texttt{sample=False}.
\item
  \emph{obs} (jnp.ndarray, optional): The observed value for this random
  variable. If provided, the sample site is conditioned on this value,
  and the function returns the observed value. If \texttt{None}, the
  site is treated as a latent (unobserved) random variable. Defaults to
  \texttt{None}.
\item
  \emph{name} (str, optional): The name of the sample site in a BF
  model. This is used to uniquely identify the random variable. Defaults
  to `x'.
\end{itemize}

\textbf{Returns} BF Levy distribution object: When \texttt{sample=False}
(for model building). JAX array: When \texttt{sample=True} (for direct
sampling). BF distribution object: When \texttt{create\_obj=True} (for
advanced use cases).

\textbf{Example Usage} from BayesForge import bf m = bf(platform=`cpu')
m.dist.levy(loc=0.0, scale=1.0, sample=True)

\textbf{Wrapper of}
https://num.pyro.ai/en/stable/distributions.html\#levy

\subsection{\texorpdfstring{\texttt{m.dist.lkj}}{m.dist.lkj}}\label{lkj}

\begin{Shaded}
\begin{Highlighting}[]
\NormalTok{m.dist.lkj(}
\NormalTok{    dimension,}
\NormalTok{    concentration}\OperatorTok{=}\FloatTok{1.0}\NormalTok{,}
\NormalTok{    sample\_method}\OperatorTok{=}\StringTok{\textquotesingle{}onion\textquotesingle{}}\NormalTok{,}
\NormalTok{    validate\_args}\OperatorTok{=}\VariableTok{None}\NormalTok{,}
\NormalTok{    name}\OperatorTok{=}\StringTok{\textquotesingle{}x\textquotesingle{}}\NormalTok{,}
\NormalTok{    obs}\OperatorTok{=}\VariableTok{None}\NormalTok{,}
\NormalTok{    mask}\OperatorTok{=}\VariableTok{None}\NormalTok{,}
\NormalTok{    infer}\OperatorTok{=}\VariableTok{None}\NormalTok{,}
\NormalTok{    sample}\OperatorTok{=}\VariableTok{False}\NormalTok{,}
\NormalTok{    seed}\OperatorTok{=}\VariableTok{None}\NormalTok{,}
\NormalTok{    shape}\OperatorTok{=}\NormalTok{(),}
\NormalTok{    event}\OperatorTok{=}\DecValTok{0}\NormalTok{,}
\NormalTok{    create\_obj}\OperatorTok{=}\VariableTok{False}\NormalTok{,}
\NormalTok{    to\_jax}\OperatorTok{=}\VariableTok{True}\NormalTok{,}
\NormalTok{)}
\end{Highlighting}
\end{Shaded}

\textbf{Lewandowski Kurowicka Joe (LKJ)} The LKJ distribution is
controlled by the concentration parameter \(\eta\) to make the
probability of the correlation matrix \(M\) proportional to
\(\det(M)^{\eta - 1}\). When \(\eta = 1\), the distribution is uniform
over correlation matrices. When \(\eta > 1\), the distribution favors
samples with large determinants. When \(\eta < 1\), the distribution
favors samples with small determinants.

\[
     P(M) \propto |\det(M)|^{\eta - 1}
  \]

\textbf{Args} - \emph{dimension} (int): The dimension of the correlation
matrices.

\begin{itemize}
\item
  \emph{concentration} (ndarray): The concentration/shape parameter of
  the distribution (often referred to as eta). Must be positive.
\item
  \emph{sample\_method} (str): Either ``cvine'' or ``onion''. Both
  methods are proposed in
  \href{https://en.wikipedia.org/wiki/Matrix_normal_distribution}{1} and
  offer the same distribution over correlation matrices. But they are
  different in how to generate samples. Defaults to ``onion''.
\item
  \emph{shape} (tuple): A multi-purpose argument for shaping. When
  \texttt{sample=False} (model building), this is used with
  \texttt{shape} to set the distribution's batch shape. When
  \texttt{sample=True} (direct sampling), this is used as
  \texttt{sample\_shape} to draw a raw JAX array of the given shape.
\item
  \emph{event} (int): The number of batch dimensions to reinterpret as
  event dimensions (used in model building).
\item
  \emph{mask} (jnp.ndarray, bool): Optional boolean array to mask
  observations.
\item
  \emph{create\_obj} (bool): If True, returns the raw BF distribution
  object instead of creating a sample site. This is essential for
  building complex distributions like \texttt{MixtureSameFamily}.
\item
  \emph{sample} (bool, optional): A control-flow argument. If
  \texttt{True}, the function will directly sample a raw JAX array from
  the distribution, bypassing the BF model context. If \texttt{False},
  it will create a \texttt{BF.sample} site within a model. Defaults to
  \texttt{False}.
\item
  \emph{obs} (jnp.ndarray, optional): The observed value for this random
  variable. If provided, the sample site is conditioned on this value,
  and the function returns the observed value. If \texttt{None}, the
  site is treated as a latent (unobserved) random variable. Defaults to
  \texttt{None}.
\item
  \emph{name} (str, optional): The name of the sample site in a BF
  model. This is used to uniquely identify the random variable. Defaults
  to `x'.
\end{itemize}

\textbf{Returns} - When \texttt{sample=False}: A BF LKJ distribution
object (for model building).

\begin{itemize}
\item
  When \texttt{sample=True}: A JAX array of samples drawn from the LKJ
  distribution (for direct sampling).
\item
  When \texttt{create\_obj=True}: The raw BF distribution object (for
  advanced use cases).

  \textbf{Example Usage} from BayesForge import bf m =
  bf(platform=`cpu') m.dist.lkj(dimension=2, concentration=1.0,
  sample=True)

  \textbf{Wrapper of}
  https://num.pyro.ai/en/stable/distributions.html\#lkj
  \textbf{References}
  https://en.wikipedia.org/wiki/Kumaraswamy\_distribution
\end{itemize}

\subsection{\texorpdfstring{\texttt{m.dist.lkj\_cholesky}}{m.dist.lkj\_cholesky}}\label{lkj_cholesky}

\begin{Shaded}
\begin{Highlighting}[]
\NormalTok{m.dist.lkj\_cholesky(}
\NormalTok{    dimension,}
\NormalTok{    concentration}\OperatorTok{=}\FloatTok{1.0}\NormalTok{,}
\NormalTok{    sample\_method}\OperatorTok{=}\StringTok{\textquotesingle{}onion\textquotesingle{}}\NormalTok{,}
\NormalTok{    validate\_args}\OperatorTok{=}\VariableTok{None}\NormalTok{,}
\NormalTok{    name}\OperatorTok{=}\StringTok{\textquotesingle{}x\textquotesingle{}}\NormalTok{,}
\NormalTok{    obs}\OperatorTok{=}\VariableTok{None}\NormalTok{,}
\NormalTok{    mask}\OperatorTok{=}\VariableTok{None}\NormalTok{,}
\NormalTok{    infer}\OperatorTok{=}\VariableTok{None}\NormalTok{,}
\NormalTok{    sample}\OperatorTok{=}\VariableTok{False}\NormalTok{,}
\NormalTok{    seed}\OperatorTok{=}\VariableTok{None}\NormalTok{,}
\NormalTok{    shape}\OperatorTok{=}\NormalTok{(),}
\NormalTok{    event}\OperatorTok{=}\DecValTok{0}\NormalTok{,}
\NormalTok{    create\_obj}\OperatorTok{=}\VariableTok{False}\NormalTok{,}
\NormalTok{    to\_jax}\OperatorTok{=}\VariableTok{True}\NormalTok{,}
\NormalTok{)}
\end{Highlighting}
\end{Shaded}

\textbf{LKJ Cholesky} The LKJ (Leonard-Kjærgaard-Jørgensen) Cholesky
distribution is a family of distributions on symmetric matrices, often
used as a prior for the Cholesky decomposition of a symmetric matrix. It
is particularly useful in BayesForge for models with covariance
structure. Given a lower triangular matrix (L) with unit diagonal
entries, the PDF is: \[
p(L | \eta) \propto \prod_{k=2}^d L_{kk}^{d-k+2\eta-2}
\]

where: * (L\_\{kk\}) is the diagonal element of (L). * (\eta) is the
shape parameter. * (d) is the dimension of the matrix.

\begin{itemize}
\tightlist
\item
  (\eta = 1): Uniform prior over correlation matrices.
\item
  (\eta \textgreater{} 1): Prior favors correlation matrices close to
  the identity matrix (i.e., weak correlations).
\item
  (\eta \textless{} 1): Prior allows stronger correlations.
\end{itemize}

\textbf{Args} - *dimension (int): The dimension of the correlation
matrices.

\begin{itemize}
\item
  *concentration (float): A parameter controlling the concentration of
  the distribution around the identity matrix. Higher values indicate
  greater concentration. Must be greater than 1.
\item
  \emph{shape} (tuple): A multi-purpose argument for shaping. When
  \texttt{sample=False} (model building), this is used with
  \texttt{shape} to set the distribution's batch shape. When
  \texttt{sample=True} (direct sampling), this is used as
  \texttt{sample\_shape} to draw a raw JAX array of the given shape.
\item
  \emph{event} (int): The number of batch dimensions to reinterpret as
  event dimensions (used in model building).
\item
  \emph{mask} (jnp.ndarray, bool): Optional boolean array to mask
  observations.
\item
  \emph{create\_obj} (bool): If True, returns the raw BF distribution
  object instead of creating a sample site. This is essential for
  building complex distributions like \texttt{MixtureSameFamily}.
\item
  \emph{sample} (bool, optional): A control-flow argument. If
  \texttt{True}, the function will directly sample a raw JAX array from
  the distribution, bypassing the BF model context. If \texttt{False},
  it will create a \texttt{BF.sample} site within a model. Defaults to
  \texttt{False}.
\item
  \emph{seed} (int, optional): An integer seed for this draw when
  \texttt{sample=True}. By default (\texttt{None}) the draw comes from
  the session's random stream, set by \texttt{bf(rand\_seed=...)}, so
  consecutive draws differ but the whole sequence is reproducible. No
  effect when \texttt{sample=False}.
\item
  \emph{obs} (jnp.ndarray, optional): The observed value for this random
  variable. If provided, the sample site is conditioned on this value,
  and the function returns the observed value. If \texttt{None}, the
  site is treated as a latent (unobserved) random variable. Defaults to
  \texttt{None}.
\item
  \emph{name} (str, optional): The name of the sample site in a BF
  model. This is used to uniquely identify the random variable. Defaults
  to `x'.

  Attributes: concentration (float): The concentration parameter.
\end{itemize}

\textbf{References}
https://mc-stan.org/docs/2\_21/functions-reference/cholesky-lkj-correlation-distribution.html

\subsection{\texorpdfstring{\texttt{m.dist.log\_normal}}{m.dist.log\_normal}}\label{log_normal}

\begin{Shaded}
\begin{Highlighting}[]
\NormalTok{m.dist.log\_normal(}
\NormalTok{    loc}\OperatorTok{=}\FloatTok{0.0}\NormalTok{,}
\NormalTok{    scale}\OperatorTok{=}\FloatTok{1.0}\NormalTok{,}
\NormalTok{    validate\_args}\OperatorTok{=}\VariableTok{None}\NormalTok{,}
\NormalTok{    name}\OperatorTok{=}\StringTok{\textquotesingle{}x\textquotesingle{}}\NormalTok{,}
\NormalTok{    obs}\OperatorTok{=}\VariableTok{None}\NormalTok{,}
\NormalTok{    mask}\OperatorTok{=}\VariableTok{None}\NormalTok{,}
\NormalTok{    infer}\OperatorTok{=}\VariableTok{None}\NormalTok{,}
\NormalTok{    sample}\OperatorTok{=}\VariableTok{False}\NormalTok{,}
\NormalTok{    seed}\OperatorTok{=}\VariableTok{None}\NormalTok{,}
\NormalTok{    shape}\OperatorTok{=}\NormalTok{(),}
\NormalTok{    event}\OperatorTok{=}\DecValTok{0}\NormalTok{,}
\NormalTok{    create\_obj}\OperatorTok{=}\VariableTok{False}\NormalTok{,}
\NormalTok{    to\_jax}\OperatorTok{=}\VariableTok{True}\NormalTok{,}
\NormalTok{)}
\end{Highlighting}
\end{Shaded}

\textbf{Log-Normal} The Log-Normal distribution is a probability
distribution defined for positive real-valued random variables,
parameterized by a location parameter (loc : \(\mu\)) and a scale
parameter (scale). It arises when the logarithm of a random variable is
normally distributed. * A random variable (X) is \textbf{log-normal} if
its natural logarithm (\ln X) is (approximately) normally distributed.
Equivalently, one can write: \[
X = \exp(Y), \quad \text{where } Y \sim \mathcal{N}(\mu, \sigma^2).
\]

Because (X) is the exponential of a normal, it is always positive
(support (x\textgreater0)). It is useful in contexts where growth,
multiplicative effects, or compounding factors dominate (e.g.~stock
prices, income, sizes of biological organisms).

\textbf{Parameters} * \(\mu\): the mean of \(\ln X\) (i.e.~the location
parameter in log-space) * \(\sigma > 0\): the standard deviation of
\(\ln X\) (i.e.~the ``scale'' in log-space)

Sometimes another parameterization uses a threshold (shift) \(\theta\),
i.e.~\(X = \theta + \exp(Y)\). For the two-parameter (no shift) case:

\begin{itemize}
\item
  \textbf{PDF} \[
  f(x; \mu, \sigma) = \frac{1}{x,\sigma\sqrt{2\pi}} ; \exp!\Bigl( -\frac{(\ln x - \mu)^2}{2\sigma^2} \Bigr), \quad x > 0.
  \]
\item
  \textbf{CDF} \[
  F(x; \mu, \sigma) = \Pr(X \le x) = \Phi!\bigl( \frac{\ln x - \mu}{\sigma} \bigr),
  \] where \(\Phi\) is the standard normal CDF.
\end{itemize}

\begin{center}\rule{0.5\linewidth}{0.5pt}\end{center}

\textbf{Args} - \emph{loc} (float): Location parameter.

\begin{itemize}
\item
  \emph{scale} (float): Scale parameter.
\item
  \emph{shape} (tuple): A multi-purpose argument for shaping. When
  \texttt{sample=False} (model building), this is used with
  \texttt{shape} to set the distribution's batch shape. When
  \texttt{sample=True} (direct sampling), this is used as
  \texttt{sample\_shape} to draw a raw JAX array of the given shape.
\item
  \emph{event} (int): The number of batch dimensions to reinterpret as
  event dimensions (used in model building).
\item
  \emph{mask} (jnp.ndarray, bool): Optional boolean array to mask
  observations.
\item
  \emph{create\_obj} (bool): If True, returns the raw BF distribution
  object instead of creating a sample site. This is essential for
  building complex distributions like \texttt{MixtureSameFamily}.
\item
  \emph{sample} (bool, optional): A control-flow argument. If
  \texttt{True}, the function will directly sample a raw JAX array from
  the distribution, bypassing the BF model context. If \texttt{False},
  it will create a \texttt{BF.sample} site within a model. Defaults to
  \texttt{False}.
\item
  \emph{seed} (int, optional): An integer seed for this draw when
  \texttt{sample=True}. By default (\texttt{None}) the draw comes from
  the session's random stream, set by \texttt{bf(rand\_seed=...)}, so
  consecutive draws differ but the whole sequence is reproducible. No
  effect when \texttt{sample=False}.
\item
  \emph{obs} (jnp.ndarray, optional): The observed value for this random
  variable. If provided, the sample site is conditioned on this value,
  and the function returns the observed value. If \texttt{None}, the
  site is treated as a latent (unobserved) random variable. Defaults to
  \texttt{None}.
\item
  \emph{name} (str, optional): The name of the sample site in a BF
  model. This is used to uniquely identify the random variable. Defaults
  to `x'.
\end{itemize}

\textbf{Returns} BF LogNormal distribution object (for model building).
JAX array of samples drawn from the LogNormal distribution (for direct
sampling). The raw BF distribution object (for advanced use cases).

\textbf{Example Usage} from BayesForge import bf m = bf(platform=`cpu')
m.dist.log\_normal(loc=0.0, scale=1.0, sample=True)

\textbf{Wrapper of:
https://num.pyro.ai/en/stable/distributions.html\#lognormal}
\textbf{References}
https://en.wikipedia.org/wiki/Log-normal\_distribution

\subsection{\texorpdfstring{\texttt{m.dist.log\_uniform}}{m.dist.log\_uniform}}\label{log_uniform}

\begin{Shaded}
\begin{Highlighting}[]
\NormalTok{m.dist.log\_uniform(}
\NormalTok{    low,}
\NormalTok{    high,}
\NormalTok{    validate\_args}\OperatorTok{=}\VariableTok{None}\NormalTok{,}
\NormalTok{    name}\OperatorTok{=}\StringTok{\textquotesingle{}x\textquotesingle{}}\NormalTok{,}
\NormalTok{    obs}\OperatorTok{=}\VariableTok{None}\NormalTok{,}
\NormalTok{    mask}\OperatorTok{=}\VariableTok{None}\NormalTok{,}
\NormalTok{    infer}\OperatorTok{=}\VariableTok{None}\NormalTok{,}
\NormalTok{    sample}\OperatorTok{=}\VariableTok{False}\NormalTok{,}
\NormalTok{    seed}\OperatorTok{=}\VariableTok{None}\NormalTok{,}
\NormalTok{    shape}\OperatorTok{=}\NormalTok{(),}
\NormalTok{    event}\OperatorTok{=}\DecValTok{0}\NormalTok{,}
\NormalTok{    create\_obj}\OperatorTok{=}\VariableTok{False}\NormalTok{,}
\NormalTok{    to\_jax}\OperatorTok{=}\VariableTok{True}\NormalTok{,}
\NormalTok{)}
\end{Highlighting}
\end{Shaded}

\textbf{Log-Uniform} Samples from a Log Uniform distribution.

The Log Uniform distribution is defined over the positive real numbers
and is the result of applying an exponential transformation to a uniform
distribution over the interval {[}low, high{]}. It is often used when
modeling parameters that must be positive.

A random variable \(X\) is \textbf{log-uniform} on \([a, b]\), with
\(0 < a < b\), if \(\ln X\) is uniformly distributed on
\([\ln a, \ln b]\). Equivalently, the density of \(X\) is proportional
to \(1/x\) over that interval. This distribution is sometimes called the
\emph{reciprocal distribution}. It is useful in modeling scales spanning
several orders of magnitude, where you want every decade (or
log-interval) to have equal weight.

For \(x \in [a, b]\):

\begin{itemize}
\item
  \textbf{PDF} \[
  f(x) = \frac{1}{x , [\ln(b) - \ln(a)]}, \quad a \le x \le b.
  \] Outside \([a, b]\), \(f(x) = 0\).
\item
  \textbf{CDF} \[
  F(x) = \Pr(X \le x) = \frac{\ln(x) - \ln(a)}{\ln(b) - \ln(a)}, \quad a \le x \le b.
  \] (For \(x < a\), \(F(x) = 0\); for \(x > b\), \(F(x) = 1\).
\end{itemize}

\textbf{Args} - \emph{low} (jnp.ndarray): The lower bound of the uniform
distribution's log-space. Must be positive.

\begin{itemize}
\item
  \emph{high} (jnp.ndarray): The upper bound of the uniform
  distribution's log-space. Must be positive.
\item
  \emph{shape} (tuple): A multi-purpose argument for shaping. When
  \texttt{sample=False} (model building), this is used with
  \texttt{shape} to set the distribution's batch shape. When
  \texttt{sample=True} (direct sampling), this is used as
  \texttt{sample\_shape} to draw a raw JAX array of the given shape.
\item
  \emph{event} (int): The number of batch dimensions to reinterpret as
  event dimensions (used in model building).
\item
  \emph{mask} (jnp.ndarray, bool): Optional boolean array to mask
  observations.
\item
  \emph{create\_obj} (bool): If True, returns the raw BF distribution
  object instead of creating a sample site. This is essential for
  building complex distributions like \texttt{MixtureSameFamily}.
\item
  \emph{sample} (bool, optional): A control-flow argument. If
  \texttt{True}, the function will directly sample a raw JAX array from
  the distribution, bypassing the BF model context. If \texttt{False},
  it will create a \texttt{BF.sample} site within a model. Defaults to
  \texttt{False}.
\item
  \emph{seed} (int, optional): An integer seed for this draw when
  \texttt{sample=True}. By default (\texttt{None}) the draw comes from
  the session's random stream, set by \texttt{bf(rand\_seed=...)}, so
  consecutive draws differ but the whole sequence is reproducible. No
  effect when \texttt{sample=False}.
\item
  \emph{obs} (jnp.ndarray, optional): The observed value for this random
  variable. If provided, the sample site is conditioned on this value,
  and the function returns the observed value. If \texttt{None}, the
  site is treated as a latent (unobserved) random variable. Defaults to
  \texttt{None}.
\item
  \emph{name} (str, optional): The name of the sample site in a BF
  model. This is used to uniquely identify the random variable. Defaults
  to `x'.
\end{itemize}

\textbf{Returns} BF LogUniform distribution object (for model building)
when \texttt{sample=False}.

\begin{verbatim}
JAX array of samples drawn from the LogUniform distribution (for direct sampling) when `sample=True`.

The raw BF distribution object (for advanced use cases) when `create_obj=True`.
\end{verbatim}

\textbf{Example Usage} from BayesForge import bf m = bf(platform=`cpu')
m.dist.log\_uniform(low=0.1, high=10.0, sample=True)

\textbf{Wrapper of}
https://num.pyro.ai/en/stable/distributions.html\#loguniform

\textbf{References}
https://en.wikipedia.org/wiki/Reciprocal\_distribution
https://docs.scipy.org/doc/scipy/reference/generated/scipy.stats.loguniform.html
https://docs.scipy.org/doc/scipy/tutorial/stats/continuous\_loguniform.html

\subsection{\texorpdfstring{\texttt{m.dist.logistic}}{m.dist.logistic}}\label{logistic}

\begin{Shaded}
\begin{Highlighting}[]
\NormalTok{m.dist.logistic(}
\NormalTok{    loc}\OperatorTok{=}\FloatTok{0.0}\NormalTok{,}
\NormalTok{    scale}\OperatorTok{=}\FloatTok{1.0}\NormalTok{,}
\NormalTok{    validate\_args}\OperatorTok{=}\VariableTok{None}\NormalTok{,}
\NormalTok{    name}\OperatorTok{=}\StringTok{\textquotesingle{}x\textquotesingle{}}\NormalTok{,}
\NormalTok{    obs}\OperatorTok{=}\VariableTok{None}\NormalTok{,}
\NormalTok{    mask}\OperatorTok{=}\VariableTok{None}\NormalTok{,}
\NormalTok{    infer}\OperatorTok{=}\VariableTok{None}\NormalTok{,}
\NormalTok{    sample}\OperatorTok{=}\VariableTok{False}\NormalTok{,}
\NormalTok{    seed}\OperatorTok{=}\VariableTok{None}\NormalTok{,}
\NormalTok{    shape}\OperatorTok{=}\NormalTok{(),}
\NormalTok{    event}\OperatorTok{=}\DecValTok{0}\NormalTok{,}
\NormalTok{    create\_obj}\OperatorTok{=}\VariableTok{False}\NormalTok{,}
\NormalTok{    to\_jax}\OperatorTok{=}\VariableTok{True}\NormalTok{,}
\NormalTok{)}
\end{Highlighting}
\end{Shaded}

\textbf{Logistic} Samples from a Logistic distribution.

The Logistic distribution is a continuous probability distribution
defined by two parameters: location and scale. It is often used to model
growth processes and is closely related to the normal distribution.Its
CDF is the logistic (sigmoid) function, which makes it appealing in
modeling probabilities, logistic regression, and various growth models.
It resembles the normal distribution in shape (bell‐shaped, symmetric)
but has \textbf{heavier tails} (i.e.~more probability in the extremes)
and simpler closed‐form expressions for the CDF.

\begin{itemize}
\tightlist
\item
  \textbf{Location parameter}: \(\mu \in \mathbb{R}\)~--- the center or
  ``mean'' of the distribution.
\item
  \textbf{Scale parameter}: \(\s > 0\)~--- controls the spread (similar
  role to standard deviation).
\item
  \textbf{Support}: \$~x \in (-\infty, +\infty) \$.
\end{itemize}

Let (X \sim \mathrm{Logistic}(\mu, s)). Then:

\[
F(x; \mu, s) = \frac{1}{1 + \exp\bigl(-\frac{x - \mu}{s}\bigr)}.
\]

Equivalently: \[
F(x) = \frac12 + \frac12 \tanh!\bigl(\frac{x - \mu}{2s}\bigr).
\]

\textbf{Args} - \emph{loc} (jnp.ndarray or float): The location
parameter, specifying the median of the distribution. Defaults to 0.0.

\begin{itemize}
\item
  \emph{scale} (jnp.ndarray or float): The scale parameter, which
  determines the spread of the distribution. Must be positive. Defaults
  to 1.0.
\item
  \emph{shape} (tuple): A multi-purpose argument for shaping. When
  \texttt{sample=False} (model building), this is used with
  \texttt{shape} to set the distribution's batch shape. When
  \texttt{sample=True} (direct sampling), this is used as
  \texttt{sample\_shape} to draw a raw JAX array of the given shape.
\item
  \emph{event} (int): The number of batch dimensions to reinterpret as
  event dimensions (used in model building).
\item
  \emph{mask} (jnp.ndarray, bool): Optional boolean array to mask
  observations.
\item
  \emph{create\_obj} (bool): If True, returns the raw BF distribution
  object instead of creating a sample site. This is essential for
  building complex distributions like \texttt{MixtureSameFamily}.
\item
  \emph{sample} (bool, optional): A control-flow argument. If
  \texttt{True}, the function will directly sample a raw JAX array from
  the distribution, bypassing the BF model context. If \texttt{False},
  it will create a \texttt{BF.sample} site within a model. Defaults to
  \texttt{False}.
\item
  \emph{seed} (int, optional): An integer seed for this draw when
  \texttt{sample=True}. By default (\texttt{None}) the draw comes from
  the session's random stream, set by \texttt{bf(rand\_seed=...)}, so
  consecutive draws differ but the whole sequence is reproducible. No
  effect when \texttt{sample=False}.
\item
  \emph{obs} (jnp.ndarray, optional): The observed value for this random
  variable. If provided, the sample site is conditioned on this value,
  and the function returns the observed value. If \texttt{None}, the
  site is treated as a latent (unobserved) random variable. Defaults to
  \texttt{None}.
\item
  \emph{name} (str, optional): The name of the sample site in a BF
  model. This is used to uniquely identify the random variable. Defaults
  to `x'.
\end{itemize}

\textbf{Returns} BF Logistic distribution object (for model building)
when \texttt{sample=False}. JAX array of samples drawn from the Logistic
distribution (for direct sampling) when \texttt{sample=True}. The raw BF
distribution object (for advanced use cases) when
\texttt{create\_obj=True}.

\textbf{Example Usage} from BayesForge import bf m = bf(platform=`cpu')
m.dist.logistic(loc=0.0, scale=1.0, sample=True)

\textbf{Wrapper of}
https://num.pyro.ai/en/stable/distributions.html\#logistic

\textbf{Reference} https://en.wikipedia.org/wiki/Logistic\_distribution

\subsection{\texorpdfstring{\texttt{m.dist.low\_rank\_multivariate\_normal}}{m.dist.low\_rank\_multivariate\_normal}}\label{low_rank_multivariate_normal}

\begin{Shaded}
\begin{Highlighting}[]
\NormalTok{m.dist.low\_rank\_multivariate\_normal(}
\NormalTok{    loc,}
\NormalTok{    cov\_factor,}
\NormalTok{    cov\_diag,}
\NormalTok{    validate\_args}\OperatorTok{=}\VariableTok{None}\NormalTok{,}
\NormalTok{    name}\OperatorTok{=}\StringTok{\textquotesingle{}x\textquotesingle{}}\NormalTok{,}
\NormalTok{    obs}\OperatorTok{=}\VariableTok{None}\NormalTok{,}
\NormalTok{    mask}\OperatorTok{=}\VariableTok{None}\NormalTok{,}
\NormalTok{    infer}\OperatorTok{=}\VariableTok{None}\NormalTok{,}
\NormalTok{    sample}\OperatorTok{=}\VariableTok{False}\NormalTok{,}
\NormalTok{    seed}\OperatorTok{=}\VariableTok{None}\NormalTok{,}
\NormalTok{    shape}\OperatorTok{=}\NormalTok{(),}
\NormalTok{    event}\OperatorTok{=}\DecValTok{0}\NormalTok{,}
\NormalTok{    create\_obj}\OperatorTok{=}\VariableTok{False}\NormalTok{,}
\NormalTok{    to\_jax}\OperatorTok{=}\VariableTok{True}\NormalTok{,}
\NormalTok{)}
\end{Highlighting}
\end{Shaded}

\textbf{Low Rank Multivariate Normal} The \emph{Low-Rank Multivariate
Normal} (LRMVN) distribution is a parameterization of the multivariate
normal distribution where the covariance matrix is expressed as a
low-rank plus diagonal decomposition: \[
\Sigma = F F^\top + D
\] where \(F\) is a low-rank matrix (capturing correlations) and \(D\)
is a diagonal matrix (capturing independent noise). This representation
is often used in probabilistic modeling and variational inference to
efficiently handle high-dimensional Gaussian distributions with
structured covariance.

For a (d)-dimensional random vector ( x \in \mathbb{R}\^{}d ), \[
p(x \mid \mu, F, D) = \frac{1}{(2\pi)^{d/2} |\Sigma|^{1/2}}
\exp\left(-\frac{1}{2}(x - \mu)^\top \Sigma^{-1} (x - \mu)\right),
\] where ( \Sigma = F F\^{}\top + D ).

Parameters: - \emph{loc} (jnp.ndarray): Mean vector.

\begin{verbatim}
cov_factor (jnp.ndarray): Matrix used to construct the covariance matrix.

cov_diag (jnp.ndarray): Diagonal elements of the covariance matrix.
    
\end{verbatim}

\begin{itemize}
\item
  \emph{sample} (bool, optional): A control-flow argument. If
  \texttt{True}, the function will directly sample a raw JAX array from
  the distribution, bypassing the BF model context. If \texttt{False},
  it will create a \texttt{BF.sample} site within a model. Defaults to
  \texttt{False}.
\item
  \emph{shape} (tuple): A multi-purpose argument for shaping. When
  \texttt{sample=False} (model building), this is used with
  \texttt{shape} to set the distribution's batch shape. When
  \texttt{sample=True} (direct sampling), this is used as
  \texttt{sample\_shape} to draw a raw JAX array of the given shape.
\item
  \emph{seed} (int, optional): An integer seed for this draw when
  \texttt{sample=True}. By default (\texttt{None}) the draw comes from
  the session's random stream, set by \texttt{bf(rand\_seed=...)}, so
  consecutive draws differ but the whole sequence is reproducible. No
  effect when \texttt{sample=False}.
\item
  \emph{obs} (jnp.ndarray, optional): The observed value for this random
  variable. If provided, the sample site is conditioned on this value,
  and the function returns the observed value. If \texttt{None}, the
  site is treated as a latent (unobserved) random variable. Defaults to
  \texttt{None}.
\item
  \emph{name} (str, optional): The name of the sample site in a BF
  model. This is used to uniquely identify the random variable. Defaults
  to `x'.
\end{itemize}

\textbf{Example Usage} from BayesForge import bf import jax.numpy as jnp
m = bf(platform=`cpu') event\_size, rank = 100, 5
m.dist.low\_rank\_multivariate\_normal( loc=m.dist.normal(0, 1,
shape=(event\_size,), sample=True) * 2, cov\_factor=m.dist.normal(0, 1,
shape=(event\_size, rank), sample=True),
cov\_diag=jnp.exp(m.dist.normal(0, 1, shape=(event\_size,),
sample=True)) * 0.1, sample=True, )

\textbf{Wrapper of}
https://num.pyro.ai/en/stable/distributions.html\#lowrankmultivariatenormal

\textbf{Reference} *
\href{https://en.wikipedia.org/wiki/Multivariate_normal_distribution}{Multivariate
normal distribution -- Wikipedia}

\begin{itemize}
\tightlist
\item
  \href{https://www.tensorflow.org/probability/api_docs/python/tfp/distributions/MultivariateNormalDiagPlusLowRank}{TensorFlow
  Probability: LowRankMultivariateNormal}
\end{itemize}

\subsection{\texorpdfstring{\texttt{m.dist.lower\_truncated\_power\_law}}{m.dist.lower\_truncated\_power\_law}}\label{lower_truncated_power_law}

\begin{Shaded}
\begin{Highlighting}[]
\NormalTok{m.dist.lower\_truncated\_power\_law(}
\NormalTok{    alpha,}
\NormalTok{    low,}
\NormalTok{    validate\_args}\OperatorTok{=}\VariableTok{None}\NormalTok{,}
\NormalTok{    name}\OperatorTok{=}\StringTok{\textquotesingle{}x\textquotesingle{}}\NormalTok{,}
\NormalTok{    obs}\OperatorTok{=}\VariableTok{None}\NormalTok{,}
\NormalTok{    mask}\OperatorTok{=}\VariableTok{None}\NormalTok{,}
\NormalTok{    infer}\OperatorTok{=}\VariableTok{None}\NormalTok{,}
\NormalTok{    sample}\OperatorTok{=}\VariableTok{False}\NormalTok{,}
\NormalTok{    seed}\OperatorTok{=}\VariableTok{None}\NormalTok{,}
\NormalTok{    shape}\OperatorTok{=}\NormalTok{(),}
\NormalTok{    event}\OperatorTok{=}\DecValTok{0}\NormalTok{,}
\NormalTok{    create\_obj}\OperatorTok{=}\VariableTok{False}\NormalTok{,}
\NormalTok{    to\_jax}\OperatorTok{=}\VariableTok{True}\NormalTok{,}
\NormalTok{)}
\end{Highlighting}
\end{Shaded}

\textbf{Lower Truncated Power Law} The \emph{Lower-Truncated Power-Law}
distribution (also known as the \emph{Pareto Type I} or \emph{power-law
with a lower bound}) models quantities that follow a heavy-tailed
power-law behavior but are bounded below by a minimum value (
x\_\{\min\} ). It is commonly used to describe phenomena such as wealth
distributions, city sizes, and biological scaling laws.

The probability density function (PDF) is given by:

For ( x \ge x\_\{\min\} \textgreater{} 0 ): \[
  f(x \mid \alpha, x_{\min}) = \frac{\alpha - 1}{x_{\min}} \left( \frac{x_{\min}}{x} \right)^{\alpha}
  \]

where \(\alpha > 1\) is the shape (scaling) parameter.

\textbf{Args} - \texttt{alpha} (jnp.ndarray): index of the power law
distribution. Must be less than -1.

\begin{itemize}
\item
  \texttt{low} (jnp.ndarray): lower bound of the distribution. Must be
  greater than 0.
\item
  \texttt{shape} (tuple): A multi-purpose argument for shaping. When
  \texttt{sample=False} (model building), this is used with
  \texttt{shape} to set the distribution's batch shape. When
  \texttt{sample=True} (direct sampling), this is used as
  \texttt{sample\_shape} to draw a raw JAX array of the given shape.
\item
  \texttt{event} (int): The number of batch dimensions to reinterpret as
  event dimensions (used in model building).
\item
  \texttt{mask} (jnp.ndarray, bool): Optional boolean array to mask
  observations.
\item
  \texttt{create\_obj} (bool): If True, returns the raw BF distribution
  object instead of creating a sample site.This is essential for
  building complex distributions like \texttt{MixtureSameFamily}.
\item
  \texttt{sample} (bool, optional): A control-flow argument. If
  \texttt{True}, the function will directly sample a raw JAX array from
  the distribution, bypassing the BF model context. If \texttt{False},
  it will create a \texttt{BF.sample} site within a model. Defaults to
  \texttt{False}.
\item
  \texttt{seed} (int, optional): An integer seed for this draw when
  \texttt{sample=True}. By default (\texttt{None}) the draw comes from
  the session's random stream, set by \texttt{bf(rand\_seed=...)}, so
  consecutive draws differ but the whole sequence is reproducible. No
  effect when \texttt{sample=False}.
\item
  \texttt{obs} (jnp.ndarray, optional): The observed value for this
  random variable. If provided, the sample site is conditioned on this
  value, and the function returns the observed value. If \texttt{None},
  the site is treated as a latent (unobserved) random variable. Defaults
  to \texttt{None}.
\item
  \texttt{name} (str, optional): The name of the sample site in a BF
  model. This is used to uniquely identify the random variable. Defaults
  to `x'.
\end{itemize}

\textbf{Returns} - When \texttt{sample=False}: A BF
LowerTruncatedPowerLaw distribution object (for model building). - When
\texttt{sample=True}: A JAX array of samples drawn from the
LowerTruncatedPowerLaw distribution (for direct sampling). - When
\texttt{create\_obj=True}: The raw BF distribution object (for advanced
use cases).

\textbf{Example Usage} from BayesForge import bf m = bf(platform=`cpu')
m.dist.lower\_truncated\_power\_law(alpha=-2.0, low=1.0, sample=True)

\textbf{Wrapper of}
https://num.pyro.ai/en/stable/distributions.html\#lowertruncatedpowerlaw

\textbf{Reference} *
\href{https://en.wikipedia.org/wiki/Power_law\#Continuous_power-law_distribution}{Power-law
distribution -- Wikipedia}

\begin{itemize}
\item
  \href{https://en.wikipedia.org/wiki/Pareto_distribution}{Pareto
  distribution -- Wikipedia}
\item
  \href{https://arxiv.org/abs/0706.1062}{Truncated Pareto distribution}
\end{itemize}

\subsection{\texorpdfstring{\texttt{m.dist.mask}}{m.dist.mask}}\label{mask}

\begin{Shaded}
\begin{Highlighting}[]
\NormalTok{m.dist.mask(mask)}
\end{Highlighting}
\end{Shaded}

Context manager that excludes observations from the likelihood.

Inside \texttt{with\ m.dist.mask(keep):}, the log-probability of every
sample site is multiplied by the boolean array \texttt{keep}, so entries
where it is False (e.g.~missing data) do not contribute.

\begin{longtable}[]{@{}lll@{}}
\toprule\noalign{}
Argument & Type & Description \\
\midrule\noalign{}
\endhead
\bottomrule\noalign{}
\endlastfoot
\texttt{mask} & jax.Array of bool & True for the observations to
keep. \\
\end{longtable}

\textbf{Returns}

\begin{itemize}
\tightlist
\item
  A NumPyro \texttt{mask} handler, used in a \texttt{with} statement.
\end{itemize}

\subsection{\texorpdfstring{\texttt{m.dist.matrix\_normal}}{m.dist.matrix\_normal}}\label{matrix_normal}

\begin{Shaded}
\begin{Highlighting}[]
\NormalTok{m.dist.matrix\_normal(}
\NormalTok{    loc,}
\NormalTok{    scale\_tril\_row,}
\NormalTok{    scale\_tril\_column,}
\NormalTok{    validate\_args}\OperatorTok{=}\VariableTok{None}\NormalTok{,}
\NormalTok{    name}\OperatorTok{=}\StringTok{\textquotesingle{}x\textquotesingle{}}\NormalTok{,}
\NormalTok{    obs}\OperatorTok{=}\VariableTok{None}\NormalTok{,}
\NormalTok{    mask}\OperatorTok{=}\VariableTok{None}\NormalTok{,}
\NormalTok{    infer}\OperatorTok{=}\VariableTok{None}\NormalTok{,}
\NormalTok{    sample}\OperatorTok{=}\VariableTok{False}\NormalTok{,}
\NormalTok{    seed}\OperatorTok{=}\VariableTok{None}\NormalTok{,}
\NormalTok{    shape}\OperatorTok{=}\NormalTok{(),}
\NormalTok{    event}\OperatorTok{=}\DecValTok{0}\NormalTok{,}
\NormalTok{    create\_obj}\OperatorTok{=}\VariableTok{False}\NormalTok{,}
\NormalTok{    to\_jax}\OperatorTok{=}\VariableTok{True}\NormalTok{,}
\NormalTok{)}
\end{Highlighting}
\end{Shaded}

\textbf{Matrix Normal} Samples from a Matrix Normal distribution, which
is a multivariate normal distribution over matrices. The distribution is
characterized by a location matrix and two lower triangular matrices
that define the correlation structure. The distribution is related to
the multivariate normal distribution in the following way. If
\(X ~ MN(loc,U,V)\) then \(vec(X) ~ MVN(vec(loc), kron(V,U) )\).

Let (X) be (n \times p), with mean matrix \[M\in\mathbb R^{n\times p}\],
row-covariance matrix \(U\in\mathbb R^{n\times n}\) (positive-definite),
and column-covariance matrix
\(V\in\mathbb R^{p\times p}\)(positive-definite). Then the PDF is:

\[
  f(X;M,U,V) = \frac{1}{(2\pi)^{\tfrac{np}{2}},|V|^{\tfrac n2},|U|^{\tfrac p2}},
  \exp!\Big[-\tfrac12,\mathrm{tr}!\big( V^{-1}(X-M)^\top,U^{-1}(X-M)\big)\Big].
  \]

Also the equivalence to vec-form is given: \[
  \mathrm{vec}(X) \sim \mathcal N_{np}\big(\mathrm{vec}(M),,V\otimes U\big). ; ; \text{([The Book of Statistical Proofs][3])}
  \]

\textbf{Args} - \texttt{loc} (array\_like): Location of the
distribution.

\begin{itemize}
\item
  \texttt{scale\_tril\_row} (array\_like): Lower cholesky of rows
  correlation matrix.
\item
  \texttt{scale\_tril\_column} (array\_like): Lower cholesky of columns
  correlation matrix.
\item
  \texttt{shape} (tuple): A multi-purpose argument for shaping. When
  \texttt{sample=False} (model building), this is used with
  \texttt{shape} to set the distribution's batch shape. When
  \texttt{sample=True} (direct sampling), this is used as
  \texttt{sample\_shape} to draw a raw JAX array of the given shape.
\item
  \texttt{event} (int): The number of batch dimensions to reinterpret as
  event dimensions (used in model building).
\item
  \texttt{mask} (jnp.ndarray, bool): Optional boolean array to mask
  observations.
\item
  \texttt{create\_obj} (bool): If True, returns the raw BF distribution
  object instead of creating a sample site.This is essential for
  building complex distributions like \texttt{MixtureSameFamily}.
\item
  \texttt{sample} (bool, optional): A control-flow argument. If
  \texttt{True}, the function will directly sample a raw JAX array from
  the distribution, bypassing the BF model context. If \texttt{False},
  it will create a \texttt{BF.sample} site within a model. Defaults to
  \texttt{False}.
\item
  \texttt{seed} (int, optional): An integer seed for this draw when
  \texttt{sample=True}. By default (\texttt{None}) the draw comes from
  the session's random stream, set by \texttt{bf(rand\_seed=...)}, so
  consecutive draws differ but the whole sequence is reproducible. No
  effect when \texttt{sample=False}.
\item
  \texttt{obs} (jnp.ndarray, optional): The observed value for this
  random variable. If provided, the sample site is conditioned on this
  value, and the function returns the observed value. If \texttt{None},
  the site is treated as a latent (unobserved) random variable. Defaults
  to \texttt{None}.
\item
  \texttt{name} (str, optional): The name of the sample site in a BF
  model. This is used to uniquely identify the random variable. Defaults
  to `x'.
\end{itemize}

\textbf{Returns} - When \texttt{sample=False}: A BF MatrixNormal
distribution object (for model building).

\begin{itemize}
\item
  When \texttt{sample=True}: A JAX array of samples drawn from the
  MatrixNormal distribution (for direct sampling).
\item
  When \texttt{create\_obj=True}: The raw BF distribution object (for
  advanced use cases).

  \textbf{Example Usage} from BayesForge import bf import jax.numpy as
  jnp m = bf(platform=`cpu') n\_rows, n\_cols = 3, 4 loc =
  jnp.zeros((n\_rows, n\_cols)) U\_row\_cov = jnp.array({[}{[}1.0, 0.5,
  0.2{]}, {[}0.5, 1.0, 0.3{]}, {[}0.2, 0.3, 1.0{]}{]}) V\_col\_cov =
  jnp.array({[}{[}2.0, -0.8, 0.1, 0.4{]}, {[}-0.8, 2.0, 0.2, -0.2{]},
  {[}0.1, 0.2, 2.0, 0.0{]}, {[}0.4, -0.2, 0.0, 2.0{]}{]}) \textbf{The
  distribution takes the Cholesky factors of both covariances}
  m.dist.matrix\_normal( loc=loc,
  scale\_tril\_row=jnp.linalg.cholesky(U\_row\_cov),
  scale\_tril\_column=jnp.linalg.cholesky(V\_col\_cov), sample=True, )

  \textbf{Wrapper of:
  https://num.pyro.ai/en/stable/distributions.html\#matrixnormal\_lowercase}
  \textbf{Reference}

  \begin{itemize}
  \tightlist
  \item
  \item
  \end{itemize}
\end{itemize}

\subsection{\texorpdfstring{\texttt{m.dist.mixture}}{m.dist.mixture}}\label{mixture}

\begin{Shaded}
\begin{Highlighting}[]
\NormalTok{m.dist.mixture(}
\NormalTok{    mixing\_distribution,}
\NormalTok{    component\_distributions,}
\NormalTok{    validate\_args}\OperatorTok{=}\VariableTok{None}\NormalTok{,}
\NormalTok{    name}\OperatorTok{=}\StringTok{\textquotesingle{}x\textquotesingle{}}\NormalTok{,}
\NormalTok{    obs}\OperatorTok{=}\VariableTok{None}\NormalTok{,}
\NormalTok{    mask}\OperatorTok{=}\VariableTok{None}\NormalTok{,}
\NormalTok{    infer}\OperatorTok{=}\VariableTok{None}\NormalTok{,}
\NormalTok{    sample}\OperatorTok{=}\VariableTok{False}\NormalTok{,}
\NormalTok{    seed}\OperatorTok{=}\VariableTok{None}\NormalTok{,}
\NormalTok{    shape}\OperatorTok{=}\NormalTok{(),}
\NormalTok{    event}\OperatorTok{=}\DecValTok{0}\NormalTok{,}
\NormalTok{    create\_obj}\OperatorTok{=}\VariableTok{False}\NormalTok{,}
\NormalTok{    to\_jax}\OperatorTok{=}\VariableTok{True}\NormalTok{,}
\NormalTok{)}
\end{Highlighting}
\end{Shaded}

\textbf{A marginalized finite mixture of component distributions.} A
\textbf{marginalised finite mixture of component distributions} refers
to a probability model in which you have a finite number \(K\) of
component distributions (e.g., normals, Poissons, etc.), each weighted
by a mixing probability from a Categorical distribution. We
\textbf{marginalise out} the latent assignment variable (i.e., you don't
explicitly model which component each data point came from). The
resulting distribution can be either a \texttt{mixture\_general} (when
component distributions are a list) or a \texttt{mixture\_same\_Family}
(when component distributions are a single distribution).

Suppose you have:

\begin{itemize}
\tightlist
\item
  Mixing weights \(\pi_1, \dots, \pi_K\), with \(\pi_j \ge 0\) and
  \(\sum_{j=1}^K \pi_j = 1\).
\item
  Component distributions (densities or PMFs) \(f_j(x \mid \theta_j)\),
  for \(j = 1,\dots,K\).
\end{itemize}

Then the marginalised mixture distribution for an observed \(X\) is:

\[
  p(X = x) = \sum_{j=1}^K \pi_j ; f_j(x \mid \theta_j).
  \]

This is exactly the mixture density/PMF you get when you marginalise the
latent component‑assignment variable \(Z \in {1,\dots,K}\) out of the
joint distribution:

\[
  p(x) = \sum_{z=1}^K p(z = j); p(x \mid z = j) = \sum_{j=1}^K \pi_j ; f_j(x \mid \theta_j).
  \]

\textbf{Args} - \emph{mixing\_distribution} (BF.distribution): The
distribution used to determine the mixing weights.

\begin{itemize}
\item
  \emph{component\_distributions} (list{[}BF.distribution{]}): The list
  of component distributions.
\item
  \emph{shape} (tuple): A multi-purpose argument for shaping. When
  \texttt{sample=False} (model building), this is used with
  \texttt{shape} to set the distribution's batch shape. When
  \texttt{sample=True} (direct sampling), this is used as
  \texttt{sample\_shape} to draw a raw JAX array of the given shape.
\item
  \emph{event} (int): The number of batch dimensions to reinterpret as
  event dimensions (used in model building).
\item
  \emph{mask} (jnp.ndarray, bool): Optional boolean array to mask
  observations.
\item
  \emph{create\_obj} (bool): If True, returns the raw BF distribution
  object instead of creating a sample site. This is essential for
  building complex distributions like \texttt{MixtureSameFamily}.
\item
  \emph{sample} (bool, optional): A control-flow argument. If
  \texttt{True}, the function will directly sample a raw JAX array from
  the distribution, bypassing the BF model context. If \texttt{False},
  it will create a \texttt{BF.sample} site within a model. Defaults to
  \texttt{False}.
\item
  \emph{seed} (int, optional): An integer seed for this draw when
  \texttt{sample=True}. By default (\texttt{None}) the draw comes from
  the session's random stream, set by \texttt{bf(rand\_seed=...)}, so
  consecutive draws differ but the whole sequence is reproducible. No
  effect when \texttt{sample=False}.
\item
  \emph{obs} (jnp.ndarray, optional): The observed value for this random
  variable. If provided, the sample site is conditioned on this value,
  and the function returns the observed value. If \texttt{None}, the
  site is treated as a latent (unobserved) random variable. Defaults to
  \texttt{None}.
\item
  \emph{name} (str, optional): The name of the sample site in a BF
  model. This is used to uniquely identify the random variable. Defaults
  to `x'.
\end{itemize}

\textbf{Returns} - When \texttt{sample=False}: A BF Mixture distribution
object (for model building).

\begin{itemize}
\item
  When \texttt{sample=True}: A JAX array of samples drawn from the
  Mixture distribution (for direct sampling).
\item
  When \texttt{create\_obj=True}: The raw BF distribution object (for
  advanced use cases).

  \textbf{Example Usage} from BayesForge import bf import jax.numpy as
  jnp m = bf(platform=`cpu') m.dist.mixture(
  mixing\_distribution=m.dist.categorical(probs=jnp.array({[}0.3,
  0.7{]}), create\_obj=True),
  component\_distributions={[}m.dist.normal(0, 1, create\_obj=True),
  m.dist.normal(2, 1, create\_obj=True){]}, sample=True, shape=(10,), )

  \textbf{Wrapper of:
  https://num.pyro.ai/en/stable/distributions.html\#mixture}
  \textbf{References}

  \begin{itemize}
  \tightlist
  \item
    https://en.wikipedia.org/wiki/Mixture\_distribution ``Mixture
    distribution''
  \item
    https://pymc3-testing.readthedocs.io/en/rtd-docs/notebooks/marginalized\_gaussian\_mixture\_model.html
    ``Marginalized Gaussian Mixture Model --- PyMC3 3.1rc3
    documentation''
  \end{itemize}
\end{itemize}

\subsection{\texorpdfstring{\texttt{m.dist.mixture\_general}}{m.dist.mixture\_general}}\label{mixture_general}

\begin{Shaded}
\begin{Highlighting}[]
\NormalTok{m.dist.mixture\_general(}
\NormalTok{    mixing\_distribution,}
\NormalTok{    component\_distributions,}
\NormalTok{    support}\OperatorTok{=}\VariableTok{None}\NormalTok{,}
\NormalTok{    validate\_args}\OperatorTok{=}\VariableTok{None}\NormalTok{,}
\NormalTok{    name}\OperatorTok{=}\StringTok{\textquotesingle{}x\textquotesingle{}}\NormalTok{,}
\NormalTok{    obs}\OperatorTok{=}\VariableTok{None}\NormalTok{,}
\NormalTok{    mask}\OperatorTok{=}\VariableTok{None}\NormalTok{,}
\NormalTok{    infer}\OperatorTok{=}\VariableTok{None}\NormalTok{,}
\NormalTok{    sample}\OperatorTok{=}\VariableTok{False}\NormalTok{,}
\NormalTok{    seed}\OperatorTok{=}\VariableTok{None}\NormalTok{,}
\NormalTok{    shape}\OperatorTok{=}\NormalTok{(),}
\NormalTok{    event}\OperatorTok{=}\DecValTok{0}\NormalTok{,}
\NormalTok{    create\_obj}\OperatorTok{=}\VariableTok{False}\NormalTok{,}
\NormalTok{    to\_jax}\OperatorTok{=}\VariableTok{True}\NormalTok{,}
\NormalTok{)}
\end{Highlighting}
\end{Shaded}

\textbf{Mixture General} A mixture distribution is a probability
distribution constructed by selecting one of several component
distributions according to specified weights, and then drawing a sample
from the chosen component. It allows modelling of heterogeneous
populations and multimodal data.

Let the mixture have (K) component distributions with densities
(f\_i(x)) for (i = 1,\dots,K). Let weights (w\_i\ge0) satisfy
(\sum\_\{i=1\}\^{}K w\_i = 1). Then the PDF is \[
f(x) = \sum_{i=1}^K w_i ; f_i(x).
\]

\textbf{Parameters:}

\begin{itemize}
\item
  \textbf{mixing\_distribution}: A \texttt{Categorical} distribution
  representing the mixing weights.
\item
  \textbf{component\_distributions}: A list of distributions
  representing the components of the mixture.
\item
  **sample (bool, optional): A control-flow argument. If \texttt{True},
  the function will directly sample a raw JAX array from the
  distribution, bypassing the BF model context. If \texttt{False}, it
  will create a \texttt{BF.sample} site within a model. Defaults to
  \texttt{False}.
\item
  \texttt{shape} (tuple): A multi-purpose argument for shaping. When
  \texttt{sample=False} (model building), this is used with
  \texttt{shape} to set the distribution's batch shape. When
  \texttt{sample=True} (direct sampling), this is used as
  \texttt{sample\_shape} to draw a raw JAX array of the given shape.
\item
  \textbf{seed (int, optional)}: An integer seed for this draw when
  \texttt{sample=True}. By default (\texttt{None}) the draw comes from
  the session's random stream, set by \texttt{bf(rand\_seed=...)}, so
  consecutive draws differ but the whole sequence is reproducible. No
  effect when \texttt{sample=False}.
\item
  \textbf{obs (jnp.ndarray, optional)}: The observed value for this
  random variable. If provided, the sample site is conditioned on this
  value, and the function returns the observed value. If \texttt{None},
  the site is treated as a latent (unobserved) random variable. Defaults
  to \texttt{None}.
\item
  \textbf{name (str, optional)}: The name of the sample site in a BF
  model. This is used to uniquely identify the random variable. Defaults
  to `x'.
\end{itemize}

\textbf{\#\#\#\# Returns:}

\begin{itemize}
\tightlist
\item
  \textbf{When \texttt{sample=False}}: A BF MixtureGeneral distribution
  object (for model building).
\item
  \textbf{When \texttt{sample=True}}: A JAX array of samples drawn from
  the MixtureGeneral distribution (for direct sampling).
\item
  \textbf{When \texttt{create\_obj=True}}: The raw BF distribution
  object (for advanced use cases).
\end{itemize}

\textbf{\#\#\#\# Example Usage:}

\begin{verbatim}
from BayesForge import bf
m = bf(platform='cpu')
m.dist.mixture_general(
    mixing_distribution=m.dist.categorical(probs=jnp.array([0.3, 0.7]), create_obj = True),     
    component_distributions=[m.dist.normal(loc=0.0, scale=1.0, create_obj=True),m.dist.normal  (loc=0.0, scale=1.0,  create_obj=True)],
    sample = True
)           
\end{verbatim}

\textbf{Wrapper of:
\url{https://num.pyro.ai/en/stable/distributions.html\#mixturegeneral}}
**Reference: *\emph{ }
\href{https://en.wikipedia.org/wiki/Matrix_normal_distribution}{1}:
https://preliz.readthedocs.io/en/latest/distributions/gallery/mixture.html
``Mixture Distribution --- PreliZ 0.22.0 documentation'' *
\href{https://statproofbook.github.io/P/matn-pdf.html}{2}:
https://reliability.readthedocs.io/en/latest/Mixture\%20models.html
``Mixture models --- reliability 0.9.0 documentation'' * {[}3{]}:
https://docs.scipy.org/doc/scipy/reference/generated/scipy.stats.Mixture.html
``Mixture --- SciPy v1.16.2 Manual''

\subsection{\texorpdfstring{\texttt{m.dist.mixture\_same\_family}}{m.dist.mixture\_same\_family}}\label{mixture_same_family}

\begin{Shaded}
\begin{Highlighting}[]
\NormalTok{m.dist.mixture\_same\_family(}
\NormalTok{    mixing\_distribution,}
\NormalTok{    component\_distribution,}
\NormalTok{    validate\_args}\OperatorTok{=}\VariableTok{None}\NormalTok{,}
\NormalTok{    name}\OperatorTok{=}\StringTok{\textquotesingle{}x\textquotesingle{}}\NormalTok{,}
\NormalTok{    obs}\OperatorTok{=}\VariableTok{None}\NormalTok{,}
\NormalTok{    mask}\OperatorTok{=}\VariableTok{None}\NormalTok{,}
\NormalTok{    infer}\OperatorTok{=}\VariableTok{None}\NormalTok{,}
\NormalTok{    sample}\OperatorTok{=}\VariableTok{False}\NormalTok{,}
\NormalTok{    seed}\OperatorTok{=}\VariableTok{None}\NormalTok{,}
\NormalTok{    shape}\OperatorTok{=}\NormalTok{(),}
\NormalTok{    event}\OperatorTok{=}\DecValTok{0}\NormalTok{,}
\NormalTok{    create\_obj}\OperatorTok{=}\VariableTok{False}\NormalTok{,}
\NormalTok{    to\_jax}\OperatorTok{=}\VariableTok{True}\NormalTok{,}
\NormalTok{)}
\end{Highlighting}
\end{Shaded}

\textbf{Finite mixture of component distributions from the same family.}
A \emph{Mixture (Same-Family)} distribution is a finite mixture in which
\textbf{all components come from the same parametric family} (for
example, all Normal distributions but with different parameters), and
are combined via mixing weights. The class is typically denoted as:

Let the mixture have (K) components. Let weights (w\_i\ge0),
(\sum\_\{i=1\}\^{}K w\_i = 1). Let the component family have density
(f(x \mid \theta\_i)) for each component (i). Then the mixture's PDF is
\[
f_X(x) = \sum_{i=1}^K w_i ; f(x \mid \theta_i).
\]

where each \(f(x \mid \theta_i)\) is from the same family with parameter
\(\theta_i\).

\textbf{Args} * \textbf{Distribution Args}: *
\textbf{mixing\_distribution}: A \texttt{Categorical} distribution
representing the mixing weights.

\begin{verbatim}
    *   **component_distributions**: A list of distributions representing the components of the mixture.

*   **Sampling / Modeling Args**:
    * `shape` (tuple): A multi-purpose argument for shaping. When `sample=False` (model building), this is used with `shape` to set the distribution's batch shape. When `sample=True` (direct sampling), this is used as `sample_shape` to draw a raw JAX array of the given shape.

    *   `event` (int): The number of batch dimensions to reinterpret as event dimensions (used in model building).
    *   `mask` (jnp.ndarray, bool): Optional boolean array to mask observations.

    *   `create_obj` (bool): If True, returns the raw BF distribution object instead of creating a sample
        site. This is essential for building complex distributions like `MixtureSameFamily`.

    *   **sample (bool, optional): A control-flow argument. If `True`, the function will directly sample a raw JAX array from the distribution, bypassing the BF model context. If `False`, it will create a `BF.sample` site within a model. Defaults to `False`.

    *   **seed (int, optional)**:  An integer seed for this draw when `sample=True`. By default (`None`) the draw comes from the session's random stream, set by `bf(rand_seed=...)`, so consecutive draws differ but the whole sequence is reproducible. No effect when `sample=False`.

    *   **obs (jnp.ndarray, optional)**:  The observed value for this random variable. If provided, the sample site is conditioned on this value, and the function returns the observed value. If `None`, the site is treated as a latent (unobserved) random variable. Defaults to `None`.

    *   **name (str, optional)**:  The name of the sample site in a BF model. This is used to uniquely identify the random variable. Defaults to 'x'.
\end{verbatim}

\textbf{Returns} * \textbf{When \texttt{sample=False}}: A BF
MixtureSameFamily distribution object (for model building). *
\textbf{When \texttt{sample=True}}: A JAX array of samples drawn from
the MixtureSameFamily distribution (for direct sampling). * \textbf{When
\texttt{create\_obj=True}}: The raw BF distribution object (for advanced
use cases).

\textbf{Example Usage} from BayesForge import bf import jax.numpy as jnp
m = bf(platform=`cpu') m.dist.mixture\_same\_family(
mixing\_distribution=m.dist.categorical(probs=jnp.array({[}0.3, 0.7{]}),
create\_obj = True), component\_distribution=m.dist.normal(loc=0.0,
scale=1.0, shape = (2,), create\_obj=True), sample = True )

\textbf{Wrapper of:
https://num.pyro.ai/en/stable/distributions.html\#mixture-, to\_jax =
True)same-family} \textbf{Reference} *
\href{https://cseweb.ucsd.edu/~elkan/250Bwinter2011/mixturemodels.pdf}{Mixture
distributions}

\subsection{\texorpdfstring{\texttt{m.dist.multinomial}}{m.dist.multinomial}}\label{multinomial}

\begin{Shaded}
\begin{Highlighting}[]
\NormalTok{m.dist.multinomial(}
\NormalTok{    total\_count}\OperatorTok{=}\DecValTok{1}\NormalTok{,}
\NormalTok{    probs}\OperatorTok{=}\VariableTok{None}\NormalTok{,}
\NormalTok{    logits}\OperatorTok{=}\VariableTok{None}\NormalTok{,}
\NormalTok{    total\_count\_max}\OperatorTok{=}\VariableTok{None}\NormalTok{,}
\NormalTok{    validate\_args}\OperatorTok{=}\VariableTok{None}\NormalTok{,}
\NormalTok{    name}\OperatorTok{=}\StringTok{\textquotesingle{}x\textquotesingle{}}\NormalTok{,}
\NormalTok{    obs}\OperatorTok{=}\VariableTok{None}\NormalTok{,}
\NormalTok{    mask}\OperatorTok{=}\VariableTok{None}\NormalTok{,}
\NormalTok{    infer}\OperatorTok{=}\VariableTok{None}\NormalTok{,}
\NormalTok{    sample}\OperatorTok{=}\VariableTok{False}\NormalTok{,}
\NormalTok{    seed}\OperatorTok{=}\VariableTok{None}\NormalTok{,}
\NormalTok{    shape}\OperatorTok{=}\NormalTok{(),}
\NormalTok{    event}\OperatorTok{=}\DecValTok{0}\NormalTok{,}
\NormalTok{    create\_obj}\OperatorTok{=}\VariableTok{False}\NormalTok{,}
\NormalTok{    to\_jax}\OperatorTok{=}\VariableTok{True}\NormalTok{,}
\NormalTok{)}
\end{Highlighting}
\end{Shaded}

\textbf{Multinomial} The multinomial distribution is a discrete
probability distribution that generalizes the Binomial distribution to
the case of \textbf{more than two possible outcomes per trial}.
Specifically:

\begin{itemize}
\tightlist
\item
  You perform a fixed number \(n\) of independent trials.
\item
  On each trial there are \(k\) possible outcomes, labeled
  (1,2,\dots,k).
\item
  The probability of outcome \(i\) on any trial is (p\_i), where (p\_1 +
  p\_2 + \cdots + p\_k = 1).
\item
  Let \(X_i\) be the number of trials (out of \(n\)) that produced
  outcome \(i\). Then the vector \(\mathbf X = (X_1, \dots, X_k)\)
  follows a multinomial distribution.
\end{itemize}

It reduces to the binomial when (k=2).

PMF (Probability Mass Function):

For \(\mathbf x = (x_1,\dots,x_k)\) with each \(x_i \ge 0\) integer and
\(\sum_{i=1}^k x_i = n\), \[
  \Pr(X_1 = x_1, \dots, X_k = x_k)
  = \frac{n!}{x_1!,x_2!\cdots x_k!}
  ; p_1^{,x_1}; p_2^{,x_2} \cdots p_k^{,x_k}.
  \] If \(\sum_i x_i \neq n\), the probability is 0.

\textbf{Args} - \emph{total\_count} (int or jnp.ndarray): The number of
trials.

\begin{itemize}
\item
  \emph{probs} (jnp.ndarray, optional): Event probabilities. Must sum to
  1.
\item
  \emph{logits} (jnp.ndarray, optional): Event log probabilities.
\item
  *total\_count\_max (int, optional): An optional integer providing an
  upper bound on \texttt{total\_count}. This is used for performance
  optimization with \texttt{lax.scan} when \texttt{total\_count} is a
  dynamic JAX tracer, helping to avoid recompilation.
\item
  \emph{shape} (tuple): A multi-purpose argument for shaping. When
  \texttt{sample=False} (model building), this is used with
  \texttt{shape} to set the distribution's batch shape. When
  \texttt{sample=True} (direct sampling), this is used as
  \texttt{sample\_shape} to draw a raw JAX array of the given shape.
\item
  \emph{event} (int): The number of batch dimensions to reinterpret as
  event dimensions (used in model building).
\item
  \emph{mask} (jnp.ndarray, bool, optional): Optional boolean array to
  mask observations.
\item
  \emph{create\_obj} (bool, optional): If True, returns the raw BF
  distribution object instead of creating a sample site. This is
  essential for building complex distributions like
  \texttt{MixtureSameFamily}.
\item
  \emph{sample} (bool, optional): A control-flow argument. If
  \texttt{True}, the function will directly sample a raw JAX array from
  the distribution, bypassing the BF model context. If \texttt{False},
  it will create a \texttt{BF.sample} site within a model. Defaults to
  \texttt{False}.
\item
  \emph{seed} (int, optional): An integer seed for this draw when
  \texttt{sample=True}. By default (\texttt{None}) the draw comes from
  the session's random stream, set by \texttt{bf(rand\_seed=...)}, so
  consecutive draws differ but the whole sequence is reproducible. No
  effect when \texttt{sample=False}.
\item
  \emph{obs} (jnp.ndarray, optional): The observed value for this random
  variable. If provided, the sample site is conditioned on this value,
  and the function returns the observed value. If \texttt{None}, the
  site is treated as a latent (unobserved) random variable. Defaults to
  \texttt{None}.
\item
  \emph{name} (str, optional): The name of the sample site in a BF
  model. This is used to uniquely identify the random variable. Defaults
  to `x'.
\end{itemize}

\textbf{Returns} - When \texttt{sample=False}: A BF Multinomial
distribution object (for model building). - When \texttt{sample=True}: A
JAX array of samples drawn from the Multinomial distribution (for direct
sampling). - When \texttt{create\_obj=True}: The raw BF distribution
object (for advanced use cases).

\textbf{Example Usage} from BayesForge import bf import jax.numpy as jnp
m = bf(platform=`cpu') m.dist.multinomial(total\_count=10,
probs=jnp.array({[}0.2, 0.3, 0.5{]}), sample=True)

\textbf{Wrapper of:
https://num.pyro.ai/en/stable/distributions.html\#multinomial}
\textbf{References} *
https://online.stat.psu.edu/stat504/book/export/html/667 ``2.3 - The
Multinomial Distribution - STAT ONLINE'' *
https://en.wikipedia.org/wiki/Multinomial\_distribution ``Multinomial
distribution'' *
https://faculty.washington.edu/yenchic/20A\_stat512/Lec7\_Multinomial.pdf
``{[}PDF{]} Lecture 7: Multinomial distribution''

\subsection{\texorpdfstring{\texttt{m.dist.multinomial\_logits}}{m.dist.multinomial\_logits}}\label{multinomial_logits}

\begin{Shaded}
\begin{Highlighting}[]
\NormalTok{m.dist.multinomial\_logits(}
\NormalTok{    logits,}
\NormalTok{    total\_count}\OperatorTok{=}\DecValTok{1}\NormalTok{,}
\NormalTok{    total\_count\_max}\OperatorTok{=}\VariableTok{None}\NormalTok{,}
\NormalTok{    validate\_args}\OperatorTok{=}\VariableTok{None}\NormalTok{,}
\NormalTok{    name}\OperatorTok{=}\StringTok{\textquotesingle{}x\textquotesingle{}}\NormalTok{,}
\NormalTok{    obs}\OperatorTok{=}\VariableTok{None}\NormalTok{,}
\NormalTok{    mask}\OperatorTok{=}\VariableTok{None}\NormalTok{,}
\NormalTok{    infer}\OperatorTok{=}\VariableTok{None}\NormalTok{,}
\NormalTok{    sample}\OperatorTok{=}\VariableTok{False}\NormalTok{,}
\NormalTok{    seed}\OperatorTok{=}\VariableTok{None}\NormalTok{,}
\NormalTok{    shape}\OperatorTok{=}\NormalTok{(),}
\NormalTok{    event}\OperatorTok{=}\DecValTok{0}\NormalTok{,}
\NormalTok{    create\_obj}\OperatorTok{=}\VariableTok{False}\NormalTok{,}
\NormalTok{    to\_jax}\OperatorTok{=}\VariableTok{True}\NormalTok{,}
\NormalTok{)}
\end{Highlighting}
\end{Shaded}

\textbf{Multinomial Logits} A \emph{multinomial logits} distribution
refers to a categorical (or more generally multinomial) distribution
over (K) classes whose probabilities are given via the softmax of a
vector of logits. That is, given a vector of real-valued logits
\(\ell = (\ell_1, …, \ell_K)\), the class probabilities are: \[
  p_k = \frac{\exp(\ell_k)}{\sum_{j=1}^K \exp(\ell_j)}.
  \] Then a single draw from the distribution yields one of the (K)
classes (or for a multinomial count version, counts over the classes)
with those probabilities.

\textbf{Categorical version (one draw):} Let (X \in {1,2,\dots,K}) be
the class. Then \[
  \Pr(X=k \mid \ell) ;=; p_k = \frac{\exp(\ell_k)}{\sum_{j=1}^K \exp(\ell_j)}.
  \] This directly uses the logits \(\ell\).

\textbf{Multinomial version (n draws):} If you draw (n) independent
draws (or equivalently count vector) with probabilities (p =
(p\_1,\ldots,p\_K)), then for a count vector (x = (x\_1,\ldots,x\_K))
with (\sum\_k x\_k = n), \[
  \Pr(X = x \mid n, p) = \frac{n!}{x_1! \cdots x_K!} ; \prod_{k=1}^K p_k^{x_k}
  \] and \(p_k\) is given by the softmax of logits.

\textbf{Args} - \emph{logits} (jnp.ndarray): Logits for each outcome.
Must be at least one-dimensional.

\begin{itemize}
\item
  \emph{total\_count} (jnp.ndarray): The total number of trials.
\item
  \emph{shape} (tuple): A multi-purpose argument for shaping. When
  \texttt{sample=False} (model building), this is used with
  \texttt{shape} to set the distribution's batch shape. When
  \texttt{sample=True} (direct sampling), this is used as
  \texttt{sample\_shape} to draw a raw JAX array of the given shape.
\item
  \emph{event} (int): The number of batch dimensions to reinterpret as
  event dimensions (used in model building).
\item
  \emph{mask} (jnp.ndarray, bool): Optional boolean array to mask
  observations.
\item
  \emph{create\_obj} (bool): If True, returns the raw BF distribution
  object instead of creating a sample site. This is essential for
  building complex distributions like \texttt{MixtureSameFamily}.
\item
  \emph{sample} (bool, optional): A control-flow argument. If
  \texttt{True}, the function will directly sample a raw JAX array from
  the distribution, bypassing the BF model context. If \texttt{False},
  it will create a \texttt{BF.sample} site within a model. Defaults to
  \texttt{False}.
\item
  \emph{seed} (int, optional): An integer seed for this draw when
  \texttt{sample=True}. By default (\texttt{None}) the draw comes from
  the session's random stream, set by \texttt{bf(rand\_seed=...)}, so
  consecutive draws differ but the whole sequence is reproducible. No
  effect when \texttt{sample=False}.
\item
  \emph{obs} (jnp.ndarray, optional): The observed value for this random
  variable. If provided, the sample site is conditioned on this value,
  and the function returns the observed value. If \texttt{None}, the
  site is treated as a latent (unobserved) random variable. Defaults to
  \texttt{None}.
\item
  \emph{name} (str, optional): The name of the sample site in a BF
  model. This is used to uniquely identify the random variable. Defaults
  to `x'.
\end{itemize}

\textbf{Returns} - When \texttt{sample=False}: A BF MultinomialLogits
distribution object (for model building).

\begin{itemize}
\item
  When \texttt{sample=True}: A JAX array of samples drawn from the
  MultinomialLogits distribution (for direct sampling).
\item
  When \texttt{create\_obj=True}: The raw BF distribution object (for
  advanced use cases).

  \textbf{Example Usage} from BayesForge import bf import jax.numpy as
  jnp m = bf(platform=`cpu')
  m.dist.multinomial\_logits(logits=jnp.array({[}1.0, 0.5{]},
  dtype=jnp.float32), total\_count=jnp.array(5, dtype=jnp.int32),
  sample=True)

  \textbf{Wrapper of}
  https://num.pyro.ai/en/stable/distributions.html\#multinomiallogits

  \textbf{Reference}

  \begin{itemize}
  \tightlist
  \item
    https://en.wikipedia.org/wiki/Multinomial\_distribution
  \end{itemize}
\end{itemize}

\subsection{\texorpdfstring{\texttt{m.dist.multinomial\_probs}}{m.dist.multinomial\_probs}}\label{multinomial_probs}

\begin{Shaded}
\begin{Highlighting}[]
\NormalTok{m.dist.multinomial\_probs(}
\NormalTok{    probs,}
\NormalTok{    total\_count}\OperatorTok{=}\DecValTok{1}\NormalTok{,}
\NormalTok{    total\_count\_max}\OperatorTok{=}\VariableTok{None}\NormalTok{,}
\NormalTok{    validate\_args}\OperatorTok{=}\VariableTok{None}\NormalTok{,}
\NormalTok{    name}\OperatorTok{=}\StringTok{\textquotesingle{}x\textquotesingle{}}\NormalTok{,}
\NormalTok{    obs}\OperatorTok{=}\VariableTok{None}\NormalTok{,}
\NormalTok{    mask}\OperatorTok{=}\VariableTok{None}\NormalTok{,}
\NormalTok{    infer}\OperatorTok{=}\VariableTok{None}\NormalTok{,}
\NormalTok{    sample}\OperatorTok{=}\VariableTok{False}\NormalTok{,}
\NormalTok{    seed}\OperatorTok{=}\VariableTok{None}\NormalTok{,}
\NormalTok{    shape}\OperatorTok{=}\NormalTok{(),}
\NormalTok{    event}\OperatorTok{=}\DecValTok{0}\NormalTok{,}
\NormalTok{    create\_obj}\OperatorTok{=}\VariableTok{False}\NormalTok{,}
\NormalTok{    to\_jax}\OperatorTok{=}\VariableTok{True}\NormalTok{,}
\NormalTok{)}
\end{Highlighting}
\end{Shaded}

\textbf{Multinomial Probs} Samples from a Multinomial distribution.

The Multinomial distribution models the number of times each of several
discrete outcomes occurs in a fixed number of trials. Each trial
independently results in one of several outcomes, and each outcome has a
probability of occurring.

\textbf{Args} - \emph{probs} (jnp.ndarray): Vector of probabilities for
each outcome. Must sum to 1.

\begin{itemize}
\item
  \emph{total\_count} (jnp.ndarray): The number of trials.
\item
  \emph{shape} (tuple): A multi-purpose argument for shaping. When
  \texttt{sample=False} (model building), this is used with
  \texttt{shape} to set the distribution's batch shape. When
  \texttt{sample=True} (direct sampling), this is used as
  \texttt{sample\_shape} to draw a raw JAX array of the given shape.
\item
  \emph{event} (int): The number of batch dimensions to reinterpret as
  event dimensions (used in model building).
\item
  \emph{mask} (jnp.ndarray, bool): Optional boolean array to mask
  observations.
\item
  \emph{create\_obj} (bool): If True, returns the raw BF distribution
  object instead of creating a sample site. This is essential for
  building complex distributions like \texttt{MixtureSameFamily}.
\item
  \emph{sample} (bool, optional): A control-flow argument. If
  \texttt{True}, the function will directly sample a raw JAX array from
  the distribution, bypassing the BF model context. If \texttt{False},
  it will create a \texttt{BF.sample} site within a model. Defaults to
  \texttt{False}.
\item
  \emph{seed} (int, optional): An integer seed for this draw when
  \texttt{sample=True}. By default (\texttt{None}) the draw comes from
  the session's random stream, set by \texttt{bf(rand\_seed=...)}, so
  consecutive draws differ but the whole sequence is reproducible. No
  effect when \texttt{sample=False}.
\item
  \emph{obs} (jnp.ndarray, optional): The observed value for this random
  variable. If provided, the sample site is conditioned on this value,
  and the function returns the observed value. If \texttt{None}, the
  site is treated as a latent (unobserved) random variable. Defaults to
  \texttt{None}.
\item
  \emph{name} (str, optional): The name of the sample site in a BF
  model. This is used to uniquely identify the random variable. Defaults
  to `x'.
\end{itemize}

\textbf{Returns} BF MultinomialProbs distribution object (for model
building). JAX array of samples drawn from the MultinomialProbs
distribution (for direct sampling). The raw BF distribution object (for
advanced use cases).

\textbf{Example Usage} from BayesForge import bf import jax.numpy as jnp
m = bf(platform=`cpu') m.dist.multinomial\_probs(probs=jnp.array({[}0.2,
0.3, 0.5{]}), total\_count=10, sample=True)

\textbf{Wrapper of}
https://num.pyro.ai/en/stable/distributions.html\#multinomialprobs

\subsection{\texorpdfstring{\texttt{m.dist.multivariate\_normal}}{m.dist.multivariate\_normal}}\label{multivariate_normal}

\begin{Shaded}
\begin{Highlighting}[]
\NormalTok{m.dist.multivariate\_normal(}
\NormalTok{    loc}\OperatorTok{=}\FloatTok{0.0}\NormalTok{,}
\NormalTok{    covariance\_matrix}\OperatorTok{=}\VariableTok{None}\NormalTok{,}
\NormalTok{    precision\_matrix}\OperatorTok{=}\VariableTok{None}\NormalTok{,}
\NormalTok{    scale\_tril}\OperatorTok{=}\VariableTok{None}\NormalTok{,}
\NormalTok{    validate\_args}\OperatorTok{=}\VariableTok{None}\NormalTok{,}
\NormalTok{    name}\OperatorTok{=}\StringTok{\textquotesingle{}x\textquotesingle{}}\NormalTok{,}
\NormalTok{    obs}\OperatorTok{=}\VariableTok{None}\NormalTok{,}
\NormalTok{    mask}\OperatorTok{=}\VariableTok{None}\NormalTok{,}
\NormalTok{    infer}\OperatorTok{=}\VariableTok{None}\NormalTok{,}
\NormalTok{    sample}\OperatorTok{=}\VariableTok{False}\NormalTok{,}
\NormalTok{    seed}\OperatorTok{=}\VariableTok{None}\NormalTok{,}
\NormalTok{    shape}\OperatorTok{=}\NormalTok{(),}
\NormalTok{    event}\OperatorTok{=}\DecValTok{0}\NormalTok{,}
\NormalTok{    create\_obj}\OperatorTok{=}\VariableTok{False}\NormalTok{,}
\NormalTok{    to\_jax}\OperatorTok{=}\VariableTok{True}\NormalTok{,}
\NormalTok{)}
\end{Highlighting}
\end{Shaded}

\textbf{Multivariate Normal} The Multivariate Normal distribution, also
known as the Gaussian distribution in multiple dimensions, is a
probability distribution that arises frequently in statistics and
machine learning. It is defined by its mean vector and covariance
matrix, which describe the central tendency and spread of the
distribution, respectively.

\[
     p(x) = \frac{1}{\sqrt{(2\pi)^n |\Sigma|}} \exp\left(-\frac{1}{2}(x - \mu)^T \Sigma^{-1} (x - \mu)\right)
  \]

where: - \(x\) is a \(n\)-dimensional vector of random variables. -
\(\mu\) is the mean vector. - \(\Sigma\) is the covariance matrix.

\textbf{Args} - \emph{loc} (tuple): The mean vector of the distribution.

\begin{itemize}
\item
  \emph{covariance\_matrix} (jnp.ndarray, optional): The covariance
  matrix of the distribution. Must be positive definite.
\item
  \emph{precision\_matrix} (jnp.ndarray, optional): The precision matrix
  (inverse of the covariance matrix) of the distribution. Must be
  positive definite.
\item
  \emph{scale\_tril} (jnp.ndarray, optional): The lower triangular
  Cholesky decomposition of the covariance matrix.
\item
  \texttt{shape} (tuple): A multi-purpose argument for shaping. When
  \texttt{sample=False} (model building), this is used with
  \texttt{shape} to set the distribution's batch shape. When
  \texttt{sample=True} (direct sampling), this is used as
  \texttt{sample\_shape} to draw a raw JAX array of the given shape.
\item
  \emph{event} (int): The number of batch dimensions to reinterpret as
  event dimensions (used in model building).
\item
  \emph{mask} (jnp.ndarray, bool): Optional boolean array to mask
  observations.
\item
  \emph{create\_obj} (bool): If True, returns the raw BF distribution
  object instead of creating a sample site. This is essential for
  building complex distributions like \texttt{MixtureSameFamily}.
\item
  \emph{sample} (bool, optional): A control-flow argument. If
  \texttt{True}, the function will directly sample a raw JAX array from
  the distribution, bypassing the BF model context. If \texttt{False},
  it will create a \texttt{BF.sample} site within a model. Defaults to
  \texttt{False}.
\item
  \emph{seed} (int, optional): An integer seed for this draw when
  \texttt{sample=True}. By default (\texttt{None}) the draw comes from
  the session's random stream, set by \texttt{bf(rand\_seed=...)}, so
  consecutive draws differ but the whole sequence is reproducible. No
  effect when \texttt{sample=False}.
\item
  \emph{obs} (jnp.ndarray, optional): The observed value for this random
  variable. If provided, the sample site is conditioned on this value,
  and the function returns the observed value. If \texttt{None}, the
  site is treated as a latent (unobserved) random variable. Defaults to
  \texttt{None}.
\item
  \emph{name} (str, optional): The name of the sample site in a BF
  model. This is used to uniquely identify the random variable. Defaults
  to `x'.
\end{itemize}

\textbf{Returns} - When \texttt{sample=False}: A BF MultivariateNormal
distribution object (for model building).

\begin{itemize}
\item
  When \texttt{sample=True}: A JAX array of samples drawn from the
  MultivariateNormal distribution (for direct sampling).
\item
  When \texttt{create\_obj=True}: The raw BF distribution object (for
  advanced use cases).

  \textbf{Example Usage} from BayesForge import bf import jax.numpy as
  jnp m = bf(platform=`cpu') m.dist.multivariate\_normal(
  loc=jnp.array({[}1.0, 0.0, -2.0{]}),
  covariance\_matrix=jnp.array({[}{[} 2.0, 0.7, -0.3{]}, {[} 0.7, 1.0,
  0.5{]}, {[}-0.3, 0.5, 1.5{]}{]}), sample=True, )

  \textbf{Wrapper of:
  https://num.pyro.ai/en/stable/distributions.html\#multivariate-normal}
  \textbf{Reference}

  \begin{itemize}
  \tightlist
  \item
    \href{https://en.wikipedia.org/wiki/Multivariate_normal_distribution}{Multivariate
    normal distribution -- Wikipedia}
  \end{itemize}
\end{itemize}

\subsection{\texorpdfstring{\texttt{m.dist.multivariate\_student\_t}}{m.dist.multivariate\_student\_t}}\label{multivariate_student_t}

\begin{Shaded}
\begin{Highlighting}[]
\NormalTok{m.dist.multivariate\_student\_t(}
\NormalTok{    df,}
\NormalTok{    loc}\OperatorTok{=}\FloatTok{0.0}\NormalTok{,}
\NormalTok{    scale\_tril}\OperatorTok{=}\VariableTok{None}\NormalTok{,}
\NormalTok{    validate\_args}\OperatorTok{=}\VariableTok{None}\NormalTok{,}
\NormalTok{    name}\OperatorTok{=}\StringTok{\textquotesingle{}x\textquotesingle{}}\NormalTok{,}
\NormalTok{    obs}\OperatorTok{=}\VariableTok{None}\NormalTok{,}
\NormalTok{    mask}\OperatorTok{=}\VariableTok{None}\NormalTok{,}
\NormalTok{    infer}\OperatorTok{=}\VariableTok{None}\NormalTok{,}
\NormalTok{    sample}\OperatorTok{=}\VariableTok{False}\NormalTok{,}
\NormalTok{    seed}\OperatorTok{=}\VariableTok{None}\NormalTok{,}
\NormalTok{    shape}\OperatorTok{=}\NormalTok{(),}
\NormalTok{    event}\OperatorTok{=}\DecValTok{0}\NormalTok{,}
\NormalTok{    create\_obj}\OperatorTok{=}\VariableTok{False}\NormalTok{,}
\NormalTok{    to\_jax}\OperatorTok{=}\VariableTok{True}\NormalTok{,}
\NormalTok{)}
\end{Highlighting}
\end{Shaded}

\textbf{Multivariate Student's t} The Multivariate Student's t
distribution is a generalization of the Student's t distribution to
multiple dimensions. It is a heavy-tailed distribution that is often
used to model data that is not normally distributed.

The PDF of the multivariate Student's t-distribution for a random vector
( \mathbf{x} \in \mathbb{R}\^{}d ) is given by:

\[
  f(\mathbf{x}) = \frac{\Gamma\left(\frac{\nu + d}{2}\right)}{\Gamma\left(\frac{\nu}{2}\right) \nu^{d/2} \pi^{d/2} | \Sigma|^{1/2}} \left(1 + \frac{1}{\nu} (\mathbf{x} - \mu)^T \Sigma^{-1} (\mathbf{x} - \mu)\right)^{-(\nu + d)/2}
  \]

where:

\begin{itemize}
\tightlist
\item
  \(\Gamma(\cdot)\) is the Gamma function.
\item
  \(\mu\) is the mean vector.
\item
  \(\Sigma\) is the scale (covariance) matrix.
\item
  \(\nu\) is the degrees of freedom.
\item
  \$ d \(is the dimensionality of\) \mathbf{x} \$.
\end{itemize}

\textbf{Args} - \emph{df} (jnp.ndarray): Degrees of freedom, must be
positive.

\begin{itemize}
\item
  \emph{loc} (jnp.ndarray): Location vector, representing the mean of
  the distribution. scale\_tril (jnp.ndarray): Lower triangular matrix
  defining the scale.
\item
  \texttt{shape} (tuple): A multi-purpose argument for shaping. When
  \texttt{sample=False} (model building), this is used with
  \texttt{shape} to set the distribution's batch shape. When
  \texttt{sample=True} (direct sampling), this is used as
  \texttt{sample\_shape} to draw a raw JAX array of the given shape.
\item
  \emph{event} (int): The number of batch dimensions to reinterpret as
  event dimensions (used in model building).
\item
  \emph{mask} (jnp.ndarray, bool): Optional boolean array to mask
  observations.
\item
  \emph{create\_obj} (bool): If True, returns the raw BF distribution
  object instead of creating a sample site. This is essential for
  building complex distributions like \texttt{MixtureSameFamily}.
\item
  \emph{sample} (bool, optional): A control-flow argument. If
  \texttt{True}, the function will directly sample a raw JAX array from
  the distribution, bypassing the BF model context. If \texttt{False},
  it will create a \texttt{BF.sample} site within a model. Defaults to
  \texttt{False}.
\item
  \emph{seed} (int, optional): An integer seed for this draw when
  \texttt{sample=True}. By default (\texttt{None}) the draw comes from
  the session's random stream, set by \texttt{bf(rand\_seed=...)}, so
  consecutive draws differ but the whole sequence is reproducible. No
  effect when \texttt{sample=False}.
\item
  \emph{obs} (jnp.ndarray, optional): The observed value for this random
  variable. If provided, the sample site is conditioned on this value,
  and the function returns the observed value. If \texttt{None}, the
  site is treated as a latent (unobserved) random variable. Defaults to
  \texttt{None}.
\item
  \emph{name} (str, optional): The name of the sample site in a BF
  model. This is used to uniquely identify the random variable. Defaults
  to `x'.
\end{itemize}

\textbf{Returns} - When \texttt{sample=False}: A BF MultivariateStudentT
distribution object (for model building). - When \texttt{sample=True}: A
JAX array of samples drawn from the MultivariateStudentT distribution
(for direct sampling). - When \texttt{create\_obj=True}: The raw BF
distribution object (for advanced use cases).

\textbf{Example Usage} from BayesForge import bf import jax.numpy as jnp
m = bf(platform=`cpu') m.dist.multivariate\_student\_t( df=2,
loc=jnp.array({[}1.0, 0.0, -2.0{]}),
scale\_tril=jnp.linalg.cholesky(jnp.array({[}{[} 2.0, 0.7, -0.3{]}, {[}
0.7, 1.0, 0.5{]}, {[}-0.3, 0.5, 1.5{]}{]})), sample=True, )

\textbf{Wrapper of}
https://num.pyro.ai/en/stable/distributions.html\#multivariatestudentt

\textbf{Reference} *
\href{https://en.wikipedia.org/wiki/Multivariate_Student\%27s_t-distribution}{Multivariate
Student-t distribution -- Wikipedia}

\subsection{\texorpdfstring{\texttt{m.dist.negative\_binomial}}{m.dist.negative\_binomial}}\label{negative_binomial}

\begin{Shaded}
\begin{Highlighting}[]
\NormalTok{m.dist.negative\_binomial(}
\NormalTok{    total\_count,}
\NormalTok{    probs}\OperatorTok{=}\VariableTok{None}\NormalTok{,}
\NormalTok{    logits}\OperatorTok{=}\VariableTok{None}\NormalTok{,}
\NormalTok{    validate\_args}\OperatorTok{=}\VariableTok{None}\NormalTok{,}
\NormalTok{    name}\OperatorTok{=}\StringTok{\textquotesingle{}x\textquotesingle{}}\NormalTok{,}
\NormalTok{    obs}\OperatorTok{=}\VariableTok{None}\NormalTok{,}
\NormalTok{    mask}\OperatorTok{=}\VariableTok{None}\NormalTok{,}
\NormalTok{    infer}\OperatorTok{=}\VariableTok{None}\NormalTok{,}
\NormalTok{    sample}\OperatorTok{=}\VariableTok{False}\NormalTok{,}
\NormalTok{    seed}\OperatorTok{=}\VariableTok{None}\NormalTok{,}
\NormalTok{    shape}\OperatorTok{=}\NormalTok{(),}
\NormalTok{    event}\OperatorTok{=}\DecValTok{0}\NormalTok{,}
\NormalTok{    create\_obj}\OperatorTok{=}\VariableTok{False}\NormalTok{,}
\NormalTok{    to\_jax}\OperatorTok{=}\VariableTok{True}\NormalTok{,}
\NormalTok{)}
\end{Highlighting}
\end{Shaded}

\textbf{Negative Binomial} The Negative Binomial distribution models the
number of failures (or the total number of trials) in a sequence of
independent Bernoulli trials with success probability (p), until a
specified number \texttt{total\_count} (r) of successes is achieved. It
is often used as a count-data model when the variance exceeds the mean
(``overdispersion'') relative to a Poisson.

PMF (Probability Mass Function): Given (X) = number of failures before
the (r)th success (so (X = 0,1,2,\dots)), with (r) successes fixed,
success probability (p), failure probability (q = 1 - p):

\[
\Pr(X = k) ;=; \binom{k + r - 1}{k} ; p^r ; q^k, \quad k = 0,1,2,\dots
\]

\textbf{Args} - \emph{total\_count} (jnp.ndarray): The total number of
events.

\begin{itemize}
\item
  \emph{shape} (tuple): A multi-purpose argument for shaping. When
  \texttt{sample=False} (model building), this is used with
  \texttt{shape} to set the distribution's batch shape. When
  \texttt{sample=True} (direct sampling), this is used as
  \texttt{sample\_shape} to draw a raw JAX array of the given shape.
\item
  \emph{event} (int): The number of batch dimensions to reinterpret as
  event dimensions (used in model building).
\item
  \emph{mask} (jnp.ndarray, bool): Optional boolean array to mask
  observations.
\item
  \emph{create\_obj} (bool): If True, returns the raw BF distribution
  object instead of creating a sample site. This is essential for
  building complex distributions like \texttt{MixtureSameFamily}.
\item
  \emph{sample} (bool, optional): A control-flow argument. If
  \texttt{True}, the function will directly sample a raw JAX array from
  the distribution, bypassing the BF model context. If \texttt{False},
  it will create a \texttt{BF.sample} site within a model. Defaults to
  \texttt{False}.
\item
  \emph{seed} (int, optional): An integer seed for this draw when
  \texttt{sample=True}. By default (\texttt{None}) the draw comes from
  the session's random stream, set by \texttt{bf(rand\_seed=...)}, so
  consecutive draws differ but the whole sequence is reproducible. No
  effect when \texttt{sample=False}.
\item
  \emph{obs} (jnp.ndarray, optional): The observed value for this random
  variable. If provided, the sample site is conditioned on this value,
  and the function returns the observed value. If \texttt{None}, the
  site is treated as a latent (unobserved) random variable. Defaults to
  \texttt{None}.
\item
  \emph{name} (str, optional): The name of the sample site in a BF
  model. This is used to uniquely identify the random variable. Defaults
  to `x'.
\end{itemize}

\textbf{Returns} BF NegativeBinomial distribution object (for model
building). JAX array of samples drawn from the NegativeBinomial
distribution (for direct sampling). The raw BF distribution object (for
advanced use cases).

\textbf{Example Usage} from BayesForge import bf import jax.numpy as jnp
m = bf(platform=`cpu') m.dist.negative\_binomial(total\_count=5.0,probs
= jnp.array({[}0.2, 0.3, 0.5{]}), sample=True)

\textbf{Wrapper of:
https://num.pyro.ai/en/stable/distributions.html\#negativebinomial}
\textbf{References} *
https://en.wikipedia.org/wiki/Negative\_binomial\_distribution *
https://www.statology.org/negative-binomial-distribution/

\subsection{\texorpdfstring{\texttt{m.dist.negative\_binomial2}}{m.dist.negative\_binomial2}}\label{negative_binomial2}

\begin{Shaded}
\begin{Highlighting}[]
\NormalTok{m.dist.negative\_binomial2(}
\NormalTok{    mean,}
\NormalTok{    concentration,}
\NormalTok{    validate\_args}\OperatorTok{=}\VariableTok{None}\NormalTok{,}
\NormalTok{    name}\OperatorTok{=}\StringTok{\textquotesingle{}x\textquotesingle{}}\NormalTok{,}
\NormalTok{    obs}\OperatorTok{=}\VariableTok{None}\NormalTok{,}
\NormalTok{    mask}\OperatorTok{=}\VariableTok{None}\NormalTok{,}
\NormalTok{    infer}\OperatorTok{=}\VariableTok{None}\NormalTok{,}
\NormalTok{    sample}\OperatorTok{=}\VariableTok{False}\NormalTok{,}
\NormalTok{    seed}\OperatorTok{=}\VariableTok{None}\NormalTok{,}
\NormalTok{    shape}\OperatorTok{=}\NormalTok{(),}
\NormalTok{    event}\OperatorTok{=}\DecValTok{0}\NormalTok{,}
\NormalTok{    create\_obj}\OperatorTok{=}\VariableTok{False}\NormalTok{,}
\NormalTok{    to\_jax}\OperatorTok{=}\VariableTok{True}\NormalTok{,}
\NormalTok{)}
\end{Highlighting}
\end{Shaded}

\textbf{Samples from a Negative Binomial} The NB2 parameterisation of
the negative-binomial distribution is a count distribution used for
modelling over-dispersed count data (variance \textgreater{} mean). It
is defined such that the variance grows \textbf{quadratically} in the
mean:

\[
mathrm{Var}(Y) = \mu + \alpha,\mu^2,
\] where (\mu = \mathbb{E}{[}Y{]}) and (\alpha\textgreater0) is the
dispersion (heterogeneity) parameter. Because of this quadratic variance
growth, it is called the NB2 family.

\[
P(k) = \frac{\Gamma(k + \alpha)}{\Gamma(k + 1) \Gamma(\alpha)} \left(\frac{\beta}{\alpha + \beta}\right)^k \left(1 - \frac{\beta}{\alpha + \beta}\right)^k
\]

One commonly used form of the NB2 parameterisation is obtained via a
Poisson‐Gamma mixture: \[
Y \mid \lambda \sim \text{Poisson}(\lambda), \quad \lambda \sim \text{Gamma}!\Big(\frac{1}{\alpha},,\frac{\mu,\alpha}{1}\Big)
\]

which marginalises to:

\[
\Pr(Y = y) = \frac{\Gamma(y + \tfrac1\alpha)}{\Gamma(\tfrac1\alpha); y!}; \bigg(\frac{1}{1 + \alpha,\mu}\bigg)^{1/\alpha}; \bigg(\frac{\alpha,\mu}{1 + \alpha,\mu}\bigg)^{y}, \quad y = 0,1,2,\dots
\] This parameterisation matches the NB2 form with mean \(\mu\) and
dispersion \(\alpha\). (See Hilbe (2011) forderivation.)

\textbf{Args} - \emph{mean} (jnp.ndarray or float): The mean of the
distribution. This is equivalent to the \texttt{mu} parameter.

\begin{itemize}
\item
  \emph{concentration} (jnp.ndarray or float): The concentration
  parameter. This is equivalent to the \(alpha\) parameter.
\item
  \emph{shape} (tuple): A multi-purpose argument for shaping. When
  \texttt{sample=False} (model building), this is used with
  \texttt{shape} to set the distribution's batch shape. When
  \texttt{sample=True} (direct sampling), this is used as
  \texttt{sample\_shape} to draw a raw JAX array of the given shape.
\item
  \emph{event} (int): The number of batch dimensions to reinterpret as
  event dimensions (used in model building).
\item
  \emph{mask} (jnp.ndarray, bool): Optional boolean array to mask
  observations.
\item
  \emph{create\_obj} (bool): If True, returns the raw BF distribution
  object instead of creating a sample site. This is essential for
  building complex distributions like \texttt{MixtureSameFamily}.
\item
  \emph{sample} (bool, optional): A control-flow argument. If
  \texttt{True}, the function will directly sample a raw JAX array from
  the distribution, bypassing the BF model context. If \texttt{False},
  it will create a \texttt{BF.sample} site within a model. Defaults to
  \texttt{False}.
\item
  \emph{seed} (int, optional): An integer seed for this draw when
  \texttt{sample=True}. By default (\texttt{None}) the draw comes from
  the session's random stream, set by \texttt{bf(rand\_seed=...)}, so
  consecutive draws differ but the whole sequence is reproducible. No
  effect when \texttt{sample=False}.
\item
  \emph{obs} (jnp.ndarray, optional): The observed value for this random
  variable. If provided, the sample site is conditioned on this value,
  and the function returns the observed value. If \texttt{None}, the
  site is treated as a latent (unobserved) random variable. Defaults to
  \texttt{None}.
\item
  \emph{name} (str, optional): The name of the sample site in a BF
  model. This is used to uniquely identify the random variable. Defaults
  to `x'.
\end{itemize}

\textbf{Returns} BF NegativeBinomial distribution object: When
\texttt{sample=False} (for model building). jnp.ndarray: A JAX array of
samples drawn from the NegativeBinomial distribution (for direct
sampling). BF NegativeBinomial distribution object: When
\texttt{create\_obj=True} (for advanced use cases).

\textbf{Example Usage} from BayesForge import bf m = bf(platform=`cpu')
m.dist.negative\_binomial2(mean=2.0, concentration=3.0, sample=True)

\textbf{Wrapper of:
https://num.pyro.ai/en/stable/distributions.html\#negativebinomial}
\textbf{Reference} * {[}Negative Binomial
Regression{]}9https://www.cambridge.org/core/books/negative-binomial-regression/12D6281A46B9A980DC6021080C9419E7)
*
\href{https://assets.cambridge.org/97805211/98158/excerpt/9780521198158_excerpt.pdf}{Negative
Binomial Regression: Second Edition}

\subsection{\texorpdfstring{\texttt{m.dist.negative\_binomial\_logits}}{m.dist.negative\_binomial\_logits}}\label{negative_binomial_logits}

\begin{Shaded}
\begin{Highlighting}[]
\NormalTok{m.dist.negative\_binomial\_logits(}
\NormalTok{    total\_count,}
\NormalTok{    logits,}
\NormalTok{    validate\_args}\OperatorTok{=}\VariableTok{None}\NormalTok{,}
\NormalTok{    name}\OperatorTok{=}\StringTok{\textquotesingle{}x\textquotesingle{}}\NormalTok{,}
\NormalTok{    obs}\OperatorTok{=}\VariableTok{None}\NormalTok{,}
\NormalTok{    mask}\OperatorTok{=}\VariableTok{None}\NormalTok{,}
\NormalTok{    infer}\OperatorTok{=}\VariableTok{None}\NormalTok{,}
\NormalTok{    sample}\OperatorTok{=}\VariableTok{False}\NormalTok{,}
\NormalTok{    seed}\OperatorTok{=}\VariableTok{None}\NormalTok{,}
\NormalTok{    shape}\OperatorTok{=}\NormalTok{(),}
\NormalTok{    event}\OperatorTok{=}\DecValTok{0}\NormalTok{,}
\NormalTok{    create\_obj}\OperatorTok{=}\VariableTok{False}\NormalTok{,}
\NormalTok{    to\_jax}\OperatorTok{=}\VariableTok{True}\NormalTok{,}
\NormalTok{)}
\end{Highlighting}
\end{Shaded}

\textbf{Negative Binomial Logits} Samples from a Negative Binomial
Logits distribution.

The Negative Binomial Logits distribution is a generalization of the
Negative Binomial distribution where the parameter `r' (number of
successes) is expressed as a function of a logit parameter. This allows
for more flexible modeling of count data.

\textbf{Args} - \emph{total\_count} (jnp.ndarray): The parameter
controlling the shape of the distribution. Represents the total number
of trials.

\begin{itemize}
\item
  \emph{logits} (jnp.ndarray): The log-odds parameter. Related to the
  probability of success.
\item
  \emph{shape} (tuple): A multi-purpose argument for shaping. When
  \texttt{sample=False} (model building), this is used with
  \texttt{shape} to set the distribution's batch shape. When
  \texttt{sample=True} (direct sampling), this is used as
  \texttt{sample\_shape} to draw a raw JAX array of the given shape.
\item
  \emph{event} (int): The number of batch dimensions to reinterpret as
  event dimensions (used in model building).
\item
  \emph{mask} (jnp.ndarray, bool): Optional boolean array to mask
  observations.
\item
  \emph{create\_obj} (bool): If True, returns the raw BF distribution
  object instead of creating a sample site. This is essential for
  building complex distributions like \texttt{MixtureSameFamily}.
\item
  \emph{sample} (bool, optional): A control-flow argument. If
  \texttt{True}, the function will directly sample a raw JAX array from
  the distribution, bypassing the BF model context. If \texttt{False},
  it will create a \texttt{BF.sample} site within a model. Defaults to
  \texttt{False}.
\item
  \emph{seed} (int, optional): An integer seed for this draw when
  \texttt{sample=True}. By default (\texttt{None}) the draw comes from
  the session's random stream, set by \texttt{bf(rand\_seed=...)}, so
  consecutive draws differ but the whole sequence is reproducible. No
  effect when \texttt{sample=False}.
\item
  \emph{obs} (jnp.ndarray, optional): The observed value for this random
  variable. If provided, the sample site is conditioned on this value,
  and the function returns the observed value. If \texttt{None}, the
  site is treated as a latent (unobserved) random variable. Defaults to
  \texttt{None}.
\item
  \emph{name} (str, optional): The name of the sample site in a BF
  model. This is used to uniquely identify the random variable. Defaults
  to `x'.
\end{itemize}

\textbf{Returns} Negative Binomial Logits: A BF Negative Binomial Logits
distribution object (for model building).

\begin{verbatim}
jnp.ndarray: A JAX array of samples drawn from the Negative Binomial Logits distribution (for direct sampling).

Negative Binomial Logits: The raw BF distribution object (for advanced use cases).
\end{verbatim}

\textbf{Example Usage} from BayesForge import bf m = bf(platform=`cpu')
m.dist.negative\_binomial\_logits(total\_count=5.0, logits=0.0,
sample=True)

\textbf{Wrapper of}
https://num.pyro.ai/en/stable/distributions.html\#Negative Binomial
Logits

\subsection{\texorpdfstring{\texttt{m.dist.negative\_binomial\_probs}}{m.dist.negative\_binomial\_probs}}\label{negative_binomial_probs}

\begin{Shaded}
\begin{Highlighting}[]
\NormalTok{m.dist.negative\_binomial\_probs(}
\NormalTok{    total\_count,}
\NormalTok{    probs,}
\NormalTok{    validate\_args}\OperatorTok{=}\VariableTok{None}\NormalTok{,}
\NormalTok{    name}\OperatorTok{=}\StringTok{\textquotesingle{}x\textquotesingle{}}\NormalTok{,}
\NormalTok{    obs}\OperatorTok{=}\VariableTok{None}\NormalTok{,}
\NormalTok{    mask}\OperatorTok{=}\VariableTok{None}\NormalTok{,}
\NormalTok{    infer}\OperatorTok{=}\VariableTok{None}\NormalTok{,}
\NormalTok{    sample}\OperatorTok{=}\VariableTok{False}\NormalTok{,}
\NormalTok{    seed}\OperatorTok{=}\VariableTok{None}\NormalTok{,}
\NormalTok{    shape}\OperatorTok{=}\NormalTok{(),}
\NormalTok{    event}\OperatorTok{=}\DecValTok{0}\NormalTok{,}
\NormalTok{    create\_obj}\OperatorTok{=}\VariableTok{False}\NormalTok{,}
\NormalTok{    to\_jax}\OperatorTok{=}\VariableTok{True}\NormalTok{,}
\NormalTok{)}
\end{Highlighting}
\end{Shaded}

\textbf{Negative Binomial with probabilities.} The Negative Binomial
distribution models the number of failures before the first success in a
sequence of independent Bernoulli trials. It is characterized by two
parameters: `concentration' (r) and `rate' (p). In this implementation,
the `concentration' parameter is derived from `total\_count' and the
`rate' parameter is derived from `probs'.

\textbf{Args} - \emph{total\_count} (jnp.ndarray): A numeric vector,
matrix, or array representing the parameter.

\begin{itemize}
\item
  \emph{probs} (jnp.ndarray): A numeric vector representing event
  probabilities. Must sum to 1.
\item
  \emph{shape} (tuple): A multi-purpose argument for shaping. When
  \texttt{sample=False} (model building), this is used with
  \texttt{shape} to set the distribution's batch shape. When
  \texttt{sample=True} (direct sampling), this is used as
  \texttt{sample\_shape} to draw a raw JAX array of the given shape.
\item
  \emph{event} (int): The number of batch dimensions to reinterpret as
  event dimensions (used in model building).
\item
  \emph{mask} (jnp.ndarray, bool): Optional boolean array to mask
  observations.
\item
  \emph{create\_obj} (bool): If True, returns the raw BF distribution
  object instead of creating a sample site. This is essential for
  building complex distributions like \texttt{MixtureSameFamily}.
\item
  \emph{sample} (bool, optional): A control-flow argument. If
  \texttt{True}, the function will directly sample a raw JAX array from
  the distribution, bypassing the BF model context. If \texttt{False},
  it will create a \texttt{BF.sample} site within a model. Defaults to
  \texttt{False}.
\item
  \emph{seed} (int, optional): An integer seed for this draw when
  \texttt{sample=True}. By default (\texttt{None}) the draw comes from
  the session's random stream, set by \texttt{bf(rand\_seed=...)}, so
  consecutive draws differ but the whole sequence is reproducible. No
  effect when \texttt{sample=False}.
\item
  \emph{obs} (jnp.ndarray, optional): The observed value for this random
  variable. If provided, the sample site is conditioned on this value,
  and the function returns the observed value. If \texttt{None}, the
  site is treated as a latent (unobserved) random variable. Defaults to
  \texttt{None}.
\item
  \emph{name} (str, optional): The name of the sample site in a BF
  model. This is used to uniquely identify the random variable. Defaults
  to `x'.
\end{itemize}

\textbf{Returns} BF NegativeBinomialProbs distribution object (for model
building). JAX array of samples drawn from the NegativeBinomialProbs
distribution (for direct sampling). The raw BF distribution object (for
advanced use cases).

\textbf{Example Usage} from BayesForge import bf import jax.numpy as jnp
m = bf(platform=`cpu')
m.dist.negative\_binomial\_probs(total\_count=10.0, probs =
jnp.array({[}0.2, 0.3, 0.5{]}), sample=True)

\textbf{Wrapper of:
https://num.pyro.ai/en/stable/distributions.html\#negativebinomialprobs}

\subsection{\texorpdfstring{\texttt{m.dist.normal}}{m.dist.normal}}\label{normal}

\begin{Shaded}
\begin{Highlighting}[]
\NormalTok{m.dist.normal(}
\NormalTok{    loc}\OperatorTok{=}\FloatTok{0.0}\NormalTok{,}
\NormalTok{    scale}\OperatorTok{=}\FloatTok{1.0}\NormalTok{,}
\NormalTok{    validate\_args}\OperatorTok{=}\VariableTok{None}\NormalTok{,}
\NormalTok{    name}\OperatorTok{=}\StringTok{\textquotesingle{}x\textquotesingle{}}\NormalTok{,}
\NormalTok{    obs}\OperatorTok{=}\VariableTok{None}\NormalTok{,}
\NormalTok{    mask}\OperatorTok{=}\VariableTok{None}\NormalTok{,}
\NormalTok{    infer}\OperatorTok{=}\VariableTok{None}\NormalTok{,}
\NormalTok{    sample}\OperatorTok{=}\VariableTok{False}\NormalTok{,}
\NormalTok{    seed}\OperatorTok{=}\VariableTok{None}\NormalTok{,}
\NormalTok{    shape}\OperatorTok{=}\NormalTok{(),}
\NormalTok{    event}\OperatorTok{=}\DecValTok{0}\NormalTok{,}
\NormalTok{    create\_obj}\OperatorTok{=}\VariableTok{False}\NormalTok{,}
\NormalTok{    to\_jax}\OperatorTok{=}\VariableTok{True}\NormalTok{,}
\NormalTok{)}
\end{Highlighting}
\end{Shaded}

\textbf{Normal} Samples from a Normal (Gaussian) distribution.

The Normal distribution is characterized by its mean (loc) and standard
deviation (scale). It's a continuous probability distribution that
arises frequently in statistics and probability theory.

\[   
  f(x \mid \mu, \sigma) = \frac{1}{\sqrt{2\pi}\sigma}\exp!\left(-\frac{(x - \mu)^2}{2\sigma^2}\right)
  \]

\textbf{Args} - \emph{loc} (jnp.ndarray): The mean of the distribution.

\begin{itemize}
\item
  \emph{scale} (jnp.ndarray): The standard deviation of the
  distribution.
\item
  \emph{shape} (tuple): A multi-purpose argument for shaping. When
  \texttt{sample=False} (model building), this is used with
  \texttt{shape} to set the distribution's batch shape. When
  \texttt{sample=True} (direct sampling), this is used as
  \texttt{sample\_shape} to draw a raw JAX array of the given shape.
\item
  \emph{event} (int): The number of batch dimensions to reinterpret as
  event dimensions (used in model building).
\item
  \emph{mask} (jnp.ndarray, bool): Optional boolean array to mask
  observations.
\item
  \emph{create\_obj} (bool): If True, returns the raw BF distribution
  object instead of creating a sample site. This is essential for
  building complex distributions like \texttt{MixtureSameFamily}.
\item
  \emph{sample} (bool, optional): A control-flow argument. If
  \texttt{True}, the function will directly sample a raw JAX array from
  the distribution, bypassing the BF model context. If \texttt{False},
  it will create a \texttt{BF.sample} site within a model. Defaults to
  \texttt{False}.
\item
  \emph{seed} (int, optional): An integer seed for this draw when
  \texttt{sample=True}. By default (\texttt{None}) the draw comes from
  the session's random stream, set by \texttt{bf(rand\_seed=...)}, so
  consecutive draws differ but the whole sequence is reproducible. No
  effect when \texttt{sample=False}.
\item
  \emph{obs} (jnp.ndarray, optional): The observed value for this random
  variable. If provided, the sample site is conditioned on this value,
  and the function returns the observed value. If \texttt{None}, the
  site is treated as a latent (unobserved) random variable. Defaults to
  \texttt{None}.
\item
  \emph{name} (str, optional): The name of the sample site in a BF
  model. This is used to uniquely identify the random variable. Defaults
  to `x'.
\end{itemize}

\textbf{Returns} - When \texttt{sample=False}: A BF Normal distribution
object (for model building).

\begin{itemize}
\item
  When \texttt{sample=True}: A JAX array of samples drawn from the
  Normal distribution (for direct sampling).
\item
  When \texttt{create\_obj=True}: The raw BF distribution object (for
  advanced use cases).

  \textbf{Example Usage} from BayesForge import bf m =
  bf(platform=`cpu') m.dist.normal(loc=0.0, scale=1.0, sample=True)

  \textbf{Wrapper of}
  https://num.pyro.ai/en/stable/distributions.html\#normal

  \textbf{References}

  \begin{itemize}
  \tightlist
  \item
    \href{https://en.wikipedia.org/wiki/Normal_distribution\#Probability_density_function}{Wikipedia:
    Normal distribution -- PDF}
  \item
    \href{https://mathworld.wolfram.com/NormalDistribution.html}{Wolfram
    MathWorld: Normal Distribution}
  \end{itemize}
\end{itemize}

\subsection{\texorpdfstring{\texttt{m.dist.ordered\_logistic}}{m.dist.ordered\_logistic}}\label{ordered_logistic}

\begin{Shaded}
\begin{Highlighting}[]
\NormalTok{m.dist.ordered\_logistic(}
\NormalTok{    predictor,}
\NormalTok{    cutpoints,}
\NormalTok{    validate\_args}\OperatorTok{=}\VariableTok{None}\NormalTok{,}
\NormalTok{    name}\OperatorTok{=}\StringTok{\textquotesingle{}x\textquotesingle{}}\NormalTok{,}
\NormalTok{    obs}\OperatorTok{=}\VariableTok{None}\NormalTok{,}
\NormalTok{    mask}\OperatorTok{=}\VariableTok{None}\NormalTok{,}
\NormalTok{    infer}\OperatorTok{=}\VariableTok{None}\NormalTok{,}
\NormalTok{    sample}\OperatorTok{=}\VariableTok{False}\NormalTok{,}
\NormalTok{    seed}\OperatorTok{=}\VariableTok{None}\NormalTok{,}
\NormalTok{    shape}\OperatorTok{=}\NormalTok{(),}
\NormalTok{    event}\OperatorTok{=}\DecValTok{0}\NormalTok{,}
\NormalTok{    create\_obj}\OperatorTok{=}\VariableTok{False}\NormalTok{,}
\NormalTok{    to\_jax}\OperatorTok{=}\VariableTok{True}\NormalTok{,}
\NormalTok{)}
\end{Highlighting}
\end{Shaded}

\textbf{Ordered Logistic} The ordered logistic distribution is used for
modeling \textbf{ordinal} outcome variables \(Y \in {1,2,\dots,K}\)
(i.e., categories with a natural order) via a latent continuous
predictor \(\eta\) and a set of increasing \emph{cut-points}
\(c_1 < c_2 < \cdots < c_{K-1}\). When \(\eta\) crosses thresholds, the
observed \(Y\)

\textbf{PMF (Probability Mass Function)} Let \(\eta \in \mathbb R\) be
the linear predictor (for example \(\eta = x^\top \beta\)), let
cutpoints \(c = (c_1, \dots,c_{K-1})\) satisfy
\(c_1 < c_2 < \cdots < c_{K-1}\). Define also \(c_0 = -\infty\),
\(c_K = +\infty\). Then for (k \in {1,\dots,K}), \[
  \Pr(Y = k \mid \eta, c) =
  \begin{cases}
  1 - \mathrm{logit}^{-1}(\eta - c_1), & k = 1, [6pt]
  \mathrm{logit}^{-1}(\eta - c_{k-1}) - \mathrm{logit}^{-1}(\eta - c_{k}), & 1 < k < K, [6pt]
  \mathrm{logit}^{-1}(\eta - c_{K-1}), & k = K.
  \end{cases}
  \] Here \(\mathrm{logit}^{-1}(z) = \frac{1}{1 + e^{-z}}\).
(\href{https://statproofbook.github.io/P/matn-pdf.html}{Stan}) In other
words: \[
  \Pr(Y \le k \mid \eta, c) = \mathrm{logit}^{-1}(\eta - c_k),
  \] and so \[
  \Pr(Y = k) = \Pr(Y \le k) - \Pr(Y \le k-1).
  \] This formulation appears in e.g.~the CDF‐link setup of the
proportional odds model.

\textbf{Args} - \emph{predictor} (jnp.ndarray): Prediction in real
domain; typically this is output of a linear model.

\begin{itemize}
\item
  \emph{cutpoints} (jnp.ndarray): Positions in real domain to separate
  categories.
\item
  \emph{shape} (tuple): A multi-purpose argument for shaping. When
  \texttt{sample=False} (model building), this is used with
  \texttt{shape} to set the distribution's batch shape. When
  \texttt{sample=True} (direct sampling), this is used as
  \texttt{sample\_shape} to draw a raw JAX array of the given shape.
\item
  \emph{event} (int): The number of batch dimensions to reinterpret as
  event dimensions (used in model building).
\item
  \emph{mask} (jnp.ndarray, bool): Optional boolean array to mask
  observations.
\item
  \emph{create\_obj} (bool): If True, returns the raw BF distribution
  object instead of creating a sample site. This is essential for
  building complex distributions like \texttt{MixtureSameFamily}.
\item
  \emph{sample} (bool, optional): A control-flow argument. If
  \texttt{True}, the function will directly sample a raw JAX array from
  the distribution, bypassing the BF model context. If \texttt{False},
  it will create a \texttt{BF.sample} site within a model. Defaults to
  \texttt{False}.
\item
  \emph{seed} (int, optional): An integer seed for this draw when
  \texttt{sample=True}. By default (\texttt{None}) the draw comes from
  the session's random stream, set by \texttt{bf(rand\_seed=...)}, so
  consecutive draws differ but the whole sequence is reproducible. No
  effect when \texttt{sample=False}.
\item
  \emph{obs} (jnp.ndarray, optional): The observed value for this random
  variable. If provided, the sample site is conditioned on this value,
  and the function returns the observed value. If \texttt{None}, the
  site is treated as a latent (unobserved) random variable. Defaults to
  \texttt{None}.
\item
  \emph{name} (str, optional): The name of the sample site in a BF
  model. This is used to uniquely identify the random variable. Defaults
  to `x'.
\end{itemize}

\textbf{Returns} - When \texttt{sample=False}: A BF OrderedLogistic
distribution object (for model building). - When \texttt{sample=True}: A
JAX array of samples drawn from the OrderedLogistic distribution (for
direct sampling). - When \texttt{create\_obj=True}: The raw BF
distribution object (for advanced use cases).

\textbf{Example Usage} from BayesForge import bf import jax.numpy as jnp
m = bf(platform=`cpu')
m.dist.ordered\_logistic(predictor=jnp.array({[}0.2, 0.5, 0.8{]}),
cutpoints=jnp.array({[}-1.0, 0.0, 1.0{]}), sample=True)

\textbf{Wrapper of}
https://num.pyro.ai/en/stable/distributions.html\#orderedlogistic

\textbf{Reference} - \url{https://en.wikipedia.org/wiki/Ordered_logit} -
\url{https://mc-stan.org/docs/2_21/functions-reference/ordered-logistic-distribution.html}
-
\href{https://labdisia.disia.unifi.it/grilli/files/Papers/Ordered_Logit.pdf\%22}{https://labdisia.disia.unifi.it/grilli/files/Papers/Ordered\_Logit.pdf''}

\subsection{\texorpdfstring{\texttt{m.dist.pareto}}{m.dist.pareto}}\label{pareto}

\begin{Shaded}
\begin{Highlighting}[]
\NormalTok{m.dist.pareto(}
\NormalTok{    scale,}
\NormalTok{    alpha,}
\NormalTok{    validate\_args}\OperatorTok{=}\VariableTok{None}\NormalTok{,}
\NormalTok{    name}\OperatorTok{=}\StringTok{\textquotesingle{}x\textquotesingle{}}\NormalTok{,}
\NormalTok{    obs}\OperatorTok{=}\VariableTok{None}\NormalTok{,}
\NormalTok{    mask}\OperatorTok{=}\VariableTok{None}\NormalTok{,}
\NormalTok{    infer}\OperatorTok{=}\VariableTok{None}\NormalTok{,}
\NormalTok{    sample}\OperatorTok{=}\VariableTok{False}\NormalTok{,}
\NormalTok{    seed}\OperatorTok{=}\VariableTok{None}\NormalTok{,}
\NormalTok{    shape}\OperatorTok{=}\NormalTok{(),}
\NormalTok{    event}\OperatorTok{=}\DecValTok{0}\NormalTok{,}
\NormalTok{    create\_obj}\OperatorTok{=}\VariableTok{False}\NormalTok{,}
\NormalTok{    to\_jax}\OperatorTok{=}\VariableTok{True}\NormalTok{,}
\NormalTok{)}
\end{Highlighting}
\end{Shaded}

\textbf{Pareto} Samples from a Pareto distribution.

The \textbf{Pareto distribution}, named after economist Vilfredo Pareto,
is a \textbf{power-law} probability distribution used to describe
phenomena with ``rich-get-richer'' or ``heavy-tail'' properties --- for
example, income distribution, city sizes, or wealth concentration. It is
characterized by:

\begin{itemize}
\tightlist
\item
  a \textbf{scale parameter} \(x_m > 0\) (the minimum possible value),
  and
\item
  a \textbf{shape parameter} \(\alpha > 0\) (which controls the tail
  heaviness).
\end{itemize}

A random variable \(X \sim \text{Pareto}(\alpha, x_m)\) takes values
\(x \ge x_m\).

\[
  f(x \mid \alpha, x_m) =
  \begin{cases}
  \dfrac{\alpha x_m^\alpha}{x^{\alpha + 1}}, & x \ge x_m, \\
  0, & x < x_m.
  \end{cases}
  \]

\textbf{Args} - \emph{scale} (jnp.ndarray or float): Scale parameter of
the Pareto distribution. Must be positive.

\begin{itemize}
\item
  \emph{alpha} (jnp.ndarray or float): Shape parameter of the Pareto
  distribution. Must be positive.
\item
  \emph{shape} (tuple): A multi-purpose argument for shaping. When
  \texttt{sample=False} (model building), this is used with
  \texttt{shape} to set the distribution's batch shape. When
  \texttt{sample=True} (direct sampling), this is used as
  \texttt{sample\_shape} to draw a raw JAX array of the given shape.
\item
  \emph{event} (int): The number of batch dimensions to reinterpret as
  event dimensions (used in model building).
\item
  \emph{mask} (jnp.ndarray, bool): Optional boolean array to mask
  observations.
\item
  \emph{create\_obj} (bool, optional): If True, returns the raw BF
  distribution object instead of creating a sample site. This is
  essential for building complex distributions like
  \texttt{MixtureSameFamily}. Defaults to False.
\item
  \emph{sample} (bool, optional): A control-flow argument. If
  \texttt{True}, the function will directly sample a raw JAX array from
  the distribution, bypassing the BF model context. If \texttt{False},
  it will create a \texttt{BF.sample} site within a model. Defaults to
  \texttt{False}.
\item
  \emph{seed} (int, optional): An integer seed for this draw when
  \texttt{sample=True}. By default (\texttt{None}) the draw comes from
  the session's random stream, set by \texttt{bf(rand\_seed=...)}, so
  consecutive draws differ but the whole sequence is reproducible. No
  effect when \texttt{sample=False}.
\item
  \emph{obs} (jnp.ndarray, optional): The observed value for this random
  variable. If provided, the sample site is conditioned on this value,
  and the function returns the observed value. If \texttt{None}, the
  site is treated as a latent (unobserved) random variable. Defaults to
  \texttt{None}.
\item
  \emph{name} (str, optional): The name of the sample site in a BF
  model. This is used to uniquely identify the random variable. Defaults
  to `x'.
\end{itemize}

\textbf{Returns} - When \texttt{sample=False}: A BF Pareto distribution
object (for model building).

\begin{itemize}
\item
  When \texttt{sample=True}: A JAX array of samples drawn from the
  Pareto distribution (for direct sampling).
\item
  When \texttt{create\_obj=True}: The raw BF distribution object (for
  advanced use cases).

  \textbf{Example Usage} from BayesForge import bf m =
  bf(platform=`cpu') m.dist.pareto(scale=2.0, alpha=3.0, sample=True)

  \textbf{Wrapper of}
  https://num.pyro.ai/en/stable/distributions.html\#pareto

  \textbf{Reference}

  \begin{itemize}
  \tightlist
  \item
    \href{https://en.wikipedia.org/wiki/Pareto_distribution\#Probability_density_function}{Wikipedia:
    Pareto distribution -- PDF}
  \item
    Arnold, B. C. (2015). \emph{Pareto Distributions}, Second Edition.
    CRC Press.
  \end{itemize}
\end{itemize}

\subsection{\texorpdfstring{\texttt{m.dist.plate}}{m.dist.plate}}\label{plate}

\begin{Shaded}
\begin{Highlighting}[]
\NormalTok{m.dist.plate(name, shape, }\OperatorTok{**}\NormalTok{kwargs)}
\end{Highlighting}
\end{Shaded}

Context manager declaring conditionally independent replicates.

Sample sites inside \texttt{with\ m.dist.plate("obs",\ N):} are
vectorised over a batch dimension of size \texttt{N} (e.g.~one
likelihood term per observation).

\begin{longtable}[]{@{}
  >{\raggedright\arraybackslash}p{(\columnwidth - 4\tabcolsep) * \real{0.3333}}
  >{\raggedright\arraybackslash}p{(\columnwidth - 4\tabcolsep) * \real{0.3333}}
  >{\raggedright\arraybackslash}p{(\columnwidth - 4\tabcolsep) * \real{0.3333}}@{}}
\toprule\noalign{}
\begin{minipage}[b]{\linewidth}\raggedright
Argument
\end{minipage} & \begin{minipage}[b]{\linewidth}\raggedright
Type
\end{minipage} & \begin{minipage}[b]{\linewidth}\raggedright
Description
\end{minipage} \\
\midrule\noalign{}
\endhead
\bottomrule\noalign{}
\endlastfoot
\texttt{name} & str & Name of the plate. \\
\texttt{shape} & int & Size of the batch dimension. \\
\texttt{**kwargs} & & Passed to \texttt{numpyro.plate}
(e.g.~\texttt{subsample\_size}, \texttt{dim}). \\
\end{longtable}

\textbf{Returns}

\begin{itemize}
\tightlist
\item
  A NumPyro \texttt{plate}, used in a \texttt{with} statement.
\end{itemize}

\subsection{\texorpdfstring{\texttt{m.dist.poisson}}{m.dist.poisson}}\label{poisson}

\begin{Shaded}
\begin{Highlighting}[]
\NormalTok{m.dist.poisson(}
\NormalTok{    rate,}
\NormalTok{    is\_sparse}\OperatorTok{=}\VariableTok{False}\NormalTok{,}
\NormalTok{    validate\_args}\OperatorTok{=}\VariableTok{None}\NormalTok{,}
\NormalTok{    name}\OperatorTok{=}\StringTok{\textquotesingle{}x\textquotesingle{}}\NormalTok{,}
\NormalTok{    obs}\OperatorTok{=}\VariableTok{None}\NormalTok{,}
\NormalTok{    mask}\OperatorTok{=}\VariableTok{None}\NormalTok{,}
\NormalTok{    infer}\OperatorTok{=}\VariableTok{None}\NormalTok{,}
\NormalTok{    sample}\OperatorTok{=}\VariableTok{False}\NormalTok{,}
\NormalTok{    seed}\OperatorTok{=}\VariableTok{None}\NormalTok{,}
\NormalTok{    shape}\OperatorTok{=}\NormalTok{(),}
\NormalTok{    event}\OperatorTok{=}\DecValTok{0}\NormalTok{,}
\NormalTok{    create\_obj}\OperatorTok{=}\VariableTok{False}\NormalTok{,}
\NormalTok{    to\_jax}\OperatorTok{=}\VariableTok{True}\NormalTok{,}
\NormalTok{)}
\end{Highlighting}
\end{Shaded}

\textbf{Poisson} The \textbf{Poisson distribution} models the
probability of observing a given number of events \(k\) occurring in a
fixed interval of time or space when these events happen independently
and at a constant average rate \(\lambda > 0\). It is widely used for
modeling \textbf{count data}, such as the number of emails received per
hour or mutations in a DNA strand per unit length.

Formally, \[
  K \sim \text{Poisson}(\lambda)
  \] where \(\lambda\) is both the \textbf{mean} and \textbf{variance}
of the distribution.

The probability mass function (PMF) of the Poisson distribution is given
by: \[
  P(K = k \mid \lambda) =
  \frac{e^{-\lambda} \lambda^k}{k!},
  \quad k = 0, 1, 2, \dots
  \]

\textbf{Args} - \emph{rate} (jnp.ndarray): The rate parameter,
representing the average number of events.

\begin{itemize}
\item
  \emph{is\_sparse} (bool, optional): Indicates whether the
  \texttt{rate} parameter is sparse. If \texttt{True}, a specialized
  sparse sampling implementation is used, which can be more efficient
  for models with many zero-rate components (e.g., zero-inflated
  models). Defaults to \texttt{False}.
\item
  \emph{shape} (tuple): A multi-purpose argument for shaping. When
  \texttt{sample=False} (model building), this is used with
  \texttt{shape} to set the distribution's batch shape. When
  \texttt{sample=True} (direct sampling), this is used as
  \texttt{sample\_shape} to draw a raw JAX array of the given shape.
\item
  \emph{event} (int): The number of batch dimensions to reinterpret as
  event dimensions (used in model building).
\item
  \emph{mask} (jnp.ndarray, bool): Optional boolean array to mask
  observations.
\item
  \emph{create\_obj} (bool): If True, returns the raw BF distribution
  object instead of creating a sample site. This is essential for
  building complex distributions like \texttt{MixtureSameFamily}.
\item
  \emph{sample} (bool, optional): A control-flow argument. If
  \texttt{True}, the function will directly sample a raw JAX array from
  the distribution, bypassing the BF model context. If \texttt{False},
  it will create a \texttt{BF.sample} site within a model. Defaults to
  \texttt{False}.
\item
  \emph{seed} (int, optional): An integer seed for this draw when
  \texttt{sample=True}. By default (\texttt{None}) the draw comes from
  the session's random stream, set by \texttt{bf(rand\_seed=...)}, so
  consecutive draws differ but the whole sequence is reproducible. No
  effect when \texttt{sample=False}.
\item
  \emph{obs} (jnp.ndarray, optional): The observed value for this random
  variable. If provided, the sample site is conditioned on this value,
  and the function returns the observed value. If \texttt{None}, the
  site is treated as a latent (unobserved) random variable. Defaults to
  \texttt{None}.
\item
  \emph{name} (str, optional): The name of the sample site in a BF
  model. This is used to uniquely identify the random variable. Defaults
  to `x'.
\end{itemize}

\textbf{Returns} - When \texttt{sample=False}: A BF Poisson distribution
object (for model building). - When \texttt{sample=True}: A JAX array of
samples drawn from the Poisson distribution (for direct sampling). -
When \texttt{create\_obj=True}: The raw BF distribution object (for
advanced use cases).

\textbf{Example Usage} from BayesForge import bf m = bf(platform=`cpu')
m.dist.poisson(rate=2.0, sample=True)

\textbf{Wrapper of:
https://num.pyro.ai/en/stable/distributions.html\#poisson}
\textbf{References} *
\href{https://en.wikipedia.org/wiki/Poisson_distribution\#Probability_mass_function}{Wikipedia:
Poisson distribution -- PMF} * Johnson, N. L., Kotz, S., \& Kemp, A. W.
(1992). \emph{Univariate Discrete Distributions} (2nd ed.). Wiley.

\subsection{\texorpdfstring{\texttt{m.dist.projected\_normal}}{m.dist.projected\_normal}}\label{projected_normal}

\begin{Shaded}
\begin{Highlighting}[]
\NormalTok{m.dist.projected\_normal(}
\NormalTok{    concentration,}
\NormalTok{    validate\_args}\OperatorTok{=}\VariableTok{None}\NormalTok{,}
\NormalTok{    name}\OperatorTok{=}\StringTok{\textquotesingle{}x\textquotesingle{}}\NormalTok{,}
\NormalTok{    obs}\OperatorTok{=}\VariableTok{None}\NormalTok{,}
\NormalTok{    mask}\OperatorTok{=}\VariableTok{None}\NormalTok{,}
\NormalTok{    infer}\OperatorTok{=}\VariableTok{None}\NormalTok{,}
\NormalTok{    sample}\OperatorTok{=}\VariableTok{False}\NormalTok{,}
\NormalTok{    seed}\OperatorTok{=}\VariableTok{None}\NormalTok{,}
\NormalTok{    shape}\OperatorTok{=}\NormalTok{(),}
\NormalTok{    event}\OperatorTok{=}\DecValTok{0}\NormalTok{,}
\NormalTok{    create\_obj}\OperatorTok{=}\VariableTok{False}\NormalTok{,}
\NormalTok{    to\_jax}\OperatorTok{=}\VariableTok{True}\NormalTok{,}
\NormalTok{)}
\end{Highlighting}
\end{Shaded}

\textbf{Projected Normal} The projected normal distribution arises by
taking a multivariate normal vector
\(\mathbf X \sim \mathcal N_n (\mu, \Sigma)\) in \(\mathbb R^n\) and
projecting it to the unit sphere. This distribution is commonly used in
directional statistics (data on circles or spheres) and supports
asymmetric and even multimodal behaviours depending on parameters.

\textbf{Args} - \emph{concentration} (jnp.ndarray): The concentration
parameter, representing the direction towards which the samples are
concentrated. Must be a JAX array with at least one dimension.

\begin{itemize}
\item
  \emph{shape} (tuple): A multi-purpose argument for shaping. When
  \texttt{sample=False} (model building), this is used with
  \texttt{shape} to set the distribution's batch shape. When
  \texttt{sample=True} (direct sampling), this is used as
  \texttt{sample\_shape} to draw a raw JAX array of the given shape.
\item
  \emph{event} (int): The number of batch dimensions to reinterpret as
  event dimensions (used in model building).
\item
  \emph{mask} (jnp.ndarray, bool): Optional boolean array to mask
  observations.
\item
  \emph{create\_obj} (bool, optional): If True, returns the raw BF
  distribution object instead of creating a sample site. This is
  essential for building complex distributions like
  \texttt{MixtureSameFamily}. Defaults to False.
\item
  \emph{sample} (bool, optional): A control-flow argument. If
  \texttt{True}, the function will directly sample a raw JAX array from
  the distribution, bypassing the BF model context. If \texttt{False},
  it will create a \texttt{BF.sample} site within a model. Defaults to
  \texttt{False}.
\item
  \emph{seed} (int, optional): An integer seed for this draw when
  \texttt{sample=True}. By default (\texttt{None}) the draw comes from
  the session's random stream, set by \texttt{bf(rand\_seed=...)}, so
  consecutive draws differ but the whole sequence is reproducible. No
  effect when \texttt{sample=False}.
\item
  \emph{obs} (jnp.ndarray, optional): The observed value for this random
  variable. If provided, the sample site is conditioned on this value,
  and the function returns the observed value. If \texttt{None}, the
  site is treated as a latent (unobserved) random variable. Defaults to
  \texttt{None}.
\item
  \emph{name} (str, optional): The name of the sample site in a BF
  model. This is used to uniquely identify the random variable. Defaults
  to `x'.
\end{itemize}

\textbf{Returns} - When \texttt{sample=False}: A BF ProjectedNormal
distribution object (for model building).

\begin{itemize}
\item
  When \texttt{sample=True}: A JAX array of samples drawn from the
  ProjectedNormal distribution (for direct sampling).
\item
  When \texttt{create\_obj=True}: The raw BF distribution object (for
  advanced use cases).

  \textbf{Example Usage} from BayesForge import bf import jax.numpy as
  jnp m = bf(platform=`cpu')
  m.dist.projected\_normal(concentration=jnp.array({[}1.0, 3.0, 2.0{]}),
  sample=True)

  \textbf{Wrapper of}
  https://num.pyro.ai/en/stable/distributions.html\#projectednormal
\end{itemize}

\subsection{\texorpdfstring{\texttt{m.dist.relaxed\_bernoulli}}{m.dist.relaxed\_bernoulli}}\label{relaxed_bernoulli}

\begin{Shaded}
\begin{Highlighting}[]
\NormalTok{m.dist.relaxed\_bernoulli(}
\NormalTok{    temperature,}
\NormalTok{    probs}\OperatorTok{=}\VariableTok{None}\NormalTok{,}
\NormalTok{    logits}\OperatorTok{=}\VariableTok{None}\NormalTok{,}
\NormalTok{    validate\_args}\OperatorTok{=}\VariableTok{None}\NormalTok{,}
\NormalTok{    name}\OperatorTok{=}\StringTok{\textquotesingle{}x\textquotesingle{}}\NormalTok{,}
\NormalTok{    obs}\OperatorTok{=}\VariableTok{None}\NormalTok{,}
\NormalTok{    mask}\OperatorTok{=}\VariableTok{None}\NormalTok{,}
\NormalTok{    infer}\OperatorTok{=}\VariableTok{None}\NormalTok{,}
\NormalTok{    sample}\OperatorTok{=}\VariableTok{False}\NormalTok{,}
\NormalTok{    seed}\OperatorTok{=}\VariableTok{None}\NormalTok{,}
\NormalTok{    shape}\OperatorTok{=}\NormalTok{(),}
\NormalTok{    event}\OperatorTok{=}\DecValTok{0}\NormalTok{,}
\NormalTok{    create\_obj}\OperatorTok{=}\VariableTok{False}\NormalTok{,}
\NormalTok{    to\_jax}\OperatorTok{=}\VariableTok{True}\NormalTok{,}
\NormalTok{)}
\end{Highlighting}
\end{Shaded}

\textbf{Relaxed Bernoulli} The Relaxed\,Bernoulli is a continuous
distribution on the interval \((0,1)\) that smoothly approximates the
discrete Bernoulli distribution (which has support \({0,1}\)). It was
introduced to allow for \emph{differentiable} sampling of approximate
binary random variables, which is useful in variational inference and
other gradient‐based optimization settings.

The probability density function (PDF) is defined as:

\begin{itemize}
\tightlist
\item
  It has a \emph{temperature} parameter \(\tau > 0\) controlling the
  smoothness: as \(\tau \to 0\), the distribution concentrates near
  \({0, 1}\) (thus approximating the discrete Bernoulli).
\item
  For larger \(\tau\), the sample becomes more ``soft'' and less binary,
  e.g., converging toward 0.5 for very high \(\tau\).
\item
  It is parameterised by a probability \(p\) (or equivalently logits
  \(\ell = \log\frac{p}{1-p}\)).
\end{itemize}

Let \(X \in (0,1)\) follow a Relaxed\,Bernoulli (Binary\,Concrete)
distribution with location parameter \(\alpha > 0\) (sometimes expressed
via logits) and temperature parameter \(\lambda > 0\). Then the PDF is:

\[
p_{,\alpha,;\lambda}(x) =
\frac{;\lambda,\alpha,x^{-\lambda - 1},(1 - x)^{-\lambda - 1};}
{\Bigl(\alpha,x^{-\lambda} + (1 - x)^{-\lambda}\Bigr)^{2}},
\quad \text{for } 0 < x < 1.
\]

\textbf{Args} - \emph{temperature} (jnp.ndarray): The temperature
parameter. Must be greater than 0.

\begin{itemize}
\item
  \emph{probs} (jnp.ndarray, optional): The probability of success. Must
  be in the interval \texttt{{[}0,\ 1{]}}. Only one of \texttt{probs} or
  \texttt{logits} can be specified.
\item
  \emph{logits} (jnp.ndarray, optional): The log-odds of success. Must
  be in the interval \texttt{{[}-inf,\ inf{]}}. Only one of
  \texttt{probs} or \texttt{logits} can be specified.
\item
  \emph{shape} (tuple): A multi-purpose argument for shaping. When
  \texttt{sample=False} (model building), this is used with
  \texttt{shape} to set the distribution's batch shape. When
  \texttt{sample=True} (direct sampling), this is used as
  \texttt{sample\_shape} to draw a raw JAX array of the given shape.
\item
  \emph{event} (int): The number of batch dimensions to reinterpret as
  event dimensions (used in model building).
\item
  \emph{mask} (jnp.ndarray, bool): Optional boolean array to mask
  observations.
\item
  \emph{create\_obj} (bool): If True, returns the raw BF distribution
  object instead of creating a sample site. This is essential for
  building complex distributions like \texttt{MixtureSameFamily}.
\item
  \emph{sample} (bool, optional): A control-flow argument. If
  \texttt{True}, the function will directly sample a raw JAX array from
  the distribution, bypassing the BF model context. If \texttt{False},
  it will create a \texttt{BF.sample} site within a model. Defaults to
  \texttt{False}.
\item
  \emph{seed} (int, optional): An integer seed for this draw when
  \texttt{sample=True}. By default (\texttt{None}) the draw comes from
  the session's random stream, set by \texttt{bf(rand\_seed=...)}, so
  consecutive draws differ but the whole sequence is reproducible. No
  effect when \texttt{sample=False}.
\item
  \emph{obs} (jnp.ndarray, optional): The observed value for this random
  variable. If provided, the sample site is conditioned on this value,
  and the function returns the observed value. If \texttt{None}, the
  site is treated as a latent (unobserved) random variable. Defaults to
  \texttt{None}.
\item
  \emph{name} (str, optional): The name of the sample site in a BF
  model. This is used to uniquely identify the random variable. Defaults
  to `x'.
\end{itemize}

\textbf{Returns} BF RelaxedBernoulli distribution object (for model
building) when \texttt{sample=False}. A JAX array of samples drawn from
the RelaxedBernoulli distribution (for direct sampling) when
\texttt{sample=True}. The raw BF distribution object (for advanced use
cases) when \texttt{create\_obj=True}.

\textbf{Example Usage} from BayesForge import bf import jax.numpy as jnp
m = bf(platform=`cpu') m.dist.relaxed\_bernoulli(temperature=1.0, probs
= jnp.array({[}0.2, 0.3, 0.5{]}), sample=True)

\textbf{Wrapper of:
https://num.pyro.ai/en/stable/distributions.html\#relaxedbernoulli}
\textbf{References} -
\href{https://www.tensorflow.org/probability/api_docs/python/tfp/distributions/RelaxedBernoulli}{tfp.distributions.RelaxedBernoulli
\textbar{} TensorFlow Probability} -
\href{https://rstudio.github.io/tfprobability/reference/tfd_relaxed_bernoulli.html}{RelaxedBernoulli
distribution with temperature and logits \ldots{}}

\subsection{\texorpdfstring{\texttt{m.dist.relaxed\_bernoulli\_logits}}{m.dist.relaxed\_bernoulli\_logits}}\label{relaxed_bernoulli_logits}

\begin{Shaded}
\begin{Highlighting}[]
\NormalTok{m.dist.relaxed\_bernoulli\_logits(}
\NormalTok{    temperature,}
\NormalTok{    logits,}
\NormalTok{    validate\_args}\OperatorTok{=}\VariableTok{None}\NormalTok{,}
\NormalTok{    name}\OperatorTok{=}\StringTok{\textquotesingle{}x\textquotesingle{}}\NormalTok{,}
\NormalTok{    obs}\OperatorTok{=}\VariableTok{None}\NormalTok{,}
\NormalTok{    mask}\OperatorTok{=}\VariableTok{None}\NormalTok{,}
\NormalTok{    infer}\OperatorTok{=}\VariableTok{None}\NormalTok{,}
\NormalTok{    sample}\OperatorTok{=}\VariableTok{False}\NormalTok{,}
\NormalTok{    seed}\OperatorTok{=}\VariableTok{None}\NormalTok{,}
\NormalTok{    shape}\OperatorTok{=}\NormalTok{(),}
\NormalTok{    event}\OperatorTok{=}\DecValTok{0}\NormalTok{,}
\NormalTok{    create\_obj}\OperatorTok{=}\VariableTok{False}\NormalTok{,}
\NormalTok{    to\_jax}\OperatorTok{=}\VariableTok{True}\NormalTok{,}
\NormalTok{)}
\end{Highlighting}
\end{Shaded}

\textbf{Relaxed Bernoulli Logits} The Relaxed Bernoulli (logits) is a
\textbf{continuous} relaxation of the standard Bernoulli distribution,
parameterised by logits (or probabilities) and a \textbf{temperature}
parameter. Rather than producing strictly 0 or 1, it produces values in
the continuous interval (0, 1). As the temperature → 0 the distribution
approximates a Bernoulli; as temperature → \(\infty\) the distribution
approximates a uniform distribution.

It is used in variational inference and deep-learning contexts to allow
gradient-based optimisation through otherwise discrete Bernoulli draws
(see the ``Concrete'' or ``Gumbel-Softmax'' literature: The Concrete
Distribution: A Continuous Relaxation of Discrete Random Variables and
Categorical Reparameterization with Gumbel‑Softmax ).

Because this is a continuous relaxation, the ``density'' is defined over
the unit interval (0,1) rather than a pmf at \{0,1\}. Using logistic or
Gumbel‐softmax style construction one can express it as:

If logits = ( \ell ), and temperature = ( \tau \textgreater{} 0 ), then
one generates \[
u \sim \mathrm{Uniform}(0,1),\quad
g = -\log(-\log(u)),\quad
x = \sigma\big((\ell + g)/\tau\big)
\] where \(\sigma(\cdot)\) is the logistic (sigmoid) function. Then
\(x \in (0,1)\) has the RelaxedBernoulli distribution.

The exact density formula can be derived via change of variables from
logistic/Gumbel, but is somewhat involved (including Jacobian of sigmoid
transform).

\textbf{Args} - \emph{temperature} (jnp.ndarray): The temperature
parameter, must be positive.

\begin{itemize}
\item
  \emph{logits} (jnp.ndarray): The logits parameter.
\item
  \emph{shape} (tuple): A multi-purpose argument for shaping. When
  \texttt{sample=False} (model building), this is used with
  \texttt{shape} to set the distribution's batch shape. When
  \texttt{sample=True} (direct sampling), this is used as
  \texttt{sample\_shape} to draw a raw JAX array of the given shape.
\item
  \emph{event} (int): The number of batch dimensions to reinterpret as
  event dimensions (used in model building).
\item
  \emph{mask} (jnp.ndarray, bool): Optional boolean array to mask
  observations.
\item
  \emph{create\_obj} (bool): If True, returns the raw BF distribution
  object instead of creating a sample site. This is essential for
  building complex distributions like \texttt{MixtureSameFamily}.
\item
  \emph{sample} (bool, optional): A control-flow argument. If
  \texttt{True}, the function will directly sample a raw JAX array from
  the distribution, bypassing the BF model context. If \texttt{False},
  it will create a \texttt{BF.sample} site within a model. Defaults to
  \texttt{False}.
\item
  \emph{seed} (int, optional): An integer seed for this draw when
  \texttt{sample=True}. By default (\texttt{None}) the draw comes from
  the session's random stream, set by \texttt{bf(rand\_seed=...)}, so
  consecutive draws differ but the whole sequence is reproducible. No
  effect when \texttt{sample=False}.
\item
  \emph{obs} (jnp.ndarray, optional): The observed value for this random
  variable. If provided, the sample site is conditioned on this value,
  and the function returns the observed value. If \texttt{None}, the
  site is treated as a latent (unobserved) random variable. Defaults to
  \texttt{None}.
\item
  \emph{name} (str, optional): The name of the sample site in a BF
  model. This is used to uniquely identify the random variable. Defaults
  to `x'.
\end{itemize}

\textbf{Returns} RelaxedBernoulliLogits: A BF RelaxedBernoulliLogits
distribution object (for model building). jnp.ndarray: A JAX array of
samples drawn from the RelaxedBernoulliLogits distribution (for direct
sampling). RelaxedBernoulliLogits: The raw BF distribution object (for
advanced use cases).

\textbf{Example Usage} from BayesForge import bf m = bf(platform=`cpu')
m.dist.relaxed\_bernoulli\_logits(temperature=1.0, logits=0.0,
sample=True)

\textbf{Wrapper of}
https://num.pyro.ai/en/stable/distributions.html\#relaxed-bernoulli-logits

\textbf{Reference} * https://arxiv.org/pdf/1611.00712

\subsection{\texorpdfstring{\texttt{m.dist.right\_truncated\_distribution}}{m.dist.right\_truncated\_distribution}}\label{right_truncated_distribution}

\begin{Shaded}
\begin{Highlighting}[]
\NormalTok{m.dist.right\_truncated\_distribution(}
\NormalTok{    base\_dist,}
\NormalTok{    high}\OperatorTok{=}\FloatTok{0.0}\NormalTok{,}
\NormalTok{    validate\_args}\OperatorTok{=}\VariableTok{None}\NormalTok{,}
\NormalTok{    name}\OperatorTok{=}\StringTok{\textquotesingle{}x\textquotesingle{}}\NormalTok{,}
\NormalTok{    obs}\OperatorTok{=}\VariableTok{None}\NormalTok{,}
\NormalTok{    mask}\OperatorTok{=}\VariableTok{None}\NormalTok{,}
\NormalTok{    infer}\OperatorTok{=}\VariableTok{None}\NormalTok{,}
\NormalTok{    sample}\OperatorTok{=}\VariableTok{False}\NormalTok{,}
\NormalTok{    seed}\OperatorTok{=}\VariableTok{None}\NormalTok{,}
\NormalTok{    shape}\OperatorTok{=}\NormalTok{(),}
\NormalTok{    event}\OperatorTok{=}\DecValTok{0}\NormalTok{,}
\NormalTok{    create\_obj}\OperatorTok{=}\VariableTok{False}\NormalTok{,}
\NormalTok{    to\_jax}\OperatorTok{=}\VariableTok{True}\NormalTok{,}
\NormalTok{)}
\end{Highlighting}
\end{Shaded}

\textbf{Right Truncated} A right-truncated distribution is a statistical
distribution that arises when the possible values of a random variable
are restricted to be below a certain specified value \texttt{high}. In
essence, the right tail of the original distribution is ``cut off'' at a
particular point, and the remaining probability is redistributed among
the allowable values. This type of distribution is common in various
fields where there are inherent upper limits or observational
constraints.

The probability density function (PDF) for a continuous right-truncated
distribution or the probability mass function (PMF) for a discrete one
is derived from the original distribution by normalizing it over the
restricted range.

\textbf{For a continuous random variable X:}

If the original probability density function is \texttt{f(x)} and the
cumulative distribution function is \texttt{F(x)}, and the distribution
is right-truncated at a value \texttt{b} (meaning \texttt{x\ ≤\ b}), the
new PDF, \texttt{f\_T(x)}, is: \[
  f_{\rm trunc}(x \mid X \le b) =
  \begin{cases}
  \displaystyle \frac{,f(x),}{,F(b),}, & x \le b, \\
  0, & x > b.
  \end{cases}
  \]

\textbf{For a discrete random variable X:}

If the original probability mass function is \texttt{P(X\ =\ k)} and the
cumulative distribution function is \texttt{F(k)}, and the distribution
is right-truncated at a value \texttt{b} (meaning \texttt{k\ ≤\ b}), the
new PMF, \texttt{P\_T(X\ =\ k)}, is:

\[P_T(X=k) = \begin{cases} \frac{P(X=k)}{F(b)} & \text{if } k \le b \\ 0 & \text{if } k > b \end{cases}\]

\textbf{Args} - \emph{base\_dist}: The base distribution to truncate.
Must be a univariate distribution with real support.

\begin{itemize}
\item
  \emph{high} (float, jnp.ndarray, optional): The upper truncation
  point. The support of the new distribution is
  \((-\infty, \text{high}]\). Defaults to 0.0.
\item
  \emph{shape} (tuple): A multi-purpose argument for shaping. When
  \texttt{sample=False} (model building), this is used with
  \texttt{shape} to set the distribution's batch shape. When
  \texttt{sample=True} (direct sampling), this is used as
  \texttt{sample\_shape} to draw a raw JAX array of the given shape.
\item
  \emph{event} (int): The number of batch dimensions to reinterpret as
  event dimensions (used in model building).
\item
  \emph{mask} (jnp.ndarray, bool): Optional boolean array to mask
  observations.
\item
  \emph{create\_obj} (bool): If True, returns the raw BF distribution
  object instead of creating a sample site. This is essential for
  building complex distributions like \texttt{MixtureSameFamily}.
\item
  \emph{sample} (bool, optional): A control-flow argument. If
  \texttt{True}, the function will directly sample a raw JAX array from
  the distribution, bypassing the BF model context. If \texttt{False},
  it will create a \texttt{BF.sample} site within a model. Defaults to
  \texttt{False}.
\item
  \emph{seed} (int, optional): An integer seed for this draw when
  \texttt{sample=True}. By default (\texttt{None}) the draw comes from
  the session's random stream, set by \texttt{bf(rand\_seed=...)}, so
  consecutive draws differ but the whole sequence is reproducible. No
  effect when \texttt{sample=False}.
\item
  \emph{obs} (jnp.ndarray, optional): The observed value for this random
  variable. If provided, the sample site is conditioned on this value,
  and the function returns the observed value. If \texttt{None}, the
  site is treated as a latent (unobserved) random variable. Defaults to
  \texttt{None}.
\item
  \emph{name} (str, optional): The name of the sample site in a BF
  model. This is used to uniquely identify the random variable. Defaults
  to `x'.
\end{itemize}

\textbf{Returns} - When \texttt{sample=False}: A BF
RightTruncatedDistribution distribution object (for model building).

\begin{itemize}
\item
  When \texttt{sample=True}: A JAX array of samples drawn from the
  RightTruncatedDistribution distribution (for direct sampling).
\item
  When \texttt{create\_obj=True}: The raw BF distribution object (for
  advanced use cases).

  \textbf{Example Usage} from BayesForge import bf m =
  bf(platform=`cpu') m.dist.right\_truncated\_distribution(base\_dist =
  m.dist.normal(0,1, create\_obj = True), high=0, sample=True)

  \textbf{Wrapper of}
  https://num.pyro.ai/en/stable/distributions.html\#righttruncateddistribution

  \textbf{References}

  \begin{itemize}
  \tightlist
  \item
    https://en.wikipedia.org/wiki/Truncated\_normal\_distribution
  \end{itemize}
\end{itemize}

\subsection{\texorpdfstring{\texttt{m.dist.sine\_bivariate\_vonmises}}{m.dist.sine\_bivariate\_vonmises}}\label{sine_bivariate_vonmises}

\begin{Shaded}
\begin{Highlighting}[]
\NormalTok{m.dist.sine\_bivariate\_vonmises(}
\NormalTok{    phi\_loc,}
\NormalTok{    psi\_loc,}
\NormalTok{    phi\_concentration,}
\NormalTok{    psi\_concentration,}
\NormalTok{    correlation}\OperatorTok{=}\VariableTok{None}\NormalTok{,}
\NormalTok{    weighted\_correlation}\OperatorTok{=}\VariableTok{None}\NormalTok{,}
\NormalTok{    validate\_args}\OperatorTok{=}\VariableTok{None}\NormalTok{,}
\NormalTok{    name}\OperatorTok{=}\StringTok{\textquotesingle{}x\textquotesingle{}}\NormalTok{,}
\NormalTok{    obs}\OperatorTok{=}\VariableTok{None}\NormalTok{,}
\NormalTok{    mask}\OperatorTok{=}\VariableTok{None}\NormalTok{,}
\NormalTok{    infer}\OperatorTok{=}\VariableTok{None}\NormalTok{,}
\NormalTok{    sample}\OperatorTok{=}\VariableTok{False}\NormalTok{,}
\NormalTok{    seed}\OperatorTok{=}\VariableTok{None}\NormalTok{,}
\NormalTok{    shape}\OperatorTok{=}\NormalTok{(),}
\NormalTok{    event}\OperatorTok{=}\DecValTok{0}\NormalTok{,}
\NormalTok{    create\_obj}\OperatorTok{=}\VariableTok{False}\NormalTok{,}
\NormalTok{)}
\end{Highlighting}
\end{Shaded}

\textbf{Sine Bivariate Von Mises}
\href{https://en.wikipedia.org/wiki/Bivariate_von_Mises_distribution}{In
probability theory and statistics, the bivariate von Mises distribution
is a probability distribution describing values on a torus. It may be
thought of as an analogue on the torus of the bivariate normal
distribution. The distribution belongs to the field of directional
statistics. The general bivariate von Mises distribution was first
proposed by Kanti Mardia in 1975. One of its variants is today used in
the field of bioinformatics to formulate a probabilistic model of
protein structure in atomic detail, such as backbone-dependent rotamer
libraries.}

\textbf{Args} - \emph{phi\_loc} (jnp.ndarray): The location parameter
for the first angle (phi).

\begin{itemize}
\item
  \emph{psi\_loc} (jnp.ndarray): The location parameter for the second
  angle (psi).
\item
  \emph{phi\_concentration} (jnp.ndarray): The concentration parameter
  for the first angle (phi). Must be positive.
\item
  \emph{psi\_concentration} (jnp.ndarray): The concentration parameter
  for the second angle (psi). Must be positive.
\item
  \emph{correlation} (jnp.ndarray, optional): The correlation parameter
  between the two angles. One of \texttt{correlation} or
  \texttt{weighted\_correlation} must be specified.
\item
  \emph{weighted\_correlation} (jnp.ndarray, optional): An alternative
  correlation parameter. One of \texttt{correlation} or
  \texttt{weighted\_correlation} must be specified.
\item
  \emph{validate\_args} (bool, optional): Whether to enable validation
  of distribution parameters. Defaults to \texttt{None}.
\item
  \emph{name} (str, optional): The name of the sample site in a BF
  model. This is used to uniquely identify the random variable. Defaults
  to `x'.
\item
  \emph{obs} (jnp.ndarray, optional): The observed value for this random
  variable. If provided, the sample site is conditioned on this value,
  and the function returns the observed value. If \texttt{None}, the
  site is treated as a latent (unobserved) random variable. Defaults to
  \texttt{None}.
\item
  \emph{mask} (jnp.ndarray, bool, optional): Optional boolean array to
  mask observations. If provided, events with a \texttt{True} mask will
  be conditioned on \texttt{obs}, while the remaining events will be
  treated as latent variables. Defaults to \texttt{None}.
\item
  \emph{sample} (bool, optional): A control-flow argument. If
  \texttt{True}, the function will directly sample a raw JAX array from
  the distribution, bypassing the BF model context. If \texttt{False},
  it will create a \texttt{BF.sample} site within a model. Defaults to
  \texttt{False}.
\item
  \emph{seed} (int, optional): An integer seed for this draw when
  \texttt{sample=True}. By default (\texttt{None}) the draw comes from
  the session's random stream, set by \texttt{bf(rand\_seed=...)}, so
  consecutive draws differ but the whole sequence is reproducible. No
  effect when \texttt{sample=False}.
\item
  \emph{shape} (tuple, optional): A multi-purpose argument for shaping.
  When \texttt{sample=False} (model building), this is used with
  \texttt{.expand(shape)} to set the distribution's batch shape. When
  \texttt{sample=True} (direct sampling), this is used as
  \texttt{sample\_shape} to draw a raw JAX array of the given shape.
\item
  \emph{event} (int, optional): The number of batch dimensions to
  reinterpret as event dimensions (used in model building).
\item
  \emph{create\_obj} (bool, optional): If True, returns the raw BF
  distribution object instead of creating a sample site. This is
  essential for building complex distributions like
  \texttt{MixtureSameFamily}. Defaults to \texttt{False}.
\end{itemize}

\textbf{Returns} BF.primitives.Messenger: A BF sample site object when
used in a model context (\texttt{sample=False}). jnp.ndarray: A JAX
array of samples drawn from the SineBivariateVonMises distribution (for
direct sampling, \texttt{sample=True}).
numpyro.distributions.Distribution: The raw BF distribution object (if
\texttt{create\_obj=True}).

\textbf{Example Usage} from BayesForge import bf import jax.numpy as jnp
m = bf(platform=`cpu')

\begin{verbatim}
**Direct sampling**
samples = m.dist.sine_bivariate_vonmises(
    phi_loc=0.0,
    psi_loc=jnp.pi,
    phi_concentration=1.0,
    psi_concentration=1.0,
    correlation=0.5,
    sample=True,
    shape=(10,)
)

**Usage within a model**
def my_model():
    angles = m.dist.sine_bivariate_vonmises(
        phi_loc=0.0,
        psi_loc=0.0,
        phi_concentration=2.0,
        psi_concentration=2.0,
        weighted_correlation=0.9,
        name='angles'
    )
    **... rest of the model**
\end{verbatim}

\textbf{Wrapper of}
https://num.pyro.ai/en/stable/distributions.html\#sinebivariatevonmises

\subsection{\texorpdfstring{\texttt{m.dist.sine\_skewed}}{m.dist.sine\_skewed}}\label{sine_skewed}

\begin{Shaded}
\begin{Highlighting}[]
\NormalTok{m.dist.sine\_skewed(}
\NormalTok{    base\_dist: numpyro.distributions.distribution.Distribution,}
\NormalTok{    skewness,}
\NormalTok{    validate\_args}\OperatorTok{=}\VariableTok{None}\NormalTok{,}
\NormalTok{    name}\OperatorTok{=}\StringTok{\textquotesingle{}x\textquotesingle{}}\NormalTok{,}
\NormalTok{    obs}\OperatorTok{=}\VariableTok{None}\NormalTok{,}
\NormalTok{    mask}\OperatorTok{=}\VariableTok{None}\NormalTok{,}
\NormalTok{    infer}\OperatorTok{=}\VariableTok{None}\NormalTok{,}
\NormalTok{    sample}\OperatorTok{=}\VariableTok{False}\NormalTok{,}
\NormalTok{    seed}\OperatorTok{=}\VariableTok{None}\NormalTok{,}
\NormalTok{    shape}\OperatorTok{=}\NormalTok{(),}
\NormalTok{    event}\OperatorTok{=}\DecValTok{0}\NormalTok{,}
\NormalTok{    create\_obj}\OperatorTok{=}\VariableTok{False}\NormalTok{,}
\NormalTok{    to\_jax}\OperatorTok{=}\VariableTok{True}\NormalTok{,}
\NormalTok{)}
\end{Highlighting}
\end{Shaded}

\textbf{Sine-skewing} The sine-skewed von Mises distribution is an
extension of the symmetric circular (or toroidal) von Mises (or
bivariate von Mises) distribution to allow for skewness (asymmetry) via
a sine-based skewing function. It is used to model directional data on
the circle (or torus) that depart from
symmetry.\href{https://arxiv.org/abs/1902.02579}{1}

\textbf{Parameters:}

\begin{itemize}
\item
  \textbf{base\_dist:} Base density on a d-dimensional torus. Supported
  base distributions include: 1D
  \texttt{numpyro.distributions.VonMises},
  \texttt{numnumpyro.distributions.SineBivariateVonMises}, 1D
  \texttt{numpyro.distributions.ProjectedNormal}, and
  \texttt{numpyro.distributions.Uniform} (-pi, pi).
\item
  \textbf{skewness:} Skewness of the distribution.
\item
  \textbf{sample (bool, optional):} A control-flow argument. If
  \texttt{True}, the function will directly sample a raw JAX array from
  the distribution, bypassing the BF model context. If \texttt{False},
  it will create a \texttt{BF.sample} site within a model. Defaults to
  \texttt{False}.
\item
  \textbf{seed (int, optional):} An integer seed for this draw when
  \texttt{sample=True}. By default (\texttt{None}) the draw comes from
  the session's random stream, set by \texttt{bf(rand\_seed=...)}, so
  consecutive draws differ but the whole sequence is reproducible. No
  effect when \texttt{sample=False}.
\item
  \textbf{obs (jnp.ndarray, optional):} The observed value for this
  random variable. If provided, the sample site is conditioned on this
  value, and the function returns the observed value. If \texttt{None},
  the site is treated as a latent (unobserved) random variable. Defaults
  to \texttt{None}.
\item
  \textbf{name (str, optional):} The name of the sample site in a BF
  model. This is used to uniquely identify the random variable. Defaults
  to `x'.
\end{itemize}

\textbf{PDF:}

The probability density function (PDF) of the Sine Skewed X distribution
is not explicitly defined here, but it is derived from the base
distribution and the skewness parameter.

\textbf{Example Usage} from BayesForge import bf import jax.numpy as jnp
m = bf(platform=`cpu') \textbf{base distribution on the torus (two
angles)} base = m.dist.sine\_bivariate\_vonmises(phi\_loc=0.0,
psi\_loc=1.0, phi\_concentration=2.0, psi\_concentration=2.0,
correlation=0.5, create\_obj=True) m.dist.sine\_skewed(base\_dist=base,
skewness=jnp.array({[}0.3, -0.2{]}), sample=True, shape=(1000,))

\subsection{\texorpdfstring{\texttt{m.dist.soft\_laplace}}{m.dist.soft\_laplace}}\label{soft_laplace}

\begin{Shaded}
\begin{Highlighting}[]
\NormalTok{m.dist.soft\_laplace(}
\NormalTok{    loc,}
\NormalTok{    scale,}
\NormalTok{    validate\_args}\OperatorTok{=}\VariableTok{None}\NormalTok{,}
\NormalTok{    name}\OperatorTok{=}\StringTok{\textquotesingle{}x\textquotesingle{}}\NormalTok{,}
\NormalTok{    obs}\OperatorTok{=}\VariableTok{None}\NormalTok{,}
\NormalTok{    mask}\OperatorTok{=}\VariableTok{None}\NormalTok{,}
\NormalTok{    infer}\OperatorTok{=}\VariableTok{None}\NormalTok{,}
\NormalTok{    sample}\OperatorTok{=}\VariableTok{False}\NormalTok{,}
\NormalTok{    seed}\OperatorTok{=}\VariableTok{None}\NormalTok{,}
\NormalTok{    shape}\OperatorTok{=}\NormalTok{(),}
\NormalTok{    event}\OperatorTok{=}\DecValTok{0}\NormalTok{,}
\NormalTok{    create\_obj}\OperatorTok{=}\VariableTok{False}\NormalTok{,}
\NormalTok{    to\_jax}\OperatorTok{=}\VariableTok{True}\NormalTok{,}
\NormalTok{)}
\end{Highlighting}
\end{Shaded}

\textbf{SoftLaplace} Smooth distribution with Laplace-like tail
behavior.

This distribution corresponds to the log-convex density:

\[
  z = (value - loc) / scale
  log_prob = log(2 / pi) - log(scale) - logaddexp(z, -z)
  \]

Like the Laplace density, this density has the heaviest possible tails
(asymptotically) while still being log-convex. Unlike the Laplace
distribution, this distribution is infinitely differentiable everywhere,
and is thus suitable for HMC and Laplace approximation.

\textbf{Args} - \emph{loc}: Location parameter.

\begin{itemize}
\item
  \emph{scale}: Scale parameter.
\item
  \emph{shape} (tuple): A multi-purpose argument for shaping. When
  \texttt{sample=False} (model building), this is used with
  \texttt{shape} to set the distribution's batch shape. When
  \texttt{sample=True} (direct sampling), this is used as
  \texttt{sample\_shape} to draw a raw JAX array of the given shape.
\item
  \emph{event} (int): The number of batch dimensions to reinterpret as
  event dimensions (used in model building).
\item
  \emph{mask} (jnp.ndarray, bool): Optional boolean array to mask
  observations.
\item
  \emph{create\_obj} (bool): If True, returns the raw BF distribution
  object instead of creating a sample site. This is essential for
  building complex distributions like \texttt{MixtureSameFamily}.
\item
  \emph{sample} (bool, optional): A control-flow argument. If
  \texttt{True}, the function will directly sample a raw JAX array from
  the distribution, bypassing the BF model context. If \texttt{False},
  it will create a \texttt{BF.sample} site within a model. Defaults to
  \texttt{False}.
\item
  \emph{seed} (int, optional): An integer seed for this draw when
  \texttt{sample=True}. By default (\texttt{None}) the draw comes from
  the session's random stream, set by \texttt{bf(rand\_seed=...)}, so
  consecutive draws differ but the whole sequence is reproducible. No
  effect when \texttt{sample=False}.
\item
  \emph{obs} (jnp.ndarray, optional): The observed value for this random
  variable. If provided, the sample site is conditioned on this value,
  and the function returns the observed value. If \texttt{None}, the
  site is treated as a latent (unobserved) random variable. Defaults to
  \texttt{None}.
\item
  \emph{name} (str, optional): The name of the sample site in a BF
  model. This is used to uniquely identify the random variable. Defaults
  to `x'.
\end{itemize}

\textbf{Returns} - When \texttt{sample=False}: A BF SoftLaplace
distribution object (for model building). - When \texttt{sample=True}: A
JAX array of samples drawn from the SoftLaplace distribution (for direct
sampling). - When \texttt{create\_obj=True}: The raw BF distribution
object (for advanced use cases).

\textbf{Example Usage} from BayesForge import bf m = bf(platform=`cpu')
m.dist.soft\_laplace(loc=0.0, scale=1.0, sample=True)

\textbf{Wrapper of}
https://num.pyro.ai/en/stable/distributions.html\#softlaplace

\subsection{\texorpdfstring{\texttt{m.dist.student\_t}}{m.dist.student\_t}}\label{student_t}

\begin{Shaded}
\begin{Highlighting}[]
\NormalTok{m.dist.student\_t(}
\NormalTok{    df,}
\NormalTok{    loc}\OperatorTok{=}\FloatTok{0.0}\NormalTok{,}
\NormalTok{    scale}\OperatorTok{=}\FloatTok{1.0}\NormalTok{,}
\NormalTok{    validate\_args}\OperatorTok{=}\VariableTok{None}\NormalTok{,}
\NormalTok{    name}\OperatorTok{=}\StringTok{\textquotesingle{}x\textquotesingle{}}\NormalTok{,}
\NormalTok{    obs}\OperatorTok{=}\VariableTok{None}\NormalTok{,}
\NormalTok{    mask}\OperatorTok{=}\VariableTok{None}\NormalTok{,}
\NormalTok{    infer}\OperatorTok{=}\VariableTok{None}\NormalTok{,}
\NormalTok{    sample}\OperatorTok{=}\VariableTok{False}\NormalTok{,}
\NormalTok{    seed}\OperatorTok{=}\VariableTok{None}\NormalTok{,}
\NormalTok{    shape}\OperatorTok{=}\NormalTok{(),}
\NormalTok{    event}\OperatorTok{=}\DecValTok{0}\NormalTok{,}
\NormalTok{    create\_obj}\OperatorTok{=}\VariableTok{False}\NormalTok{,}
\NormalTok{    to\_jax}\OperatorTok{=}\VariableTok{True}\NormalTok{,}
\NormalTok{)}
\end{Highlighting}
\end{Shaded}

\textbf{Student's t} The Student's t-distribution is a probability
distribution that arises in hypothesis testing involving the mean of a
normally distributed population when the population standard deviation
is unknown. It is similar to the normal distribution, but has heavier
tails, making it more robust to outliers. For large \(\nu\), it
converges to the Normal distribution. \[
  X \sim t_\nu(\mu, \sigma)
  \] where:

\begin{itemize}
\tightlist
\item
  \(\mu\) is the \textbf{location (mean)} parameter
\item
  \(\sigma > 0\) is the \textbf{scale} parameter
\item
  \(\nu > 0\) is the \textbf{degrees of freedom} controlling the tail
  heaviness
\end{itemize}

The probability density function (pdf) of the Student's t distribution
is: \[
  f(x \mid \nu, \mu, \sigma) = \frac{\Gamma!\left(\frac{\nu+1}{2}\right)}
  {\Gamma!\left(\frac{\nu}{2}\right)\sqrt{\pi\nu}\sigma}
  \left[1 + \frac{1}{\nu}
  \left(\frac{x - \mu}{\sigma}\right)^2
  \right]^{-\frac{\nu+1}{2}}
  \]

\begin{itemize}
\tightlist
\item
  The heavier tails (for small \(ν\)) allow for larger outliers.
\item
  For \(\nu \to \infty\), this approaches
  \(\mathcal{N}(\mu, \sigma^2)\).
\end{itemize}

\textbf{Args} - \emph{df} (jnp.ndarray): Degrees of freedom, must be
positive.

\begin{itemize}
\item
  \emph{loc} (jnp.ndarray): Location parameter, defaults to 0.0.
\item
  \emph{scale} (jnp.ndarray): Scale parameter, defaults to 1.0.
\item
  \emph{shape} (tuple): A multi-purpose argument for shaping. When
  \texttt{sample=False} (model building), this is used with
  \texttt{.expand(shape)} to set the distribution's batch shape. When
  \texttt{sample=True} (direct sampling), this is used as
  \texttt{sample\_shape} to draw a raw JAX array of the given shape.
\item
  \emph{event} (int): The number of batch dimensions to reinterpret as
  event dimensions (used in model building).
\item
  \emph{mask} (jnp.ndarray, bool): Optional boolean array to mask
  observations.
\item
  \emph{create\_obj} (bool): If True, returns the raw BF distribution
  object instead of creating a sample site. This is essential for
  building complex distributions like \texttt{MixtureSameFamily}.
\item
  \emph{sample} (bool, optional): A control-flow argument. If
  \texttt{True}, the function will directly sample a raw JAX array from
  the distribution, bypassing the BF model context. If \texttt{False},
  it will create a \texttt{BF.sample} site within a model. Defaults to
  \texttt{False}.
\item
  \emph{seed} (int, optional): An integer seed for this draw when
  \texttt{sample=True}. By default (\texttt{None}) the draw comes from
  the session's random stream, set by \texttt{bf(rand\_seed=...)}, so
  consecutive draws differ but the whole sequence is reproducible. No
  effect when \texttt{sample=False}.
\item
  \emph{obs} (jnp.ndarray, optional): The observed value for this random
  variable. If provided, the sample site is conditioned on this value,
  and the function returns the observed value. If \texttt{None}, the
  site is treated as a latent (unobserved) random variable. Defaults to
  \texttt{None}.
\item
  \emph{name} (str, optional): The name of the sample site in a BF
  model. This is used to uniquely identify the random variable. Defaults
  to `x'.
\end{itemize}

\textbf{Returns} - When \texttt{sample=False}: A BF StudentT
distribution object (for model building). - When \texttt{sample=True}: A
JAX array of samples drawn from the StudentT distribution (for direct
sampling). - When \texttt{create\_obj=True}: The raw BF distribution
object (for advanced use cases).

\textbf{Example Usage} from BayesForge import bf m = bf(platform=`cpu')
m.dist.student\_t(df = 2, loc=0.0, scale=1.0, sample=True)

\textbf{Wrapper of}
https://num.pyro.ai/en/stable/distributions.html\#studentt

\textbf{References} *
\href{https://en.wikipedia.org/wiki/Student\%27s_t-distribution\#Probability_density_function}{Wikipedia
-- Student's t distribution: PDF}

\subsection{\texorpdfstring{\texttt{m.dist.track}}{m.dist.track}}\label{track}

\begin{Shaded}
\begin{Highlighting}[]
\NormalTok{m.dist.track(name, x)}
\end{Highlighting}
\end{Shaded}

Record a derived quantity so it is saved with the posterior draws.

\begin{longtable}[]{@{}
  >{\raggedright\arraybackslash}p{(\columnwidth - 4\tabcolsep) * \real{0.3333}}
  >{\raggedright\arraybackslash}p{(\columnwidth - 4\tabcolsep) * \real{0.3333}}
  >{\raggedright\arraybackslash}p{(\columnwidth - 4\tabcolsep) * \real{0.3333}}@{}}
\toprule\noalign{}
\begin{minipage}[b]{\linewidth}\raggedright
Argument
\end{minipage} & \begin{minipage}[b]{\linewidth}\raggedright
Type
\end{minipage} & \begin{minipage}[b]{\linewidth}\raggedright
Description
\end{minipage} \\
\midrule\noalign{}
\endhead
\bottomrule\noalign{}
\endlastfoot
\texttt{name} & str & Name of the quantity in \texttt{m.posteriors}. \\
\texttt{x} & jax.Array & The value, computed from the parameters inside
the model. \\
\end{longtable}

\textbf{Returns}

\begin{itemize}
\tightlist
\item
  jax.Array: \texttt{x} unchanged.
\end{itemize}

\textbf{Example Usage}

sigma2 = m.dist.track(``sigma2'', s ** 2) \# inside a model

\subsection{\texorpdfstring{\texttt{m.dist.truncated\_cauchy}}{m.dist.truncated\_cauchy}}\label{truncated_cauchy}

\begin{Shaded}
\begin{Highlighting}[]
\NormalTok{m.dist.truncated\_cauchy(}
\NormalTok{    loc}\OperatorTok{=}\FloatTok{0.0}\NormalTok{,}
\NormalTok{    scale}\OperatorTok{=}\FloatTok{1.0}\NormalTok{,}
\NormalTok{    low}\OperatorTok{=}\VariableTok{None}\NormalTok{,}
\NormalTok{    high}\OperatorTok{=}\VariableTok{None}\NormalTok{,}
\NormalTok{    validate\_args}\OperatorTok{=}\VariableTok{None}\NormalTok{,}
\NormalTok{    name}\OperatorTok{=}\StringTok{\textquotesingle{}x\textquotesingle{}}\NormalTok{,}
\NormalTok{    obs}\OperatorTok{=}\VariableTok{None}\NormalTok{,}
\NormalTok{    mask}\OperatorTok{=}\VariableTok{None}\NormalTok{,}
\NormalTok{    infer}\OperatorTok{=}\VariableTok{None}\NormalTok{,}
\NormalTok{    sample}\OperatorTok{=}\VariableTok{False}\NormalTok{,}
\NormalTok{    seed}\OperatorTok{=}\VariableTok{None}\NormalTok{,}
\NormalTok{    shape}\OperatorTok{=}\NormalTok{(),}
\NormalTok{    event}\OperatorTok{=}\DecValTok{0}\NormalTok{,}
\NormalTok{    create\_obj}\OperatorTok{=}\VariableTok{False}\NormalTok{,}
\NormalTok{    to\_jax}\OperatorTok{=}\VariableTok{True}\NormalTok{,}
\NormalTok{)}
\end{Highlighting}
\end{Shaded}

\textbf{Truncated Cauchy} The Cauchy distribution, also known as the
Lorentz distribution, is a continuous probability distribution that
appears frequently in various areas of mathematics and physics. It is
characterized by its heavy tails, which extend to infinity. The
truncated version limits the support of the Cauchy distribution to a
specified interval.

Let the base Cauchy have parameters \(\mu\) (location) and
\(\gamma > 0\) (scale). Let the truncation bounds be \(a\)
(\texttt{low}) and \(b\) (\texttt{high}) (with \(a < b\)). Then the PDF
of the truncated Cauchy at \(x \in [a,b]\) is: \[
f_{T}(x) = \frac{1}{\pi,\gamma};\frac{1}{1 + \bigl(\tfrac{x - \mu}{\gamma}\bigr)^2} ;\Bigg/ ; \left[F_{\mathrm{Cauchy}}!\left(\tfrac{b - \mu}{\gamma}\right) - F_{\mathrm{Cauchy}}!\left(\tfrac{a - \mu}{\gamma}\right)\right],
\] and (f\_\{T\}(x)=0) for (x \textless{} a) or (x \textgreater{} b).

Here \(F_{\mathrm{Cauchy}}\) is the standard Cauchy CDF \$with
parameters \(\mu,\gamma)\).

\textbf{Args} - \emph{loc} (float): Location parameter of the Cauchy
distribution.

\begin{itemize}
\item
  \emph{scale} (float): Scale parameter of the Cauchy distribution.
\item
  \emph{shape} (tuple): A multi-purpose argument for shaping. When
  \texttt{sample=False} (model building), this is used with
  \texttt{shape} to set the distribution's batch shape. When
  \texttt{sample=True} (direct sampling), this is used as
  \texttt{sample\_shape} to draw a raw JAX array of the given shape.
\item
  \emph{event} (int): The number of batch dimensions to reinterpret as
  event dimensions (used in model building).
\item
  \emph{mask} (jnp.ndarray, bool): Optional boolean array to mask
  observations.
\item
  \emph{create\_obj} (bool): If True, returns the raw BF distribution
  object instead of creating a sample site. This is essential for
  building complex distributions like \texttt{MixtureSameFamily}.
\item
  \emph{sample} (bool, optional): A control-flow argument. If
  \texttt{True}, the function will directly sample a raw JAX array from
  the distribution, bypassing the BF model context. If \texttt{False},
  it will create a \texttt{BF.sample} site within a model. Defaults to
  \texttt{False}.
\item
  \emph{seed} (int, optional): An integer seed for this draw when
  \texttt{sample=True}. By default (\texttt{None}) the draw comes from
  the session's random stream, set by \texttt{bf(rand\_seed=...)}, so
  consecutive draws differ but the whole sequence is reproducible. No
  effect when \texttt{sample=False}.
\item
  \emph{obs} (jnp.ndarray, optional): The observed value for this random
  variable. If provided, the sample site is conditioned on this value,
  and the function returns the observed value. If \texttt{None}, the
  site is treated as a latent (unobserved) random variable. Defaults to
  \texttt{None}.
\item
  \emph{name} (str, optional): The name of the sample site in a BF
  model. This is used to uniquely identify the random variable. Defaults
  to `x'.
\end{itemize}

\textbf{Returns} BF TruncatedCauchy distribution object (for model
building) when \texttt{sample=False}.

\begin{verbatim}
JAX array of samples drawn from the TruncatedCauchy distribution (for direct sampling) when `sample=True`.

The raw BF distribution object (for advanced use cases) when `create_obj=True`.
\end{verbatim}

\textbf{Example Usage} from BayesForge import bf m = bf(platform=`cpu')
m.dist.truncated\_cauchy(loc=0.0, scale=1.0, sample=True)

\textbf{Wrapper of:
https://num.pyro.ai/en/stable/distributions.html\#truncatedcauchy}
\textbf{References} -
\href{https://en.wikipedia.org/wiki/Cauchy_distribution}{Cauchy
distribution} -
\href{https://www.tandfonline.com/doi/full/10.1080/00207390600595223}{Full
article: A truncated Cauchy distribution} -
\href{https://rstudio.github.io/tfprobability/reference/tfd_truncated_cauchy.html}{The
Truncated Cauchy distribution. --- tfd\_truncated\_cauchy}

\subsection{\texorpdfstring{\texttt{m.dist.truncated\_distribution}}{m.dist.truncated\_distribution}}\label{truncated_distribution}

\begin{Shaded}
\begin{Highlighting}[]
\NormalTok{m.dist.truncated\_distribution(}
\NormalTok{    base\_dist,}
\NormalTok{    low}\OperatorTok{=}\VariableTok{None}\NormalTok{,}
\NormalTok{    high}\OperatorTok{=}\VariableTok{None}\NormalTok{,}
\NormalTok{    validate\_args}\OperatorTok{=}\VariableTok{None}\NormalTok{,}
\NormalTok{    name}\OperatorTok{=}\StringTok{\textquotesingle{}x\textquotesingle{}}\NormalTok{,}
\NormalTok{    obs}\OperatorTok{=}\VariableTok{None}\NormalTok{,}
\NormalTok{    mask}\OperatorTok{=}\VariableTok{None}\NormalTok{,}
\NormalTok{    infer}\OperatorTok{=}\VariableTok{None}\NormalTok{,}
\NormalTok{    sample}\OperatorTok{=}\VariableTok{False}\NormalTok{,}
\NormalTok{    seed}\OperatorTok{=}\VariableTok{None}\NormalTok{,}
\NormalTok{    shape}\OperatorTok{=}\NormalTok{(),}
\NormalTok{    event}\OperatorTok{=}\DecValTok{0}\NormalTok{,}
\NormalTok{    create\_obj}\OperatorTok{=}\VariableTok{False}\NormalTok{,}
\NormalTok{    to\_jax}\OperatorTok{=}\VariableTok{True}\NormalTok{,}
\NormalTok{)}
\end{Highlighting}
\end{Shaded}

\textbf{Truncated} Samples from a Truncated Distribution.

A \textbf{truncated distribution} arises when you take a random variable
\(X\) that originally has some distribution (with PDF \(f_X(x)\) and CDF
\(F_X(x)\)) and you restrict attention only to those values of \(X\)
that are \emph{above} a given truncation point \(a\). In other words you
only observe \(X\) when \(X > a\). All the ``mass'' below (or equal to)
\(a\) is \textbf{excluded} (not just unobserved, but removed from the
sample/analysis). This differs from \emph{censoring}, where values below
a threshold might be known (for example ``\textless{} a''), but here
they are entirely excluded from the domain. Left truncation is common in
many applied fields --- for instance:

\begin{itemize}
\tightlist
\item
  In survival analysis: subjects whose event time happens before the
  study start are not included.
\item
  In insurance or losses: only losses above a deductible (threshold) are
  recorded, so the loss distribution is left-truncated at that
  deductible.
\item
  In industrial/life-data: items used before data collection start, so
  only those with lifetime \textgreater{} some lower bound are observed.
\end{itemize}

Let (X) be a random variable with PDF. Choose truncation points
\(a < X \le b\). Then the truncated distribution (Y = X \mid (a
\textless{} X \le b)) has:\\
\[
  f_Y(y) = \frac{f_X(y)}{F_X(b) - F_X(a)}, \quad \text{for } a < y \le b,
  \] and (0) outside that interval. * The lower bound (a) can be
(-\infty) (so only an upper truncation), and the upper bound (b) can be
(+\infty) (so only a lower truncation).
(\href{https://statproofbook.github.io/P/matn-pdf.html}{Statistics How
To}) * If both bounds are finite (i.e., doubly truncated), that is a
valid and often‐used scenario.

\textbf{Args} - \emph{base\_dist}: The base distribution to be
truncated. This should be a univariate distribution. Currently, only the
following distributions are supported: Cauchy, Laplace, Logistic,
Normal, and StudentT.

\begin{itemize}
\item
  \emph{shape} (tuple): A multi-purpose argument for shaping. When
  \texttt{sample=False} (model building), this is used with
  \texttt{shape} to set the distribution's batch shape. When
  \texttt{sample=True} (direct sampling), this is used as
  \texttt{sample\_shape} to draw a raw JAX array of the given shape.
\item
  \emph{event} (int): The number of batch dimensions to reinterpret as
  event dimensions (used in model building).
\item
  \emph{mask} (jnp.ndarray, bool): Optional boolean array to mask
  observations.
\item
  \emph{create\_obj} (bool): If True, returns the raw BF distribution
  object instead of creating a sample site. This is essential for
  building complex distributions like \texttt{MixtureSameFamily}.
\item
  \emph{sample} (bool, optional): A control-flow argument. If
  \texttt{True}, the function will directly sample a raw JAX array from
  the distribution, bypassing the BF model context. If \texttt{False},
  it will create a \texttt{BF.sample} site within a model. Defaults to
  \texttt{False}.
\item
  \emph{seed} (int, optional): An integer seed for this draw when
  \texttt{sample=True}. By default (\texttt{None}) the draw comes from
  the session's random stream, set by \texttt{bf(rand\_seed=...)}, so
  consecutive draws differ but the whole sequence is reproducible. No
  effect when \texttt{sample=False}.
\item
  \emph{obs} (jnp.ndarray, optional): The observed value for this random
  variable. If provided, the sample site is conditioned on this value,
  and the function returns the observed value. If \texttt{None}, the
  site is treated as a latent (unobserved) random variable. Defaults to
  \texttt{None}.
\item
  \emph{name} (str, optional): The name of the sample site in a BF
  model. This is used to uniquely identify the random variable. Defaults
  to `x'.
\end{itemize}

\textbf{Returns} - When \texttt{sample=False}: A BF
TruncatedDistribution distribution object (for model building). - When
\texttt{sample=True}: A JAX array of samples drawn from the
TruncatedDistribution distribution (for direct sampling). - When
\texttt{create\_obj=True}: The raw BF distribution object (for advanced
use cases).

\textbf{Example Usage} from BayesForge import bf m = bf(platform=`cpu')
m.dist.truncated\_distribution(base\_dist = m.dist.normal(0,1,
create\_obj = True), high=1, low = 0, sample=True)

\textbf{Wrapper of:
https://num.pyro.ai/en/stable/distributions.html\#truncateddistribution}
\textbf{References} -
\href{https://en.wikipedia.org/wiki/Truncated_distribution}{Truncated
distribution} -
\href{https://pages.stern.nyu.edu/~wgreene/Text/Edition8/PDF/M19_GREE1366_08_SE_C19.pdf}{19}
-
\href{https://search.r-project.org/CRAN/refmans/LaplacesDemon/html/dist.Truncated.html}{R:
Truncated Distributions} -
\href{https://dspace.cuni.cz/bitstream/handle/20.500.11956/127921/130308341.pdf?isAllowed=y&sequence=1&utm_source.com}{Truncated
data thesis}

\subsection{\texorpdfstring{\texttt{m.dist.truncated\_normal}}{m.dist.truncated\_normal}}\label{truncated_normal}

\begin{Shaded}
\begin{Highlighting}[]
\NormalTok{m.dist.truncated\_normal(}
\NormalTok{    loc}\OperatorTok{=}\FloatTok{0.0}\NormalTok{,}
\NormalTok{    scale}\OperatorTok{=}\FloatTok{1.0}\NormalTok{,}
\NormalTok{    low}\OperatorTok{=}\VariableTok{None}\NormalTok{,}
\NormalTok{    high}\OperatorTok{=}\VariableTok{None}\NormalTok{,}
\NormalTok{    validate\_args}\OperatorTok{=}\VariableTok{None}\NormalTok{,}
\NormalTok{    name}\OperatorTok{=}\StringTok{\textquotesingle{}x\textquotesingle{}}\NormalTok{,}
\NormalTok{    obs}\OperatorTok{=}\VariableTok{None}\NormalTok{,}
\NormalTok{    mask}\OperatorTok{=}\VariableTok{None}\NormalTok{,}
\NormalTok{    infer}\OperatorTok{=}\VariableTok{None}\NormalTok{,}
\NormalTok{    sample}\OperatorTok{=}\VariableTok{False}\NormalTok{,}
\NormalTok{    seed}\OperatorTok{=}\VariableTok{None}\NormalTok{,}
\NormalTok{    shape}\OperatorTok{=}\NormalTok{(),}
\NormalTok{    event}\OperatorTok{=}\DecValTok{0}\NormalTok{,}
\NormalTok{    create\_obj}\OperatorTok{=}\VariableTok{False}\NormalTok{,}
\NormalTok{    to\_jax}\OperatorTok{=}\VariableTok{True}\NormalTok{,}
\NormalTok{)}
\end{Highlighting}
\end{Shaded}

\textbf{Truncated Normal} A truncated normal distribution is derived
from a normal (Gaussian) random variable by restricting (truncating) its
domain to an interval \([a, b]\) (which could be one‐sided, e.g., (a)
only or (b) only). It is defined by its location (\texttt{loc}), scale
(\texttt{scale}), lower bound \(a\) (\texttt{low}), and upper bound
\(b\) (\texttt{high}). In effect: if
\(X \sim \mathcal N(\mu, \sigma^2)\), then the truncated version
\(Y = X | (a \le X \le b)\) has the same ``shape'' but only supports
values in \([a,b]\). This is used when you know that values outside a
range are impossible or not observed (e.g., measurement limits, natural
bounds).

Let \(X \sim \mathcal N(\mu, \sigma^2)\). Define truncation bounds
\(a < b\) (could be \(a = -\infty\) or \(b = +\infty\) for one‐sided
truncation). Then for \(y \in [a,b]\),

\[
f_Y(y) = \frac{1}{\sigma} ; \frac{\varphi!\bigl(\frac{y - \mu}{\sigma}\bigr)}{\Phi!\bigl(\frac{b - \mu}{\sigma}\bigr) - \Phi!\bigl(\frac{a - \mu}{\sigma}\bigr)} ,
\]

where:

\begin{itemize}
\tightlist
\item
  \(\varphi(z)\) is the standard normal PDF,
  \(\varphi(z) = \frac1{\sqrt{2\pi}} e^{-z^2/2}\).
\item
  \(\Phi(z)\) is the standard normal CDF.
\item
  The denominator
  \(\Phi\bigl(\tfrac{b - \mu}{\sigma}\bigr) - \Phi\bigl(\tfrac{a - \mu}{\sigma}\bigr)\)
  normalises the density so the total probability over ({[}a,b{]}) is 1.
\end{itemize}

For values (y \textless{} a) or (y \textgreater{} b), (f\_Y(y) = 0).

\textbf{Args} - \emph{loc} (float): The location parameter of the normal
distribution.

\begin{itemize}
\item
  \emph{sample} (float): The scale parameter of the normal distribution.
\item
  \emph{shape} (tuple): A multi-purpose argument for shaping. When
  \texttt{sample=False} (model building), this is used with
  \texttt{shape} to set the distribution's batch shape. When
  \texttt{sample=True} (direct sampling), this is used as
  \texttt{sample\_shape} to draw a raw JAX array of the given shape.
\item
  \emph{event} (int): The number of batch dimensions to reinterpret as
  event dimensions (used in model building).
\item
  \emph{mask} (jnp.ndarray, bool): Optional boolean array to mask
  observations.
\item
  \emph{create\_obj} (bool): If True, returns the raw BF distribution
  object instead of creating a sample site. This is essential for
  building complex distributions like \texttt{MixtureSameFamily}.
\item
  \emph{sample} (bool, optional): A control-flow argument. If
  \texttt{True}, the function will directly sample a raw JAX array from
  the distribution, bypassing the BF model context. If \texttt{False},
  it will create a \texttt{BF.sample} site within a model. Defaults to
  \texttt{False}.
\item
  \emph{seed} (int, optional): An integer seed for this draw when
  \texttt{sample=True}. By default (\texttt{None}) the draw comes from
  the session's random stream, set by \texttt{bf(rand\_seed=...)}, so
  consecutive draws differ but the whole sequence is reproducible. No
  effect when \texttt{sample=False}.
\item
  \emph{obs} (jnp.ndarray, optional): The observed value for this random
  variable. If provided, the sample site is conditioned on this value,
  and the function returns the observed value. If \texttt{None}, the
  site is treated as a latent (unobserved) random variable. Defaults to
  \texttt{None}.
\item
  \emph{name} (str, optional): The name of the sample site in a BF
  model. This is used to uniquely identify the random variable. Defaults
  to `x'.
\end{itemize}

\textbf{Returns} BF TruncatedNormal distribution object (for model
building). JAX array of samples drawn from the TruncatedNormal
distribution (for direct sampling). The raw BF distribution object (for
advanced use cases).

\textbf{Example Usage} from BayesForge import bf m = bf(platform=`cpu')
m.dist.truncated\_normal(loc=0.0, scale=1.0, sample=True)

\textbf{Wrapper of:
https://num.pyro.ai/en/stable/distributions.html\#truncatednormal\_lowercase}
\textbf{References} -
\href{https://en.wikipedia.org/wiki/Truncated_normal_distribution}{Truncated
normal distribution} -
\href{https://people.math.sc.edu/burkardt/m_src/truncated_normal/truncated_normal.html}{TRUNCATED\_NORMAL
- The Truncated Normal Distribution} -
\href{https://www.statisticshowto.com/truncated-normal-distribution/}{Truncated
Distribution / Truncated Normal Distribution - Statistics \ldots{}} -
\href{https://real-statistics.com/normal-distribution/truncated-normal-distribution/}{Truncated
Normal Distribution \textbar{} Real Statistics Using Excel} -
\href{https://stats.stackexchange.com/questions/525894/understanding-the-pdf-of-a-truncated-normal-distribution}{Understanding
the pdf of a truncated normal distribution}

\subsection{\texorpdfstring{\texttt{m.dist.truncated\_polya\_gamma}}{m.dist.truncated\_polya\_gamma}}\label{truncated_polya_gamma}

\begin{Shaded}
\begin{Highlighting}[]
\NormalTok{m.dist.truncated\_polya\_gamma(}
\NormalTok{    batch\_shape}\OperatorTok{=}\NormalTok{(),}
\NormalTok{    validate\_args}\OperatorTok{=}\VariableTok{None}\NormalTok{,}
\NormalTok{    name}\OperatorTok{=}\StringTok{\textquotesingle{}x\textquotesingle{}}\NormalTok{,}
\NormalTok{    obs}\OperatorTok{=}\VariableTok{None}\NormalTok{,}
\NormalTok{    mask}\OperatorTok{=}\VariableTok{None}\NormalTok{,}
\NormalTok{    infer}\OperatorTok{=}\VariableTok{None}\NormalTok{,}
\NormalTok{    sample}\OperatorTok{=}\VariableTok{False}\NormalTok{,}
\NormalTok{    seed}\OperatorTok{=}\VariableTok{None}\NormalTok{,}
\NormalTok{    shape}\OperatorTok{=}\NormalTok{(),}
\NormalTok{    event}\OperatorTok{=}\DecValTok{0}\NormalTok{,}
\NormalTok{    create\_obj}\OperatorTok{=}\VariableTok{False}\NormalTok{,}
\NormalTok{    to\_jax}\OperatorTok{=}\VariableTok{True}\NormalTok{,}
\NormalTok{)}
\end{Highlighting}
\end{Shaded}

\textbf{Truncated PolyaGamma} Samples from a Truncated PolyaGamma
distribution.

This distribution is a truncated version of the PolyaGamma distribution,
defined over the interval {[}0, truncation\_point{]}. It is often used
in Bayesian non-parametric models.

\textbf{Args} - \emph{batch\_shape} (tuple): The shape of the batch
dimension.

\begin{itemize}
\item
  \emph{event} (int): The number of batch dimensions to reinterpret as
  event dimensions.
\item
  \emph{mask} (jnp.ndarray, bool): Optional boolean array to mask
  observations.
\item
  \emph{create\_obj} (bool): If True, returns the raw BF distribution
  object instead of creating a sample site.
\item
  \emph{sample} (bool, optional): A control-flow argument. If
  \texttt{True}, the function will directly sample a raw JAX array from
  the distribution, bypassing the BF model context. If \texttt{False},
  it will create a \texttt{BF.sample} site within a model. Defaults to
  \texttt{False}.
\item
  \emph{seed} (int, optional): An integer seed for this draw when
  \texttt{sample=True}. By default (\texttt{None}) the draw comes from
  the session's random stream, set by \texttt{bf(rand\_seed=...)}, so
  consecutive draws differ but the whole sequence is reproducible. No
  effect when \texttt{sample=False}.
\item
  \emph{obs} (jnp.ndarray, optional): The observed value for this random
  variable. If provided, the sample site is conditioned on this value,
  and the function returns the observed value. If \texttt{None}, the
  site is treated as a latent (unobserved) random variable. Defaults to
  \texttt{None}.
\item
  \emph{name} (str, optional): The name of the sample site in a BF
  model. This is used to uniquely identify the random variable. Defaults
  to `x'.
\end{itemize}

\textbf{Returns} - When \texttt{sample=False}: A BF Truncated PolyaGamma
distribution object (for model building). - When \texttt{sample=True}: A
JAX array of samples drawn from the Truncated PolyaGamma distribution
(for direct sampling). - When \texttt{create\_obj=True}: The raw BF
distribution object (for advanced use cases).

\textbf{Example Usage} from BayesForge import bf m = bf(platform=`cpu')
m.dist.truncated\_polya\_gamma(batch\_shape=(), sample=True)

\textbf{Wrapper of}
https://num.pyro.ai/en/stable/distributions.html\#truncatedpolygammadistribution

\subsection{\texorpdfstring{\texttt{m.dist.two\_sided\_truncated\_distribution}}{m.dist.two\_sided\_truncated\_distribution}}\label{two_sided_truncated_distribution}

\begin{Shaded}
\begin{Highlighting}[]
\NormalTok{m.dist.two\_sided\_truncated\_distribution(}
\NormalTok{    base\_dist,}
\NormalTok{    low}\OperatorTok{=}\FloatTok{0.0}\NormalTok{,}
\NormalTok{    high}\OperatorTok{=}\FloatTok{1.0}\NormalTok{,}
\NormalTok{    validate\_args}\OperatorTok{=}\VariableTok{None}\NormalTok{,}
\NormalTok{    name}\OperatorTok{=}\StringTok{\textquotesingle{}x\textquotesingle{}}\NormalTok{,}
\NormalTok{    obs}\OperatorTok{=}\VariableTok{None}\NormalTok{,}
\NormalTok{    mask}\OperatorTok{=}\VariableTok{None}\NormalTok{,}
\NormalTok{    infer}\OperatorTok{=}\VariableTok{None}\NormalTok{,}
\NormalTok{    sample}\OperatorTok{=}\VariableTok{False}\NormalTok{,}
\NormalTok{    seed}\OperatorTok{=}\VariableTok{None}\NormalTok{,}
\NormalTok{    shape}\OperatorTok{=}\NormalTok{(),}
\NormalTok{    event}\OperatorTok{=}\DecValTok{0}\NormalTok{,}
\NormalTok{    create\_obj}\OperatorTok{=}\VariableTok{False}\NormalTok{,}
\NormalTok{    to\_jax}\OperatorTok{=}\VariableTok{True}\NormalTok{,}
\NormalTok{)}
\end{Highlighting}
\end{Shaded}

\textbf{Two Sided Truncated} A ``two-sided truncated distribution'' is a
general concept: you take a base continuous distribution and
\textbf{restrict it} to an interval ({[}\texttt{low}, \texttt{high}{]}),
discarding all mass outside, then \textbf{renormalize} so the inner
portion integrates to 1. I'll spell out the general formulas, caveats,
sampling strategies, and special cases (e.g.~truncated normal) to
illustrate.

\[
  f(x) =
  \begin{cases}
  \dfrac{p(x)}{P(\text{low} \le X \le \text{high})}, & \text{if } \text{low} \le x \le \text{high}, \\[6pt]
  0, & \text{otherwise}.
  \end{cases}
  \]

where \(p(x)\) is the probability density function of the base
distribution.

\textbf{Args} - \emph{base\_dist}: The base distribution to truncate.

\begin{itemize}
\item
  \emph{low}: The lower bound for truncation.
\item
  \emph{high}: The upper bound for truncation.
\item
  \emph{sample} (bool, optional): A control-flow argument. If
  \texttt{True}, the function will directly sample a raw JAX array from
  the distribution, bypassing the BF model context. If \texttt{False},
  it will create a \texttt{BF.sample} site within a model. Defaults to
  \texttt{False}.
\item
  \emph{seed} (int, optional): An integer seed for this draw when
  \texttt{sample=True}. By default (\texttt{None}) the draw comes from
  the session's random stream, set by \texttt{bf(rand\_seed=...)}, so
  consecutive draws differ but the whole sequence is reproducible. No
  effect when \texttt{sample=False}.
\item
  \emph{obs} (jnp.ndarray, optional): The observed value for this random
  variable. If provided, the sample site is conditioned on this value,
  and the function returns the observed value. If \texttt{None}, the
  site is treated as a latent (unobserved) random variable. Defaults to
  \texttt{None}.
\item
  \emph{name} (str, optional): The name of the sample site in a BF
  model. This is used to uniquely identify the random variable. Defaults
  to `x'.
\end{itemize}

\textbf{Returns} - When \texttt{sample=False}: A BF
TwoSidedTruncatedDistribution distribution object (for model building).
- When \texttt{sample=True}: A JAX array of samples drawn from the
TwoSidedTruncatedDistribution distribution (for direct sampling). - When
\texttt{create\_obj=True}: The raw BF distribution object (for advanced
use cases).

\textbf{Wrapper of:
https://num.pyro.ai/en/stable/distributions.html\#twosidedtruncateddistribution}

\subsection{\texorpdfstring{\texttt{m.dist.uniform}}{m.dist.uniform}}\label{uniform}

\begin{Shaded}
\begin{Highlighting}[]
\NormalTok{m.dist.uniform(}
\NormalTok{    low}\OperatorTok{=}\FloatTok{0.0}\NormalTok{,}
\NormalTok{    high}\OperatorTok{=}\FloatTok{1.0}\NormalTok{,}
\NormalTok{    validate\_args}\OperatorTok{=}\VariableTok{None}\NormalTok{,}
\NormalTok{    name}\OperatorTok{=}\StringTok{\textquotesingle{}x\textquotesingle{}}\NormalTok{,}
\NormalTok{    obs}\OperatorTok{=}\VariableTok{None}\NormalTok{,}
\NormalTok{    mask}\OperatorTok{=}\VariableTok{None}\NormalTok{,}
\NormalTok{    infer}\OperatorTok{=}\VariableTok{None}\NormalTok{,}
\NormalTok{    sample}\OperatorTok{=}\VariableTok{False}\NormalTok{,}
\NormalTok{    seed}\OperatorTok{=}\VariableTok{None}\NormalTok{,}
\NormalTok{    shape}\OperatorTok{=}\NormalTok{(),}
\NormalTok{    event}\OperatorTok{=}\DecValTok{0}\NormalTok{,}
\NormalTok{    create\_obj}\OperatorTok{=}\VariableTok{False}\NormalTok{,}
\NormalTok{    to\_jax}\OperatorTok{=}\VariableTok{True}\NormalTok{,}
\NormalTok{)}
\end{Highlighting}
\end{Shaded}

\textbf{Uniform} The Uniform distribution is the simplest continuous
distribution: every value in the interval ({[}a, b{]}) is
\textbf{equally likely}. It is widely used for modeling complete
randomness within a fixed range, random sampling, and as a building
block for other distributions.

\[   
f(x) =
\begin{cases}
\dfrac{1}{b - a}, & a \le x \le b, [2mm]
0, & \text{otherwise}.
\end{cases}
\]

\textbf{Args} - \emph{low} (jnp.ndarray): The lower bound of the uniform
interval.

\begin{itemize}
\item
  \emph{high} (jnp.ndarray): The upper bound of the uniform interval.
\item
  \emph{shape} (tuple): A multi-purpose argument for shaping. When
  \texttt{sample=False} (model building), this is used with
  \texttt{shape} to set the distribution's batch shape. When
  \texttt{sample=True} (direct sampling), this is used as
  \texttt{sample\_shape} to draw a raw JAX array of the given shape.
\item
  \emph{event} (int): The number of batch dimensions to reinterpret as
  event dimensions (used in model building).
\item
  \emph{mask} (jnp.ndarray, bool): Optional boolean array to mask
  observations.
\item
  \emph{create\_obj} (bool, optional): If True, returns the raw BF
  distribution object instead of creating a sample site. This is
  essential for building complex distributions like
  \texttt{MixtureSameFamily}. Defaults to False.
\item
  \emph{sample} (bool, optional): A control-flow argument. If
  \texttt{True}, the function will directly sample a raw JAX array from
  the distribution, bypassing the BF model context. If \texttt{False},
  it will create a \texttt{BF.sample} site within a model. Defaults to
  \texttt{False}.
\item
  \emph{seed} (int, optional): An integer seed for this draw when
  \texttt{sample=True}. By default (\texttt{None}) the draw comes from
  the session's random stream, set by \texttt{bf(rand\_seed=...)}, so
  consecutive draws differ but the whole sequence is reproducible. No
  effect when \texttt{sample=False}.
\item
  \emph{obs} (jnp.ndarray, optional): The observed value for this random
  variable. If provided, the sample site is conditioned on this value,
  and the function returns the observed value. If \texttt{None}, the
  site is treated as a latent (unobserved) random variable. Defaults to
  \texttt{None}.
\item
  \emph{name} (str, optional): The name of the sample site in a BF
  model. This is used to uniquely identify the random variable. Defaults
  to `x'.
\end{itemize}

\textbf{Returns} BF Uniform distribution object (for model building)
when \texttt{sample=False}.

\begin{verbatim}
JAX array of samples drawn from the Uniform distribution (for direct
sampling) when `sample=True`.

The raw BF distribution object (for advanced use cases) when
`create_obj=True`.
\end{verbatim}

\textbf{Example Usage} from BayesForge import bf m = bf(platform=`cpu')
m.dist.uniform(low=0.0, high=1.0, sample=True)

\textbf{Wrapper of}
https://num.pyro.ai/en/stable/distributions.html\#uniform

\textbf{References} *
\href{https://en.wikipedia.org/wiki/Uniform_distribution_\%28continuous\%29}{Wikipedia
-- Uniform distribution (continuous)}

\subsection{\texorpdfstring{\texttt{m.dist.unit}}{m.dist.unit}}\label{unit}

\begin{Shaded}
\begin{Highlighting}[]
\NormalTok{m.dist.unit(}
\NormalTok{    log\_factor,}
\NormalTok{    validate\_args}\OperatorTok{=}\VariableTok{None}\NormalTok{,}
\NormalTok{    name}\OperatorTok{=}\StringTok{\textquotesingle{}x\textquotesingle{}}\NormalTok{,}
\NormalTok{    obs}\OperatorTok{=}\VariableTok{None}\NormalTok{,}
\NormalTok{    mask}\OperatorTok{=}\VariableTok{None}\NormalTok{,}
\NormalTok{    infer}\OperatorTok{=}\VariableTok{None}\NormalTok{,}
\NormalTok{    sample}\OperatorTok{=}\VariableTok{False}\NormalTok{,}
\NormalTok{    seed}\OperatorTok{=}\VariableTok{None}\NormalTok{,}
\NormalTok{    shape}\OperatorTok{=}\NormalTok{(),}
\NormalTok{    event}\OperatorTok{=}\DecValTok{0}\NormalTok{,}
\NormalTok{    create\_obj}\OperatorTok{=}\VariableTok{False}\NormalTok{,}
\NormalTok{    to\_jax}\OperatorTok{=}\VariableTok{True}\NormalTok{,}
\NormalTok{)}
\end{Highlighting}
\end{Shaded}

\textbf{Unit} * The Unit distribution is the simplest continuous
probability distribution. * Every number in ({[}0,1{]}) is equally
likely. * Often used in simulation, Monte Carlo methods, and as a
building block for generating other distributions via the inverse-CDF
method. \[
f(x) =
\begin{cases}
1, & 0 \le x \le 1, [1mm]
0, & \text{otherwise}.
\end{cases}
\]

\textbf{Args} - \emph{log\_factor} (jnp.ndarray): Log factor for the
unit distribution. This parameter determines the \emph{shape} and batch
size of the distribution.

\begin{itemize}
\item
  \emph{shape} (tuple): A multi-purpose argument for shaping. When
  \texttt{sample=False} (model building), this is used with
  \texttt{shape} to set the distribution's batch shape. When
  \texttt{sample=True} (direct sampling), this is used as
  \texttt{sample\_shape} to draw a raw JAX array of the given shape.
\item
  \emph{event} (int): The number of batch dimensions to reinterpret as
  event dimensions (used in model building).
\item
  \emph{mask} (jnp.ndarray, bool): Optional boolean array to mask
  observations.
\item
  \emph{create\_obj} (bool): If True, returns the raw BF distribution
  object instead of creating a sample site. This is essential for
  building complex distributions like \texttt{MixtureSameFamily}.
\item
  \emph{sample} (bool, optional): A control-flow argument. If
  \texttt{True}, the function will directly sample a raw JAX array from
  the distribution, bypassing the BF model context. If \texttt{False},
  it will create a \texttt{BF.sample} site within a model. Defaults to
  \texttt{False}.
\item
  \emph{seed} (int, optional): An integer seed for this draw when
  \texttt{sample=True}. By default (\texttt{None}) the draw comes from
  the session's random stream, set by \texttt{bf(rand\_seed=...)}, so
  consecutive draws differ but the whole sequence is reproducible. No
  effect when \texttt{sample=False}.
\item
  \emph{obs} (jnp.ndarray, optional): The observed value for this random
  variable. If provided, the sample site is conditioned on this value,
  and the function returns the observed value. If \texttt{None}, the
  site is treated as a latent (unobserved) random variable. Defaults to
  \texttt{None}.
\item
  \emph{name} (str, optional): The name of the sample site in a BF
  model. This is used to uniquely identify the random variable. Defaults
  to `x'.
\end{itemize}

\textbf{Returns} BF Unit distribution object: When \texttt{sample=False}
(for model building). jnp.ndarray: A JAX array of samples drawn from the
Unit distribution (for direct sampling). BF Unit distribution object:
When \texttt{create\_obj=True} (for advanced use cases).

\textbf{Example Usage} from BayesForge import bf import jax.numpy as jnp
m = bf(platform=`cpu') m.dist.unit(log\_factor=jnp.ones(5), sample=True)

\textbf{Wrapper of}
https://num.pyro.ai/en/stable/distributions.html\#unit

\subsection{\texorpdfstring{\texttt{m.dist.weibull}}{m.dist.weibull}}\label{weibull}

\begin{Shaded}
\begin{Highlighting}[]
\NormalTok{m.dist.weibull(}
\NormalTok{    scale,}
\NormalTok{    concentration,}
\NormalTok{    validate\_args}\OperatorTok{=}\VariableTok{None}\NormalTok{,}
\NormalTok{    name}\OperatorTok{=}\StringTok{\textquotesingle{}x\textquotesingle{}}\NormalTok{,}
\NormalTok{    obs}\OperatorTok{=}\VariableTok{None}\NormalTok{,}
\NormalTok{    mask}\OperatorTok{=}\VariableTok{None}\NormalTok{,}
\NormalTok{    infer}\OperatorTok{=}\VariableTok{None}\NormalTok{,}
\NormalTok{    sample}\OperatorTok{=}\VariableTok{False}\NormalTok{,}
\NormalTok{    seed}\OperatorTok{=}\VariableTok{None}\NormalTok{,}
\NormalTok{    shape}\OperatorTok{=}\NormalTok{(),}
\NormalTok{    event}\OperatorTok{=}\DecValTok{0}\NormalTok{,}
\NormalTok{    create\_obj}\OperatorTok{=}\VariableTok{False}\NormalTok{,}
\NormalTok{    to\_jax}\OperatorTok{=}\VariableTok{True}\NormalTok{,}
\NormalTok{)}
\end{Highlighting}
\end{Shaded}

\textbf{Weibull} Samples from a Weibull distribution.

The Weibull distribution is widely used for modeling \textbf{lifetime or
reliability data}. Its shape parameter (k) controls the hazard function:

\begin{itemize}
\tightlist
\item
  (k \textless{} 1): decreasing hazard (infant mortality)
\item
  (k = 1): constant hazard → reduces to \textbf{Exponential
  distribution}
\item
  (k \textgreater{} 1): increasing hazard (aging/failure over time)
\end{itemize}

\[
f(x \mid \lambda, k) =
\begin{cases}
\dfrac{k}{\lambda} \left(\dfrac{x}{\lambda}\right)^{k-1} e^{-(x/\lambda)^k}, & x \ge 0, [2mm]
0, & x < 0
\end{cases}
\]

\textbf{Args} - \emph{scale} (jnp.ndarray): The scale parameter of the
Weibull distribution. Must be positive.

\begin{itemize}
\item
  \emph{concentration} (jnp.ndarray): The concentration parameter of the
  Weibull distribution. Must be positive.
\item
  \emph{shape} (tuple): A multi-purpose argument for shaping. When
  \texttt{sample=False} (model building), this is used with
  \texttt{shape} to set the distribution's batch shape. When
  \texttt{sample=True} (direct sampling), this is used as
  \texttt{sample\_shape} to draw a raw JAX array of the given shape.
\item
  \emph{event} (int): The number of batch dimensions to reinterpret as
  event dimensions (used in model building).
\item
  \emph{mask} (jnp.ndarray, bool): Optional boolean array to mask
  observations.
\item
  \emph{create\_obj} (bool): If True, returns the raw BF distribution
  object instead of creating a sample site. This is essential for
  building complex distributions like \texttt{MixtureSameFamily}.
\item
  \emph{sample} (bool, optional): A control-flow argument. If
  \texttt{True}, the function will directly sample a raw JAX array from
  the distribution, bypassing the BF model context. If \texttt{False},
  it will create a \texttt{BF.sample} site within a model. Defaults to
  \texttt{False}.
\item
  \emph{seed} (int, optional): An integer seed for this draw when
  \texttt{sample=True}. By default (\texttt{None}) the draw comes from
  the session's random stream, set by \texttt{bf(rand\_seed=...)}, so
  consecutive draws differ but the whole sequence is reproducible. No
  effect when \texttt{sample=False}.
\item
  \emph{obs} (jnp.ndarray, optional): The observed value for this random
  variable. If provided, the sample site is conditioned on this value,
  and the function returns the observed value. If \texttt{None}, the
  site is treated as a latent (unobserved) random variable. Defaults to
  \texttt{None}.
\item
  \emph{name} (str, optional): The name of the sample site in a BF
  model. This is used to uniquely identify the random variable. Defaults
  to `x'.
\end{itemize}

\textbf{Returns} BF Weibull distribution object (for model building)
when \texttt{sample=False}. JAX array of samples drawn from the Weibull
distribution (for direct sampling) when \texttt{sample=True}. The raw BF
distribution object (for advanced use cases) when
\texttt{create\_obj=True}.

\textbf{Example Usage} from BayesForge import bf m = bf(platform=`cpu')
m.dist.weibull(scale=1.0, concentration=2.0, sample=True)

\textbf{Wrapper of}
https://num.pyro.ai/en/stable/distributions.html\#weibull

\textbf{References} *
\href{https://en.wikipedia.org/wiki/Weibull_distribution\#Probability_density_function}{Wikipedia
-- Weibull distribution}

\begin{itemize}
\tightlist
\item
  Weibull, W. (1951). \emph{A Statistical Distribution Function of Wide
  Applicability.} Journal of Applied Mechanics.
\end{itemize}

\subsection{\texorpdfstring{\texttt{m.dist.wishart}}{m.dist.wishart}}\label{wishart}

\begin{Shaded}
\begin{Highlighting}[]
\NormalTok{m.dist.wishart(}
\NormalTok{    concentration,}
\NormalTok{    scale\_matrix}\OperatorTok{=}\VariableTok{None}\NormalTok{,}
\NormalTok{    rate\_matrix}\OperatorTok{=}\VariableTok{None}\NormalTok{,}
\NormalTok{    scale\_tril}\OperatorTok{=}\VariableTok{None}\NormalTok{,}
\NormalTok{    validate\_args}\OperatorTok{=}\VariableTok{None}\NormalTok{,}
\NormalTok{    name}\OperatorTok{=}\StringTok{\textquotesingle{}x\textquotesingle{}}\NormalTok{,}
\NormalTok{    obs}\OperatorTok{=}\VariableTok{None}\NormalTok{,}
\NormalTok{    mask}\OperatorTok{=}\VariableTok{None}\NormalTok{,}
\NormalTok{    infer}\OperatorTok{=}\VariableTok{None}\NormalTok{,}
\NormalTok{    sample}\OperatorTok{=}\VariableTok{False}\NormalTok{,}
\NormalTok{    seed}\OperatorTok{=}\VariableTok{None}\NormalTok{,}
\NormalTok{    shape}\OperatorTok{=}\NormalTok{(),}
\NormalTok{    event}\OperatorTok{=}\DecValTok{0}\NormalTok{,}
\NormalTok{    create\_obj}\OperatorTok{=}\VariableTok{False}\NormalTok{,}
\NormalTok{    to\_jax}\OperatorTok{=}\VariableTok{True}\NormalTok{,}
\NormalTok{)}
\end{Highlighting}
\end{Shaded}

\textbf{Wishart} The Wishart distribution is a multivariate distribution
used to model positive definite matrices, often representing covariance
matrices. It's commonly used in Bayesian statistics and machine
learning, particularly in models involving covariance estimation.

\[
  \mathbf{X} \sim \text{Wishart}(\text{scale_matrix}=\mathbf{V}, \text{concentration}=n)
  \]

where:

\begin{itemize}
\tightlist
\item
  \begin{enumerate}
  \def\labelenumi{(\alph{enumi})}
  \setcounter{enumi}{15}
  \tightlist
  \item
    = dimension of the square matrices
  \end{enumerate}
\item
  (\mathbf{V} \in \mathbb{R}\^{}\{p \times p\}) = positive definite
  \textbf{scale matrix} (\texttt{scale\_matrix})
\item
  (n \ge p) = \textbf{degrees of freedom} (\texttt{concentration})
\end{itemize}

The Wishart distribution is the \textbf{multivariate generalization of
the chi-squared distribution}. It is commonly used as the distribution
of the \textbf{sample covariance matrix} for multivariate normal
samples.

\[
  f(\mathbf{X} \mid \mathbf{V}, n) =
  \frac{ |\mathbf{X}|^{(n-p-1)/2} \exp\big(-\frac{1}{2} \mathrm{tr}(\mathbf{V}^{-1} \mathbf{X}) \big) }
  { 2^{np/2} |\mathbf{V}|^{n/2} \Gamma_p(n/2) }, \quad \mathbf{X} \succ 0
  \] where:

\begin{itemize}
\tightlist
\item
  \(|\mathbf{X}|\) = determinant of \(\mathbf{X}\)
\item
  \(\mathrm{tr}(\cdot)\) = trace
\item
  \(\Gamma_p(\cdot)\) = multivariate Gamma function:
\end{itemize}

\textbf{Args} - \emph{concentration} (jnp.ndarray): Positive
concentration parameter analogous to the concentration of a Gamma
distribution. The concentration must be larger than the dimensionality
of the scale matrix.

\begin{itemize}
\item
  \emph{scale\_matrix} (jnp.ndarray, optional): Scale matrix analogous
  to the inverse rate of a Gamma distribution.
\item
  \emph{rate\_matrix} (jnp.ndarray, optional): Rate matrix anaologous to
  the rate of a Gamma distribution.
\item
  \emph{scale\_tril} (jnp.ndarray, optional): Cholesky decomposition of
  the :code:\texttt{scale\_matrix}.
\item
  \emph{shape} (tuple): A multi-purpose argument for shaping. When
  \texttt{sample=False} (model building), this is used with
  \texttt{shape} to set the distribution's batch shape. When
  \texttt{sample=True} (direct sampling), this is used as
  \texttt{sample\_shape} to draw a raw JAX array of the given shape.
\item
  \emph{event} (int): The number of batch dimensions to reinterpret as
  event dimensions (used in model building).
\item
  \emph{mask} (jnp.ndarray, bool, optional): Optional boolean array to
  mask observations.
\item
  \emph{create\_obj} (bool, optional): If True, returns the raw BF
  distribution object instead of creating a sample site. This is
  essential for building complex distributions like
  \texttt{MixtureSameFamily}.
\item
  \emph{sample} (bool, optional): A control-flow argument. If
  \texttt{True}, the function will directly sample a raw JAX array from
  the distribution, bypassing the BF model context. If \texttt{False},
  it will create a \texttt{BF.sample} site within a model. Defaults to
  \texttt{False}.
\item
  \emph{seed} (int, optional): An integer seed for this draw when
  \texttt{sample=True}. By default (\texttt{None}) the draw comes from
  the session's random stream, set by \texttt{bf(rand\_seed=...)}, so
  consecutive draws differ but the whole sequence is reproducible. No
  effect when \texttt{sample=False}.
\item
  \emph{obs} (jnp.ndarray, optional): The observed value for this random
  variable. If provided, the sample site is conditioned on this value,
  and the function returns the observed value. If \texttt{None}, the
  site is treated as a latent (unobserved) random variable. Defaults to
  \texttt{None}.
\item
  \emph{name} (str, optional): The name of the sample site in a BF
  model. This is used to uniquely identify the random variable. Defaults
  to `x'.
\end{itemize}

\textbf{Returns} - When \texttt{sample=False}: A BF Wishart distribution
object (for model building). - When \texttt{sample=True}: A JAX array of
samples drawn from the Wishart distribution (for direct sampling). -
When \texttt{create\_obj=True}: The raw BF distribution object (for
advanced use cases).

\textbf{Example Usage} from BayesForge import bf import jax.numpy as jnp
m = bf(platform=`cpu') m.dist.wishart(concentration=5.0,
scale\_matrix=jnp.eye(2), sample=True)

\textbf{Wrapper of}
https://num.pyro.ai/en/stable/distributions.html\#wishart

\textbf{References} *
\href{https://en.wikipedia.org/wiki/Wishart_distribution}{Wikipedia --
Wishart distribution}

\subsection{\texorpdfstring{\texttt{m.dist.wishart\_cholesky}}{m.dist.wishart\_cholesky}}\label{wishart_cholesky}

\begin{Shaded}
\begin{Highlighting}[]
\NormalTok{m.dist.wishart\_cholesky(}
\NormalTok{    concentration,}
\NormalTok{    scale\_matrix}\OperatorTok{=}\VariableTok{None}\NormalTok{,}
\NormalTok{    rate\_matrix}\OperatorTok{=}\VariableTok{None}\NormalTok{,}
\NormalTok{    scale\_tril}\OperatorTok{=}\VariableTok{None}\NormalTok{,}
\NormalTok{    validate\_args}\OperatorTok{=}\VariableTok{None}\NormalTok{,}
\NormalTok{    name}\OperatorTok{=}\StringTok{\textquotesingle{}x\textquotesingle{}}\NormalTok{,}
\NormalTok{    obs}\OperatorTok{=}\VariableTok{None}\NormalTok{,}
\NormalTok{    mask}\OperatorTok{=}\VariableTok{None}\NormalTok{,}
\NormalTok{    infer}\OperatorTok{=}\VariableTok{None}\NormalTok{,}
\NormalTok{    sample}\OperatorTok{=}\VariableTok{False}\NormalTok{,}
\NormalTok{    seed}\OperatorTok{=}\VariableTok{None}\NormalTok{,}
\NormalTok{    shape}\OperatorTok{=}\NormalTok{(),}
\NormalTok{    event}\OperatorTok{=}\DecValTok{0}\NormalTok{,}
\NormalTok{    create\_obj}\OperatorTok{=}\VariableTok{False}\NormalTok{,}
\NormalTok{    to\_jax}\OperatorTok{=}\VariableTok{True}\NormalTok{,}
\NormalTok{)}
\end{Highlighting}
\end{Shaded}

\textbf{Wishart Cholesky} * The \textbf{Wishart} distribution is a
distribution over positive definite matrices, often used as a prior for
covariance or precision matrices in multivariate normal models. * The
\textbf{Cholesky parameterization} of the Wishart (called
``wishart\_cholesky'' in Stan, for example) reparameterizes the Wishart
over its \textbf{lower (or upper) triangular Cholesky factor}. This is
useful for numerical stability and unconstrained parameterization in
Bayesian sampling frameworks. * In this parameterization, one works with
a lower‐triangular matrix \(L_W\) such that \[
\Sigma = L_W L_W^\top
\] and imposes a density over \(L_W\) corresponding to the induced
Wishart density on (\Sigma).

The probability density function (PDF) is given by:

If the dimension is \(K\), degrees of freedom \(\nu > K - 1\), and
\(L_S\) is the lower‐triangular Cholesky factor of the scale matrix
\(S\), then the log density for \(L_W\) under the
\textbf{wishart\_cholesky} distribution is:

\[
\log p(L_W \mid \nu, L_S)
= \log \bigl[ \text{Wishart}(L_W L_W^\top \mid \nu, L_S L_S^\top) \bigr] + \log \bigl| J_{f^{-1}} \bigr|
\]

where \(J_{f^{-1}}\) is the Jacobian of the transformation from \(L_W\)
to \(\Sigma = L_W L_W^\top\).

The Jacobian term (absolute log determinant) is:

\[
\log |J_{f^{-1}}| = K \log 2 ; \sum_{k=1}^K (K - k + 1),\log L_{W_{k,k}}.
\]

Thus the density includes both the usual Wishart density on \(\Sigma\)
plus this extra Jacobian (which penalizes or weights the diagonal
entries of \(L_W\)).

\textbf{Parameters} - concentration: (Tensor) Positive concentration
parameter analogous to the concentration of a \texttt{Gamma}
distribution. The concentration must be larger than the dimensionality
of the scale matrix.

\begin{itemize}
\item
  scale\_matrix: (Tensor, optional) Scale matrix analogous to the
  inverse rate of a \texttt{Gamma} distribution. If not provided,
  \texttt{rate\_matrix} or \texttt{scale\_tril} must be.
\item
  rate\_matrix: (Tensor, optional) Rate matrix anaologous to the rate of
  a \texttt{Gamma} distribution. If not provided, \texttt{scale\_matrix}
  or \texttt{scale\_tril} must be.
\item
  scale\_tril: (Tensor, optional) Cholesky decomposition of the
  :code:\texttt{scale\_matrix}. If not provided, \texttt{scale\_matrix}
  or \texttt{rate\_matrix} must be.
\item
  sample (bool, optional): A control-flow argument. If \texttt{True},
  the function will directly sample a raw JAX array from the
  distribution, bypassing the BF model context. If \texttt{False}, it
  will create a \texttt{BF.sample} site within a model. Defaults to
  \texttt{False}.
\item
  seed (int, optional): An integer seed for this draw when
  \texttt{sample=True}. By default (\texttt{None}) the draw comes from
  the session's random stream, set by \texttt{bf(rand\_seed=...)}, so
  consecutive draws differ but the whole sequence is reproducible. No
  effect when \texttt{sample=False}.
\item
  obs (jnp.ndarray, optional): The observed value for this random
  variable. If provided, the sample site is conditioned on this value,
  and the function returns the observed value. If \texttt{None}, the
  site is treated as a latent (unobserved) random variable. Defaults to
  \texttt{None}.
\item
  name (str, optional): The name of the sample site in a BF model. This
  is used to uniquely identify the random variable. Defaults to `x'.
\end{itemize}

\textbf{References} -
\href{https://en.wikipedia.org/wiki/Wishart_distribution}{Wishart
distribution} \href{https://statproofbook.github.io/P/matn-pdf.html}{2}:
https://mc-stan.org/docs/2\_32/functions-reference/wishart-cholesky-distribution.html''28.2
Wishart distribution, Cholesky Parameterization - Stan''

\subsection{\texorpdfstring{\texttt{m.dist.zero\_inflated\_distribution}}{m.dist.zero\_inflated\_distribution}}\label{zero_inflated_distribution}

\begin{Shaded}
\begin{Highlighting}[]
\NormalTok{m.dist.zero\_inflated\_distribution(}
\NormalTok{    base\_dist,}
\NormalTok{    gate}\OperatorTok{=}\VariableTok{None}\NormalTok{,}
\NormalTok{    gate\_logits}\OperatorTok{=}\VariableTok{None}\NormalTok{,}
\NormalTok{    validate\_args}\OperatorTok{=}\VariableTok{None}\NormalTok{,}
\NormalTok{    name}\OperatorTok{=}\StringTok{\textquotesingle{}x\textquotesingle{}}\NormalTok{,}
\NormalTok{    obs}\OperatorTok{=}\VariableTok{None}\NormalTok{,}
\NormalTok{    mask}\OperatorTok{=}\VariableTok{None}\NormalTok{,}
\NormalTok{    infer}\OperatorTok{=}\VariableTok{None}\NormalTok{,}
\NormalTok{    sample}\OperatorTok{=}\VariableTok{False}\NormalTok{,}
\NormalTok{    seed}\OperatorTok{=}\VariableTok{None}\NormalTok{,}
\NormalTok{    shape}\OperatorTok{=}\NormalTok{(),}
\NormalTok{    event}\OperatorTok{=}\DecValTok{0}\NormalTok{,}
\NormalTok{    create\_obj}\OperatorTok{=}\VariableTok{False}\NormalTok{,}
\NormalTok{    to\_jax}\OperatorTok{=}\VariableTok{True}\NormalTok{,}
\NormalTok{)}
\end{Highlighting}
\end{Shaded}

\textbf{Generic Zero Inflated} A Zero-Inflated distribution combines a
base distribution with a Bernoulli distribution to model data with an
excess of zero values. It assumes that each observation is either drawn
from the base distribution or is a zero with probability determined by
the Bernoulli distribution (the ``gate''). A \textbf{zero-inflated
distribution} arises when you take a random variable \(X\) that
originally has some distribution (with PMF \(f_X(x)\) ) \emph{and} you
add extra mass at zero. In other words you only observe \(X\) when
\(X \ge 0\) (if count) and you assume there are more zeros than what the
original distribution predicts --- so you mix in a point-mass at zero. A
zero-inflated model assumes two processes: a ``structural zero'' process
with probability \(\pi\), and another process (the base
count/distribution) with probability \(1 - \pi\), which itself may
generate zeros \emph{or} non-zeros.

\begin{itemize}
\tightlist
\item
  When / Why use it:

  \begin{itemize}
  \tightlist
  \item
    \textbf{Data collection constraints}: You can only observe values in
    a process when some condition is met; others are automatically zero
    (for example, when the event cannot happen for some units).
  \item
    \textbf{Mixtures of processes}: Some observations are ``structural
    zeros'' (no risk) and others follow a regular count process (with
    risk).
  \item
    \textbf{Overdispersion \& excess zeros}: If you try a standard count
    distribution (Poisson, Negative Binomial) and you observe many more
    zeros than predicted (given the non-zero counts), a zero-inflated
    alternative may fit better.
  \end{itemize}
\item
  The general PMF form is (for discrete count models): \[
    P(X=0) = \pi + (1 - \pi),P_{\text{base}}(0),
    \quad
    P(X=k) = (1 - \pi),P_{\text{base}}(k), ; k > 0.
    \]
\end{itemize}

\textbf{Args} - \emph{base\_dist} (Distribution): The base distribution
to be zero-inflated (e.g., Poisson, NegativeBinomial).

\begin{itemize}
\item
  \emph{gate} (jnp.ndarray, optional): Probability of extra zeros
  (between 0 and 1).
\item
  \emph{gate\_logits} (jnp.ndarray, optional): Log-odds of extra zeros.
\item
  \emph{validate\_args} (bool, optional): Whether to validate parameter
  values. Defaults to None.
\item
  \emph{shape} (tuple): A multi-purpose argument for shaping. When
  \texttt{sample=False} (model building), this is used with
  \texttt{shape} to set the distribution's batch shape. When
  \texttt{sample=True} (direct sampling), this is used as
  \texttt{sample\_shape} to draw a raw JAX array of the given shape.
\item
  \emph{event} (int): The number of batch dimensions to reinterpret as
  event dimensions (used in model building).
\item
  \emph{mask} (jnp.ndarray, bool): Optional boolean array to mask
  observations.
\item
  \emph{create\_obj} (bool): If True, returns the raw BF distribution
  object instead of creating a sample site. This is essential for
  building complex distributions like \texttt{MixtureSameFamily}.
\item
  \emph{sample} (bool, optional): A control-flow argument. If
  \texttt{True}, the function will directly sample a raw JAX array from
  the distribution, bypassing the BF model context. If \texttt{False},
  it will create a \texttt{BF.sample} site within a model. Defaults to
  \texttt{False}.
\item
  \emph{seed} (int, optional): An integer seed for this draw when
  \texttt{sample=True}. By default (\texttt{None}) the draw comes from
  the session's random stream, set by \texttt{bf(rand\_seed=...)}, so
  consecutive draws differ but the whole sequence is reproducible. No
  effect when \texttt{sample=False}.
\item
  \emph{obs} (jnp.ndarray, optional): The observed value for this random
  variable. If provided, the sample site is conditioned on this value,
  and the function returns the observed value. If \texttt{None}, the
  site is treated as a latent (unobserved) random variable. Defaults to
  \texttt{None}.
\item
  \emph{name} (str, optional): The name of the sample site in a BF
  model. This is used to uniquely identify the random variable. Defaults
  to `x'.
\end{itemize}

\textbf{Returns} - When \texttt{sample=False}: A BF
ZeroInflatedDistribution distribution object (for model building). -
When \texttt{sample=True}: A JAX array of samples drawn from the
ZeroInflatedDistribution distribution (for direct sampling). - When
\texttt{create\_obj=True}: The raw BF distribution object (for advanced
use cases).

\textbf{Example Usage} from BayesForge import bf m = bf(platform=`cpu')
m.dist.zero\_inflated\_distribution(base\_dist=m.dist.poisson(rate=5,
create\_obj = True), gate = 0.3, sample=True)

\textbf{Wrapper of:
https://num.pyro.ai/en/stable/distributions.html\#zeroinflateddistribution}
\textbf{References} -
\href{https://en.wikipedia.org/wiki/Zero-inflated_model}{Zero-inflated
model} -
\href{https://larmarange.github.io/guide-R/analyses_avancees/modeles-zero-inflated.html}{48
Modèles de comptage zero-inflated et hurdle -- guide-R}

\subsection{\texorpdfstring{\texttt{m.dist.zero\_inflated\_negative\_binomial2}}{m.dist.zero\_inflated\_negative\_binomial2}}\label{zero_inflated_negative_binomial2}

\begin{Shaded}
\begin{Highlighting}[]
\NormalTok{m.dist.zero\_inflated\_negative\_binomial2(}
\NormalTok{    mean,}
\NormalTok{    concentration,}
\NormalTok{    gate}\OperatorTok{=}\VariableTok{None}\NormalTok{,}
\NormalTok{    gate\_logits}\OperatorTok{=}\VariableTok{None}\NormalTok{,}
\NormalTok{    validate\_args}\OperatorTok{=}\VariableTok{None}\NormalTok{,}
\NormalTok{    name}\OperatorTok{=}\StringTok{\textquotesingle{}x\textquotesingle{}}\NormalTok{,}
\NormalTok{    obs}\OperatorTok{=}\VariableTok{None}\NormalTok{,}
\NormalTok{    mask}\OperatorTok{=}\VariableTok{None}\NormalTok{,}
\NormalTok{    infer}\OperatorTok{=}\VariableTok{None}\NormalTok{,}
\NormalTok{    sample}\OperatorTok{=}\VariableTok{False}\NormalTok{,}
\NormalTok{    seed}\OperatorTok{=}\VariableTok{None}\NormalTok{,}
\NormalTok{    shape}\OperatorTok{=}\NormalTok{(),}
\NormalTok{    event}\OperatorTok{=}\DecValTok{0}\NormalTok{,}
\NormalTok{    create\_obj}\OperatorTok{=}\VariableTok{False}\NormalTok{,}
\NormalTok{    to\_jax}\OperatorTok{=}\VariableTok{True}\NormalTok{,}
\NormalTok{)}
\end{Highlighting}
\end{Shaded}

\textbf{Zero-Inflated Negative Binomial} A Zero-Inflated Negative
Binomial distribution is used for count data that exhibit \textbf{both}
(a) over-dispersion relative to a Poisson (i.e., variance \textgreater{}
mean) \emph{and} (b) an excess of zero counts beyond what a standard
Negative Binomial would predict. It assumes two latent processes:

\begin{enumerate}
\def\labelenumi{\arabic{enumi}.}
\tightlist
\item
  With probability \(\pi\) (sometimes denoted \(\psi\) or
  ``zero-inflation probability'') you are in a ``structural zero'' state
  → you observe a zero.
\item
  With probability \(1 - \pi\), you come from a regular Negative
  Binomial distribution (with parameters e.g.~mean \(\mu\) and
  dispersion parameter \(\alpha\) or size/r parameter) and then you
  might observe zero or a positive count.
\end{enumerate}

Thus the model is a mixture of a point‐mass at zero + a Negative
Binomial for counts.

This distribution combines a Negative Binomial distribution with a
binary gate variable. Observations are either drawn from the Negative
Binomial distribution with probability (1 - gate) or are treated as zero
with probability `gate'. This models data with excess zeros compared to
what a standard Negative Binomial distribution would predict.

Let \(X\) denote the count random variable, support \({0,1,2,\dots}\).
Denote:

\begin{itemize}
\tightlist
\item
  \(\pi\) = probability of being in the ``always zero'' (structural
  zero) process, with \$ 0 \le \pi \le 1\$.
\item
  The count process is NB with parameters: mean \(\mu>0\) and
  dispersion/shape parameter \(\alpha>0\) (or equivalently size
  parameter (r)). Then:
\end{itemize}

\[
  P(X = 0) = \pi ;+; (1 - \pi); P_{\text{NB}}(0;\mu,\alpha),
  \] \[
  P(X = k) = (1 - \pi); P_{\text{NB}}(k;\mu,\alpha), \quad k = 1,2,3,\dots
  \]

Here \(P_{\text{NB}}(k;\mu,\alpha)\) is the PMF of the negative binomial
distribution for count \(k\).

\textbf{Args} - \emph{mean} (jnp.ndarray or float): The mean of the
Negative Binomial 2 distribution.

\begin{itemize}
\item
  \emph{concentration} (jnp.ndarray or float): The concentration
  parameter of the Negative Binomial 2 distribution.
\item
  \emph{shape} (tuple): A multi-purpose argument for shaping. When
  \texttt{sample=False} (model building), this is used with
  \texttt{shape} to set the distribution's batch shape. When
  \texttt{sample=True} (direct sampling), this is used as
  \texttt{sample\_shape} to draw a raw JAX array of the given shape.
\item
  \emph{event} (int): The number of batch dimensions to reinterpret as
  event dimensions (used in model building).
\item
  \emph{mask} (jnp.ndarray, bool): Optional boolean array to mask
  observations.
\item
  \emph{create\_obj} (bool): If True, returns the raw BF distribution
  object instead of creating a sample site. This is essential for
  building complex distributions like \texttt{MixtureSameFamily}.
\item
  \emph{sample} (bool, optional): A control-flow argument. If
  \texttt{True}, the function will directly sample a raw JAX array from
  the distribution, bypassing the BF model context. If \texttt{False},
  it will create a \texttt{BF.sample} site within a model. Defaults to
  \texttt{False}.
\item
  \emph{seed} (int, optional): An integer seed for this draw when
  \texttt{sample=True}. By default (\texttt{None}) the draw comes from
  the session's random stream, set by \texttt{bf(rand\_seed=...)}, so
  consecutive draws differ but the whole sequence is reproducible. No
  effect when \texttt{sample=False}.
\item
  \emph{obs} (jnp.ndarray, optional): The observed value for this random
  variable. If provided, the sample site is conditioned on this value,
  and the function returns the observed value. If \texttt{None}, the
  site is treated as a latent (unobserved) random variable. Defaults to
  \texttt{None}.
\item
  \emph{name} (str, optional): The name of the sample site in a BF
  model. This is used to uniquely identify the random variable. Defaults
  to `x'.
\end{itemize}

\textbf{Returns} - When \texttt{sample=False}: A BF
ZeroInflatedNegativeBinomial2 distribution object (for model building).
- When \texttt{sample=True}: A JAX array of samples drawn from the
ZeroInflatedNegativeBinomial2 distribution (for direct sampling). - When
\texttt{create\_obj=True}: The raw BF distribution object (for advanced
use cases).

\textbf{Example Usage} from BayesForge import bf m = bf(platform=`cpu')
m.dist.zero\_inflated\_negative\_binomial2(mean=2.0, concentration=3.0,
gate = 0.3, sample=True)

\textbf{Wrapper of:
https://num.pyro.ai/en/stable/distributions.html\#zeroinflatednegativebinomial2}
\textbf{References} -
\href{https://www.ewadirect.com/proceedings/tns/article/view/27624}{Social
Networks Count Data: Negative Binomial Distributions and Extensions
Versus Poisson and Bernoulli Models} -
\href{https://pmc.ncbi.nlm.nih.gov/articles/PMC7880198/}{Bayesian
Zero-Inflated Negative Binomial Regression Based on Pólya-Gamma Mixtures
- PMC} -
\href{https://search.r-project.org/CRAN/refmans/emdbook/html/dzinbinom.html}{R:
Zero-inflated negative binomial distribution}

\subsection{\texorpdfstring{\texttt{m.dist.zero\_inflated\_poisson}}{m.dist.zero\_inflated\_poisson}}\label{zero_inflated_poisson}

\begin{Shaded}
\begin{Highlighting}[]
\NormalTok{m.dist.zero\_inflated\_poisson(}
\NormalTok{    gate,}
\NormalTok{    rate}\OperatorTok{=}\FloatTok{1.0}\NormalTok{,}
\NormalTok{    validate\_args}\OperatorTok{=}\VariableTok{None}\NormalTok{,}
\NormalTok{    name}\OperatorTok{=}\StringTok{\textquotesingle{}x\textquotesingle{}}\NormalTok{,}
\NormalTok{    obs}\OperatorTok{=}\VariableTok{None}\NormalTok{,}
\NormalTok{    mask}\OperatorTok{=}\VariableTok{None}\NormalTok{,}
\NormalTok{    infer}\OperatorTok{=}\VariableTok{None}\NormalTok{,}
\NormalTok{    sample}\OperatorTok{=}\VariableTok{False}\NormalTok{,}
\NormalTok{    seed}\OperatorTok{=}\VariableTok{None}\NormalTok{,}
\NormalTok{    shape}\OperatorTok{=}\NormalTok{(),}
\NormalTok{    event}\OperatorTok{=}\DecValTok{0}\NormalTok{,}
\NormalTok{    create\_obj}\OperatorTok{=}\VariableTok{False}\NormalTok{,}
\NormalTok{    to\_jax}\OperatorTok{=}\VariableTok{True}\NormalTok{,}
\NormalTok{)}
\end{Highlighting}
\end{Shaded}

\textbf{A Zero Inflated Poisson} The Zero-Inflated Poisson distribution
is a discrete count‐distribution designed for data with \emph{more
zeros} than would be expected under a standard Poisson. Essentially, it
assumes two underlying processes:

\begin{itemize}
\tightlist
\item
  With probability \(\pi\) you are in a ``structural zero'' state (i.e.,
  you automatically get a zero count).
\item
  With probability \(1 - \pi\) you draw from a standard Poisson
  distribution with parameter \(\lambda\).
\end{itemize}

This results in a mixture distribution that places more mass at zero
than a Poisson alone would. It's widely used in, for instance, ecology
(species counts with many zeros), insurance/claims problems, and any
count‐data setting with excess zeros.

The probability density function (PDF) is given by:

Let \(X\) be a random variable following a Zero-Inflated Poisson with
parameters

\begin{itemize}
\tightlist
\item
  \(\pi\in[0,1]\): zero-inflation probability (chance of being
  structural zero)
\item
  \(\lambda>0\): mean of the Poisson component
\end{itemize}

Then: \[
P(X = k) =
\begin{cases}
\pi ;+; (1 - \pi),e^{-\lambda}, & \text{if } k = 0, [6pt]
(1 - \pi),\dfrac{\lambda^k e^{-\lambda}}{k!}, & \text{if } k = 1,2,3,\dots
\end{cases}
\]\\
(Here the Poisson component contributes also a mass at zero of
\((1 - \pi)e^{-\lambda}\), plus the extra \(\pi\) mass.)

\textbf{Args} - \emph{gate} (jnp.ndarray): The probability of observing
a zero.

\begin{itemize}
\item
  \emph{rate} (jnp.ndarray): The rate parameter of the underlying
  Poisson distribution.
\item
  \emph{shape} (tuple): A multi-purpose argument for shaping. When
  \texttt{sample=False} (model building), this is used with
  \texttt{shape} to set the distribution's batch shape. When
  \texttt{sample=True} (direct sampling), this is used as
  \texttt{sample\_shape} to draw a raw JAX array of the given shape.
\item
  \emph{event} (int): The number of batch dimensions to reinterpret as
  event dimensions (used in model building).
\item
  \emph{mask} (jnp.ndarray, bool): Optional boolean array to mask
  observations.
\item
  \emph{create\_obj} (bool): If True, returns the raw BF distribution
  object instead of creating a sample site. This is essential for
  building complex distributions like \texttt{MixtureSameFamily}.
\item
  \emph{sample} (bool, optional): A control-flow argument. If
  \texttt{True}, the function will directly sample a raw JAX array from
  the distribution, bypassing the BF model context. If \texttt{False},
  it will create a \texttt{BF.sample} site within a model. Defaults to
  \texttt{False}.
\item
  \emph{seed} (int, optional): An integer seed for this draw when
  \texttt{sample=True}. By default (\texttt{None}) the draw comes from
  the session's random stream, set by \texttt{bf(rand\_seed=...)}, so
  consecutive draws differ but the whole sequence is reproducible. No
  effect when \texttt{sample=False}.
\item
  \emph{obs} (jnp.ndarray, optional): The observed value for this random
  variable. If provided, the sample site is conditioned on this value,
  and the function returns the observed value. If \texttt{None}, the
  site is treated as a latent (unobserved) random variable. Defaults to
  \texttt{None}.
\item
  \emph{name} (str, optional): The name of the sample site in a BF
  model. This is used to uniquely identify the random variable. Defaults
  to `x'.
\end{itemize}

\textbf{Returns} BF ZeroInflatedPoisson distribution object (when
\texttt{sample=False}). JAX array of samples drawn from the
ZeroInflatedPoisson distribution (when \texttt{sample=True}). The raw BF
distribution object (when \texttt{create\_obj=True}).

\textbf{Example Usage} from BayesForge import bf m = bf(platform=`cpu')
m.dist.zero\_inflated\_poisson(gate = 0.3, rate=2.0, sample=True)

\textbf{Wrapper of:
https://num.pyro.ai/en/stable/distributions.html\#zeroinflatedpoisson}
\textbf{References} * https://en.wikipedia.org/wiki/Zero-inflated\_model
``Zero-inflated model'' *
https://www.statisticshowto.com/zero-inflated-poisson-distribution
``Zero-Inflated Poisson distribution - Statistics How To''

\subsection{\texorpdfstring{\texttt{m.dist.zero\_sum\_normal}}{m.dist.zero\_sum\_normal}}\label{zero_sum_normal}

\begin{Shaded}
\begin{Highlighting}[]
\NormalTok{m.dist.zero\_sum\_normal(}
\NormalTok{    scale,}
\NormalTok{    event\_shape,}
\NormalTok{    validate\_args}\OperatorTok{=}\VariableTok{None}\NormalTok{,}
\NormalTok{    name}\OperatorTok{=}\StringTok{\textquotesingle{}x\textquotesingle{}}\NormalTok{,}
\NormalTok{    obs}\OperatorTok{=}\VariableTok{None}\NormalTok{,}
\NormalTok{    mask}\OperatorTok{=}\VariableTok{None}\NormalTok{,}
\NormalTok{    infer}\OperatorTok{=}\VariableTok{None}\NormalTok{,}
\NormalTok{    sample}\OperatorTok{=}\VariableTok{False}\NormalTok{,}
\NormalTok{    seed}\OperatorTok{=}\VariableTok{None}\NormalTok{,}
\NormalTok{    shape}\OperatorTok{=}\NormalTok{(),}
\NormalTok{    event}\OperatorTok{=}\DecValTok{0}\NormalTok{,}
\NormalTok{    create\_obj}\OperatorTok{=}\VariableTok{False}\NormalTok{,}
\NormalTok{    to\_jax}\OperatorTok{=}\VariableTok{True}\NormalTok{,}
\NormalTok{)}
\end{Highlighting}
\end{Shaded}

\textbf{Zero Sum Normal} A \textbf{zero-sum normal} is a variant of a
multivariate normal in which one (or more) linear constraint(s) force
certain components to \textbf{sum to zero}. In practice, it's used to
model vectors of random effects (e.g.~in hierarchical models) where the
effects are constrained to sum to zero (to avoid overparameterization or
enforce identifiability).

\[
  ZSN(\sigma) = N(0, \sigma^2 (I - \tfrac{1}{n}J)) \\
  \text{where} \ ~ J_{ij} = 1 \ ~ \text{and} \\
  n = \text{number of zero-sum axes}
  \]

where \(J\)is the all-ones matrix and \(n\) is the number of constrained
elements. So the covariance is \(\sigma^2 (I - \tfrac{1}{n} J)\), which
ensures that the sum of the components is 0 (since
\(J \mathbf{x} = \mathbf{1} \sum_i x_i\)).

\textbf{Args} - \emph{scale} (array\_like): Standard deviation of the
underlying normal distribution before the zerosum constraint is
enforced.

\begin{itemize}
\item
  \emph{event\_shape} (tuple): Shape of the event dimension of the
  distribution.
\item
  \emph{shape} (tuple): A multi-purpose argument for shaping. When
  \texttt{sample=False} (model building), this is used with
  \texttt{shape} to set the distribution's batch shape. When
  \texttt{sample=True} (direct sampling), this is used as
  \texttt{sample\_shape} to draw a raw JAX array of the given shape.
\item
  \emph{event} (int): The number of batch dimensions to reinterpret as
  event dimensions (used in model building).
\item
  \emph{mask} (jnp.ndarray, bool): Optional boolean array to mask
  observations.
\item
  \emph{create\_obj} (bool): If True, returns the raw BF distribution
  object instead of creating a sample site. This is essential for
  building complex distributions like \texttt{MixtureSameFamily}.
\item
  \emph{sample} (bool, optional): A control-flow argument. If
  \texttt{True}, the function will directly sample a raw JAX array from
  the distribution, bypassing the BF model context. If \texttt{False},
  it will create a \texttt{BF.sample} site within a model. Defaults to
  \texttt{False}.
\item
  \emph{seed} (int, optional): An integer seed for this draw when
  \texttt{sample=True}. By default (\texttt{None}) the draw comes from
  the session's random stream, set by \texttt{bf(rand\_seed=...)}, so
  consecutive draws differ but the whole sequence is reproducible. No
  effect when \texttt{sample=False}.
\item
  \emph{obs} (jnp.ndarray, optional): The observed value for this random
  variable. If provided, the sample site is conditioned on this value,
  and the function returns the observed value. If \texttt{None}, the
  site is treated as a latent (unobserved) random variable. Defaults to
  \texttt{None}.
\item
  \emph{name} (str, optional): The name of the sample site in a BF
  model. This is used to uniquely identify the random variable. Defaults
  to `x'.
\end{itemize}

\textbf{Returns} - When \texttt{sample=False}: A BF ZeroSumNormal
distribution object (for model building).

\begin{itemize}
\item
  When \texttt{sample=True}: A JAX array of samples drawn from the
  ZeroSumNormal distribution (for direct sampling).
\item
  When \texttt{create\_obj=True}: The raw BF distribution object (for
  advanced use cases).

  \textbf{Example Usage} from BayesForge import bf m =
  bf(platform=`cpu') m.dist.zero\_sum\_normal(scale=1.0, event\_shape =
  (2,), sample = True)

  \textbf{Wrapper of}
  https://num.pyro.ai/en/stable/distributions.html\#zerosumnormal

  \textbf{References}
  https://www.pymc.io/projects/docs/en/latest/api/distributions/generated/pymc.ZeroSumNormal.html
\end{itemize}




\end{document}
